% !TeX program = xelatex
% Standalone English volume. Use XeLaTeX until the references stabilize.
\documentclass[a4paper,11pt]{article}
\usepackage[T1]{fontenc}
\usepackage{lmodern}
\usepackage[english]{babel}
\usepackage[left=26mm,right=26mm,top=25mm,bottom=25mm,headheight=15pt,headsep=7mm,footskip=11mm]{geometry}
\usepackage{amsmath,amssymb,mathtools,amsthm}
\newtheorem{proposition}{Proposition}[section]
\newtheorem{lemma}[proposition]{Lemma}
\newtheorem{theorem}[proposition]{Theorem}
\newtheorem{corollary}[proposition]{Corollary}
\newtheorem{maintheorem}[proposition]{Main Theorem}
\theoremstyle{definition}
\newtheorem{definition}[proposition]{Definition}
\theoremstyle{remark}
\newtheorem{remark}[proposition]{Remark}
\makeatletter
\renewcommand*\l@section{\@dottedtocline{1}{0em}{2.4em}}
\renewcommand*\l@subsection{\@dottedtocline{2}{1.45em}{3.45em}}
\renewcommand*\l@subsubsection{\@dottedtocline{3}{3.8em}{4.55em}}
\makeatother
\usepackage{xcolor,fancyhdr,titlesec,needspace,xurl}
\usepackage{longtable,multicol}
\usepackage{hyperref}
\definecolor{ink}{HTML}{1A3542}
\definecolor{muted}{HTML}{4F606A}
\hypersetup{unicode=true,colorlinks=true,linkcolor=ink,filecolor=ink,urlcolor=ink,citecolor=ink,bookmarksnumbered=true,
 pdfdisplaydoctitle=true,pdflang={en},
 pdftitle={The generalized Riemann hypothesis for Dirichlet L-functions. Edition 930, Volume XVII},
 pdfsubject={A natural-language reconstruction of the proposed proof},
 pdfauthor={A mathematical reconstruction}}
\titleformat{\section}{\Large\bfseries\color{ink}}{\thesection}{.65em}{}
\titlespacing*{\section}{0pt}{.8em}{.65em}
\setlength{\parindent}{1.2em}
\linespread{1.045}
\setlength{\parskip}{.22em plus .05em minus .03em}
\setlength{\emergencystretch}{2em}
\setcounter{topnumber}{2}
\setcounter{tocdepth}{3}\setcounter{secnumdepth}{3}
\widowpenalty=4500
\clubpenalty=4500
\pagestyle{fancy}\fancyhf{}
\fancyhead[L]{\small\color{muted}Dirichlet GRH \textperiodcentered\ Volume I}
\fancyhead[R]{\small\color{muted}Dirichlet zeros and power intervals}
\fancyfoot[C]{\small\thepage}
\renewcommand{\headrulewidth}{.3pt}
\newcommand{\lead}[1]{\par\smallskip\noindent{\bfseries\color{ink}#1.}\enspace}
\newcommand{\sour}[2]{\unskip\penalty150\hspace{.35em}{\small\hyperlink{ref#1}{[#2]}}}
\titleformat{\subsection}{\large\bfseries}{\thesubsection}{.6em}{}
\titlespacing*{\subsection}{0pt}{1.15em plus .15em}{.4em}
\setlength{\abovedisplayskip}{9pt plus 2pt minus 3pt}
\setlength{\belowdisplayskip}{9pt plus 2pt minus 3pt}
\newcommand{\RR}{\mathbb R}\newcommand{\CC}{\mathbb C}\newcommand{\QQ}{\mathbb Q}
\newcommand{\ZZ}{\mathbb Z}\newcommand{\ord}{\operatorname{ord}}
\newcommand{\Frob}{\operatorname{Frob}}\newcommand{\Gal}{\operatorname{Gal}}
\newcommand{\len}{\operatorname{length}}\newcommand{\Res}{\operatorname*{Res}}
\newcommand{\wmean}{\operatorname{av}}\newcommand{\Ker}{\operatorname{ker}}
\newcommand{\Lh}{\mathcal L}\newcommand{\Gg}{\mathcal G}
\newcommand{\real}{\operatorname{Re}}\newcommand{\imag}{\operatorname{Im}}
\newcommand{\abs}[1]{\left|#1\right|}\newcommand{\eps}{\varepsilon}
\newcommand{\anchoritem}[2]{\par\Needspace{2.8\baselineskip}\hypertarget{ref#1}{}\noindent{\bfseries\color{ink}#2.}\enspace}
\newcommand{\srcloc}[2]{{\small\nolinkurl{#1.lean}}, line~#2}
\numberwithin{equation}{section}
\raggedbottom

\newcommand{\Sel}{\mathcal S}
\newcommand{\Cl}{\operatorname{Cl}}
\newcommand{\Aut}{\operatorname{Aut}}
\newcommand{\GL}{\operatorname{GL}}
\newcommand{\OO}{\mathcal O}
\hypersetup{bookmarksdepth=3}


\addto\captionsenglish{\renewcommand{\partname}{Volume}}


\makeatletter
\expandafter\def\csname r@cor:ang153\endcsname{{106.18}{70}{}{proposition.106.18}{Dirichlet-GRH-v930-Volume-XX.pdf}}
\expandafter\def\csname r@cor:cost128\endcsname{{100.17}{79}{}{proposition.100.17}{Dirichlet-GRH-v930-Volume-XIX.pdf}}
\expandafter\def\csname r@cor:fibres89\endcsname{{92.10}{9}{}{proposition.92.10}{Dirichlet-GRH-v930-Volume-XVIII.pdf}}
\expandafter\def\csname r@cor:normaldegree347\endcsname{{127.2}{27}{}{proposition.127.2}{Dirichlet-GRH-v930-Volume-XXIV.pdf}}
\expandafter\def\csname r@cor:original200\endcsname{{83.7}{57}{}{proposition.83.7}{Dirichlet-GRH-v930-Volume-XVI.pdf}}
\expandafter\def\csname r@cor:real91\endcsname{{94.10}{45}{}{proposition.94.10}{Dirichlet-GRH-v930-Volume-XVIII.pdf}}
\expandafter\def\csname r@cor:rel155\endcsname{{106.21}{73}{}{proposition.106.21}{Dirichlet-GRH-v930-Volume-XX.pdf}}
\expandafter\def\csname r@cor:trace235\endcsname{{110.7}{43}{}{proposition.110.7}{Dirichlet-GRH-v930-Volume-XXI.pdf}}
\expandafter\def\csname r@eq:gap\endcsname{{6.3}{26}{}{equation.6.3}{Dirichlet-GRH-v930-Volume-I.pdf}}
\expandafter\def\csname r@eq:invar148\endcsname{{20.10}{32}{}{equation.20.10}{Dirichlet-GRH-v930-Volume-IV.pdf}}
\expandafter\def\csname r@eq:letters111\endcsname{{9.1}{117}{}{equation.9.1}{Dirichlet-GRH-v930-Volume-I.pdf}}
\expandafter\def\csname r@eq:main127\endcsname{{92.17}{20}{}{equation.92.17}{Dirichlet-GRH-v930-Volume-XVIII.pdf}}
\expandafter\def\csname r@eq:pair55\endcsname{{65.1}{21}{}{equation.65.1}{Dirichlet-GRH-v930-Volume-XIII.pdf}}
\expandafter\def\csname r@eq:period94\endcsname{{121.7}{51}{}{equation.121.7}{Dirichlet-GRH-v930-Volume-XXIII.pdf}}
\expandafter\def\csname r@eq:real127\endcsname{{92.15}{19}{}{equation.92.15}{Dirichlet-GRH-v930-Volume-XVIII.pdf}}
\expandafter\def\csname r@eq:region198\endcsname{{20.12}{33}{}{equation.20.12}{Dirichlet-GRH-v930-Volume-IV.pdf}}
\expandafter\def\csname r@eq:relevant133\endcsname{{98.2}{3}{}{equation.98.2}{Dirichlet-GRH-v930-Volume-XIX.pdf}}
\expandafter\def\csname r@eq:seed127\endcsname{{92.16}{20}{}{equation.92.16}{Dirichlet-GRH-v930-Volume-XVIII.pdf}}
\expandafter\def\csname r@eq:tail114\endcsname{{19.15}{17}{}{equation.19.15}{Dirichlet-GRH-v930-Volume-IV.pdf}}
\expandafter\def\csname r@eq:window\endcsname{{{W}}{2}{}{AMS.2}{Dirichlet-GRH-v930-Volume-I.pdf}}
\expandafter\def\csname r@lem:auxiliaryroots339\endcsname{{99.16}{57}{}{proposition.99.16}{Dirichlet-GRH-v930-Volume-XIX.pdf}}
\expandafter\def\csname r@lem:branch243\endcsname{{95.1}{65}{}{proposition.95.1}{Dirichlet-GRH-v930-Volume-XVIII.pdf}}
\expandafter\def\csname r@lem:cas96\endcsname{{121.25}{55}{}{proposition.121.25}{Dirichlet-GRH-v930-Volume-XXIII.pdf}}
\expandafter\def\csname r@lem:coefficient199\endcsname{{83.1}{52}{}{proposition.83.1}{Dirichlet-GRH-v930-Volume-XVI.pdf}}
\expandafter\def\csname r@lem:covol127\endcsname{{92.23}{19}{}{proposition.92.23}{Dirichlet-GRH-v930-Volume-XVIII.pdf}}
\expandafter\def\csname r@lem:dimension239\endcsname{{94.22}{55}{}{proposition.94.22}{Dirichlet-GRH-v930-Volume-XVIII.pdf}}
\expandafter\def\csname r@lem:flag272\endcsname{{94.26}{58}{}{proposition.94.26}{Dirichlet-GRH-v930-Volume-XVIII.pdf}}
\expandafter\def\csname r@lem:flagmarked264\endcsname{{102.6}{94}{}{proposition.102.6}{Dirichlet-GRH-v930-Volume-XIX.pdf}}
\expandafter\def\csname r@lem:flagstep338\endcsname{{71.23}{39}{}{proposition.71.23}{Dirichlet-GRH-v930-Volume-XIV.pdf}}
\expandafter\def\csname r@lem:lattices247\endcsname{{106.1}{52}{}{proposition.106.1}{Dirichlet-GRH-v930-Volume-XX.pdf}}
\expandafter\def\csname r@lem:mass209\endcsname{{93.7}{26}{}{proposition.93.7}{Dirichlet-GRH-v930-Volume-XVIII.pdf}}
\expandafter\def\csname r@lem:mass32\endcsname{{121.17}{48}{}{proposition.121.17}{Dirichlet-GRH-v930-Volume-XXIII.pdf}}
\expandafter\def\csname r@lem:mass90\endcsname{{94.7}{42}{}{proposition.94.7}{Dirichlet-GRH-v930-Volume-XVIII.pdf}}
\expandafter\def\csname r@lem:moment271\endcsname{{106.8}{61}{}{proposition.106.8}{Dirichlet-GRH-v930-Volume-XX.pdf}}
\expandafter\def\csname r@lem:oneweight355\endcsname{{120.10}{21}{}{proposition.120.10}{Dirichlet-GRH-v930-Volume-XXIII.pdf}}
\expandafter\def\csname r@lem:project271\endcsname{{71.5}{19}{}{proposition.71.5}{Dirichlet-GRH-v930-Volume-XIV.pdf}}
\expandafter\def\csname r@lem:realflag265\endcsname{{105.1}{32}{}{proposition.105.1}{Dirichlet-GRH-v930-Volume-XX.pdf}}
\expandafter\def\csname r@lem:root94\endcsname{{121.20}{51}{}{proposition.121.20}{Dirichlet-GRH-v930-Volume-XXIII.pdf}}
\expandafter\def\csname r@lem:shuffle107\endcsname{{9.10}{115}{}{proposition.9.10}{Dirichlet-GRH-v930-Volume-I.pdf}}
\expandafter\def\csname r@lem:spread111\endcsname{{9.12}{117}{}{proposition.9.12}{Dirichlet-GRH-v930-Volume-I.pdf}}
\expandafter\def\csname r@lem:terminalclass259\endcsname{{111.2}{58}{}{proposition.111.2}{Dirichlet-GRH-v930-Volume-XXI.pdf}}
\expandafter\def\csname r@lem:zero133\endcsname{{98.1}{3}{}{proposition.98.1}{Dirichlet-GRH-v930-Volume-XIX.pdf}}
\expandafter\def\csname r@nav:I-28\endcsname{{6.6}{32}{}{subsection.6.6}{Dirichlet-GRH-v930-Volume-I.pdf}}
\expandafter\def\csname r@nav:VI-08\endcsname{{121.16}{51}{}{subsection.121.16}{Dirichlet-GRH-v930-Volume-XXIII.pdf}}
\expandafter\def\csname r@nav:VI-15\endcsname{{92.8}{18}{}{subsection.92.8}{Dirichlet-GRH-v930-Volume-XVIII.pdf}}
\expandafter\def\csname r@prop:cancelsquare256\endcsname{{74.9}{111}{}{proposition.74.9}{Dirichlet-GRH-v930-Volume-XIV.pdf}}
\expandafter\def\csname r@prop:cofactor286\endcsname{{105.12}{43}{}{proposition.105.12}{Dirichlet-GRH-v930-Volume-XX.pdf}}
\expandafter\def\csname r@prop:depth245\endcsname{{95.3}{67}{}{proposition.95.3}{Dirichlet-GRH-v930-Volume-XVIII.pdf}}
\expandafter\def\csname r@prop:der95\endcsname{{93.16}{33}{}{proposition.93.16}{Dirichlet-GRH-v930-Volume-XVIII.pdf}}
\expandafter\def\csname r@prop:flagunmarked264\endcsname{{102.1}{87}{}{proposition.102.1}{Dirichlet-GRH-v930-Volume-XIX.pdf}}
\expandafter\def\csname r@prop:gradient248\endcsname{{106.5}{57}{}{proposition.106.5}{Dirichlet-GRH-v930-Volume-XX.pdf}}
\expandafter\def\csname r@prop:index90\endcsname{{94.6}{41}{}{proposition.94.6}{Dirichlet-GRH-v930-Volume-XVIII.pdf}}
\expandafter\def\csname r@prop:jacorbit266\endcsname{{75.6}{118}{}{proposition.75.6}{Dirichlet-GRH-v930-Volume-XIV.pdf}}
\expandafter\def\csname r@prop:kernel270\endcsname{{106.7}{60}{}{proposition.106.7}{Dirichlet-GRH-v930-Volume-XX.pdf}}
\expandafter\def\csname r@prop:linkedline265\endcsname{{74.4}{104}{}{proposition.74.4}{Dirichlet-GRH-v930-Volume-XIV.pdf}}
\expandafter\def\csname r@prop:oppline256\endcsname{{102.8}{96}{}{proposition.102.8}{Dirichlet-GRH-v930-Volume-XIX.pdf}}
\expandafter\def\csname r@prop:order268\endcsname{{106.6}{59}{}{proposition.106.6}{Dirichlet-GRH-v930-Volume-XX.pdf}}
\expandafter\def\csname r@prop:parity262\endcsname{{71.18}{35}{}{proposition.71.18}{Dirichlet-GRH-v930-Volume-XIV.pdf}}
\expandafter\def\csname r@prop:poincarenorm261\endcsname{{106.3}{55}{}{proposition.106.3}{Dirichlet-GRH-v930-Volume-XX.pdf}}
\expandafter\def\csname r@prop:positions264\endcsname{{121.4}{33}{}{proposition.121.4}{Dirichlet-GRH-v930-Volume-XXIII.pdf}}
\expandafter\def\csname r@prop:projecttests260\endcsname{{71.4}{17}{}{proposition.71.4}{Dirichlet-GRH-v930-Volume-XIV.pdf}}
\expandafter\def\csname r@prop:res23\endcsname{{6.11}{32}{}{proposition.6.11}{Dirichlet-GRH-v930-Volume-I.pdf}}
\expandafter\def\csname r@prop:res80\endcsname{{7.7}{94}{}{proposition.7.7}{Dirichlet-GRH-v930-Volume-I.pdf}}
\expandafter\def\csname r@prop:scale151\endcsname{{106.13}{67}{}{proposition.106.13}{Dirichlet-GRH-v930-Volume-XX.pdf}}
\expandafter\def\csname r@prop:schurframe259\endcsname{{75.7}{120}{}{proposition.75.7}{Dirichlet-GRH-v930-Volume-XIV.pdf}}
\expandafter\def\csname r@prop:shell263\endcsname{{75.5}{117}{}{proposition.75.5}{Dirichlet-GRH-v930-Volume-XIV.pdf}}
\expandafter\def\csname r@prop:short107\endcsname{{9.9}{114}{}{proposition.9.9}{Dirichlet-GRH-v930-Volume-I.pdf}}
\expandafter\def\csname r@prop:spacing276\endcsname{{94.30}{63}{}{proposition.94.30}{Dirichlet-GRH-v930-Volume-XVIII.pdf}}
\expandafter\def\csname r@prop:terminalcore255\endcsname{{95.4}{68}{}{proposition.95.4}{Dirichlet-GRH-v930-Volume-XVIII.pdf}}
\expandafter\def\csname r@prop:tracelat253\endcsname{{122.7}{63}{}{proposition.122.7}{Dirichlet-GRH-v930-Volume-XXIII.pdf}}
\expandafter\def\csname r@prop:vacuumcore254\endcsname{{122.8}{64}{}{proposition.122.8}{Dirichlet-GRH-v930-Volume-XXIII.pdf}}
\expandafter\def\csname r@prop:vol127\endcsname{{92.25}{20}{}{proposition.92.25}{Dirichlet-GRH-v930-Volume-XVIII.pdf}}
\expandafter\def\csname r@sec:char94\endcsname{{121.16}{51}{}{subsection.121.16}{Dirichlet-GRH-v930-Volume-XXIII.pdf}}
\expandafter\def\csname r@sec:graphtransport337\endcsname{{71.20}{38}{}{subsection.71.20}{Dirichlet-GRH-v930-Volume-XIV.pdf}}
\expandafter\def\csname r@sec:missing\endcsname{{6.6}{32}{}{subsection.6.6}{Dirichlet-GRH-v930-Volume-I.pdf}}
\expandafter\def\csname r@sec:orbital133\endcsname{{98.1}{3}{}{subsection.98.1}{Dirichlet-GRH-v930-Volume-XIX.pdf}}
\expandafter\def\csname r@sec:refined112\endcsname{{92.5}{11}{}{subsection.92.5}{Dirichlet-GRH-v930-Volume-XVIII.pdf}}
\expandafter\def\csname r@seriesguide\endcsname{{1.1}{4}{}{subsection.1.1}{Dirichlet-GRH-v930-Volume-I.pdf}}
\expandafter\def\csname r@thm:auxiliaryflag339\endcsname{{99.17}{57}{}{proposition.99.17}{Dirichlet-GRH-v930-Volume-XIX.pdf}}
\expandafter\def\csname r@thm:average269\endcsname{{8.3}{102}{}{proposition.8.3}{Dirichlet-GRH-v930-Volume-I.pdf}}
\expandafter\def\csname r@thm:column242\endcsname{{121.3}{31}{}{proposition.121.3}{Dirichlet-GRH-v930-Volume-XXIII.pdf}}
\expandafter\def\csname r@thm:common111\endcsname{{9.13}{118}{}{proposition.9.13}{Dirichlet-GRH-v930-Volume-I.pdf}}
\expandafter\def\csname r@thm:common133\endcsname{{98.2}{4}{}{proposition.98.2}{Dirichlet-GRH-v930-Volume-XIX.pdf}}
\expandafter\def\csname r@thm:commonnorm128\endcsname{{100.16}{78}{}{proposition.100.16}{Dirichlet-GRH-v930-Volume-XIX.pdf}}
\expandafter\def\csname r@thm:compact248\endcsname{{71.15}{31}{}{proposition.71.15}{Dirichlet-GRH-v930-Volume-XIV.pdf}}
\expandafter\def\csname r@thm:conductor232\endcsname{{19.15}{12}{}{proposition.19.15}{Dirichlet-GRH-v930-Volume-IV.pdf}}
\expandafter\def\csname r@thm:cusp247\endcsname{{106.4}{56}{}{proposition.106.4}{Dirichlet-GRH-v930-Volume-XX.pdf}}
\expandafter\def\csname r@thm:decay96\endcsname{{121.26}{55}{}{proposition.121.26}{Dirichlet-GRH-v930-Volume-XXIII.pdf}}
\expandafter\def\csname r@thm:direction249\endcsname{{71.6}{19}{}{proposition.71.6}{Dirichlet-GRH-v930-Volume-XIV.pdf}}
\expandafter\def\csname r@thm:dominantflag344\endcsname{{9.21}{125}{}{proposition.9.21}{Dirichlet-GRH-v930-Volume-I.pdf}}
\expandafter\def\csname r@thm:downcore254\endcsname{{71.13}{28}{}{proposition.71.13}{Dirichlet-GRH-v930-Volume-XIV.pdf}}
\expandafter\def\csname r@thm:dyad234\endcsname{{110.4}{40}{}{proposition.110.4}{Dirichlet-GRH-v930-Volume-XXI.pdf}}
\expandafter\def\csname r@thm:edge242\endcsname{{67.9}{82}{}{proposition.67.9}{Dirichlet-GRH-v930-Volume-XIII.pdf}}
\expandafter\def\csname r@thm:energy153\endcsname{{106.17}{70}{}{proposition.106.17}{Dirichlet-GRH-v930-Volume-XX.pdf}}
\expandafter\def\csname r@thm:first250\endcsname{{71.10}{25}{}{proposition.71.10}{Dirichlet-GRH-v930-Volume-XIV.pdf}}
\expandafter\def\csname r@thm:firstsquare255\endcsname{{71.3}{16}{}{proposition.71.3}{Dirichlet-GRH-v930-Volume-XIV.pdf}}
\expandafter\def\csname r@thm:flagindex336\endcsname{{99.14}{55}{}{proposition.99.14}{Dirichlet-GRH-v930-Volume-XIX.pdf}}
\expandafter\def\csname r@thm:fourierslopes354\endcsname{{120.8}{20}{}{proposition.120.8}{Dirichlet-GRH-v930-Volume-XXIII.pdf}}
\expandafter\def\csname r@thm:global167\endcsname{{72.4}{47}{}{proposition.72.4}{Dirichlet-GRH-v930-Volume-XIV.pdf}}
\expandafter\def\csname r@thm:graph274\endcsname{{94.27}{59}{}{proposition.94.27}{Dirichlet-GRH-v930-Volume-XVIII.pdf}}
\expandafter\def\csname r@thm:grh16\endcsname{{1.1}{1}{}{proposition.1.1}{Dirichlet-GRH-v930-Volume-I.pdf}}
\expandafter\def\csname r@thm:identity127\endcsname{{92.24}{20}{}{proposition.92.24}{Dirichlet-GRH-v930-Volume-XVIII.pdf}}
\expandafter\def\csname r@thm:imagedet355\endcsname{{127.3}{28}{}{proposition.127.3}{Dirichlet-GRH-v930-Volume-XXIV.pdf}}
\expandafter\def\csname r@thm:individual259\endcsname{{74.2}{101}{}{proposition.74.2}{Dirichlet-GRH-v930-Volume-XIV.pdf}}
\expandafter\def\csname r@thm:invar148\endcsname{{20.11}{31}{}{proposition.20.11}{Dirichlet-GRH-v930-Volume-IV.pdf}}
\expandafter\def\csname r@thm:levels88\endcsname{{7.11}{97}{}{proposition.7.11}{Dirichlet-GRH-v930-Volume-I.pdf}}
\expandafter\def\csname r@thm:lowesttwo261\endcsname{{75.4}{116}{}{proposition.75.4}{Dirichlet-GRH-v930-Volume-XIV.pdf}}
\expandafter\def\csname r@thm:markedcore253\endcsname{{102.7}{95}{}{proposition.102.7}{Dirichlet-GRH-v930-Volume-XIX.pdf}}
\expandafter\def\csname r@thm:maximal130\endcsname{{110.1}{37}{}{proposition.110.1}{Dirichlet-GRH-v930-Volume-XXI.pdf}}
\expandafter\def\csname r@thm:mixedfourier334\endcsname{{83.19}{68}{}{proposition.83.19}{Dirichlet-GRH-v930-Volume-XVI.pdf}}
\expandafter\def\csname r@thm:moment258\endcsname{{83.14}{63}{}{proposition.83.14}{Dirichlet-GRH-v930-Volume-XVI.pdf}}
\expandafter\def\csname r@thm:neighbormove356\endcsname{{73.3}{75}{}{proposition.73.3}{Dirichlet-GRH-v930-Volume-XIV.pdf}}
\expandafter\def\csname r@thm:neighbour93\endcsname{{94.11}{46}{}{proposition.94.11}{Dirichlet-GRH-v930-Volume-XVIII.pdf}}
\expandafter\def\csname r@thm:norm177\endcsname{{72.21}{64}{}{proposition.72.21}{Dirichlet-GRH-v930-Volume-XIV.pdf}}
\expandafter\def\csname r@thm:norm91\endcsname{{94.9}{44}{}{proposition.94.9}{Dirichlet-GRH-v930-Volume-XVIII.pdf}}
\expandafter\def\csname r@thm:packing246\endcsname{{111.17}{72}{}{proposition.111.17}{Dirichlet-GRH-v930-Volume-XXI.pdf}}
\expandafter\def\csname r@thm:pair107\endcsname{{9.11}{116}{}{proposition.9.11}{Dirichlet-GRH-v930-Volume-I.pdf}}
\expandafter\def\csname r@thm:parabolic129\endcsname{{100.15}{77}{}{proposition.100.15}{Dirichlet-GRH-v930-Volume-XIX.pdf}}
\expandafter\def\csname r@thm:project151\endcsname{{106.15}{68}{}{proposition.106.15}{Dirichlet-GRH-v930-Volume-XX.pdf}}
\expandafter\def\csname r@thm:project94\endcsname{{121.22}{52}{}{proposition.121.22}{Dirichlet-GRH-v930-Volume-XXIII.pdf}}
\expandafter\def\csname r@thm:rangenew50\endcsname{{6.29}{62}{}{proposition.6.29}{Dirichlet-GRH-v930-Volume-I.pdf}}
\expandafter\def\csname r@thm:rationalflag342\endcsname{{9.18}{122}{}{proposition.9.18}{Dirichlet-GRH-v930-Volume-I.pdf}}
\expandafter\def\csname r@thm:raw276\endcsname{{94.28}{60}{}{proposition.94.28}{Dirichlet-GRH-v930-Volume-XVIII.pdf}}
\expandafter\def\csname r@thm:realcore253\endcsname{{71.16}{32}{}{proposition.71.16}{Dirichlet-GRH-v930-Volume-XIV.pdf}}
\expandafter\def\csname r@thm:realrow257\endcsname{{74.1}{100}{}{proposition.74.1}{Dirichlet-GRH-v930-Volume-XIV.pdf}}
\expandafter\def\csname r@thm:region198\endcsname{{20.13}{33}{}{proposition.20.13}{Dirichlet-GRH-v930-Volume-IV.pdf}}
\expandafter\def\csname r@thm:rel155\endcsname{{106.20}{72}{}{proposition.106.20}{Dirichlet-GRH-v930-Volume-XX.pdf}}
\expandafter\def\csname r@thm:secondrow256\endcsname{{74.5}{105}{}{proposition.74.5}{Dirichlet-GRH-v930-Volume-XIV.pdf}}
\expandafter\def\csname r@thm:seesaw245\endcsname{{71.14}{30}{}{proposition.71.14}{Dirichlet-GRH-v930-Volume-XIV.pdf}}
\expandafter\def\csname r@thm:short246\endcsname{{111.1}{57}{}{proposition.111.1}{Dirichlet-GRH-v930-Volume-XXI.pdf}}
\expandafter\def\csname r@thm:signedreturn342\endcsname{{83.26}{73}{}{proposition.83.26}{Dirichlet-GRH-v930-Volume-XVI.pdf}}
\expandafter\def\csname r@thm:splitflag343\endcsname{{9.20}{123}{}{proposition.9.20}{Dirichlet-GRH-v930-Volume-I.pdf}}
\expandafter\def\csname r@thm:target90\endcsname{{94.8}{43}{}{proposition.94.8}{Dirichlet-GRH-v930-Volume-XVIII.pdf}}
\expandafter\def\csname r@thm:tensor249\endcsname{{106.9}{62}{}{proposition.106.9}{Dirichlet-GRH-v930-Volume-XX.pdf}}
\expandafter\def\csname r@thm:terminalscalar357\endcsname{{96.1}{72}{}{proposition.96.1}{Dirichlet-GRH-v930-Volume-XVIII.pdf}}
\expandafter\def\csname r@thm:terminalwalk338\endcsname{{71.24}{40}{}{proposition.71.24}{Dirichlet-GRH-v930-Volume-XIV.pdf}}
\expandafter\def\csname r@thm:third263\endcsname{{74.6}{107}{}{proposition.74.6}{Dirichlet-GRH-v930-Volume-XIV.pdf}}
\expandafter\def\csname r@thm:transverse250\endcsname{{106.2}{53}{}{proposition.106.2}{Dirichlet-GRH-v930-Volume-XX.pdf}}
\expandafter\def\csname r@thm:vac242\endcsname{{67.8}{81}{}{proposition.67.8}{Dirichlet-GRH-v930-Volume-XIII.pdf}}
\expandafter\def\csname r@thm:word275\endcsname{{83.10}{59}{}{proposition.83.10}{Dirichlet-GRH-v930-Volume-XVI.pdf}}
\expandafter\def\csname r@thm:wordrank134\endcsname{{98.4}{6}{}{proposition.98.4}{Dirichlet-GRH-v930-Volume-XIX.pdf}}
\makeatother
\begin{document}
\setcounter{part}{16}
\fancyhead[L]{\small\color{muted}Dirichlet GRH \textperiodcentered\ Volume XVII}
\fancyhead[R]{\small\color{muted}Projected Bessel periods}
\thispagestyle{plain}\begin{center}
{\LARGE\bfseries The generalized Riemann hypothesis\\[2pt]for Dirichlet \(L\)-functions}\\[6pt]
{\small 27 November 2026 \quad\textperiodcentered\quad Edition 930}\end{center}
\part{Bessel models and projected periods}\label{part:XVII}
\begin{abstract}\noindent We construct Bessel functionals for specified characters and compact types. At split finite places we match sections of complementary ranks. For the prescribed auxiliary type we prove that an archimedean translation close to the identity gives zero on the original spherical rows when its other valuations are at good split primes. We give a different rational translation compatible with every marked exchange and every unmarked integral root. We compute the contribution of the normalized terminal input to the determinant of the signed period image.\end{abstract}
\setcounter{section}{85}
\noindent The \hyperref[seriesguide]{guide in Volume~I} lists all twenty-seven volumes.
\begingroup\small\setlength{\parskip}{0pt}\tableofcontents\endgroup
\clearpage

\section{The global Bessel integral}\label{nav:g-VII-0}\label{nav:volume-XVII}

\subsection{An integral frame for the terminal frequency}\label{nav:VII-00}
\label{sec:middle131}

Use the common fields, intervals, words and primes of
Theorem \ref{thm:common111}.
The terminal vector in the Bessel datum has norm \(-T\). We take
it to be \(z=\varpi e\), where \(T=\varpi\bar\varpi\), with the
Witt flag and the group \(B=N\rtimes H\) of Section~\ref{sec:char94}.
At \(v\mid T\) choose the split factor in which \(\varpi\) has
valuation one as the first factor. Put \(O=\ZZ_T\), and lift the
given nonzero column covector to \(\ell_v:O^k\to O\).
In that factor set
\[
 M_v=\{(w,z_0)\in O^k\oplus O:z_0\equiv\ell_v(w)\pmod T\}.
\]
Let \(a_v^0\) act on the middle space by
\((w,z_0)\mapsto(w,\ell_v(w)+Tz_0)\), and by the identity on
all \(u_i^+,v_i^+\). The dual action in the other factor defines
a split unitary element. Set \(a^0=\prod_{v\mid T}a_v^0\), extended
by one elsewhere, and \(K^0=a^0K_\sigma(a^0)^{-1}\).

\begin{theorem}[The column level, terminal character and integral frame]
\label{thm:midlevel131}
The finite level \(K^0\) induces the full prescribed covector
level \(K_\pi(\ell)\), with
\[
       B_f\cap K^0=(N_f\cap K^0)\rtimes K_\pi(\ell).
\]
At \(v\mid T\), among locally integral terminal frequencies
\(\beta\in O\times O\), precisely those with first component
in \(TO\) have trivial character on \(N_v\cap K_v^0\).
In particular \(\chi_\varpi\) is trivial on \(N_f\cap K^0\).
The rational vectors \(u_1,\ldots,u_s,\varpi e\) span a
saturated submodule of the adelic integral lattice with this
stabilizer at every finite place.
\end{theorem}
\begin{proof}
The first lattice at \(v\) is \(O^s\oplus M_v\oplus O^s\) in
the Witt splitting. The second is its dual for the Hermitian
pairing. For \(nh\in B_v\) the induced action on the middle
quotient of the flag is \((w,z_0)\mapsto(hw,z_0)\). This
preserves \(M_v\) in both directions exactly when
\(h\in\GL_k(O)\) and \(\ell_v h=\ell_v\bmod T\). Indeed
projection of \(M_v\) onto \(O^k\) is surjective, and the
residue relation determines the condition on \(\ell_v\).
The block embedding of each such \(h\) preserves the entire
lattice pair. Thus if \(nh\) preserves it, so do \(h\) and
\(n\) separately. Away from \(T\) the levels are those of
Lemma~\ref{lem:orbit110}.

In the two split factors the first \(s-1\) simple coordinates
of an integral unipotent at this level are integral because
the lattices on both strings of isotropic lines have the original
bases. In the first factor \(Te^+\) is in \(M_v\), giving
\(d_s(n)^+\in T^{-1}O\). In the second factor \(e^-\) is
integral, since it pairs with \(M_v\) by the integral last
coordinate up to sign. Hence \(d_s(n)^-\in O\). These facts
give sufficiency by the conductor of the additive character of
\(\QQ_T\). For necessity take the elementary shear adding
\(c(z_0-\ell_v(w))u_s^+/T\) to the middle vector \((w,z_0)\)
in the first factor, with \(c\in O\). It and its inverse preserve
the lattice because the numerator is a multiple of \(T\) on
\(M_v\). Its dual fixes \(e^-\). It belongs to \(N_v\), its
earlier simple coordinates vanish and \(d_s(n)=(c/T,0)\).
Triviality for all \(c\) forces \(\beta^+\in TO\). Away from
\(T\) integrality suffices by Lemma~\ref{lem:orbit110}.

The span of \(Te^+\) splits off in \(M_v\). Indeed the vectors
\((w,\ell_v(w))\), together with the multiples of \(Te^+\),
give a direct decomposition. The first component of \(\varpi\)
is \(T\) times a unit. In the dual lattice \(e^-\) pairs to a
unit with some \((w,\ell_v(w))\), since the residue covector is
nonzero. It is therefore primitive and extends to a basis over
the discrete valuation ring. The second component of \(\varpi\)
is a unit. The \(u_i\) are separate basis lines in both factors.
Away from \(T\) the scalar \(\varpi\) is a unit and
\(u_i=e_{+,i}+e_{-,i}\) is part of an integral basis in its
coordinate plane, including at places above two. This proves
saturation. Replacing \(\varpi e\) by \(\varpi e+b u_s\) for
any \(b\in\OO_E\) gives the same saturated module.
\end{proof}

\begin{theorem}[Preservation of both row and unipotent volumes]
\label{thm:midnorm131}
The row volumes for \(K^0\) and \(K_\sigma\) agree, and
\[
       \operatorname{vol}_f(N_f\cap K^0)
                    =\operatorname{vol}_f(N_f\cap K_\sigma).
\]
For the norm-\(T\) frequency the class-coordinate constant
\(c_{\rho,K^0}\) of \eqref{eq:main127} is consequently
\(c_{\rho,K_\sigma}\) for every word. In the column probability
norm the ratio of squared-norm constants is
\begin{equation}
                \frac{C_{K^0}}{C_{K_\sigma}}=(T^k-1)^f .
                \label{eq:midratio131}
\end{equation}
Summation over the entire stabilizer \(Q(F_+)\) of the isotropic
\(s\)-space is isometric on these refined columns, including
their zero-sum subspaces, with the norms of Theorem \ref{thm:maximal130}.
\end{theorem}
\begin{proof}
The row assertion follows from conjugacy in the unimodular
group \(G_f\) of finite adelic points. At \(v\mid T\) the
conjugating element lies in the Levi of the full flag. Its
middle block has determinant \(T\). Among the independent
upper blocks of \(N_v\), a block from the middle to one of
the first \(s\) lines acquires determinant \(T^{-1}\) and
a block from one of the last \(s\) lines to the middle acquires
determinant \(T\). These occur \(s\) times each. The other
blocks have determinant one. Product additive measure in the
blocks is Haar measure as in Proposition~\ref{prop:vol127},
so the unipotent modulus is one as claimed.

Theorem~\ref{thm:midlevel131} gives the finite decomposition and
character compatibility needed in Lemma~\ref{lem:covol127} and
Theorem~\ref{thm:identity127}. The column stabilizers remain
trivial and there are \((T^k-1)^f\) times the original number
of classes. The class-coordinate constant and \eqref{eq:midratio131}
now follow from those proofs and the column probability norm of
Theorem \ref{thm:commonnorm128}.
In the proof of Theorem~\ref{thm:maximal130} the intermediate
unipotent character separations depend only on the rational flag.
The finite support at the last separation requires compact row
levels at every finite place and unipotent reduction at \(g_1,g_2\).
Both hold here, and the input seeds are defined by the finite
decomposition just proved. The same isometry follows for the
norm-\(T\) frequency.
\end{proof}

For comparison, at \(K^0\) the frequency \(1\) has zero Bessel
coefficient on right invariant rows. In the inner integral at
\(h_0\), translate its unipotent variable on the right by
\(h_0 n h_0^{-1}\) for one of the shears in the preceding proof.
The row is fixed, while \(\chi_1(h_0nh_0^{-1})=\chi_1(n)\)
can differ from one. Section~\ref{sec:refined112} gives a level
supporting both frequencies. The construction here chooses the
frequency attached to the unnormalized terminal vector and makes
that vector integral together with the entire isotropic frame.

\subsection{The Gram determinant of two terminal frames}\label{nav:VII-01}
\label{sec:gram132}

Work in the Hermitian space \(V\) over \(E/F_+\) of
Section~\ref{sec:char94}. Its signature at every real place is
\((s+k,s+1)\). A frame consists of a basis \(u_1,\ldots,u_s\) of
a totally isotropic space \(L\) and a vector \(z\in L^\perp\)
of norm \(-n\), where \(n\in F_+\) is totally positive. Let
\(Y=L+Ez\). Use primes for the data of a second frame. Write
\(I=L\cap L'\), \(d=\dim_E I\), and \(A=(Y+Y')^\perp\).
For \(y=(u_1,\ldots,u_s,z)\) and the corresponding \(y'\) put
\[
 Z_{ij}=h(y_i,y'_j),\qquad
 a=\det Z_{1:s,1:s},\quad b=\det Z,\quad
 \mathcal G=b\bar b-nn'a\bar a .
\]

\begin{theorem}[The common definite subspace]\label{thm:cross132}
The radical of \(A\) is \(I\). It has a complement \(C\) over \(E\)
that is positive definite at every real place. Its dimension is
\(k\) if \(Y\cap Y'\) contains a vector of negative norm at a
real place, and \(k-1\) otherwise. In the first case the vector
can be chosen rational over \(E\) and totally negative.
In both cases \(U(C)\), acting identically on \(C^\perp\), fixes
both frames pointwise. Its projection to each of the two definite
column groups is the usual unitary action on a subspace of the
stated dimension.

If \(d=0\), the number \(\mathcal G\) is totally positive when
\(\dim C=k-1\), and is zero when \(\dim C=k\).
For \(d>0\) these same tests hold in \(I^\perp/I\), using any
bases of \(L/I,L'/I\) and the images of \(z,z'\).
\end{theorem}
\begin{proof}
At a real place the form on \(Y\) is negative semidefinite with
radical \(L\). Split off hyperbolic planes for a basis of \(L\).
The remaining space has signature \((k,1)\), and the image of
\(z\) has negative norm. Its perpendicular there is positive
definite. Thus \(Y^\perp\) is positive semidefinite with radical
\(L\). The same assertions hold with primes. In particular the
vectors of norm zero in \(A\) lie in both \(L\) and \(L'\).
Conversely \(I\) is orthogonal to both frames and to \(A\). This
proves the assertion on the radical. Any \(E\)-linear complement
to \(I\) in \(A\) is therefore totally positive.

Inside \(Y\cap Y'\), the kernel of projection to \(Y/L\) is
exactly \(I\). Indeed any vector in that kernel has norm zero
and its membership in \(Y'\) puts it in \(L'\). Consequently
\(\dim(Y\cap Y')=d+j\) with \(j=0\) or \(1\). When \(j=1\)
choose an \(E\)-vector in this intersection outside \(I\).
Its coefficient of \(z\) is nonzero, so its norm is negative
at every place. Dimensions of the intersections are unchanged
under field extension, proving the assertion starting at a single
real place as well. Nondegeneracy of \(V\) gives
\[
    \dim C=\dim V-2(s+1)+(d+j)-d=k-1+j .
\]
Since \(Y,Y'\subset C^\perp\), the extended action fixes every
displayed frame vector. On \(Y^\perp/L\) the map from \(C\) is
isometric, and so is the map to \((Y')^\perp/L'\).
In each case the group fixes the orthogonal complement of its
image. These are the two column representations on their defining
Hermitian spaces. The projection of this group to either flag
Levi fixes the terminal line and every graded flag vector.

The Gram matrix of the concatenated list \((y,y')\) has off-diagonal
blocks \(Z,Z^*\). The two diagonal blocks are zero apart from
their final entries \(-n,-n'\). A nonzero term of the determinant
selects the same number of entries from each diagonal block by
counting its columns. Terms selecting none sum to
\((-1)^{s+1}b\bar b\), by interchanging the two blocks of columns.
Selecting both final entries leaves the two size-\(s\) off-diagonal
blocks and contributes \((-1)^snn'a\bar a\). Hence
\begin{equation}
       \det\operatorname{Gram}(y,y')=(-1)^{s+1}\mathcal G .
       \label{eq:cross132}
\end{equation}
For \(d=0,j=1\) the concatenated list is dependent. For \(d=j=0\)
it is a basis of \(A^\perp\), with \(A\) positive definite of
dimension \(k-1\). The signature of \(A^\perp\) is
\((s+1,s+1)\). The sign in \eqref{eq:cross132} now gives total
positivity. For \(d>0\) split off the hyperbolic partners of
\(I\) to see that \(I^\perp/I\) has signature
\((s-d+k,s-d+1)\). Both spaces \(Y,Y'\) lie in \(I^\perp\).
Their images intersect in \((Y\cap Y')/I\) and their isotropic
parts now have zero intersection. The argument for \(d=0\)
applies there, with the determinant of an empty minor understood
as one. A change of either isotropic basis multiplies the
reduced invariant by the norm of the determinant of the change,
which preserves its sign and its vanishing.
\end{proof}

\begin{theorem}[A divisibility bound at both auxiliary primes]
\label{thm:gap132}
Use the level \(K^0\) and \(z=\varpi e\) in
Theorem~\ref{thm:midlevel131}. For \(\gamma\in G(F_+)\) take
\(y'=\gamma y\). If \(\gamma_f\in B_fK^0B_f\), then
\[
    \frac{\mathcal G}{(g_1g_2)^{2(s+1)}}\in\mathcal O_{F_+}.
\]
In particular, for every such pair with \(\mathcal G\ne0\),
\begin{equation}
     \prod_{\tau:F_+\hookrightarrow\mathbb R}\tau(\mathcal G)
              \ \ge\ (g_1g_2)^{mf},\qquad m=2(s+1).
     \label{eq:gap132}
\end{equation}
\end{theorem}
\begin{proof}
Every \(b_0\in B_f\) acts on the frame by an upper unitriangular
matrix. Indeed the successive \(u_i\) are fixed modulo earlier
lines and the terminal vector is fixed modulo \(L\).
Both exterior products \(u_1\wedge\cdots\wedge u_s\) and
\(u_1\wedge\cdots\wedge u_s\wedge z\) are therefore fixed.
The minors \(a,b\) are the Hermitian pairings of each of these
exterior products with its translate. In a finite factorization
\(\gamma_f=b_1 t b_2\), \(t\in K^0\), invariance of the
Hermitian pairings gives the same two minors with \(t\) in
place of \(\gamma_f\). The frame is integral by
Theorem~\ref{thm:midlevel131} and \(t\) preserves its ambient
integral Hermitian lattice. At \(T\) the two lattices in the
split factors are paired duals. Thus the minors, and hence
\(\mathcal G\), are integral at every finite place.

At a place over \(g_1\) or \(g_2\), write \(g\) for the
underlying rational prime. In the first split factor take the
Witt order \(u_1,\ldots,u_s,W_0,e,v_s,\ldots,v_1\). It is an
integral basis there as \(g\) is odd. The specified row group is
upper unitriangular modulo \(g\), with the \(k\) entries of the
\(e\)-column in the \(W_0\)-rows equal to zero. Hence the
reduced matrix sends the ordered frame \((u_1,\ldots,u_s,z)\)
to itself up to upper unitriangular change of basis. Here
\(\varpi\) is a unit. In the other factor \(u_i\) is dual to
\(v_i\), up to the chosen unit pairing. The \(v_i\) occur
in reverse order, so the inverse-transpose action sends
\(u_i\) to itself modulo the earlier \(u_j\). The \(e\)-row
of the inverse upper unitriangular matrix has only its diagonal
and possible entries in the \(v\)-columns. Thus the second
factor also sends \(z\) to itself modulo the \(u_j\).

Lift these residue coefficients in each factor. Subtract from
each vector of \(t y\) the resulting linear combination of the
first frame. All \(s+1\) differences lie in \(g\) times the
local lattice. These column operations on the concatenated list
have determinant one. The new Gram determinant has \(s+1\)
factors of \(g\) from the rows and \(s+1\) from the columns
by sesquilinearity, including the entries where both occur.
Equation~\eqref{eq:cross132} gives the asserted divisibility
at every place above both primes, which split completely.
If the invariant is nonzero the full Gram list is independent,
so \(d=0\), and Theorem~\ref{thm:cross132} gives total positivity.
The norm of the resulting nonzero totally positive algebraic
integer is a positive rational integer. This proves
\eqref{eq:gap132}.
\end{proof}

The two minors, their signed difference and the flag congruence
appear in \srcloc{Formalism/SharpRestriction/CuspSpanNormBounds/BlockEnds}{53},
line~62 there, and
\srcloc{Formalism/SharpRestriction/CuspSpanNormBounds/IsPrimitiveSub}{139}.
The real sign is also the assertion at
\srcloc{Formalism/PeriodIntegrals/UnitaryGroupMassFormula/CountPositivity/LeadingMinorFullEqProd}{97}.
The bound concerns the unreduced invariant at the specified
level. For a pair sharing an isotropic direction this invariant
vanishes, even when the reduced determinant is positive.

\subsubsection{A binary unitary congruence and its real weight}\label{nav:VII-02}
\label{sub:binary230}

The normal quotient of a pair of terminal hyperplanes has
signature \((1,1)\). We give an estimate for its integral
isometries that is uniform in the CM field and in the point
at which the real weight is evaluated.

Let \(E/F\) be CM, \(f=[F:\QQ]\), and take one embedding
\(\tau:E\hookrightarrow\CC\) over each real place. For
positive integers \(Q,T\) with \(T>\max(2,Q)\), use the
lattice and group
\[
 \mathcal Z=\{(u,v)\in\mathcal O_E^2:u+v\in Q\mathcal O_E\},
 \qquad
 \Gamma=\{g\in U(1,1)(F): (g-1)\mathcal Z\subseteq T\mathcal Z\}.
\]
Here the Hermitian form is \(\operatorname{diag}(1,-1)\).
The positive lines at a real place are \(\CC(1,z)\) with
\(z\) in the open unit disk \(\mathbb D\). Matrices act on
these lines, hence on \(\mathbb D^f\) componentwise.

\begin{lemma}[Separation and stabilizers of positive lines]
\label{lem:binary230}
The set \(\Gamma\) is a group. Every element has determinant
one. No nonidentity element fixes a point of \(\mathbb D^f\).
For \(g\ne1\) its lower left entry \(c\) satisfies
\[
          \prod_\tau |\tau c|\ \ge\ (T/Q)^f.
\]
\end{lemma}
\begin{proof}
In the \(\mathcal O_E\)-basis \((1,-1),(0,Q)\) the matrix of
\(g\) has integral entries and is the identity modulo \(T\).
Its determinant has absolute value one at every embedding by
unitarity. If the determinant differs from one, their
difference belongs to \(T\mathcal O_E\) with absolute norm
at least \(T^{2f}\), but each conjugate has absolute value
at most two. Thus \(\det g=1\). The inverse has integral
entries in the same basis and is also the identity modulo
\(T\). Closure under multiplication follows in this basis.

If \(g\) fixes a tuple of positive lines, at each place it
preserves both that line and its negative perpendicular, and
acts on each by a scalar of absolute value one. Consequently
\(|\tau(2-\operatorname{tr}g)|\le4\) at every place.
In the lattice basis,
\(2-\operatorname{tr}g=\det(g-1)\in T^2\mathcal O_E\).
Since \(T^2>4\), taking norms gives \(\operatorname{tr}g=2\).
The two eigenvalues are then one, and the two invariant lines
give a basis, so \(g=1\). In particular if \(c=0\), then
\(g\) fixes the tuple of lines of slope zero and is the
identity. Finally \((Q,0)\in\mathcal Z\) gives
\(Qc/T\in\mathcal O_E\) by taking the second coordinate.
Its nonzero absolute norm gives the asserted product.
\end{proof}

Give the disk the metric \(2|dz|/(1-|z|^2)\), write its
distance as \(d_{\mathbb D}\), and put
\[
 \Phi(z)=\exp\!\left(\tfrac12\sum_\tau
                     d_{\mathbb D}(0,z_\tau)\right)
        =\prod_\tau\sqrt{\frac{1+|z_\tau|}{1-|z_\tau|}}
                \qquad(z\in\mathbb D^f).
\]
The second formula follows by integration along a radius.
Radial projection bounds any path length below by the same
radial integral. To see that the action preserves the metric, take a
matrix \(\left(\begin{smallmatrix}a&b\\ c&e\end{smallmatrix}\right)\) in the unitary group.
The new slope is \((c+ez)/(a+bz)\), its derivative has
modulus \(|a+bz|^{-2}\), and the new value of \(1-|z|^2\)
is \((1-|z|^2)/|a+bz|^2\). Every disk point can be carried
to zero by a unitary matrix, for the columns
\((1,z),(\bar z,1)\) divided by \(\sqrt{1-|z|^2}\)
give one carrying zero to \(z\).

\begin{theorem}[A uniform bound for the sum over a congruence orbit]
\label{thm:disk230}
For all \(z\in\mathbb D^f\) and real \(p\ge3\),
\[
 \sum_{g\in\Gamma}\Phi(gz)^{-3}\le 10^f,\qquad
 \sum_{g\in\Gamma,\ g\ne1}\{\Phi(z)\Phi(gz)\}^{-p}
     \le 10^f\Phi(z)^{-3}
                      \left(\frac{Q}{2T}\right)^{f(p-3)}.
\]
\end{theorem}
\begin{proof}
At each place \(g0\) has slope \(c/a\), with
\(|a|^2-|c|^2=1\), giving
\(\sinh(d_{\mathbb D}(0,g0)/2)=|c|\).
Write \(y_g=g^{-1}0\). If \(g\ne h\), transport the two
centres by \(g\) and apply Lemma~\ref{lem:binary230} to
\(gh^{-1}\). The product over the places of
\(\sinh(d_{\mathbb D}(y_{g,\tau},y_{h,\tau})/2)\) is at
least \((T/Q)^f\). The products of disks of hyperbolic
radius \(R=\log 2\) about the centres \(y_g\) are therefore
disjoint. An intersection would bound that product by
\(\sinh(R)^f=(3/4)^f\).

Hyperbolic area in radius \(u\), after angular integration,
is \(2\pi\sinh u\,du\), by substituting the Euclidean
radius \(\tanh(u/2)\) in the metric area element.
A disk of radius \(R\) has area \(\pi/2\), and the integral
of \(\exp(-3d_{\mathbb D}(z,y)/2)\) over the whole disk
is \(2\pi\int_0^\infty e^{-3u/2}\sinh u\,du=8\pi/5\),
independently of its centre \(z\).
For \(y\) in the product about \(y_g\), each distance to
\(z_\tau\) exceeds the distance at the centre by at most
\(R\). Integrating the product of the exponential weights
over these disjoint sets gives
\[
 \sum_{g\in\Gamma}\Phi(gz)^{-3}
 \le\left(e^{3R/2}\frac{8\pi/5}{\pi/2}\right)^f<10^f.
\]
One may first sum over any finite subset and take the
supremum, so the estimate also proves convergence.

For \(g\ne1\) the triangle inequality gives
\(d_{\mathbb D}(0,z_\tau)+d_{\mathbb D}(0,(gz)_\tau)
 \ge d_{\mathbb D}(0,y_{g,\tau})\). The exponential of
half the latter distance is
\(\sqrt{1+|\tau c|^2}+|\tau c|\), since
\(d_{\mathbb D}(0,g^{-1}0)=d_{\mathbb D}(0,g0)\).
Multiplying at the places yields
\(\Phi(z)\Phi(gz)\ge(2T/Q)^f\).
Apply this bound to the power \(p-3\) and sum the remaining
power by the first assertion.
\end{proof}

The orbital symmetry for these frames and the resulting rank estimate are proved together in Section~\ref{sec:orbital133} and Theorem~\ref{thm:wordrank134}.

\subsection{Zero Gram determinant and separation of isotropic subspaces}\label{nav:VII-03}
\label{sec:success135}

Keep the two rational frames of Section~\ref{sec:gram132}. Write
\(L_i=\langle u_1,\ldots,u_i\rangle\), put \(L_0=0\), and use
primed notation for the second flag. Define the \(E\)-linear
successor \(S:L\to Y\) by \(Su_i=u_{i+1}\) for \(i<s\) and
\(Su_s=z\). Define \(S'\) in the same way. The two characters
are taken with these marked terminal vectors, possibly of
different negative norms. Call the pair relevant when the two
real characters agree on the common unipotent group, as in
\eqref{eq:relevant133}.

\begin{theorem}[The common part of two relevant flags]
\label{thm:success135}
Suppose \(I=L\cap L'\) has dimension \(d>0\) and the pair is
relevant. For every nonzero \(x\in I\), the least \(p\) for
which \(x\in L_p\) is also the least \(q\) for which \(x\in L'_q\).
The set of these indices is precisely \(s-d+1,\ldots,s\).
The intersection \(Y\cap Y'\) contains a totally negative
rational vector. Moreover \(\operatorname{rank}h|_{L\times L'}
=s-d\).
\end{theorem}
\begin{proof}
Take a nonzero \(x\in I\) with the two indices \(p,q\).
For \(v\in (L_p+L'_q)^\perp\) consider the skew-adjoint operator
\[
       X(w)=x h(v,w)-v h(x,w).
\]
For \(w\in L_i\), the second term is zero. The first vanishes
for \(i\le p\), and for \(i>p\) it belongs to \(L_p\subset
L_{i-1}\). On \(L^\perp\) the image belongs to \(L\).
These relations and skew adjointness imply \(X\in\operatorname{Lie}N\).
Indeed on the flag, its middle quotient, and the dual flag its
diagonal actions are zero, making it strictly upper triangular
in the completed Witt flag. The same argument for \(L'_q\)
puts it in the other unipotent Lie algebra as well.

Write \(x=\sum_i x_i u_i=\sum_i x'_i u'_i\). Evaluating the
simple-coordinate functional on \(X\) gives
\(\sum_i x_i h(v,Su_i)=h(v,Sx)\). This includes the coefficient
on the marked terminal vector since \(h(x,z)=0\). The primed
functional is \(h(v,S'x)\). By relevance their real traces agree.
Replacing \(v\) by \(\alpha v\) for all \(\alpha\in E\) shows
that \(h(v,Sx-S'x)=0\). For example, at one complex embedding
one can test a rational value and its product with a nonzero
purely imaginary element of \(E\). The two real traces detect
both components. Taking the second perpendicular gives the
rational identity
\begin{equation}
                Sx-S'x\in L_p+L'_q .
                \label{eq:success135}
\end{equation}
Thus one may subtract an element of \(L_p\) from \(Sx\) to
obtain the same vector as \(S'x\) plus an element of \(L'_q\).
If \(p,q<s\), this vector lies in \(I\), and its two indices
are \(p+1,q+1\). The leading coefficients at these indices are
\(x_p,x'_q\), which are nonzero. If \(p=s\) and \(q<s\), its
first expression is \(x_s z\) plus a vector of \(L\), hence has
strictly negative norm. Its second expression lies in \(L'\)
and has norm zero. This contradiction also excludes \(p<s,q=s\).
Iterating the first case therefore forces \(p=q\). Starting with
any nonzero \(x\) and repeating until both indices are \(s\),
the last application of \eqref{eq:success135} gives an element
of \(Y\cap Y'\) of norm \(x_s\bar x_s h(z,z)\), with \(x_s\)
the then-current nonzero leading coefficient.

The filtration \(I\cap L_i\) can increase dimension by at most
one at each index. Each occupied index below \(s\) gives the next
one by the construction, and the total number of increases is
\(d\). This proves the asserted index set. Finally both \(L,L'\)
lie in the perpendicular of the common negative vector. At any
real place this perpendicular has negative index \(s\).
If \(v\in L'\cap L^\perp\), the totally isotropic span \(L+Ev\)
has dimension at most \(s\), forcing \(v\in I\).
The kernel of the cross pairing on \(L'\) is therefore \(I\),
which proves the rank statement.
\end{proof}

\begin{corollary}[Vanishing of the mixed terms with zero determinant]
\label{cor:zero135}
For the convergent cells of Section~\ref{sec:orbital133} and two
different words, every cell with \(\mathcal G=0\) has zero real
operator. The same conclusion holds for the two available
frequencies at the unrefined level \(K_\sigma\).
\end{corollary}
\begin{proof}
If the characters disagree, integration over their common
unipotent fibre gives zero by Lemma~\ref{lem:zero133}.
Otherwise, for \(d>0\) Theorem~\ref{thm:success135} gives a
common negative vector. For \(d=0\) the determinant test in
Theorem~\ref{thm:cross132} gives the same conclusion. In both
cases the common compact group has rank \(k\). Theorem~\ref{thm:common133}
now gives zero because the column types are inequivalent.
Its fibre and compact averaging arguments used character
compatibility and trivial finite column stabilizers, both of
which also hold for \(K_\sigma\) by Lemma~\ref{lem:covol127}.
The differential calculation above uses the actual vector
in each flag and applies to both \(z=e\) and \(z=\varpi e\).
\end{proof}

Here is a size bound that includes identical words. Choose
the standard Euclidean norm at each complex embedding over
\(\tau:F_+\hookrightarrow\mathbb R\), in the original signed
orthonormal coordinates of \(h\). Put
\[
 h_\tau(g)=\max_{1\le i\le s}
 \frac{\|g^{-1}(u_1\wedge\cdots\wedge u_i)\|}
      {\|g^{-1}(u_1\wedge\cdots\wedge u_{i-1})\|},\qquad
 H(g_\infty)=\biggl(\prod_\tau h_\tau(g_\tau)\biggr)^{1/f},
\]
using the Euclidean exterior norm and denominator one for \(i=1\).
These are positive continuous functions on the real period
quotients. The group \(B_\infty\) fixes each oriented exterior
product. At one place the \(i\)-th ratio is the length of the
orthogonal projection of \(g^{-1}u_i\) off the preceding vectors,
by Gram--Schmidt. Equivalently \(h_\tau(g)\) is the smallest
possible maximum of the Euclidean lengths after replacing each
\(u_i\) by \(u_i\) plus any combination of preceding vectors.
Each minimizing change exists, with the coefficients chosen
independently for different columns.

\begin{theorem}[A height separation for different isotropic spaces]
\label{thm:height135}
Suppose \(s\ge2\), \(K=K_\sigma\) or \(K^0\),
\(\gamma_f\in B_f K B_f\), and \(L'=\gamma L\ne L\).
If the two frames are relevant, then at every \(g_\infty\)
\begin{equation}
              H(g_\infty)H(\gamma^{-1}g_\infty)\ \ge\ g_1g_2 .
              \label{eq:height135}
\end{equation}
The assertion holds for any supported choice of terminal
frequencies at the indicated level.
\end{theorem}
\begin{proof}
First \(h|_{L\times L'}\) is nonzero. If \(d>0\), this follows
from Theorem~\ref{thm:success135} since \(d<s\). If \(d=0\)
and the pairing vanishes, \(L+L'\) is totally isotropic of
dimension \(2s>s+1\), larger than the negative index of \(V\).

Choose \(i,j\) with \(i+j\) as small as possible such that
\(\xi=h(u_i,\gamma u_j)\ne0\). This entry is unchanged by
independent upper unitriangular changes in both isotropic bases.
Indeed every other term in its expansion has strictly smaller
sum of indices. The entire zero region of smaller sums is
preserved by the changes and their inverses. As in
Theorem~\ref{thm:gap132}, factor \(\gamma_f=b_1 t b_2\),
\(t\in K\). Moving \(b_1\) to the first argument and using
these triangular actions shows that \(\xi=h(u_i,t u_j)\)
in every finite completion. The vectors \(u_i\) are integral
at both levels. Thus \(\xi\in\mathcal O_E\). At either auxiliary
prime the reduction of \(t\) preserves the isotropic flag in
both split factors, as checked in Theorem~\ref{thm:gap132}.
Consequently \(\xi\in(g_1g_2)\mathcal O_E\). Its nonzero
field norm gives
\[
                   \prod_\tau|\tau(\xi)|\ge(g_1g_2)^f ,
\]
where either complex embedding above \(\tau\) gives the same
absolute value.

Invariance under \(g_\tau^{-1}\) writes \(\tau(\xi)\) as the
pairing of a vector from each of the two transformed flags
\(g_\tau^{-1}L, g_\tau^{-1}\gamma L\).
The same smallest-index observation permits the two independent
Gram--Schmidt changes just described. In signed coordinates
\(|h(v,w)|\le\|v\|\|w\|\) because the signature matrix has
Euclidean operator norm one. It follows that
\(|\tau(\xi)|\le h_\tau(g_\tau)h_\tau(\gamma^{-1}g_\tau)\).
Multiplication proves \eqref{eq:height135}.
\end{proof}

The source defines the common Lie-algebra condition and the
skew-adjoint tests in
\srcloc{Formalism/SharpRestriction/CuspSpanNormBounds/BlockEnds}{79}
and line~83 there. The mixed zero-determinant assertion also
occurs in
\srcloc{Formalism/SharpRestriction/CuspSpanNormBounds/BruhatCellsKronecker/NormAdditiveCharacterTypePrimeEqOne}{386}.

\subsection{Summing the period seeds above the height cutoff}\label{nav:VII-04}
\label{sec:angle149}

Use the period fields, ranks and levels of
Section \ref{nav:I-28}.
Write \(f=[F_+:\QQ]\), \(r=\dim_E V\), \(s=\dim_E L\),
\(Q_0=g_1g_2\), and \(R_0=\sqrt{Q_0}\). We take \(s\ge2\),
as on all sufficiently distant intervals. At the real place
\(\tau\), use the Euclidean norm of the signed coordinates in
Section~\ref{sec:success135}. For a complex subspace \(A\)
write \(p_A\) for its Euclidean orthogonal projection and put
\[
  d_\tau(A,A')=\|p_A-p_{A'}\|_{\rm op},\qquad
  A_\tau(g)=g_\tau^{-1}L_\tau .
\]
The projection distance is at most one. Indeed the self-adjoint
operator \(p_A-p_{A'}\) lies between \(-1\) and \(1\) as a
quadratic form. Both \(A_\tau\) and the height \(H\) are
well defined on the real seed quotient. Write \(P\) initially
for ordinary summation of compactly supported seeds over
\(B(F_+)\backslash G(F_+)\).

\begin{lemma}[The size of a crossing in terms of the angles]
\label{lem:angle149}
Let the two finite supports meet and suppose the two isotropic
spaces \(L\) and \(\gamma L\) are distinct and the two unipotent
characters agree on their common real subgroup. At either of
the levels in Theorem~\ref{thm:height135},
\[
 Q_0^{\,f}\ \le\
     \{H(g)H(\gamma^{-1}g)\}^{f}
      \prod_\tau d_\tau\big(A_\tau(g),A_\tau(\gamma^{-1}g)\big).
\]
\end{lemma}
\begin{proof}
The finite support condition gives \(\gamma_f\in B_f K B_f\).
The proof of Theorem~\ref{thm:height135} produces a nonzero
entry \(\xi=h(u_i,\gamma u_j)\in Q_0\mathcal O_E\) with
\(i+j\) minimal. Its norm gives
\(\prod_\tau|\tau(\xi)|\ge Q_0^{\,f}\).
At a real place put \(A=A_\tau(g)\), \(A'=A_\tau(\gamma^{-1}g)\).
The two independent Gram--Schmidt operations from that proof
give \(v\in A,w\in A'\) with \(h(v,w)=\tau(\xi)\) and lengths
bounded respectively by \(h_\tau(g)\) and
\(h_\tau(\gamma^{-1}g)\). All earlier cross entries vanish,
so these operations preserve this particular value. The
Hermitian form vanishes on \(A\times A\). Thus
\[
 |h(v,w)|=|h(v,w-p_Aw)|
  \le \|v\|\,\|(p_{A'}-p_A)w\|
  \le \|v\|\,\|w\|\,d_\tau(A,A').
\]
Here we used \(p_{A'}w=w\) and the operator norm one of the
signature matrix. Multiplication over the real places proves
the claim.
\end{proof}

\begin{theorem}[An upper bound at arbitrary height]
\label{thm:sectors149}
Fix a real place \(\tau_0\) and \(R>0\).
For seeds supported on \(H\le R\) there is a measurable
partition based solely on \(A_{\tau_0}(g)\) into at most
\[
       \bigl(1+16\max\{1,(R^2/Q_0)^f\}\bigr)^{2rs}
\]
parts. On each part full automorphic summation is an isometry
from the seed norm. For every square-integrable seed of the
selected types supported at these heights,
\begin{equation}
     \|P\Phi\|_{L^2([G])}
       \le \bigl(1+16\max\{1,(R^2/Q_0)^f\}\bigr)^{rs}
                  \|\Phi\|_{\rm seed}.
       \label{eq:sectors149}
\end{equation}
The assertion includes finite linear combinations of words
and class columns at a fixed supported level and frequency, and
the two supported frequencies together at \(K_\sigma\).
\end{theorem}
\begin{proof}
Set \(\eta=\tfrac18\min\{1,(Q_0/R^2)^f\}\). The complex
\(r\)-by-\(s\) matrices, equipped with the operator norm,
form a real normed space of dimension \(2rs\). Choose a
maximal \(\eta\)-separated family among the orthonormal
\(s\)-frames in that space. The balls of radius \(\eta/2\)
about them have disjoint interiors and lie in the ball of
radius \(1+\eta/2\) about zero. Lebesgue volume bounds the
number by \((1+2/\eta)^{2rs}\). Such a finite maximal family
exists by successively adding points and using this bound.
Every plane has a frame \(U\) within \(\eta\) of some chosen
frame \(V\). Then
\(\|UU^*-VV^*\|\le2\|U-V\|\le2\eta\), by subtracting one
factor at a time. The closed projection balls of radius
\(2\eta\) about the chosen planes therefore cover. Assign a
plane to the first such ball that contains it. The resulting
Borel parts have diameter at most \(4\eta\).

Two indices contributing at the same real point to one part
have their subspaces at \(\tau_0\) within \(4\eta\) and their
heights at most \(R\). All other projection distances are
at most one. If the spaces are distinct and the relative
cell is relevant, Lemma~\ref{lem:angle149} would give
\[
        Q_0^{\,f}\le R^{2f}4\eta\le Q_0^{\,f}/2 .
\]
If the cell is irrelevant, its unfolded integral is zero by
Lemma~\ref{lem:zero133}. The imposed height and angle
indicators preserve this cancellation. They are left
\(B_\infty\)-invariant, and in an unfolded product they are
constant along the common unipotent fibre.

First take bounded measurable sections of compact support
modulo \(B_\infty\) in the chosen part. Their scalar envelopes
have compact support modulo \(B(F_+)\), by the finite class
description. The ordinary sums are locally finite and products
of the periodized envelopes are integrable, by the compact
lifting argument of
Theorem \ref{thm:maximal130}.
The Gram integral can therefore be unfolded. By the preceding
paragraph only the terms with the same isotropic space
survive. Their sum is the norm squared of summation over
the stabilizer of \(L\), which equals the seed norm squared
by the same theorem and by Theorem~\ref{thm:midnorm131} at
the refined level. That isometry extends to measurable
truncations: on a fixed compact support the ordinary sums are
locally finite with uniformly bounded multiplicity on the
compact image, so \(L^2\) approximation agrees with ordinary
summation there. Bounded compact truncations within a part
are dense in the \(L^2\) space of its sections. This proves
the isometric extension on each part. The finite partition is
orthogonal for the seed norm by disjoint supports. Apply the
triangle inequality to the sum of its images, then
Cauchy--Schwarz over the parts. Taking the square root of the
covering bound proves \eqref{eq:sectors149}. All extensions
agree on intersecting height ranges, by their agreement on
bounded compact sections and density.
\end{proof}

\begin{theorem}[Full summation with a weighted real tail]
\label{thm:tail149}
Set \(D_0=2frs\), fix \(\sigma>0\), and write
\[
  C(D_0,\sigma)=
        \frac{2\cdot17^{rs}}{\sqrt{1-2^{-2\sigma/D_0}}}.
\]
For a square-integrable seed \(\Phi\) let
\(\Phi^{<}=\mathbf1_{H<R_0}\Phi\). If
\[
   T(\Phi)=\big\|\mathbf1_{H\ge R_0}
               (H/R_0)^{D_0+\sigma}\Phi\big\|_{\rm seed}<\infty ,
\]
the ordinary sums of bounded compact truncations converging
to \(\Phi\) in this weighted \(L^2\) norm and in the seed norm
have a unique limit \(P\Phi\) in \(L^2([G])\). It satisfies
\begin{equation}
       \|P\Phi-P\Phi^{<}\|\le C(D_0,\sigma)T(\Phi).
       \label{eq:tail149}
\end{equation}
If \(\Psi\) is supported in \(H<S\) with \(0<S\le R_0\),
the difference between \(\langle P\Phi,P\Psi\rangle\) and
\(\langle\Phi,\Psi\rangle_{\rm seed}\) is bounded by
\[
 C(D_0,\sigma)\|\Psi\|_{\rm seed}
        \big\|\mathbf1_{H\ge Q_0/S}
              (H/R_0)^{D_0+\sigma}\Phi\big\|_{\rm seed}.
\]
\end{theorem}
\begin{proof}
Write \(\Phi_j\) for the part with
\(2^{j/D_0}R_0\le H<2^{(j+1)/D_0}R_0\), \(j\ge0\).
For \(R\ge R_0\) the factor in \eqref{eq:sectors149} is at
most \(17^{rs}(R/R_0)^{D_0}\). Hence
\[
 \sum_{j\ge0}\|P\Phi_j\|
 \le 2\cdot17^{rs}\sum_{j\ge0}
       2^{-\sigma j/D_0}\,
         2^{(1+\sigma/D_0)j}\|\Phi_j\|
 \le C(D_0,\sigma)T(\Phi).
\]
The second inequality follows by Cauchy--Schwarz and the
definition of the disjoint shells. Their images sum absolutely
in the Hilbert space. Below \(R_0\) the isometry in Theorem \ref{thm:low136}, applies. This proves the estimate and continuity
in the sum of the two stated norms. Truncation in support and
in the pointwise norm gives density in this space, so the
extension is unique.

To test against \(\Psi\), first restrict \(\Phi\) to
\(H<Q_0/S\). On a pair of supports the two heights now have
product less than \(Q_0\). The same-space calculation in
the preceding proof applies, and Theorem~\ref{thm:height135}
excludes every relevant different-space term. Compact
approximation followed by \eqref{eq:sectors149} at the finite
height \(Q_0/S\) proves equality of the pairings on this part.
The remainder has zero seed pairing with \(\Psi\) because
\(Q_0/S\ge S\). Its automorphic pairing is at most the product
of its norm from \eqref{eq:tail149} and the seed norm of
\(\Psi\), by the low-height isometry. This proves the last
assertion.
\end{proof}

\subsection{The full projected period at fixed finite level}\label{nav:VII-08}
\label{sec:adjoint157}

Fix one period field, one translate, the finite level \(K^0\),
and the terminal frequency \(\varpi\). The selected words \(c\)
give finitely many inequivalent compact column types \(\rho_c\)
for a common real discrete series \(\sigma\). Let \(U\) be its
product lowest compact type and \(\mathcal M\) its automorphic
multiplicity Hilbert space at this finite level. Thus the
\(\sigma\)-isotypic space is \(\sigma\widehat\otimes\mathcal M\).
Denote by \(\Pi\) the projection onto \(U\otimes\mathcal M\).
All projections and adjoints in this section are Hilbert space
operators.

Let \(S_c\) denote the space of square-integrable real Bessel
sections of type \(\rho_c\). Give the finite class covector space
\(D_c\) its sum inner product multiplied by \(c_{\rho_c,K^0}\)
from \eqref{eq:main127}; equivalently, this is the inner product
induced by the seed for a real section of norm one. The finite
class subspace with zero sums on scalar orbits can be used in
its place. Write
\[
     S=\bigoplus_c (S_c\widehat\otimes D_c),\qquad
     D=\bigoplus_c D_c,\qquad
     J_U:U\otimes D\longrightarrow S .
\]
On the \(c\)-th summand \(J_U\) is \(J_c|_U\otimes 1_{D_c}\),
where \(J_c:\sigma\to S_c\) is the isometry in
Corollary~\ref{cor:rel155}. The map from \(S\) to actual
period seeds preserves inner products. Indeed polarization of
Theorem~\ref{thm:identity127} in the two real sections gives
their \(S_c\) inner product times the stated class inner product,
and distinct words are orthogonal. Finite sums of pure tensors
are dense.

\begin{lemma}[The mass of a compact type in a measurable region]
\label{lem:weights157}
If \(\Omega\subset B_\infty\backslash G_\infty\) is right
invariant under the maximal compact subgroup, there are
numbers \(0\le t_c(\Omega)\le1\) such that
\[
  J_U^*\mathbf1_\Omega J_U=1_U\otimes T_\Omega,\qquad
  T_\Omega|_{D_c}=t_c(\Omega)\,1_{D_c}.
\]
For any unit \(v\in U\),
\(t_c(\Omega)=\int_\Omega\|(J_cv)(g)\|^2\,d\dot g\).
For \(\Omega_R=\{H<R\}\) one has \(t_c(\Omega_R)\uparrow1\)
as \(R\to\infty\), for all \(c\) simultaneously.
\end{lemma}
\begin{proof}
The compression of \(\mathbf1_\Omega\) to \(J_c(U)\)
commutes with the compact group. On this irreducible type it
is a scalar by Schur's lemma. The diagonal quadratic form
on a unit vector gives the displayed integral. Multiplication
by the indicator preserves each word and acts trivially on
its class factor. The height is finite at every point and
right compact-invariant. Monotone convergence gives the limit,
and the word set at the fixed field and level is finite.
\end{proof}

\begin{theorem}[A bounded extension of the projected sum]
\label{thm:adjoint157}
The ordinary sum followed by \(\Pi\), initially on compact
seeds, has a unique bounded extension \(P_U:S\to U\otimes
\mathcal M\). There is a bounded operator \(A:\mathcal M\to D\)
such that
\begin{equation}
      P_U^*=J_U(1_U\otimes A),\qquad
      P_U=(1_U\otimes A^*)J_U^*,\qquad \|P_U\|=\|A\|.
      \label{eq:adj157}
\end{equation}
For every \(R\) with \(t_*(R)=\min_c t_c(\Omega_R)>0\),
\begin{equation}
       A^* T_{\Omega_R} A\ \le\ B(R)^2\,1_{\mathcal M},
       \qquad \|A\|\le B(R)/\sqrt{t_*(R)} .
       \label{eq:bound157}
\end{equation}
The bound concerns the fixed field and level and also holds
when the automorphic multiplicity is zero.
\end{theorem}
\begin{proof}
For \(m\in\mathcal M\), \(v\in\sigma^\infty\), let
\(F=v\otimes m\) in automorphic \(L^2\). On seeds of height
less than \(R\) the functional \(\Phi\mapsto\langle P\Phi,F\rangle\)
has norm at most \(B(R)\|v\|\|m\|\) by
Theorem~\ref{thm:sectors149}. For fixed \(v,m\) the Riesz
representatives agree on overlaps. They define a locally
square-integrable real section \(W_{v,m}\) with values in the
appropriate column and finite class spaces. Right translation
in the real argument replaces \(v\) by \(\sigma(g)v\).
Indeed test this equality on smooth compact seeds, use the
right equivariance of ordinary summation, then translate in
the pairing. Consequently each weak right derivative of
\(W_{v,m}\) is represented by \(W_{d\sigma(X)v,m}\), and
the same assertion holds for arbitrary words.

In a relatively compact coordinate chart on the quotient, lift
sections to a local trivialization. Right Lie directions span
the quotient tangent space. Coordinate derivatives of each
fixed order are therefore linear combinations with locally
bounded smooth coefficients of the preceding Lie derivatives,
including lower orders from the bundle trivialization. Local
Sobolev embedding gives a smooth representative and bounds its
evaluation at the identity coset by a constant times a finite
sum of \(\|d\sigma(I)v\|\|m\|\), using one finite height for
a neighborhood of that coset. Thus at each word and class
coordinate the evaluation \(v\mapsto W_{v,m}(1)\) is a
continuous Bessel map on \(\sigma^\infty\). This follows for
the left covariance by evaluating the right-translation
identity at elements of \(B_\infty\).

For each word the space of these maps into the compact column
representation is a line. To check the product assertion, pair
with elementary tensors in the dual column and regard the
result as a functional on the exterior tensor of local row
and dual column representations. Fixing all but one tensor
factor gives a local functional on its known line. Choose in
each factor a vector on which its generator takes the value one.
Applying the one-place assertion successively gives the product
of the local functionals times the value on this chosen tuple.
Such tensors span a dense subspace in the smooth topology:
compact Fourier truncations converge there and each of the
resulting compact isotypic spaces is spanned by product
vectors by admissibility of the local discrete series. Hence
the factorization holds on all smooth vectors.

Apply this fact in an orthonormal basis of \(D_c\), and then
use right translation to recover the section at every \(g\).
There is a unique vector \(A_c m\in D_c\) with
\[
              (W_{v,m})_c=J_c v\otimes A_c m
                        \qquad (v\in\sigma^\infty).
\]
The Riesz representative depends linearly on the second
argument vector \(F\) in this adjoint construction, so \(A_c\)
is linear. Set \(Am=(A_cm)_c\). For a fixed unit \(v\in U\),
the square norm of \(W_{v,m}\) restricted to \(H<R\) is
\(\sum_c t_c(\Omega_R)\|A_c m\|^2\) by Lemma~\ref{lem:weights157}.
Its Riesz bound was \(B(R)^2\|m\|^2\) after squaring. This
proves \eqref{eq:bound157}, and some finite \(R\) has
\(t_*(R)\ge1/2\) by that lemma.

It follows that \(J_U(1_U\otimes A)\) is bounded on \(U\otimes
\mathcal M\). On elementary vectors its pairing with all compact
seeds is exactly the defining pairing for \(\Pi P\). Density
and the adjoint identity give the unique extension and
\eqref{eq:adj157}. The equality of norms follows from the
isometry of \(J_U\). In particular height truncation converges
after projection for every input in \(S\), by continuity.
\end{proof}

\begin{corollary}[Recovering the period from a truncated real vector]
\label{cor:mass157}
For every right compact-invariant measurable region,
\[
           P_U\mathbf1_\Omega J_U=1_U\otimes(A^*T_\Omega).
\]
If \(t_*(R)>0\), the operator on normalized real vectors is
recovered by
\[
       P_U J_U=(P_U\mathbf1_{\Omega_R}J_U)
                         (1_U\otimes T_{\Omega_R}^{-1}).
\]
These identities continue to hold at fixed level on the refined
zero-sum class spaces.
\end{corollary}
\begin{proof}
Substitute the first lemma into \eqref{eq:adj157}. Inversion
of the finitely many positive scalars proves the second identity.
Each acts on the entire class space of its word by a scalar
and thus preserves its zero-sum subspace.
\end{proof}

\subsection{Two terminal frequencies at a common finite level}\label{nav:VII-09}
\label{sec:adelic159}

Keep the actual integral level \(K^0\) of
Theorem~\ref{thm:midlevel131} and write \(K_0=K^0\).
Put \(a=a_\varpi\in G(F_+)\), the rational isometry multiplying
the isotropic vectors \(u_i\) by \(\varpi\), their duals by
\(\bar\varpi^{-1}\), and acting identically on the middle space.
Use \(a_f\) for its finite component, also regarded as an adele
with identity real component. Define
\[
       K_1=a_fK_0a_f^{-1},\qquad L=K_0\cap K_1,\qquad
       J=T^{\,fs(r-s)},\qquad f=[F_+:\QQ].
\]
The subscript on \(K_1\) refers to frequency one. Use the
same class covectors for both frequencies. Write
\(\Phi_{\beta,K}(W,\ell)\) for the seed \eqref{eq:seed127}
with real section \(W\) and class tuple \(\ell\).

\begin{theorem}[Rational change of terminal frequency]
\label{thm:conj159}
The class level \(K_H\) is identical for \(K_0\) and \(K_1\).
At \(K_1\) the unit-frequency character is trivial on the
finite unipotent level. Choose normalized real sections with
\(W_\varpi(g)=J^{-1/2}W_1(a_\infty g)\). Then
\begin{align}
   \Phi_{\varpi,K_0}(W_\varpi,\ell)(g)
     &=J^{-1/2}\Phi_{1,K_1}(W_1,\ell)(a g a_f^{-1}),\notag\\
   \Pi P\Phi_{\varpi,K_0}(W_\varpi,\ell)
     &=J^{-1/2}R(a_f^{-1})\,\Pi P\Phi_{1,K_1}(W_1,\ell).
   \label{eq:conj159}
\end{align}
The second formula holds for all square-integrable real
sections, with projected sums understood by continuous
extension. Their class-coordinate constants in
\eqref{eq:main127} satisfy \(c_{\rho,K_1}=Jc_{\rho,K_0}\).
\end{theorem}
\begin{proof}
Conjugation by \(a\) normalizes \(N\), centralizes \(H=U(W_0)\),
and satisfies \(\chi_1(ana^{-1})=\chi_\varpi(n)\) on the
adelic unipotent group. Indeed all simple coordinates on the
isotropic string are unchanged and the terminal coordinate
is multiplied by \(\varpi\). Thus the finite period subgroups
for the two levels are conjugate with the same \(H\)-projection.
Triviality of the character at \(K_0\) in
Theorem~\ref{thm:midlevel131} proves the assertion at \(K_1\).
At the real places Proposition~\ref{prop:scale151} gives the
stated unitary change of sections and its factor \(J^{-1/2}\).
Conjugation sends \(B_fK_0\) onto \(B_fK_1\) and leaves the
finite class argument in \(H_f\) unchanged. Substituting in
\eqref{eq:seed127} proves the first equality on and off the
finite support.

On compact seeds, left multiplication by \(a\) permutes the
cosets \(B(F_+)\backslash G(F_+)\) since \(a\) normalizes
\(B\). Ordinary summation therefore removes that left factor
and retains the right translation. Real projection commutes
with finite translation, which is unitary on the common
automorphic \(L^2\) space. This proves \eqref{eq:conj159}
on compact supports. For the extension, first compute the
finite norms. The determinant of \(\operatorname{Ad}a\) on
\(\operatorname{Lie}(N)\) is in \(F_+^\times\). Its product
of real absolute values is \(J\) by
Proposition~\ref{prop:scale151}; the compact block in that
calculation is fixed. The product formula gives finite Haar
modulus \(J^{-1}\). Hence \(N_f\cap K_1\) has volume
\(J^{-1}\operatorname{vol}_f(N_f\cap K_0)\). The two row
volumes agree by unimodularity. Formula~\eqref{eq:main127}
now gives the stated relation of class norms. The real change
of variables and this relation make the two seed norms
proportional by a fixed positive factor. Boundedness for the
projected sum at \(K_0\) by Theorem~\ref{thm:adjoint157}
therefore gives boundedness at \(K_1\) on compact seeds.
Density extends the identity.
\end{proof}

\begin{proposition}[The common level at the split prime]
\label{prop:index159}
The indices of \(L\) in \(K_0\) and in \(K_1\) coincide and
are the integer
\[
             I={r\brack s}_T^{\,f},\qquad
       {r\brack s}_T=\prod_{j=1}^{s}
                         \frac{T^{r-s+j}-1}{T^j-1}.
\]
One has \(N_f\cap K_1\subset N_f\cap K_0\) of index \(J\),
and \(B_f\cap L=(N_f\cap K_1)\rtimes K_H\). In particular
both terminal characters are trivial on \(N_f\cap L\).
Furthermore
\begin{equation}
        \frac{c_{\rho,L}}{c_{\rho,K_0}}=\frac{J}{I},
          \qquad 1-\frac f{T-1}\ \le\ \frac{J}{I}\ \le\ 1 .
        \label{eq:index159}
\end{equation}
\end{proposition}
\begin{proof}
Outside \(T\) the levels agree. At odd primes away from the
auxiliary primes, \(a\) and its inverse preserve the standard
lattice in its integral Witt basis. At a dyadic place its
matrix on each signed plane has entries
\((\varpi+\bar\varpi^{-1})/2\) and
\((\varpi-\bar\varpi^{-1})/2\) on the diagonal and off it,
respectively. They are integral, also for the inverse, since
\(\varpi\equiv1\bmod 2\mathcal O_E\). At \(g_1,g_2\) the
imposed congruence of Theorem \ref{thm:common111} gives the
identity reduction, so \(a\) normalizes those levels as
well.

Now let \(v\mid T\) and choose the split factor where
\(\varpi\) has valuation one. The paired unitary group is
described by its first factor \(\mathrm{GL}_r(\QQ_T)\).
Write \(O=\ZZ_T\). By Theorem~\ref{thm:midlevel131},
\(K_{0,v}\) stabilizes
\(\Lambda_0=O^s\oplus M_v\oplus O^s\) in the Witt splitting.
Since \(\bar\varpi\) is a unit here, \(K_{1,v}\) stabilizes
\(\Lambda_1=TO^s\oplus M_v\oplus O^s\).
This sublattice of \(\Lambda_0\) contains \(T\Lambda_0\).
Its image modulo \(T\Lambda_0\) has dimension \(r-s\).
Reduction of \(K_{0,v}\) surjects onto
\(\mathrm{GL}_r(\mathbf F_T)\), and the intersection is the
preimage of the stabilizer of this subspace. Counting the
ordered bases of an \(s\)-dimensional subspace and dividing
by \(|\mathrm{GL}_s(\mathbf F_T)|\) gives the displayed
Gaussian binomial coefficient, also the count in dimension
\(r-s\) by duality. Equal row volumes give the same index
in \(K_{1,v}\). Multiplication over the \(f\) split places
gives \(I\).

An element of \(N_v\cap K_{1,v}\) preserves the entire first
isotropic space. It therefore preserves its intersection
\(TO^s\) with \(\Lambda_1\), and so also \(O^s\). As
\(\Lambda_0=\Lambda_1+O^s\), it preserves \(\Lambda_0\) in
both directions. This proves the inclusion on the finite
unipotents. The quotient of their volumes gives index \(J\).
The assertions for \(B_f\) follow from the common compact
factor and the semidirect product at \(K_0\). Formula
\eqref{eq:main127} gives \(c_{\rho,L}/c_{\rho,K_0}=J/I\).
After removing \(T^{s(r-s)}\) from the Gaussian binomial
coefficient, its reciprocal factor is
\(\prod_{j=1}^{s}(1-T^{-j})/(1-T^{-(r-s+j)})\). It lies
between \(1-\sum_{j\ge1}T^{-j}=1-1/(T-1)\) and one.
Raising to \(f\) proves \eqref{eq:index159}.
\end{proof}

\begin{theorem}[Finite averages of the common-level seeds]
\label{thm:common159}
Let \(e_i\) be probability averaging of right translations by
\(K_i\), \(i=0,1\). For the prime frequency at \(i=0\) and
the unit frequency at \(i=1\), the ordinary seed identities are
\[
 e_0\Phi_{\varpi,L}=\frac{J}{I}\Phi_{\varpi,K_0},
       \qquad e_1\Phi_{1,L}=\frac1I\Phi_{1,K_1}.
\]
The real sections and class covectors in each equality are
unchanged. The finite averaging operator \(I e_i\) is a sum
of integral translations on functions at the common level.
\end{theorem}
\begin{proof}
On an \(L\)-fixed input \(e_i\) is the sum of the right
translations over representatives of \(K_i/L\), divided by
\(I\). At a point of \(B_fK_i\) write the finite argument as
\(bk_0'\). In the average replace \(k_0'k\) by \(k'\).
A term is supported precisely when
\(k'\in(B_f\cap K_i)L\). The number of supported cosets is
\([B_f\cap K_i:B_f\cap L]\), which is \(J\) for \(i=0\)
and one for \(i=1\). In each supported term the character
on the unipotent part of \(B_f\cap K_i\) is trivial for its
designated frequency, and the column is right \(K_H\)-fixed.
The value of every supported term is consequently the
coarse-level value at \(b\). Off \(B_fK_i\) all terms vanish.
This proves both identities. In the unaveraged sum every
coefficient is the integer one. In particular it preserves
integrality of values when those values range in a fixed
integral module.
\end{proof}

\subsection{Absolute convergence of the global Bessel Gram matrix}\label{nav:VII-11}
\label{sec:absolute163}

Put \(F=E^+\), \(f=[F:\QQ]\), \(G=U(V)\), and \(B=N\rtimes H\)
as in Section~\ref{sec:char94}. Work at one of the fields and one
of the translates of Section \ref{nav:I-28}. Write \(r=\dim_E V\)
and \(q=s+1\) for the negative index. All the selected words have
the same real row \(\sigma\). The finite level is \(K=K^0\), the
frequency is \(\varpi\), and \(U\) denotes the lowest compact
type of \(\sigma\). We use the ordinary seed of
\eqref{eq:seed127}, with its right level character when present.
The class spaces \(D_c\) and the isometry \(J_U\) have the
normalization of Section~\ref{sec:adjoint157}.

In the construction preceding Theorem \ref{thm:common111}, choose
the two free spacings as follows. Put
\[
 d=4r^2,\qquad L=8(f+1)(r+1)^3,\qquad
 M=L+2(d+2),\qquad \eta=L-4fr^2>0 .
\]
Take \(\Delta\ge M+1\) and \(H_0\ge M+1\) in addition to the
lower bound already imposed on \(H_0\). For instance one can
take the least integers meeting these conditions and the earlier
lower bounds. The three successive kinds of gap in the real row
are \(4\Delta,\ B_\tau,\ 2H_0\). The least \(B_\tau\) grows
faster than every fixed power of \(x\), while \(f,r,H_0,\Delta\)
are polynomially bounded for this choice. Every adjacent gap
therefore exceeds \(M\), at every real place of all sufficiently
large fields in the stated intervals. This choice retains the
interlacing, the common terminal prime construction and the
inequality \(4k\Delta\le B_\tau/x\) of Volume~I.

\begin{lemma}[Polynomial decay from gaps in the real row]
\label{lem:translate163}
Let \(G_0=U(p,q)\) in signed orthonormal coordinates and let
\(w(g)=\|g\|_{\rm op}\) for its defining complex representation.
If every adjacent gap of a decreasing Harish-Chandra parameter
\(\lambda\) exceeds the integer \(M\), every coefficient \(a(g)\)
of its discrete series between two compact-finite vectors satisfies
\[
                  |a(g)|\le C_{a,M} w(g)^{-M}.
\]
For \(S=G_0\) or \(S=B_\tau\), one has
\(\int_S w(h)^{-D}\,dh<\infty\) for \(D>4r^2\).
\end{lemma}
\begin{proof}
We use the coherent character theorem for discrete series. In
a fixed regular chamber the characters with the same compact
Cartan and chamber labels form the restriction of a coherent
family. Multiplying a member by a finite-dimensional character
gives the sum of the members shifted by its weights, counted
with multiplicity. This is the character identity (8.3) of
Schmid and Vilonen, applied to the open orbit for the given
discrete series. Its identification in the regular chamber is
given in their Sections 7--8.~\cite{XVII-XII-IV-sv163}

Tensor the Harish-Chandra module once with the defining
representation. Its weights are the \(r\) distinct coordinate
vectors \(e_i\). When all shifts stay in the chosen chamber,
the character identity gives the discrete series with parameters
\(\lambda+e_i\), each once. They have distinct infinitesimal
characters. Indeed two strictly decreasing lists in the same
permutation orbit coincide. Characters determine multiplicities
in a finite-length Harish-Chandra module. Splitting this module
into generalized eigenspaces for the center of the enveloping
algebra now gives a direct sum, since each nonzero block has
length one. After \(m\) shifts the difference between the
changes in two entries has absolute value at most \(m\).
Consequently we may repeat this argument through \(M\) tensor
factors. Every summand of the resulting module is unitarizable.

For any matrix entry \(g_{ij}\), the function
\(g_{ij}^{M}a(g)\) is a coefficient of this tensor product.
In the algebraic direct sum its vector has finitely many
components. Its compact-finite covector on each summand is
represented by a Hilbert vector, since it is supported on
finitely many compact types of finite multiplicity. The
original coefficient equals the resulting sum of unitary
coefficients on the whole group. Indeed both sides are real
analytic by the analytic-vector theorem for compact-finite
vectors of an irreducible unitary representation. The
intertwining of Lie algebra modules gives all derivatives
at the identity, and \(U(p,q)\) is connected. Each summand
is bounded by the product of its two vector norms. Taking
the maximum over \(i,j\) and using
\(w(g)\le r\max_{i,j}|g_{ij}|\) proves the first assertion.

In signed coordinates \(g^{-1}=Jg^*J\), with \(J\) a unitary
involution. Hence \(w(g)=w(g^{-1})\ge1\). In the ball
\(w(g)\le R\) of a closed subgroup \(S\), take a maximal
\(1/(2R)\)-separated set for the matrix norm. Packing balls
in the real vector space of complex \(r\)-by-\(r\) matrices
bounds its cardinality by \(C_rR^{4r^2}\) for \(R\ge1\).
Each point of the original ball is within \(1/(2R)\) of a
center \(h\), so \(h^{-1}g\) lies in the compact set
\(\{u\in S:\|u-1\|_{\rm op}\le1/2\}\).
Left invariance bounds the Haar volume by \(C_SR^{4r^2}\).
Summing on dyadic shells proves the integral assertion.
\end{proof}

At a real place write
\[
 \nu_\tau(g)=\inf_{b\in B_\tau}w(bg),\qquad
 V(g_\infty)=\prod_\tau\nu_\tau(g_\tau).
\]
These functions are at least one. They are continuous because
right multiplication by \(t\) changes \(\nu_\tau\) by a
factor between \(w(t^{-1})^{-1}\) and \(w(t)\). The same
letter \(V\) with a group element as argument will always
denote this product.

\begin{proposition}[Decay and reproduction of the Bessel sections]
\label{prop:weight163}
For every word \(c\) and compact-finite vector \(v\) of the
real row, \(W=J_cv\) satisfies
\[
  \sup_g V(g)^L\|W(g)\|<\infty,\qquad
  \int_{B_\infty\backslash G_\infty}\|W(g)\|V(g)^L\,d\dot g<\infty .
\]
In particular \(\|H^{\,2D_0}W\|_2<\infty\) for the height of
Section~\ref{sec:angle149} and \(D_0=2frs\). On the lowest
compact type, one compactly supported smooth right convolution
reproduces all these sections.
\end{proposition}
\begin{proof}
At one place use the scalar realization in the proof of
Theorem~\ref{thm:rel155}. For \(z\) in the dual compact column
and a compact-finite \(u\) with \(\ell(u)\ne0\), each coordinate
\(z(W(g))\), up to a fixed nonzero normalization, is
\(Qf_{v\otimes z,u}(g,1)\). Here \(Q\) integrates the coefficient
over the diagonal period group \(\mathcal H\subset
G_\tau\times H_\tau\) against its unitary character.
Projection onto its first factor identifies \(\mathcal H\)
with \(B_\tau\). The second factor is compact. Expanding \(u\)
in a basis of the compact column and applying
Lemma~\ref{lem:translate163} bounds the absolute value of the
integrand by \(Cw(bg)^{-M}\), uniformly in its compact variable.

Choose \(g\) in its left coset with \(w(g)\le2\nu_\tau(g)\).
Write \(D=d+2\). From \(w(bg)\ge\nu_\tau(g)\) and
\(w(bg)\ge w(b)/w(g^{-1})\) we obtain
\[
 w(bg)^{-L-2D}\le 2^D\nu_\tau(g)^{-L}w(b)^{-D}.
\]
Integration over \(b\) proves the pointwise bound for the
norm, which is constant on the left coset. For the integral
bound use the quotient integration formula on the nonnegative
majorant. The weight is unchanged along the left fibre and
\(\nu_\tau(g)\le w(g)\), giving an upper bound
\(C\int_{G_\tau}w(y)^{-M+L}\,dy<\infty\).
Apply these inequalities to finitely many coordinates, then
take tensor products over the real places. Compact-finite
vectors are finite sums of such tensors by admissibility.

Each Gram--Schmidt length in \(h_\tau\) is at most
\(\|g^{-1}u_i\|\le w(g)\|u_i\|\). Left invariance yields
\(H\le C V^{1/f}\). Thus \(H^{4D_0}\le C'V^{8rs}\) and
\(8rs\le L\). The two estimates just proved imply the required
square integral. For reproduction, conjugation-average a smooth
compactly supported approximate identity over the maximal
compact group. The resulting operators commute with that
group and tend strongly to the identity. The lowest compact
type occurs once. On it one of these operators is therefore a
nonzero scalar. Divide the corresponding function by that
scalar and use right equivariance of \(J_c\). The identity
holds pointwise by smoothness. In particular the exponent
\(4r^2\sqrt{2q}\) required in
\srcloc{Formalism/LocalToGlobal/SiegelLatticeTheory/HaarTransversals/WuDecay}{31}
is covered since it is at most \(L\).
\end{proof}

Let \(\Phi\) denote one of the resulting scalar seeds, including
its finite class covector. It vanishes outside \(B_fK\) in the
finite variable and is bounded there in modulus by \(C_\Phi
V(g_\infty)^{-L}\). We give a count uniform in the point of the
automorphic quotient. Its constant may depend on the number field,
the level and the chosen coordinates.

\begin{lemma}[Rational cosets in a height ball]\label{lem:cosets163}
For \(R\ge1\) and \(t\in G(\mathbb A_F)\) the number of indices
\(\gamma\in B(F)\backslash G(F)\) with \(\gamma_f t_f\in B_fK\)
and \(V(\gamma_\infty t_\infty)\le R\) is at most \(C_xR^{4fr^2}\).
Consequently the sum of \(|\Phi(\gamma t)|\) over indices of
height at least \(R\) is at most \(C_{\Phi,x}R^{-\eta}\).
\end{lemma}
\begin{proof}
The adelic quotients of \(N\) and \(H\) are compact, as proved
in Section~\ref{sec:char94}. Their compact sets of representatives
give one for \(B(F)\backslash B(\mathbb A_F)\). Reduce the
\(H\)-coordinate first, then the \(N\)-coordinate. For \(y_f\in
B_fK\), reduce \(y_f\) to \(K\) by \(B_f\) and reduce each
\(y_\tau\) to a representative of norm \(\le2\nu_\tau(y_\tau)\)
by \(B_\tau\). Removing the rational part of the combined left
multiplier gives \(\delta\in B(F)\) such that
\begin{equation}
       (\delta y)_f\in C_f,\qquad
       w((\delta y)_\tau)\le C_\tau\nu_\tau(y_\tau).
       \label{eq:lifts163}
\end{equation}
The finite compact set and all constants here are fixed.

Choose such a representative for each contributing index and
fix one of them as base point. Relative to the base point the
other rational matrices are distinct, have finite components
in \(C_fC_f^{-1}\), and have archimedean norms \(\ll_x R^2\).
Every entry belongs to a fixed fractional ideal of \(E\).
In its \(2f\)-dimensional real embedding an ideal is a lattice.
Disjoint balls of a fixed positive radius about its points in a
box of radius \(O_x(R^2)\) give \(O_x(R^{4f})\) possible entries.
There are \(r^2\) entries. This proves the count, including the
case where it is zero. Sum the bound \(C_\Phi(2^j R)^{-L}\)
times the count at \(2^{j+1}R\) over \(j\ge0\) for the tail.
\end{proof}

\begin{theorem}[The ordinary sum and its projected Gram matrix]
\label{thm:ordinary163}
On \(J_U(U\otimes D)\) the ordinary rational periodization \(P\)
converges absolutely at every point. It is a bounded function in
automorphic \(L^2\) and takes values in the prescribed real
isotypic space and lowest compact type. Its Gram matrix is the
one in Theorem~\ref{thm:adjoint157}. More explicitly, for a
common unit \(v\in U\) and finite class vectors \(\ell,\ell'\),
write \(\Phi,\Psi\) for their seeds. With \(P_{\rm fl}\) the
rational flag parabolic, one has
\begin{align}
 \langle P\Phi,P\Psi\rangle
 &=\langle\ell,\ell'\rangle_D+
     \sum_{\substack{\gamma\in B(F)\backslash G(F)/B(F)\\
                      \gamma\notin P_{\rm fl}(F)}}
        \int_{\Gamma_\gamma\backslash G(\mathbb A_F)}
          \Phi(t)\overline{\Psi(\gamma^{-1}t)}\,dt ,
       \label{eq:geom163}\\[-3pt]
 &\hspace{25mm}\Gamma_\gamma=B(F)\cap\gamma B(F)\gamma^{-1}.
 \notag
\end{align}
The sum of the integrals of absolute values is finite. For
different words all cells of zero Gram determinant in
Corollary~\ref{cor:zero135} vanish in this global formula.
\end{theorem}
\begin{proof}
The majorant in Lemma~\ref{lem:cosets163} is bounded for \(R=1\).
It also gives a bounded periodization if the norm of the real
section is used at every finite class. This latter periodization
is integrable. Indeed its integral unfolds by Tonelli to the
single seed quotient. In the class coordinates of
Lemma~\ref{lem:covol127}, the real integrals reduce to the
\(L^1\) norm of that section times the finite volume of
\(\Lambda_i\backslash B_\infty\). The compactness here follows
from the compact unipotent lattice quotient and \(H_\infty\)
compact. Proposition~\ref{prop:weight163} proves finiteness.
Thus the majorant belongs to \(L^1\cap L^\infty\), hence to
\(L^2\). Cutoffs \(V\le R\) converge by dominated convergence
on the automorphic quotient and also on the seed quotient.
Their real support is compact modulo \(B_\infty\) by the
bounded representative construction. Locally on the group the
ordinary sum is continuous. On a compact real neighborhood
right translation changes \(V\) by bounded factors, so the
tail is dominated by a common summable series at its center.
The finite variable is locally constant.

For a unit vector \(v\in U\), use an inner product linear in its
first variable, put \(c(g)=\langle\sigma(g)v,v\rangle\), and
put \(p(g)=d_\sigma\overline{c(g)}\) on \(G_\infty\), with
\(d_\sigma\) its formal degree. Lemma~\ref{lem:translate163}
on the real factors gives \(p\in L^1\), also for a finite sum
of tensor vectors \(v\). Schur orthogonality gives
\(\sigma(p)w=\langle w,v\rangle v\),
\(p=p^*=p*p\), and
\(p*\delta_t*p=c(t)p\). For the last identity insert the
translate between the two coefficients in the convolution and
integrate their common variable by the same orthogonality.
These are identities of integrable kernels. For any unitary
right representation \(R\), a vector \(z=R(p)z\) consequently
satisfies \(\langle R(t)z,z\rangle=c(t)\|z\|^2\).
If \(z\ne0\), the assignment \(\sigma(t)v\mapsto
R(t)z/\|z\|\) preserves all inner products of orbit vectors.
Its linear extension and completion give a copy of \(\sigma\)
carrying \(v\) to \(z/\|z\|\).

Right convolution by \(p\) fixes the seed made with \(v\),
by the isometry and equivariance of \(J_c\). The convolution
is pointwise defined and continuous, using boundedness of the
seed and \(p\in L^1\), so the Hilbert space identity holds
everywhere on its support. It commutes with ordinary
periodization since the sum of absolute values of the translated
integrands is bounded by the uniform majorant times
\(|p(t)|\). The resulting automorphic vector is therefore
fixed by \(R(p)\) and belongs to the required copy at \(v\),
or is zero. Linearity in pure inputs handles \(U\otimes D\).
Applying \(\Pi\) to the cutoff limits now identifies the
ordinary sum with the extension in Theorem~\ref{thm:adjoint157}.

The two periodized absolute majorants have integrable product
by Cauchy--Schwarz. Tonelli applied to that product justifies
unfolding both rational sums, with the displayed stabilizer for
each pair. The summands with \(\gamma\in P_{\rm fl}(F)\)
sum to the inner product of the two parabolic periodizations.
Theorem~\ref{thm:parabolic129} evaluates this as the seed pairing.
The same absolute bound with the two real norms gives absolute
convergence on each finite cell before contraction by the
covectors. Corollary~\ref{cor:zero135} therefore applies
directly to the full real sections.
\end{proof}

\subsection{Separation of Bessel periods at one finite place}\label{nav:VII-12}
\label{sec:depth163}

Keep the archimedean vectors, number field, and original level
\(K\) of Section~\ref{sec:absolute163}. Choose one of the already
specified places \(v\mid T\) of \(F\), with maximal ideal
\(\mathfrak p_v\) and residue cardinality \(Q\). In a faithful
matrix action on the underlying \(F_v\)-space of \(V\), choose
a lattice stabilized by \(K_v\). Such a lattice is the sum of
the translates of any one lattice under \(K_v\). Compactness
bounds their denominators, and the sum contains the original
lattice. For \(j\ge1\) let \(J_j\) be the kernel of the action
of \(K_v\) modulo \(\mathfrak p_v^j\) on the chosen lattice.
Put
\[
 J_0=K_v,\qquad L_j=(B(F_v)\cap K_v)J_j,\qquad
 \alpha_j=[K_v:L_j]^{-1}.
\]
Normality of \(J_j\) in \(K_v\) makes \(L_j\) an open subgroup.
For an original seed \(\Phi\) put
\[
  \Phi_j(t)=\Phi(t)\,\mathbf1_{B(F_v)J_j}(t_v),\qquad
  S_j\xi=P\Phi_j\quad\text{if }\Phi=J_U\xi,\quad
  \xi\in U\otimes D .
\]
We identify \(J_U\xi\) here with its scalar adelic seed.
The finite class covectors in the domain carry their original
normalization. On the new support the seed transforms on the
right under \(L_j\) by the restriction of the old level
character. Indeed \(B(F_v)J_j\) is stable on the right by
\(B(F_v)\cap K_v\), since this subgroup normalizes \(J_j\).
The real projection argument of Theorem~\ref{thm:ordinary163}
applies without change to these seeds. Its majorant works at
every \(j\) since the additional indicator has modulus at most one.

\begin{lemma}[An arithmetic lower bound for overlapping supports]
\label{lem:deep163}
Let \(h_i\) be any of the original finite class representatives
in \(H_f\). Suppose that \(\gamma\in G(F)\setminus B(F)\),
\(\gamma_f h_i\in B_fK\), and
\(\gamma_v h_{i,v}\in B(F_v)J_j\). For every \(g\in G_\infty\),
\[
                 Q^j\le C_x V(g)\,V(\gamma_\infty g).
\]
The constant is independent of \(j,\gamma,g\) and of the class
index \(i\) at these fixed data.
\end{lemma}
\begin{proof}
Apply \eqref{eq:lifts163} to \(y=(g,h_i)\) and to \(\gamma y\).
Let \(\delta_0,\delta_1\in B(F)\) be the multipliers, and set
\(a=\delta_1\gamma\delta_0^{-1}\). It lies outside \(B(F)\).
At the finite places \(a_f\in C_fC_f^{-1}\), and at each
real place
\[
        w(a_\tau)\le C_\tau^2
                     \nu_\tau(\gamma_\tau g_\tau)\nu_\tau(g_\tau).
\]
At \(v\) write \(c_0=\delta_{0,v}h_{i,v}\),
\(c_1=\delta_{1,v}\gamma_v h_{i,v}\). Both lie in the fixed
compact set \(C_v\). We have \(c_0\in B(F_v)\) and
\(c_1=b k\) with \(b\in B(F_v),\, k\in J_j\).
As \(k\in K_v\), the matrices \(b=c_1 k^{-1}\) range in a
fixed compact set. In the chosen lattice coordinates
\[
              a_v-bc_0^{-1}=b(k-1)c_0^{-1}
\]
has entries in \(\mathfrak p_v^{\,j-c}\) for a fixed integer \(c\).

We record the equations cutting out \(B\) inside \(G\).
In the Witt frame of Section~\ref{sec:char94}, with
\(A_i=\sum_{h\le i}Eu_h\) and \(A_0=0\), they are
\[
 (a-1)u_i\in A_{i-1}\quad(1\le i\le s),\qquad
                      (a-1)e\in A_s .
\]
The first conditions force an isometry to stabilize each flag
space and to act as the identity on its successive lines.
The dual lines then have identity action as well. The last
condition fixes \(e\) in \(A_s^\perp/A_s=W_0\perp Ee\).
The induced isometry of this middle quotient lies in \(U(W_0)\).
Removing that Levi factor gives precisely an element of \(N\).
The reverse inclusion follows from the same description.
Expanding the conditions in an \(F\)-basis of \(E\) gives
finitely many affine linear equations in \(F\)-matrix
coordinates.

One of their values \(z\in F\) at \(a\) is nonzero. All possible
values belong to a fixed fractional ideal by the finite bound on
\(a_f\). The congruence above gives
\(\operatorname{ord}_v(z)\ge j-O_x(1)\), since every equation
vanishes on \(bc_0^{-1}\) and its linear coefficients are fixed.
Choose a positive rational integer clearing the fractional
ideal. The absolute integral norm of its product with \(z\)
is divisible by \(Q^{\max(0,j-O_x(1))}\). At each infinite
place \(|\tau z|\le C'_\tau w(a_\tau)\), since the equations
are affine linear and \(w(a_\tau)\ge1\). The two bounds for
the same field norm give the assertion.
\end{proof}

\begin{theorem}[The projected Gram matrix at increasing depth]
\label{thm:deep163}
Use the same Haar measure on the automorphic quotient at every
level. On the entire original space \(U\otimes D\), including
the subspace of finite columns of zero sum on each scalar orbit,
\[
             \|\alpha_j^{-1}S_j^*S_j-1\|_{\rm op}
                         \le C_x Q^{-j\eta},\qquad
                         \eta=L-4fr^2>0 .
\]
Consequently, at each fixed field and choice of parameters,
\(S_j\) is injective for all sufficiently large \(j\).
\end{theorem}
\begin{proof}
Take two original seeds \(\Phi,\Psi\). After unfolding the
first rational sum in \(\langle P\Phi_j,P\Psi_j\rangle\),
reduce its finite variable to the original class cells \(h_iK\)
as in Lemma~\ref{lem:covol127}. On \(h_iK\) the first indicator
restricts the local right coordinate to \(L_j\). In fact
\(B(F_v)J_j\cap K_v=(B(F_v)\cap K_v)J_j\).
The two factors have the same unitary right character on
\(K^{(v)}L_j\), so it cancels in their product. The original
stabilizer \(\Lambda_i\) acts on this restricted finite cell
since its conjugate at \(v\) lies in \(B(F_v)\cap K_v\).
A real fundamental region for \(\Lambda_i\) and the indicated
finite cell therefore give, with the ordinary quotient measures,
\[
 \alpha_j^{-1}\langle P\Phi_j,P\Psi_j\rangle
 =\operatorname{vol}_f(K)\sum_i
    \int_{\Lambda_i\backslash G_\infty}
       \Phi(g,h_i)\overline{P\Psi_j(g,h_i)}\,dg .
\]
The bounded integrable majorants of
Theorem~\ref{thm:ordinary163} justify the unfolding. The
identity coset in \(P\Psi_j(g,h_i)\) contributes \(\Psi(g,h_i)\)
because \(h_{i,v}\in B(F_v)\). Hence its integral is exactly
\(\langle\Phi,\Psi\rangle_{\rm seed}\).

Write \(t=V(g)\). Each remaining nonzero term of \(P\Psi_j(g,h_i)\)
has height at least \(Q^j/(C_x t)\) by Lemma~\ref{lem:deep163}.
The tail estimate of Lemma~\ref{lem:cosets163} now gives
\[
 |P\Psi_j(g,h_i)-\Psi(g,h_i)|
       \le C_{\Psi,x}\,Q^{-j\eta}t^\eta .
\]
When the lower cutoff is less than one, the bound at height
one gives the same inequality after enlarging the constant.
The integral of \(|\Phi(g,h_i)|V(g)^\eta\) on the real class
quotient is finite. Disintegrate over
\(\Lambda_i\backslash B_\infty\), which is compact, and apply
Proposition~\ref{prop:weight163} with \(\eta\le L\).
Combining these estimates proves the asserted decay for each
matrix entry of the difference from the original seed pairing.
The seed map \(J_U\) is isometric. Summing the squared bounds
over an orthonormal basis of its finite-dimensional domain gives
the operator bound with a fixed constant. When the bound is
less than one, every nonzero input has positive squared image
norm. This proves injectivity for all those \(j\).
\end{proof}

\begin{proposition}[Averaging back to the prescribed level]
\label{prop:average163}
Let \(e_K\) be right averaging at \(K_v\) against the conjugate
of the old level character, with probability measure. On the
same inputs,
\[
         e_K S_j=\alpha_j S_0,\qquad
         \alpha_j^{-1} S_j^*(1-e_K)S_j
              =\alpha_j^{-1}S_j^*S_j-\alpha_j S_0^*S_0 .
\]
\end{proposition}
\begin{proof}
For \(t_v\in B(F_v)K_v\) write \(t_v=bk_0\).
The indicator in \(\Phi_j(tk)\), averaged over \(k\in K_v\),
has integral \(\alpha_j\) because \(k_0 k\) must lie in
\(L_j\). The original seed contributes its level character on
the right, canceled by the conjugate character in the integral.
Off \(B(F_v)K_v\) every right translate is zero. Thus averaging
the seed gives \(\alpha_j\Phi\). Absolute summability allows
the compact average to commute with periodization. This proves
the first equality. Since \(e_K=e_K^*=e_K^2\), taking the
Gram matrix of \(e_K S_j\) gives the second equality.
\end{proof}

The original period matrix is \(S_0^*S_0\). The estimate
in Theorem~\ref{thm:deep163} gives a small error for large \(j\)
at fixed \(x\). Its constant includes the real coefficient
norms and the covolumes of the chosen fields. The required bound
for the matrix at \(j=0\) on the original intervals, and its
pairing with the designated integral theta rows, remain open.

\subsection{Vanishing on the proposed lowest compact types}\label{nav:VI-10}
\label{sec:branch106}

Consider the Bessel integral \eqref{eq:period94} on the fields of
Section~\ref{sec:char94}. At a real place write \(P=k+s,\ Q=s+1\)
and \(r=P+Q\), where \(s=m/2-1\ge1\) and \(k\ge2\) is even.
The group \(H_\tau=U(k)\) acts on the positive subspace \(W_0\).
After ordering the positive coordinates with \(W_0\) first, it
is the upper left block in \(U(P)\) and acts trivially on the
remaining coordinates of \(K_\tau=U(P)\times U(Q)\).
We use the right actions on both spaces of automorphic forms and
the conjugation of the column in \eqref{eq:period94}.

\begin{lemma}[A necessary compact branching condition]\label{lem:branch106}
Suppose that at one real place the row belongs to a compact type
of highest weight \((A;B)\), with \(A_1\ge\cdots\ge A_P\), and
the column belongs to the type of highest weight
\(\nu_1\ge\cdots\ge\nu_k\). A nonzero value of the integral
requires
\begin{equation}
             A_i\ge\nu_i\ge A_{i+s}\qquad(1\le i\le k).
             \label{eq:branch106}
\end{equation}
These inequalities are necessary and sufficient for the occurrence
of the column in the restriction of this compact row type.
\end{lemma}
\begin{proof}
For \(y\in H_\tau\), simultaneous right translation of the two
inputs changes \(h_0\) to \(h_0y\) in \eqref{eq:period94}.
Right invariance of the measure gives
\({\cal P}_\beta(R(y)F,R(y)\phi)={\cal P}_\beta(F,\phi)\).
On compact types the resulting sesquilinear form pairs only
equivalent irreducible summands, by Schur's lemma. Multiplicity
spaces and the \(U(Q)\)-factor carry the trivial \(H_\tau\)-action.

For the restriction \(U(n)\) to \(U(n-1)\), a weight \(b\)
occurs in a row \(c\) exactly when \(c_i\ge b_i\ge c_{i+1}\).
One proof uses the tableau character after a common determinant
twist makes \(c\) a partition. The boxes containing the letter
\(n\) form a horizontal strip. Removing them leaves precisely
the shapes \(b\) satisfying the inequalities. Conversely all
tableaux of each such shape \(b\) extend by filling the strip
with \(n\). Setting the last variable equal to one gives the
restriction character. Undoing the twist subtracts the same
integer from all entries on both sides. This is also the
Gelfand--Tsetlin branching rule \cite[Theorem 2.1]{XVII-XII-IV-molev106}.

Iterating the inequalities gives \eqref{eq:branch106}. Conversely,
for \(s\ge1\) fill the next row of length \(k+s-1\) by
\(\gamma_i=\max(A_{i+1},\nu_i)\) for \(i\le k\) and
\(\gamma_i=A_{i+1}\) for \(i>k\), taking only indices in the row.
It is decreasing in the weak sense and \(A_i\ge\gamma_i\ge A_{i+1}\).
For \(1\le i\le k\), one has \(\gamma_i\ge\nu_i\) and
\(\gamma_{i+s-1}\le\nu_i\) because \(A_{i+s}\le\nu_i\) and any
entry of \(\nu\) used at that index is at most \(\nu_i\).
For \(s=1\) this row equals \(\nu\). Induction supplies an entire
branching pattern.
\end{proof}

Let \(\lambda_1>\cdots>\lambda_r\) be an ordered Harish-Chandra
parameter. A chamber of the discrete-series packet assigns \(P\)
entries to the positive subspace, in positions \(j_1<\cdots<j_P\).
In the lowest compact type the shift of the parameter is
\(\rho_{\rm n}-\rho_{\rm c}\), the difference of half the sums of
the positive noncompact and compact roots. At position \(j_i\) the
half-sum of all positive roots is \((r+1)/2-j_i\), and the half-sum of
the compact roots is \((P+1)/2-i\). Subtracting twice the latter
gives the first equality below. Subtracting the compact column's
own half-sum gives the second.
\begin{equation}
 A_i=\lambda_{j_i}-k/2+2i-j_i,\qquad
 \nu_i=u_i-(k+1)/2+i.                            \label{eq:mu106}
\end{equation}
Here \(u\) denotes the column Harish-Chandra parameter. These are
the conventions of \srcloc{Formalism/PeriodIntegrals/UpqRep}{231}.
The entries of \(A\) are dominant. In fact
\(\lambda_{j_i}-\lambda_{j_{i+1}}\ge j_{i+1}-j_i\) by regular
integrality, so \(A_i-A_{i+1}\ge 2(j_{i+1}-j_i)-2\ge0\).

\begin{theorem}[Vanishing for the explicit terminal weights]\label{thm:zero106}
Use the lists printed just before
Theorem \ref{thm:target90}.
The global integral vanishes for every lowest compact row type in
the packet and the specified column. This holds for every terminal
prime, every \(\beta\in E^\times\), and all smooth finite components.
It already holds at a place with the low column weight.
\end{theorem}
\begin{proof}
Write \(v_i=2(k-1-i)\) for \(1\le i<k\). The row is obtained by
adjoining \(-j(t_{\rm r}+4)\) for \(1\le j\le m\). The column
has \(u_i=v_i+1/2\) for \(i<k\), \(u_k=-1/2\), with an additional
\(a_\tau\) in \(u_1\), equal to \(a\) at a high place and zero
otherwise. Since \(j_1\ge1,\ j_k\ge k\), monotonicity of \(\lambda\)
in \eqref{eq:mu106} gives
\[
 \begin{aligned}
 A_1&\le v_1-k/2+1,            &\nu_1&=v_1-k/2+1+a_\tau,\\
 A_k&\le -(t_{\rm r}+4)+k/2,  &\nu_k&=k/2-1 .
 \end{aligned}
\]
The last inequality required by \(A_k\ge\nu_k\) thus fails by
at least \(t_{\rm r}+3>0\). Lemma~\ref{lem:branch106} applies
in every chamber. Finite translations and finite averages commute
with \(H_\tau\), so the same vanishing holds after either operation.
\end{proof}

\begin{proposition}[Order of a derivative that could contribute]\label{prop:jets106}
Let \(f\) be the number of real places and \(h\) the number with
\(a_\tau=a\). Starting with the lowest compact types of
Theorem~\ref{thm:zero106}, a differential operator of total order
less than
\begin{equation}
                    ha+f(t_{\rm r}+3)                 \label{eq:jets106}
\end{equation}
has zero period against the specified column. Furthermore, operators
generated solely by the terminal one-parameter subgroups \(n(t)\)
of Lemma~\ref{lem:root94} have zero period in every order.
\end{proposition}
\begin{proof}
At one place the compact span of derivatives of order at most \(d\)
is a quotient of a subspace of the sum of
\(\mathfrak{gl}_r(\CC)^{\otimes j}\otimes V_{(A;B)}\), \(0\le j\le d\).
Indeed composition of derived actions intertwines the compact
action on the input with conjugation on each Lie algebra factor.
As a \(U(P)\)-module the latter factor consists of copies of the
standard module \(V\), its dual, \(V\otimes V^*\), and trivial
modules, by the four blocks in an \(r\)-by-\(r\) matrix.

Tensoring a dominant weight \(c\) with \(V\) increases precisely
one coordinate by one in each constituent. Tensoring with \(V^*\)
decreases precisely one. For instance, in the alternant numerator
of the Weyl character formula multiply by \(\sum_i z_i\) or by
\(\sum_i z_i^{-1}\). Distributing yields the sum of alternants with
one exponent moved by one. A repeated exponent gives zero and the
surviving sequences are still strictly decreasing. The denominator
is unchanged. Consequently each Lie algebra factor adds at most
one to the total of all upward coordinate moves, also in the
\(V\otimes V^*\) block. Every resulting highest weight \(A'\)
therefore satisfies \(\sum_i\max(A'_i-A_i,0)\le d\).
For the column to occur we need \(A'_i\ge\nu_i\) for \(i\le k\).
The distinct indices \(1,k\) alone require at least
\(a_\tau+t_{\rm r}+3\) upward moves. Expanding a global differential
operator in words supported at individual places and summing these
local lower bounds proves \eqref{eq:jets106}.

Finally each \(n(t)\) acts only on the three-space spanned by
\(u_s,e,v_s\) and fixes \(W_0\) pointwise. Its derivatives commute
with \(H_\tau\), so applying any of them preserves the list of
\(H_\tau\)-types available in the row. Theorem~\ref{thm:zero106}
again gives zero.
\end{proof}

The exponential bound for the analytic integral in Theorem \ref{thm:raw97}, holds in particular for a zero matrix. Its
Corollary~\ref{cor:alg97} requires a nonzero algebraic coefficient. The
explicit weights above cannot supply that coefficient on these
lowest compact types. The scalar in the identification in
\srcloc{Formalism/SharpRestriction/GlobalTheoremInputs}{65}
cannot turn the zero pairing on the specified spaces into a
nonzero one.

\subsection{Global Bessel coefficients at the prescribed finite levels}\label{nav:VI-11}
\label{sec:global110}

Let \(E/F_+\), \(G\), \(H=U(W_0)\), \(N\) and \(\chi_\beta\)
be as in Section~\ref{sec:char94}. Put \(B=NH\) and write
\(p_H:B\to H\) for the quotient with kernel \(N\). Use the
row and column levels \(K_\sigma,K_\pi\) of
Proposition \ref{prop:index90}.
At the finite places away from \(g_1,g_2\) use the respective
standard lattice stabilizers. In particular the original
stabilizers are used at the terminal prime \(T\). We take
probability measure on \([H]\) and \([N]\) as in
\eqref{eq:period94}. All assertions below use the shortened
period field and the intervals of Proposition \ref{prop:short107}.

\begin{proposition}[Column multiplicities and their inner products]\label{prop:column110}
Let \(M_H=\mathfrak m_k I_\pi\), where \(I_\pi\) is the index
of \(K_\pi\) in the full finite stabilizer and \(\mathfrak m_k\)
is the mass of the standard definite genus. Then \(M_H\) is the
number of finite level classes. If \(\rho\) is the irreducible
unitary representation of \(H_\infty\) of a specified column
weight, put \(D_\rho=\dim\rho=\prod_\tau W_k(u_\tau)\).
Its multiplicity in \(L^2([H])^{K_\pi}\) is \(M_HD_\rho\),
the quantity \(d_{u,*}\) of Theorem \ref{thm:target90}.

Choose representatives \(y_1,\ldots,y_{M_H}\) in
\(H(\mathbb A_{F_+,f})\). On the class \(y_i\), and zero outside
that class, set \(\phi_{i,a,b}(h_\infty,y_i)
=\langle e_a,\rho(h_\infty)e_b\rangle\) for an orthonormal
basis of the representation, using an inner product linear in
the second entry. These \(M_HD_\rho^2\) functions are orthogonal
and each has squared norm \(1/(M_HD_\rho)\).
\end{proposition}
\begin{proof}
At the full level a class corresponds to a lattice \(L\) in
the genus, with finite stabilizer \(\Aut(L)\). Finiteness
follows by taking its integral action on the underlying lattice
with its positive definite trace form. Refinements of this class
are the orbits of \(\Aut(L)\) on the \(I_\pi\) cosets of the
smaller level. All the refined stabilizers are trivial, as in
the proof of Corollary~\ref{cor:fibres89}. Indeed each stabilizer element has
unipotent reduction on the underlying \(2fk\)-dimensional
rational lattice at the split prime \(g_1\).
The principal congruence kernel in \(\GL_{2fk}(\ZZ_{g_1})\) is
torsion free by the binomial formula at the odd prime \(g_1\).
Thus an element of finite order here has \(g_1\)-power order.
An element of order \(g_1\) would require a rational cyclotomic
factor of degree \(g_1-1>2fk\) in its characteristic polynomial.
Hence the action on refinements is free. It contributes
\(I_\pi/|\Aut(L)|\) classes. Summing proves the formula for
\(M_H\).

Each refined class in \([H]/K_\pi\) is one copy of \(H_\infty\).
They have equal measure because they have the same finite level
volume and trivial stabilizers. Each copy therefore has mass
\(1/M_H\). The right regular representation on one such copy
contains \(\rho\) with multiplicity \(D_\rho\) by the compact
group decomposition into matrix coefficients. Schur orthogonality
gives their stated inner products. In particular the multiplicity
space can be described by one covector in \(\rho^\vee\) for each
class, by evaluation at \((1,y_i)\).
\end{proof}

\begin{lemma}[The period domain at finite level]\label{lem:orbit110}
Let \(K_B=B(\mathbb A_{F_+,f})\cap K_\sigma\).
The natural map
\[
   Y=B(F_+)\backslash B(\mathbb A_{F_+})/K_B
      \longrightarrow
   X=G(F_+)\backslash G(\mathbb A_{F_+})/K_\sigma
\]
is a smooth embedding with compact image.
Moreover \(p_H(K_B)=K_\pi\), with \(K_\pi\) embedded in
\(K_\sigma\), and \(\chi_\beta\) extended trivially on \(H\)
is trivial on \(K_B\) for every \(\beta\in\OO_E\).
\end{lemma}
\begin{proof}
Suppose representatives \(b_1,b_2\) have the same image. Then
\(b_2=\gamma b_1 k_f\) with \(\gamma\in G(F_+)\) and
\(k_f\in K_\sigma\). The real component of \(\gamma\) belongs
to \(B_\infty\). Since \(B\) is a closed algebraic subgroup over
\(F_+\), its polynomial equations evaluated at the rational
element \(\gamma\) already vanish over \(F_+\). Thus
\(\gamma\in B(F_+)\), and \(k_f\in K_B\). This proves injectivity.
Both quotients are locally discrete group quotients of their
real Lie groups by left translation. To see discreteness, at
a fixed finite level the allowed rational matrices lie in a
fixed fractional lattice in all their real embeddings.
The local map is the immersion of \(B_\infty\) in \(G_\infty\).
The quotient for \(B\) is compact. Indeed \([N]\) is compact
by its additive filtration, and \([H]\) is compact by
definiteness. Compact sets of representatives for these two
quotients give a compact set of representatives for \([B]\),
by first moving the \(H\)-component rationally and then the
\(N\)-component. An injective immersion of a compact manifold
into a Hausdorff manifold is an embedding.

The action on the middle quotient projects \(b=nh\in K_B\) to
the \(W_0\)-matrix of \(h\). It preserves the standard lattice
there. Indeed projection to the \(W_0\)-coordinates and inclusion
of those coordinates both preserve the original lattice, and the
same argument applies to \(b^{-1}\). At \(g_1,g_2\) this block
is upper unitriangular modulo the prime, so \(h\in K_\pi\).
Conversely the block embedding of \(K_\pi\) has identity on
the other graded spaces. It satisfies all row level conditions,
including the zero entries in the column of \(e\). Thus
\(K_B=(K_\sigma\cap N(\mathbb A_{F_+,f}))\rtimes K_\pi\).
For \(n\) in this intersection, \(n u_{i+1}\) and \(ne\) are
integral at each finite place. Their relevant \(u_i\)-coefficients
are integral even at two, since they are the coefficients of
the original positive basis vectors in the respective sums of
the \(u_j\). Traces of integral elements to \(\QQ_p\) are
integral. The standard finite additive character is therefore
trivial on every term defining \(\chi_\beta(n)\) for
\(\beta\in\OO_E\). Invariance under \(H\)-conjugation extends
this character to \(B\) as asserted.
\end{proof}

For a row function \(F\) write its Bessel coefficient as
\[
       {\cal F}_\beta F(h)=
             \int_{[N]}F(nh)\overline{\chi_\beta(n)}\,dn.
\]
It is a function on \([H]/K_\pi\), by the character invariance
and the inclusion of finite levels in Lemma~\ref{lem:orbit110}.
Let \(\mathcal C\) be a sum of the full isotypic spaces of any
finite set of column types in \(L^2([H])^{K_\pi}\). It is finite
dimensional and smooth, by Proposition~\ref{prop:column110}.

\begin{theorem}[Simultaneous prescribed global coefficients]\label{thm:section110}
For any finite set \(J\) of distinct nonzero elements of \(\OO_E\)
there is an \(H_\infty\)-equivariant linear map
\[
         S:\bigoplus_{\beta\in J}\mathcal C
            \longrightarrow C_c^\infty(X)
     \quad\text{with}\quad
         {\cal F}_\alpha S((\phi_\beta)_\beta)=\phi_\alpha
                     \quad(\alpha\in J).
\]
Its image can be chosen finite under \(K_\infty\). For every
\(\eta>0\) it can also be chosen with \(L^2\) operator norm
at most \(\eta\), using the direct sum norm on the input
and the fixed quotient measure on \(X\).
\end{theorem}
\begin{proof}
First prescribe the restriction of the desired function to \(Y\)
by the formula
\[
             (nh)\longmapsto
                 \sum_{\beta\in J}\chi_\beta(n)\phi_\beta(h).
\]
It descends on the left by rational invariance of each factor
and on the right by Lemma~\ref{lem:orbit110}. Distinct
\(\chi_\alpha,\chi_\beta\) are orthogonal on \([N]\). Indeed
their ratio is nontrivial on the last one-parameter subgroup
of Lemma~\ref{lem:root94}, as is seen by varying its real
argument at one place. Right invariance of the probability
measure forces the integral of this ratio to vanish.
Consequently the prescribed restriction already has all the
stated coefficients.

Choose an \(H_\infty\)-invariant Riemannian metric on \(X\) by
averaging a metric over this compact right action. In a tubular
neighborhood of the compact embedded \(Y\), pull its prescribed
function back by normal projection and multiply by a smooth
function of squared normal distance which is one on \(Y\) and
zero near the boundary of the neighborhood. Extension by zero
gives a compactly supported smooth function. The neighborhood,
projection and cutoff can be chosen invariant under \(H_\infty\);
the construction is linear and equivariant. Denote the resulting
map by \(S_0\). Its norm can be made arbitrarily small by
shrinking the tube. In detail, the prescribed functions on \(Y\)
are uniformly bounded on the unit ball of their finite dimensional
domain, and \(Y\) has positive codimension in \(X\), since
\(NH\) is a proper Lie subgroup of smaller dimension than \(G\).
Smooth quotient measure gives \(Y\) measure zero. The volumes
of tubes of decreasing radius tend to zero by compactness.

To obtain compact-finite functions, apply to \(S_0\) finite
sums \(P_j\) of the right \(K_\infty\)-isotypic orthogonal
projections, exhausting its irreducible types. These commute
with \(H_\infty\). On smooth functions they converge uniformly
on compact sets to the identity. For example, on the compact
group apply the Peter--Weyl decomposition to
\(k\mapsto F(xk)\). Partial sums converge in every compact
Sobolev norm by the compact Laplacian. With cutoff eigenvalue
\(j\) the norm of the omitted tail in order \(t\) is at most
\((1+j)^{-1}\) times the norm in order \(t+2\). The latter norms
are bounded uniformly for \(x\) on a compact set because all
its products with \(K_\infty\) form a compact set. Sobolev
embedding then gives uniform convergence there.
The supports after projection lie in the compact \(K_\infty\)-orbit
of the original support. Thus on the finite dimensional source
the joint coefficient of \(P_jS_0\) converges in operator norm
to the identity. The coefficient maps and \(P_jS_0\) are
\(H_\infty\)-equivariant and the finite column level is \(K_\pi\).
Since \(\mathcal C\) contains the entire isotypic spaces of the
input types, the range lies in \(\bigoplus_{\beta\in J}\mathcal C\).
For some \(j\) this endomorphism
\(A_j\) satisfies \(\|A_j-1\|<1/2\). Taking
\(S=P_jS_0 A_j^{-1}\) gives exactly the required identity.
Since \(\|P_j\|\le1\) on \(L^2(X)\) and \(\|A_j^{-1}\|\le2\),
starting with \(\|S_0\|<\eta/2\) proves the norm assertion.
\end{proof}

\begin{corollary}[The two terminal frequencies on the same quotient]\label{cor:gram110}
Take \(J=\{1,\varpi\}\) for the chosen terminal generator
\(\varpi\bar\varpi=T\), and include the prescribed column
types in \(\mathcal C\). For any orthonormal basis
\((\phi_i)\) of \(\mathcal C\), there are smooth compact-finite
row functions \(F_{\beta,i}\) at the given levels with
\[
           {\cal P}_\alpha(F_{\beta,i},\phi_j)
                 =\mathbf1_{\alpha=\beta}\mathbf1_{i=j}
                      \qquad(\alpha,\beta\in J).
\]
\end{corollary}
\begin{proof}
Apply Theorem~\ref{thm:section110} to the vector with \(\phi_i\)
in component \(\beta\) and all other components zero, then
pair the resulting coefficient with \(\phi_j\) on \([H]\).
\end{proof}

The functions in this construction have a finite set of compact
types that can depend on the chosen tube. Their construction
imposes neither a particular infinitesimal character nor the
condition of lying in an iterated theta lift. In fact the
lowest-type component for the coded source row gives zero in
every entry by Corollary \ref{cor:code109}. To transfer the displayed global matrix to the
required arithmetic row one must control its spectral projection
and its values on the integral theta lattice.

\subsection{Fourier inversion between two positions of the Bessel period}\label{sec:exchange336}
Fix a word \(u\) and its field from Theorem~\ref{thm:common111}.
Write \(A\) for the definite space of the column, \(a=\dim_E A\),
\(H=U(A)\), and \(G\) for the row group. Write \(s\) for the length
of its isotropic flag. Fix the marked vector \(z=z(u)\) and the
simple-root character \(\chi\) of \(N\). For a column \(\phi\) and
a smooth automorphic row \(F\), put
\[
 L(F)=c_H\int_{[H]}\overline{\phi(h)}
        \int_{[N]} F(nh)\overline{\chi(n)}\,dn\,dh ,
 \qquad (R(g)F)(t)=F(tg).
\]
Here \(c_H\) is the scalar in the given pairing, \([N]\) has mass
one, and \(H\) acts identically on the orthogonal complement of
\(A\). The two quotients have compact sets of representatives by
Lemma~\ref{lem:orbit110}. Thus these integrals exist for each
fixed translate. For \(n_0\in N(\mathbb A)\) one has
\(L(R(n_0)F)=\chi(n_0)L(F)\). Indeed \(hn_0h^{-1}\) lies in \(N\)
and has the same simple-root coordinates as \(n_0\). Right
translation by this element in the inner quotient preserves its
measure and contributes \(\chi(n_0)\).

At position one the split factors are those over the chosen
terminal generator \(\pi_T(u)\) and the conjugate marked generators
\(\bar\pi_{g_j}\). Denote these chosen halves by \(S_u\).
At \(v\in S_u\) let \(k=F_{+,v}\), \(O\) its valuation ring,
\(\pi\) a uniformizer and \(q=|O/\pi O|\). Identifying the unitary
group with the general linear group on this half determines the
action on the other half by inverse adjoint. Use coordinates
\(e_1,\ldots,e_s,A,z,f_1,\ldots,f_s\), with \(e_i,f_i\) dual. Set
\(d_i=e_{i+1}\) for \(i<s\) and \(d_s=z\). We use \(d_i^*\) for
the corresponding coordinate functional. For \(x\in A_k\),
\(y\in A_k^*\) and \(i\le j\), put
\[
 X_i(x)=1+x d_i^*,\qquad Y_i(y)=1+e_i y,\qquad
 Z_{ij}(t)=1+e_i t d_j^* .
\]
The covectors on \(A_k\) are extended by zero on the other summands.
The exchange group at index \(i\) fixes all basis columns except
\(d_i\), and its displacement on \(d_i\) can be any vector of
\(A_k\). It is therefore exactly \(X_i(A_k)\). At position one
these are the indices \(1\le i\le s\). At the symmetric position
there are no exchanges. All the \(X_i\) commute.

\begin{lemma}[The character on the split flag]\label{lem:exchangeroot336}
The elements \(Y_i(y)\) and \(Z_{ij}(t)\) belong to \(N(k)\).
Their \(\chi\)-values are one except for
\(\chi(Z_{ii}(t))=\psi_i(t)\), a nontrivial additive character of
\(k\). With the character
\(\psi_v(\operatorname{Tr}_{E_v/k}(\theta\sum_l
 h(n d_l,f_l)))\), linearity of \(h\) in its first entry gives
\(\psi_i(t)=\psi_v(\theta_v^+t)\), where \(+\) denotes the chosen
half. For \(P=\prod_{j\ge i}X_j(x_j)\) one has
\[
 P Y_i(y)P^{-1}=Y_i(y)\prod_{j\ge i} Z_{ij}(-y(x_j)).
\]
\end{lemma}
\begin{proof}
Each indicated element sends an isotropic flag vector to itself
plus earlier flag vectors and acts trivially on the middle quotient.
On the opposite half the added map in the inverse adjoint has
input the coordinate dual to \(e_i\), so it can move only
\(f_i\). It fixes all the columns used to compute the character.
On the chosen half only \(Z_{ii}(t)\) contributes to the
sum of adjacent coordinates, and its contribution is \(t\). The
split trace adds the two components and \(\theta_v^+\ne0\).
Finally \(P e_i=e_i\), \(P^{-1}=1-\sum_{j\ge i}x_jd_j^*\), and
\(yP^{-1}=y-\sum_{j\ge i}y(x_j)d_j^*\). Multiplying the
rank-one matrices gives the asserted identity.
\end{proof}

Fix the original additive Haar measure \(\mu_i\) on each \(X_i\).
Write \(L_{>i}\) for integration of \(L(R(\prod_{j>i}X_j(x_j))F)\)
over the indicated variables at this place and \(L_{\ge i}\) for
the result after also integrating \(x_i\). Exchanges at other
places may already be present. These definitions include the
functional \(L_{>s}=L\) when there are no other exchanges.

\begin{theorem}[Finite Fourier inversion for a specified row]\label{thm:exchangefourier336}
All these exchange integrals converge absolutely. Suppose that
\(F\) is fixed by \(X_i(U)\) and \(Y_i(V)\) for full \(O\)-lattices
\(U\subset A_k,\ V\subset A_k^*\). Write \(\perp\) for annihilators
under \((x,y)\mapsto\psi_i(y(x))\). Choose \(U,V\) small enough
that \(U\subset V^\perp\). Put \(D=V^\perp/U\), \(D'=U^\perp/V\)
and \(b(x)=L_{>i}(R(X_i(x))F)\). For every \(y\in U^\perp\),
\[
 \begin{split}
 L_{\ge i}(R(Y_i(y))F)
   &=\mu_i(U)\sum_{x\in D} b(x)\,\overline{\psi_i(y(x))},\\
 L_{>i}(F)&=\frac{1}{|D|\mu_i(U)}
                  \sum_{y\in D'}L_{\ge i}(R(Y_i(y))F).
 \end{split}
\]
The first formula at \(y=0\) evaluates the exchanged functional
on the original row. If \(U=V^\perp\), its value is
\(\mu_i(U)L_{>i}(F)\).
\end{theorem}
\begin{proof}
First consider the joint integrand before any exchange. Fix a
basis of \(A_k\). Smoothness gives an integer \(m\ge1\) such
that the row is fixed by the principal congruence subgroup of
depth \(m\) in these coordinates. The character \(\psi_i\) is
the same for every \(i\) at this place. Let \(c\) be the integer
for which the annihilator of \(O\) under this character is \(\pi^cO\).
For a tail \(P\) beginning at \(i\), the lemma and
\(N\)-equivariance show that \(L(R(P)F)\) vanishes unless
\(x_i\in\pi^{c-m}O^a\). Indeed it must be unchanged by
right multiplication by every \(Y_i(y)\) with \(y\in\pi^mO^{a*}\).
The change on moving this matrix to the left of \(P\) is
\(\overline{\psi_i(y(x_i))}\).

To apply the same argument after \(x_1,\ldots,x_{i-1}\) have
been bounded, move their commuting product onto \(F\). Its open
stabilizer implies that the translated row takes only finitely
many values on this compact set. A single deeper
congruence subgroup fixes them all. The tail argument then
bounds \(x_i\) uniformly on this set. Starting at the first
index proves compact support in all variables. Finite products
of places satisfy the same argument since their groups commute.
The joint integrand is locally constant. It is therefore
absolutely integrable, and Fubini's theorem applies.

Conjugating \(Y_i(y)\) through the factors \(X_j,\ j>i\), produces
only \(Z_{ij}\) with \(j>i\), on which \(\chi=1\).
The subgroup \(Z_{ii}\) commutes with these factors. Thus
\(L_{>i}\) is \(Y_i\)-invariant and transforms by \(\psi_i\)
under \(Z_{ii}\). Applying the lemma with the additional factor
\(X_i(x)\) gives the multiplier \(\overline{\psi_i(y(x))}\)
in the integrand for \(R(Y_i(y))F\). Taking \(y\in V\) shows
that \(b\) vanishes off \(V^\perp\). It is constant on cosets of
\(U\) by right invariance. Integration on these cosets proves
the first formula. The pairing identifies \(D'\) with the full
character group of \(D\). Indeed in a basis diagonalizing the
inclusion \(U\subset V^\perp\), each pair of cyclic \(O\)-module
quotients has the same order and has zero annihilator on either
side. Summing the characters over \(D'\) therefore selects
\(x=0\) with multiplicity \(|D|\), which proves the second formula.
\end{proof}

In particular both directions are integral away from the residue
characteristic after removing the displayed Haar scalar. More
precisely, work in a coefficient ring containing the finite
character values. After division by \(\mu_i(U)\), the forward
matrix has entries in this ring and its inverse is the conjugate
transpose divided by \(|D|\), a power of the residue characteristic.
This assertion concerns the full vector of Fourier values. For
one value at \(y=0\) the formula retains the individual signs.

\begin{proposition}[The comparison on the two prescribed translates]\label{prop:exchangegram336}
Let \(\Phi\) and \(W\) be the given finite-dimensional Hilbert
spaces of columns and rows for a word with \(s>0\). Take the two
specified finite translates \(g_{\rm sym},g_1\). Inserting the
given column and row forms, denote the symmetric pairing on the
first translate by \(\langle\phi,SF\rangle\) and the position-one
pairing on the second by \(\langle\phi,JF\rangle\), with the
inner product linear in its second argument. Use the conjugate
column space if the original parametrization uses conjugated
forms. For a basis \(F_1,\ldots,F_n\) of \(W\), the two pairings
satisfy
\[
 \langle\phi,SF\rangle=\kappa\langle U^s\phi,JF\rangle
                 \quad(\phi\in\Phi,\ F\in W)
\]
for some unitary \(U\) on \(\Phi\) and some \(\kappa\ne0\) if
and only if a single \(t>0\) satisfies
\[
       \langle SF_j,SF_l\rangle=t\langle JF_j,JF_l\rangle
                        \qquad(1\le j,l\le n).
\]
Every entry on the right is a finite sum of products of symmetric
periods on the second translate. A positive scalar multiplying
\(U\) can be absorbed into \(\kappa\).
\end{proposition}
\begin{proof}
The two linear maps to \(\Phi\) are given by finite-dimensional
Riesz representation. For an orthonormal column basis \((h_k)\)
the \(k\)-th coordinate of \(JF_j\) is the finite exchange sum
of Theorem~\ref{thm:flagindex336} for the column \(h_k\) and the
row \(R(g_1)F_j\). Thus its Gram entries are obtained by summing
products of these coordinates with the first factor conjugated.
The coordinates of \(SF_j\) use \(R(g_{\rm sym})F_j\).
If the pairing identity holds, \(S=\kappa(U^s)^*J\), which
gives \(t=|\kappa|^2\). Conversely the Gram identities extend
to every pair of rows by sesquilinearity. The map
\(JF\mapsto t^{-1/2}SF\) is then well defined and preserves the
inner product. Extend it to a unitary \(V\) on \(\Phi\) by
matching orthonormal bases of the two orthogonal complements.
Diagonalizing \(V^*\) and taking an \(s\)-th root of each
eigenvalue gives a unitary \(U\) with \(U^s=V^*\). Taking
\(\kappa=\sqrt t\) proves the assertion.
\end{proof}
The finite Fourier identities determine the entries to be compared.
Their proportionality on the specified automorphic row space
remains to be proved.

\begin{corollary}[A translate annihilating the entire row space]\label{cor:zerotranslate336}
Suppose \(a,s\ge1\) and fix a split finite place and a finite
translate \(g_0\). There is a compact lattice \(C\subset k^a\)
such that for every \(x\notin C\), every column \(\phi\) and
every \(F\in W\),
\[
             L_\phi(R(X_i(x)g_0)F)=0 .
\]
If the position-one pairing on \(R(g_1)W\) is nonzero, the
proportionality in Proposition~\ref{prop:exchangegram336} fails
for each of these choices \(g_{\rm sym}=X_i(x)g_0\).
\end{corollary}
\begin{proof}
Open stabilizers of the translates of a finite basis of \(W\)
have open intersection. In particular \(Y_i(V)\) fixes all of
\(R(g_0)W\) for a single full lattice \(V\). The commutator
formula, with just the factor \(X_i(x)\), forces the symmetric
coefficient to vanish outside \(C=V^\perp\). This uses only
the character on \(N\) and is uniform in \(\phi\). The symmetric
Gram matrix is zero there. The nonzero position-one pairing
has a strictly positive diagonal entry in some row basis.
\end{proof}

The supplied source records both translates as data. In its
\(\mathrm{Tied}\) condition, the exponential-twist and
archimedean identities involve \(g_1\), and the root-of-unity
identity involves the character factors. All three are preserved
when \(g_{\rm sym}\) changes. The Hecke condition called
\(\mathrm{adm5}\) depends only on the underlying arithmetic datum.
Consequently, for any packet satisfying these conditions and
having a nonzero position-one pairing, they continue to hold
for choices of \(g_{\rm sym}\) excluded by the corollary. The
comparison for a selected packet must use a property of its
particular \(g_{\rm sym}\). The two definitions occur in
\texttt{ZetaProductPairs.lean}, lines 197--205 and 422--427.


\subsection{A Fourier operator on the refined column level}\label{sec:incidence338}
Let \(k=\mathbb Q_T\) at a split terminal place and take
\(\pi=\varpi\) on the chosen half. Write \(O=\mathbb Z_T\),
\(q=T\), \(K=\GL_a(O)\), \(V=O^a/\pi O^a\).
In \(V\) choose the last coordinate vector \(v_0\), and in its
dual choose the last coordinate covector \(\ell_0\). An integral
change of basis makes any prescribed primitive covector take this
form. Write \(P\) for the stabilizer of \(\ell_0\) in \(K\) and
\(P'\) for the stabilizer of \(v_0\). Use the twisted unitary
column action \(\rho=\widetilde R_H\) of
Section~\ref{sec:graphtransport337}. Its restriction to the finite
scalar residues on \(P\)- or \(P'\)-invariants is well defined
because lifts differ by an element congruent to one.
A superscript \(\circ\) on either fixed space means that the
sum of these scalar actions is zero on the vector.

Let \(\psi_0(\bar t)=\psi(t/\pi)\) for the simple additive character
of conductor \(O\). For \(0\ne\eta\in V^*\) choose \(t_\eta\in K\)
with \(\ell_0t_\eta^{-1}=\eta\) modulo \(\pi\). For \(v\ne0\)
choose \(k_v\in K\) with \(k_vv_0=v\) modulo \(\pi\).
Set \(h=\operatorname{diag}(1,\ldots,1,\pi^{-1})\). We use
\[
 \begin{split}
 B\phi&=\sum_{\eta\ne0}\psi_0(\eta(v_0))\rho(t_\eta)\phi ,\\
 C\xi&=\sum_{v\ne0}\overline{\psi_0(\ell_0(v))}\rho(k_v)\xi ,
 \\
 D\phi&=\rho(h^{-1}) B\phi .
 \end{split}
\]
\begin{theorem}[An invertible operator on the scalar zero-sum part]\label{thm:incidence338}
For every unitary \(H(k)\)-representation, \(D\) maps
\((\mathcal H^P)^\circ\) onto itself and
\[
        \|D\phi\|^2=q^a\|\phi\|^2,\qquad
        D^{-1}=q^{-a}C\rho(h).
\]
Both \(D\) and \(q^aD^{-1}\) are finite sums of group actions with
coefficients in \(\mathbb Z[\psi_0(\mathbb F_q)]\). They preserve
each scalar character subspace wherever that character is defined.
\end{theorem}
\begin{proof}
The choices in \(B\) differ by right multiplication by \(P\) and those
in \(C\) by right multiplication by \(P'\). Reindexing shows that
\(B\phi\) is \(P'\)-fixed, \(C\xi\) is \(P\)-fixed, and
\[
    \rho(k_v)B\phi=
       \sum_{\eta\ne0}\psi_0(\eta(v))\rho(t_\eta)\phi .
\]
The analogous identity for \(\rho(t_\eta)C\xi\) has coefficients
\(\overline{\psi_0(\eta(v))}\). The action of each scalar unit
commutes with every group element in the definitions. If
\(\phi\) has zero scalar sum, then
\(\sum_{\eta\ne0}\rho(t_\eta)\phi=0\). Indeed the summands for a
single line in \(V^*\) are one translate of the scalar sum, by
the stabilizer description. The same argument applies on vector
lines to \(\xi\). Summing the displayed identity against
\(\overline{\psi_0(\ell_0(v))}\) and including \(v=0\), whose term
is zero, gives \(CB\phi=q^a\phi\) by finite character orthogonality.
The second identity gives \(BC\xi=q^a\xi\) on the corresponding
vector space.

Put \(N=q^a-1\). In the first identity sum the squared norms over
all \(v\in V\), setting its value at zero equal to zero. The left
side is \(N\|B\phi\|^2\) by unitarity. On the right, orthogonality
leaves \(q^a\sum_{\eta\ne0}\|\rho(t_\eta)\phi\|^2=q^aN\|\phi\|^2\).
Lastly \(h^{-1}P'h=P\). In the matrix of a member of \(P'\),
the entries above the last diagonal entry in the last column
are divisible by \(\pi\) and the last diagonal entry is one
modulo \(\pi\). Conjugating by \(h^{-1}\) divides those entries by
\(\pi\) and multiplies the off-diagonal entries of the last row
by \(\pi\), giving exactly \(P\). The action of \(h^{-1}\) is
unitary with inverse \(\rho(h)\). This proves the theorem.
\end{proof}

On automorphic columns the action is by the specified finite
translations. It commutes with projection at a different place.
At the chosen place the theorem applies inside any closed subspace
stable under its full local group. In particular it acts on the
actual cuspidal column space of the specified real type with all
the other finite levels fixed.


\subsection{Bessel exchange on the original signed auxiliary rows}\label{sec:rationalperiod342}
Keep the signed auxiliary row projectors, the terminal graph and the
separately zero-sum columns on one field of Section~\ref{sec:missing}.
Choose $d_*$ from Theorem~\ref{thm:rationalflag342}. For the ordinary
global character $\chi$ of the unipotent group $N$, write
\[
                     \chi_*(n)=\chi(d_* n d_*^{-1}).
\]
It is a character of $N(F)\backslash N(\mathbb A_F)$ since $d_*$
is rational and normalizes $N$. The column group $H$ preserves it.
Use superscripts on $\mathcal B,\mathcal J$ to specify the character.

\begin{theorem}[A rational character compatible with the original projector]\leavevmode\label{thm:rationalperiod342}
For each original signed row $A$ and every indicated column $\phi$,
\[
 \mathcal J^{\chi_*}_\phi\bigl(\widetilde R(g_T)A\bigr)
    = C_T\!\prod_{v\mid g_1g_2}\prod_{i=1}^{s}\mu_{i,v}(\pi_v O_v^a)
                           \,\mathcal B^{\chi_*}_{D_T^s\phi}(A).
\]
Both sides use $A$ at the original signed auxiliary levels. The terminal
scalar and column operator are those of Theorem~\ref{thm:terminalwalk338}.
At either auxiliary place the character on the ordinary period agrees
with $\lambda_\Gamma$ throughout $N\cap\Gamma$. For any row and column
one also has the exact global identity
\[
              \mathcal B^{\chi_*}_\phi(A)
                         =\mathcal B^{\chi}_\phi(R(d_*)A),
\]
where $R(d_*)$ is the full ordinary adelic translation, including
all archimedean places. No Haar scalar occurs in this second identity.
\end{theorem}
\begin{proof}
In an auxiliary chosen half, conjugation by $d_*$ multiplies the
$j$-th string entry by the ratio of its two consecutive eigenvalues.
This ratio differs from the one in Theorem~\ref{thm:auxiliaryflag339}
by a unit congruent to one modulo $\pi_v$ and has valuation $-1$.
On $O_v$ both ratios give the same additive character, because their
difference lies in $O_v$ and the original conductor ideal is $O_v$.
The product formula on the string entries of
Lemma~\ref{lem:auxiliaryroots339} proves the compact character
agreement. The annihilator of $O_v$ for each new simple character
is now $\pi_v O_v$. The original signed row is fixed by
$Y_i(O_v^{a*})$ and $X_i(\pi_v O_v^a)$. These lattices are exact
annihilators of each other for the commutator pairing with $\chi_*$.
The finite Fourier inversion in Theorem~\ref{thm:exchangefourier336}
therefore has only its zero coset and contributes
$\mu_{i,v}(\pi_v O_v^a)$ for each integration on this row.

At $T$ every new simple-root ratio is one modulo $\pi_v$. In
particular the annihilator of $O_v$ is still $O_v$ and the
characters on $\pi_v^{-1}O_v/O_v$ are exactly the old ones. In the
single-step proof of Lemma~\ref{lem:flagstep338} the only finite
character entries are $\psi_i(\pi_v^{-1}y)$ for $y\in O_v$. Its
zero term, the transfer $D_v$ and each Haar factor are therefore
identical. The conjugation of the off-diagonal string entries still
has trivial character, since a diagonal cannot change their root
positions. This verifies the commutator calculation with the new
character separately at each simple root. The compact-support
argument of Theorem~\ref{thm:exchangefourier336} applies here because
at each step it uses only the annihilator of the current nontrivial
simple character and an open stabilizer of the smooth translated
row. It bounds the support at the finitely many marked places
successively. Thus one may perform all terminal steps and all
auxiliary integrations in the indicated order, proving the first
identity.

In the ordinary integral over $[H]$ and $[N]$, one has
$A(nh d_*)=A(d_*^{-1}n d_* h)$ by commutation with $H$ and left
invariance under $G(F)$. Make the change $n'=d_*^{-1}n d_*$.
Conjugation preserves $N(F)$, and the mass-one invariant measure
on its adelic quotient is preserved. Equivalently the local
Jacobians multiply to one by the product formula on its rational
Lie algebra. The character becomes $\chi(d_* n' d_*^{-1})$.
The column and its measure have stayed fixed. This proves the
second identity for the original Haar normalization.
\end{proof}

In the original Bessel normalization the rational translation on the
right includes its components away from the auxiliary primes. The
height bound in Theorem~\ref{thm:rationalflag342} controls the
valuations and the archimedean sizes of the rational flag coordinates.
Its lower bound forces an unmarked finite valuation on each isotropic
coordinate. Measuring the translations on the prescribed lattices
requires accounting for their integral frames. Identifying
the resulting signed value with the independent integral theta
column in \srcloc{Formalism/SharpRestriction/Divisibility/ZetaProductPairs}{107}
still requires evaluation with those components and its prescribed
endpoint projector. Theorem~\ref{thm:signedreturn342} gives exact
vanishing for the attempted comparison which simply projects a
matched auxiliary translate back to its original signed type.


\subsection{When the unmarked places annihilate the Bessel period}\label{sec:flagvanishing344}
Use the original auxiliary character, row level and terminal frame
of Theorem~\ref{thm:rationalperiod342}. Write $g=g_1g_2$ and
$\rho_j=c_j/c_{j+1}$ with $c_{s+1}=1$. At a good unmarked split
place the row has the full lattice level. The two flag halves are
integral and the absolute-trace character on every simple string
entry has conductor $O_v$. These properties hold at all but finitely
many completely split places. In fact the trace multiplier and the
coefficients of the original frame and its inverse are units or
integral as required outside finitely many primes. The standard
finite additive character of $\QQ_p$ has conductor $\ZZ_p$.
Whenever prescribing the support of the new valuations, we may
include the exceptions in the unmarked congruences in
Theorem~\ref{thm:splitflag343} or Theorem~\ref{thm:dominantflag344}.

\begin{lemma}[The integral root condition on a spherical row]\leavevmode\label{lem:rootvanish344}
At a good unmarked split place, $\chi_*(n)=\chi(d_*nd_*^{-1})$
is trivial on $N\cap K_v$ precisely when
$v_{\mathfrak p}(\rho_j)\ge0$ at both primes $\mathfrak p$ of
$E$ above $v$, for every $1\le j\le s$. A negative valuation on
any one root implies $\mathcal B^{\chi_*}_\phi(A)=0$ for every
global column $\phi$ and every original row $A$ fixed by $K_v$.
\end{lemma}
\begin{proof}
The string is $e_1,\ldots,e_s,z,f_s,\ldots,f_1$ on the chosen
half. The simple entries have multipliers $\rho_1,\ldots,\rho_s$
followed by $\bar\rho_s,\ldots,\bar\rho_1$ under conjugation.
The original additive coefficients are units. The character on
$N$ is their product on these simple entries, by the split trace
calculation of Lemmas~\ref{lem:exchangeroot336}
and~\ref{lem:auxiliaryroots339}. This last calculation of the trace
uses only the flag and the inverse adjoint on the second half.
Every element of $N\cap K_v$ has integral entries there, and each
individual simple string root subgroup $Z(O_v)$ lies in $K_v$.
The conductor therefore gives the equivalence.

If a multiplier has negative valuation, choose $u\in Z(O_v)$
with $\chi_*(u)\ne1$. The root involves only the string and fixes
the column space, so $u$ commutes with $H$. In the ordinary
integral, right translation of the unipotent variable by $u$ is a
measure-preserving change of variables on $[N]$. It multiplies the
value on replacing $A$ by $R(u)A$ by $\chi_*(u)$.
The row is fixed by $u$. Hence its period is zero for every column.
\end{proof}

\begin{theorem}[A lower bound for the finite denominators and exact vanishing]\leavevmode\label{thm:flagvanish344}
Suppose a rational diagonal matches the auxiliary residues and is
the identity in the split flag coordinates modulo $T^b$. Set
\[
       D_j^-=\sum_{\mathfrak p\nmid gT}
                \max(0,-v_{\mathfrak p}(\rho_j))\log N\mathfrak p .
\]
For every simple ratio one has the global inequality
\[
     \exp(D_j^-/n)
         \left(\prod_{\sigma:E\hookrightarrow\CC}
                    |\sigma(\rho_j)-1|\right)^{1/n}\ge T^b/g .
\]
In particular take a flag from Theorem~\ref{thm:splitflag343}
whose unmarked nonzero valuations occur at good split places
and for which $|\sigma(c_i)-1|\le\delta<1/2$ everywhere.
For $s\ge1$ one has $\mathcal B^{\chi_*}_\phi(A)=0$ for every
original row $A$ fixed at these good places and every column $\phi$.
Both sides of the
exchange identity of Theorem~\ref{thm:rationalperiod342} then
vanish on its separately zero-sum columns.
\end{theorem}
\begin{proof}
The element $y=g(\rho_j-1)/T^b$ is nonzero because
$v_{\mathfrak p}(\rho_j)=-1$ at every $\mathfrak p\mid g$.
There $g$ has valuation one and $T$ is a unit, so $y$ is a unit.
At the primes over $T$ the full congruence makes $y$ integral.
At any other prime where $\rho_j$ has negative valuation the
ultrametric inequality gives
$v_{\mathfrak p}(\rho_j-1)=v_{\mathfrak p}(\rho_j)$.
At all remaining primes $y$ is integral. Thus
$\log|N_{E/\QQ}(y)|\ge-D_j^-$. Evaluating that same norm at
the archimedean embeddings gives the displayed inequality.

For the last positive root $\rho_s=c_s$. Its archimedean geometric
mean in the formula is at most $\delta$, so
$D_s^-\ge n\log(T^b/(g\delta))>0$. Indeed the chosen terminal
prime satisfies $T\equiv1\bmod g$ and $T>g$ by
Theorem~\ref{thm:levels88}. At least one unmarked prime of
$E$ has negative valuation on this ratio. On the assumed support
it lies at a good split place where the row is spherical.
Lemma~\ref{lem:rootvanish344} forces the ordinary period to be zero.
The exchange formula expresses its first-position value as an
explicit multiple of this period evaluated at another column, so
that value is zero as well.
\end{proof}

\begin{proposition}[Exchange with integral ratios at every compensating place]\leavevmode\label{prop:dominantexchange344}
Use the diagonal of Theorem~\ref{thm:dominantflag344}, with the
unmarked exceptions included in $M$. The rational character
$\chi_*$ agrees with the prescribed signed auxiliary character on
$N\cap\Gamma$, has the original terminal residue characters and
is trivial on $N\cap K_v$ at every good unmarked place.
At each new split place there are compactly supported local row
and column tests at the original lattice levels with nonzero
ordinary local integral. For the actual global signed rows and
the separately zero-sum columns the exchange is
\[
 \mathcal J^{\chi_*}_\phi(\widetilde R(g_T)A)
    =C_T\!\prod_{v\mid g_1g_2}\prod_{i=1}^s\mu_{i,v}(\pi_v O_v^a)
            \,\mathcal B^\chi_{D_T^s\phi}(R(d_*)A).
\]
\end{proposition}
\begin{proof}
The prescribed auxiliary leading residues make the conductor
pairing on each integral string entry equal to the signed
character. At $T$ the ratios are one modulo $T^b$, so their
residual Fourier factors on $\pi_v^{-1}O_v/O_v$ are the original
ones. At a new split place apply Lemma~\ref{lem:rootvanish344}
to the integral ratios. At other good unmarked places the diagonal
and its inverse preserve the lattice and normalize $N$, so
conjugation preserves the original trivial restriction to
$N\cap K_v$.

For the local assertion take the row $\mathbf1_{K_v}$ on $G(F_v)$
and the column $\mathbf1_{H(F_v)\cap K_v}$ on $H(F_v)$. On the
column support, $nh\in K_v$ exactly when $n\in N\cap K_v$.
The character there is trivial. The integral is the product of
the positive volumes of these two compact open subgroups in $H$
and $N$. For the global identity repeat the auxiliary Fourier
inversions and terminal steps of Theorem~\ref{thm:rationalperiod342}.
They use the auxiliary and terminal characters just checked.
Rational conjugation on the ordinary period moves $d_*$ to the
row without a Haar factor. These identities hold for the specified
global automorphic vectors and require no archimedean proximity
to the identity.
\end{proof}

With these integral ratios the lower bound in
Theorem~\ref{thm:flagvanish344} applies with $D_j^-=0$ for
every $j$. Thus the required archimedean displacement is present
on each simple root. The local tests just constructed establish
compatibility of the character with the level. The signed global
theta coefficient at this displaced archimedean argument, its
arithmetic normalization and its size on the original fields
remain the period estimates required in Section~\ref{sec:missing}.

\section{Discrete series and real Bessel periods}\label{nav:g-VII-1}

\subsection{The size of the global integral on fixed compact types}\label{nav:VII-17}
Write \(F_+=E^+\) and use the signed row group and the two levels of
Corollary~\ref{cor:real91}. Fix a constant \(C_\lambda>0\). In the
interlacing lists of Theorem~\ref{thm:target90} suppose also that
the absolute values of the row parameters are at most
\(\exp(C_\lambda x^\eps)\). This includes its explicit lists for
large \(x\), since \(k,m,t_{\rm r}\) are bounded by fixed powers
of \(x\). Take the actual prime \(T\) of Corollary~\ref{cor:real91}
with the stronger generator congruence of Theorem~\ref{thm:neighbour93}.

\begin{lemma}[Bounds for the field, level and terminal prime]\label{lem:size97}
There is \(C_1\), depending only on the fixed family and parameter
bounds, such that, throughout the common intervals at large \(x\),
\[
 (16-o(1))m x^\eps\le\log T\le x^{C_1},\qquad
 \log B+\log I_\sigma\le x^{C_1}.
\]
Here \(B=(8n_0)^{2n_0}|D_E|^{r/2}\), \(n_0=2fr\). The level
index still has logarithm at most \(x^{C_1}\), after increasing
\(C_1\), when the principal subgroup of the enlarged lattice in
Theorem \ref{thm:project94} is imposed at every place above \(T\).
\end{lemma}
\begin{proof}
The degrees, \(k,m,r\), and the logarithm of the cyclotomic
conductor bound \(N\) are bounded by fixed powers of \(x\), by
Theorems~\ref{thm:rangenew50} and~\ref{thm:levels88}. The
conductor--discriminant formula gives \(\log|D_E|\le 2f\log N\).
The two auxiliary levels have \(I_\sigma\le(g_1g_2)^{fr^2}\),
by counting all residue matrices. These estimates bound \(\log B\)
and \(\log I_\sigma\).

For the upper bound on the terminal scale use the formula in
Corollary~\ref{cor:real91}. The compact mass \(\mathfrak m_r\)
is at most the number of its genus classes. The full-level
counting argument of
Lemma \ref{lem:mass32}
gives \(\log^+\mathfrak m_r\le C(2fr)^4\log(4frN)\), where
\(\log^+y=\max(0,\log y)\). In the denominator
\(d_{u,*}\ge\mathfrak m_k\ge(k!w_E^k)^{-1}\), with
\(w_E=|\mu(E)|\). Indeed the identity lattice has that automorphism
order. Each vector of norm one has a single nonzero coordinate
which is a root of unity, by the conjugate absolute-value argument
in Lemma~\ref{lem:mass90}. The columns are orthogonal. Also
\(\varphi(w_E)\le2f\), so \(w_E\le 2(2f)^2\) by the elementary
totient bound used in Theorem~\ref{thm:target90}. Each Weyl
dimension in the numerator is at most
\((2+2e^{C_\lambda x^\eps})^{r(r-1)/2}\), by its product formula.
Finally \(\Gamma\ge1\). The logarithm of the upper end of the
prime search interval is therefore at most a fixed power of \(x\).
Its lower bound is Corollary~\ref{cor:real91}.

At a place above \(T\), the principal subgroup of the enlarged
lattice contains the original principal subgroup modulo \(T^2\).
Indeed the enlargement has elementary divisors
\((T^{-1},1,\ldots,1)\) in the chosen split factor. Conjugation
back to the old coordinates can decrease the valuation of any
matrix entry by at most one. The dual factor is determined by
the first. The extra index is at most \(T^{2fr^2}\) by counting
matrices modulo \(T^2\). This proves the last assertion.
\end{proof}

At infinity use a \(K_\infty\)-finite smooth square-integrable row
form in discrete-series representations of these infinitesimal
characters. At each place take the compact type of highest weight
\(\mu_\tau=\lambda_\tau+\rho_{\rm n}-\rho_{\rm c}\), in the chamber
of its parameter. Here \(\rho_{\rm n}\) and \(\rho_{\rm c}\)
are the half sums of the positive noncompact and compact roots.
This specifies the type used below, also described in
\srcloc{Formalism/PeriodIntegrals/UpqRep}{231}. More generally the
next bound holds for any compact types and Casimir scalars with the
bound on \(L\) in its proof. Finite right translations in the
projector of Theorem \ref{thm:project94} preserve the archimedean type.

\begin{theorem}[Exponential decay in the terminal frequency]\label{thm:raw97}
For some absolute \(c>0\) and all sufficiently large \(x\)
in the common intervals, every such row form \(F\) of the
original level or the smaller level in Lemma~\ref{lem:size97}
and every smooth column form \(\phi\) satisfy
\begin{equation}
 |{\cal P}_{\varpi}(F,\phi)|\le
       \exp(-c\sqrt T/r)\,\|F\|_{L^2([G])}\,\|\phi\|_{L^2([H])}.
 \label{eq:raw97}
\end{equation}
The normalization is the analytic integral of the row and column
forms, with the character of last vector \(\varpi e\).
\end{theorem}
\begin{proof}
We first verify the derivative parameter \(L\) of Lemma \ref{lem:cas96}. Each coordinate of \(\mu_\tau\) has magnitude at most
\(r+\max_i|\lambda_{\tau,i}|\), since there are at most \(r-1\)
roots through a coordinate in the sums of positive roots.
In either compact unitary factor every weight coordinate lies
between the least and greatest highest weight entries of that
factor. Indeed weights are dominated by the highest weight and
permutation of entries preserves the weight set. Placing the
largest entry first in the dominance inequalities gives the upper
bound. Placing the least entry last and using equality of total
weights gives the lower bound. A compact Lie generator of
Frobenius norm one diagonalizes by that compact group. Its
eigenvalues on the type are scalar products of its diagonal
angles with the weights, times \(i\), so their magnitudes are
at most \(\sqrt r\max_i|\mu_{\tau,i}|\).

For the trace form the complex Casimir is \(\sum_{i,j}E_{ij}E_{ji}\).
On a highest weight vector of weight \(\nu\) the diagonal terms
sum to \(\sum_i\nu_i^2\). For \(i<j\) the two terms of that pair
sum to \(\nu_i-\nu_j\), by commuting the raising generator past
the lowering generator. The infinitesimal character \(\lambda_\tau\)
has the same central action as the compact representation with
highest weight \(\nu=\lambda_\tau-\rho\), after arranging the
entries in decreasing order. Here \(\rho_i=(r+1-2i)/2\).
The resulting scalar is
\(\omega_\tau=\sum_i(\lambda_{\tau,i}^2-\rho_i^2)\).
Thus both the compact Lie norm and \(\sqrt{|\omega_\tau|}\) are
bounded by \(L=\exp(C_2 x^\eps)\) for some fixed \(C_2\).

Apply Theorem \ref{thm:decay96}.
Its derivative orders can be chosen freely. Take
\(A=\lfloor\delta\sqrt T/r\rfloor\) at one real place, with an
absolute \(\delta>0\) small enough in terms of the absolute
constant in that theorem. Since \(r(L+fr^2+1)/\sqrt T\to0\)
by Lemma~\ref{lem:size97}, the ratio
\(Cr(L+A+fr^2+1)/(4\pi\sqrt T)\) is at most \(1/2\) for all
sufficiently large \(x\). Its \(A\)-th power accounts for the
factor \(2^{-A}\). The logarithm of the remaining multiplier
is at most
\[
 \tfrac12\log I+Cfr^2\log(2fr^2B)
                +fr^2\log\big(Cr(L+A+fr^2+1)\big).
\]
It is bounded by a fixed power of \(x\), by Lemma~\ref{lem:size97}.
On the other hand \(\sqrt T/r\) grows faster than every fixed
power of \(x\). Absorbing this multiplier into \(2^{-A}\) proves
\eqref{eq:raw97} with an absolute smaller \(c\).
\end{proof}

\begin{corollary}[A necessary arithmetic normalization]\label{cor:alg97}
Let \(p\ne0\) be an algebraic period in a number field \(J\)
whose values in its complex embeddings take the form
\[
    \sigma(p)=c_\sigma
         {\cal P}_{\varpi_\sigma}(F_\sigma,\phi_\sigma).
\]
Suppose the inputs in every embedding satisfy Theorem~\ref{thm:raw97}
with the same parameter bounds, and \(\varpi_\sigma\) is the
corresponding conjugate generator. Write
\(\mathfrak d(p)=\{a\in\mathcal O_J:ap\in\mathcal O_J\}\).
Then
\begin{equation}
 \frac{\log {\rm N}\mathfrak d(p)}{[J:\QQ]}+
 \frac1{[J:\QQ]}\sum_{\sigma:J\hookrightarrow\CC}
        \log\!\big(|c_\sigma|\,\|F_\sigma\|_2\,\|\phi_\sigma\|_2\big)
       \ \ge\ c\sqrt T/r .                                    \label{eq:alg97}
\end{equation}
In particular for integral \(p\) and unit inputs the geometric
mean of \(|c_\sigma|\) is at least \(\exp(c\sqrt T/r)\).
\end{corollary}
\begin{proof}
The fractional ideal of \(p\) is its integral numerator divided
by \(\mathfrak d(p)\), and the numerator has norm at least one.
Thus \(\sum_\sigma\log|\sigma(p)|\ge-\log {\rm N}\mathfrak d(p)\).
Use \eqref{eq:raw97} in each embedding and rearrange.
\end{proof}

For the high orders of Proposition~\ref{prop:der95}, one has
\(ha\log T=o(\sqrt T/r)\). Indeed \(h,r\) and \(\log T\) are
bounded by powers of \(x\), \(a=\lceil x\exp(x^\eps)\rceil\),
and \(\log T\ge(16-o(1))m x^\eps\), with \(m\to\infty\).
The denominator \(T^f\) of the finite projector likewise has
logarithm \(o(\sqrt T/r)\). The height sought in \eqref{eq:gap}
concerns an algebraic map between the prescribed lattices. The
source separately specifies a nonzero scalar and isometries for the
comparison of the two pairings in
\srcloc{Formalism/SharpRestriction/GlobalTheoremInputs}{65}.
Equation \eqref{eq:alg97} quantifies what this scalar, any change
in the test vectors, and any denominators must account for if the
raw fixed-type integrals represent a nonzero algebraic period.
A bound on that algebraic map and the uniform bound for
\(\mathcal R_{**}(x)\) in Proposition~\ref{prop:res80} remain open.

\subsection{A common real Bessel model for all translated columns}\label{nav:VII-18}
The Bessel subgroup at a real place is \(S=H_\tau\ltimes N_\tau\),
with \(H_\tau=U(k)\) acting on \(W_0\). Write \(\pi_u\) for the
compact column representation and \(\sigma_\lambda\) for a discrete
series in a row packet with ordered Harish-Chandra parameter
\(\lambda\). Since the period conjugates the column, its local
multiplicity is
\[
 m_\beta(\sigma_\lambda,u)=\dim\operatorname{Hom}_{S}
      (\sigma_\lambda^\infty\otimes\pi_u^\vee,\chi_\beta).
\]
Here \(\operatorname{Hom}\) means continuous equivariant linear
forms and the superscript denotes the smooth representation.
The column is finite dimensional. Existence of one such form
is a necessary condition on the entire real representation in a
nonzero global period on a factorizable pair. Fix the data at
the other places of such a pair. The resulting local form is
equivariant by change of variables in the period and continuous
by integration of smooth automorphic vectors over its compact
domain.

We use Xue's real local Gan--Gross--Prasad theorem
\cite[Theorem 1.1]{XVII-XII-I-xue108}. We describe the choices used in
applying it. Replace the Hermitian form by its negative at this
place. The unitary matrix groups and the flag subgroup stay
the same. In this convention \(W_0\) is negative definite and
the distinguished line \(\CC e\) is positive, as required in
Xue's definition for even \(k\). Pair the isotropic basis vectors
with their negatives on one side to restore pairing one.
The character used in the local theorem has a nonzero coefficient
on each of the same simple coordinates \(d_1,\ldots,d_s\).
Indeed the terms of Xue's trace pair a dual flag vector with the
image of the next flag vector, and at the last place in the flag
with the image of \(e\). These are exactly the defining coordinates
in Section \ref{nav:VI-08}.
A diagonal isometry scaling \(u_j\) by \(c_j\in\CC^\times\) and
\(v_j\) by \(\overline c_j^{-1}\), and fixing the middle space,
conjugates the coefficients by the respective factors
\(c_j/c_{j+1}\) and \(c_s\), or their inverses in the reverse
conjugation. Starting at \(s\) realizes arbitrary nonzero
coefficients. This conjugation fixes \(H_\tau\) pointwise.
The local multiplicities consequently agree for every
\(\beta\ne0\), including the prescribed terminal generator at
the given real place.

We recall the two sign computations used below. In a regular
discrete-series packet the ordered entries label the generators
of its component group, one of order two for each entry. The
restriction of a packet character to each generator is a sign.
Changing the sign at a given entry interchanges its positive and
negative positions in the chamber. More precisely, for a fixed
rank and Whittaker normalization the signature sign at position
\(j\) is its packet sign times \((-1)^{j-1}\) and a fixed sign.
This is the explicit real parametrization in
\cite[Appendix A]{XVII-XII-I-atobe108}. The calibrations with compact
groups below use only ratios, so that fixed sign cancels.

For the distinguished pair, Xue computes the label at a row
entry as the product of its epsilon factors against all column
entries. The label at a column entry is the product against all
row entries. Here the row exponents are the integers \(\lambda_j\)
and the dual column exponents are \(-u_i\), which are half-integers
outside \(\ZZ\). Thus a factor is evaluated at the angular
character with exponent \(\lambda_j-u_i\), using exponents for
\((z/\bar z)^t=(z/|z|)^{2t}\). For a half-integral exponent
outside \(\ZZ\), the real-place epsilon factor is one fixed sign
for \(t>0\) and its opposite for \(t<0\). This is
\cite[Proposition 2.1 and Corollary 2.2]{XVII-XII-I-ggp108}. Replacing the
additive character by its inverse, or using the conjugate
convention for all exponents, multiplies every such factor by
one common sign. In products of the same parity its contribution
cancels when the products are divided. The dual column entries
are indexed below by their original \(u_i\), even though their
decreasing order as dual weights is reversed.

\begin{proposition}[Crossing a row parameter]\label{prop:cross108}
Keep the row packet fixed and increase just the largest column
parameter from \(u_1^0\) to \(u_1^a\). At each row entry in the
open interval \((u_1^0,u_1^a)\) the distinguished row packet
label changes sign. Elsewhere it is unchanged. If at least one
entry is crossed, a single irreducible row representation cannot
have a local Bessel model for both columns.
\end{proposition}
\begin{proof}
All factors at a row generator cancel on division except the
one belonging to the first column entry. The sign of
\(\lambda_j-u_1\) changes precisely on the indicated interval.
The endpoints are disjoint from the integral row entries.
Xue's necessary condition then assigns two different characters
to the same row packet. A fixed irreducible member has just one
character. In particular Theorem~\ref{thm:pair107} crosses both
\(v_1+1\) and \(v_1+a\), since \(u_1^0=v_1+1/2\) and
\(u_1^a=v_1+a+1/2\). Its compact pairings across all translates
therefore cannot come from one irreducible real Bessel model
for both columns at a place where the assignment varies.
\end{proof}

\begin{theorem}[One real discrete series for both Bessel models]\label{thm:real108}
On the same fields and intervals, retain all column weights and
the shared \(v_i\). Use the decreasing row parameter
\begin{equation}
 \lambda=\big(v_1+a+2,\ v_1+a+1,\ v_1,\ldots,v_{k-1},\
        -1,\ -(t_{\rm r}+4),\ldots,-(m-3)(t_{\rm r}+4)\big).
 \label{eq:list108}
\end{equation}
In the original Hermitian form take the first slot negative,
the slots from the second through \(-1\) positive, and make the
remaining signs alternate starting negative. For the discrete
series in this chamber and each \(\beta\ne0\), one has
\[
       m_\beta(\sigma_\lambda,u^0)
       =m_\beta(\sigma_\lambda,u^a)=1 .
\]
Its lowest compact type also contains the two integral isometric
column embeddings of Lemma~\ref{lem:shuffle107}.
\end{theorem}
\begin{proof}
We determine the two packet characters separately, including
which definite space occurs. For either column the two leading
row entries lie strictly above \(u_1\), each gap
\((u_{i+1},u_i)\) contains just the common integer \(v_i\),
and exactly \(2s\) entries lie below \(u_k=-1/2\).
In particular their relative order is identical for the two
columns. The ranks satisfy \(r=k+2s+1\).

First keep this column, including its dual in the local Hom,
and take a comparison row of rank \(k+1\) on a positive definite
space that extends a positive definite \(k\)-space. Take one
integer above \(u_1\), the same \(v_i\) in the gaps, and one
integer below \(u_k\). The comparison has local multiplicity
one. Indeed for adjacent ranks the unipotent radical in the
Bessel model is trivial. Subtracting the half-sums of roots
converts strict interlacing of these integral and half-integral
parameters into the two weak highest weight inequalities of
ordinary compact branching. The compact restriction
has multiplicity one. Both compact spaces are positive and the
line added to the even-dimensional column is positive, so this
is a relevant pair in Xue's convention.

At column generator \(i\) the comparison row has \(i\) entries
above \(u_i\), whereas \eqref{eq:list108} has \(i+1\). Its
rank has increased by \(2s\), an even number. The ratio of the
epsilon products is therefore \(-1\) at every column generator,
whether one counts positive or negative differences. Relative
to the positive definite comparison column this reverses every
signature sign. Thus the required column lives on the negative
definite space for the sign-reversed form, as does exactly our
fixed \(\pi_u^\vee\).

For the row calibration keep the rank \(r\) and \(\lambda\) in
\eqref{eq:list108}. Compare with the positive definite compact
pair of ranks \(r,r-1\), taking column entries
\(\lambda_i-1/2\) for \(1\le i<r\). Strict integral spacing
makes these interlace, so compact branching again gives
multiplicity one. For the row generator at slot \(j\), the
comparison column has \(j-1\) entries above \(\lambda_j\).
Let \(b_j\) be the corresponding number in the actual column.
Since the two column ranks differ by \(2s\), the ratio at this
generator is \((-1)^{b_j-(j-1)}\). In order along the row,
\[
 b_1=b_2=0,\qquad b_{i+2}=i\ (1\le i<k),\qquad
 b_j=k\quad(j\ge k+2).
\]
The comparison row was everywhere positive. For the sign-reversed
form the ratio is therefore its required signature sign. It is
positive in the first slot, negative from the second slot through
\(-1\), and then alternates starting positive. Returning to the
original form gives the chamber in the statement. Its positive
and negative counts are \(k+s,s+1\), respectively. Both packet
characters now satisfy Xue's sufficiency condition, proving the
two equalities.

In this chamber \(j_i=i+1\) for the first \(k+1\) positive
slots. The first such slot has parameter \(v_1+a+1\), and hence
compact highest weight \(v_1+a+1-k/2=\nu_1^{\,a}\). The next
\(k\) positive slots give \(\nu^{\,0}_1,\ldots,\nu^{\,0}_k\)
by the calculation of Theorem~\ref{thm:pair107}. Its two
coordinate embeddings, with the same index sets \(I_0,I_a\),
apply to this lowest compact type.
\end{proof}

\begin{corollary}[A common terminal prime with every real model]\label{cor:realplace108}
For a common admissible \(\Gamma\), the row \eqref{eq:list108}
gives a single prime at its own formal-degree scale such that
the finite congruences and the analytic estimate of
Theorem~\ref{thm:raw97} hold, and every translated column admits
a nonzero continuous Bessel functional on the real row
representations. The same prime and row work for all placements
of the prescribed number of high weights.
\end{corollary}
\begin{proof}
The common integers \(v_i\) still interlace both column lists.
The largest absolute parameter is at most
\(a+2k+2+m(t_{\rm r}+4)\). The products of row formal degrees
and of column dimensions are independent of the placement.
The prime-window and doubled-modulus argument in the proof of
Theorem~\ref{thm:pair107} thus applies with the new row. At each
real place take a nonzero functional given by
Theorem~\ref{thm:real108}. Their tensor product at this finite
set of places is nonzero, as seen by evaluating at a tensor of
vectors with nonzero individual values. Each functional takes
a nonzero value on some compact-finite smooth vectors, since
the compact-finite subspace is dense in the smooth representation
by the Fourier expansion for the maximal compact group.
\end{proof}

Local multiplicity one determines a line of continuous
functionals. Two further facts would be needed to estimate the
global integral on the prescribed lattices. First, its restriction
to the specified vectors, including the lowest compact vectors
and the finite test data, needs a nonzero global value. Second,
the scalar relating that analytic value to an algebraic coefficient
must satisfy the arithmetic comparison of Corollary~\ref{cor:alg97}.
The compact pairing of value one in Lemma~\ref{lem:shuffle107}
is a pairing for \(H_\tau\). The local Bessel condition also
prescribes the action of \(N_\tau\).

\subsection{The prescribed row and the interior column weights}\label{nav:VII-19}
The row in \eqref{eq:list108} establishes existence when one can
choose its entries. The source fixes another row in
\srcloc{Formalism/SharpRestriction/CuspSpanNormBounds/RelevantZeroGramTerms}{165}.
We now use that row itself. At each real place its decreasing
Harish-Chandra parameter and its chamber are
\begin{equation}
      \lambda=(v_1,\ldots,v_{k-1},-b_1,\ldots,-b_m),
      \qquad (+,\ldots,+\,;\,+,-,+,-,\ldots,+,-).
      \label{eq:src109}
\end{equation}
There are \(k-1\) plus signs before the semicolon and \(m=2s+2\)
alternating signs after it. The \(v_i,b_j\) are integral.
The column entries are half-integers outside \(\ZZ\), and
\[
         u_1>v_1>u_2>\cdots>v_{k-1}>u_k>-b_1>
                         -b_2>\cdots>-b_m .
\]
These inequalities follow directly from the interlacing and
the positive gaps imposed in
\srcloc{Formalism/LocalRingCardinality/Params}{56}.
All words at a given place share \(v,b,e,u_1,u_k\) by
\srcloc{Formalism/LocalRingCardinality/RhoDetPrimeEqMuTildeMulChiPow}{96}.
Thus the interior entries distinguish words, while \(u_1\) may
also depend on which real place is under consideration.

\begin{theorem}[The real Bessel model on the prescribed row]\label{thm:src109}
For \eqref{eq:src109} and every column satisfying the displayed
interlacing, \(m_\beta(\sigma_\lambda,u)=1\) for every
\(\beta\in\CC^\times\). On the lowest compact row type, occurrence
of this column on restriction to \(H_\tau\) is equivalent to
the vanishing of every entry in the nonnegative integral vector
\[
 d_i=u_i-v_i-\tfrac12\quad(1\le i<k),\qquad
 d_k=u_k+b_1-\tfrac12 .
\]
If they all vanish, that compact type has the integral column
embedding of Lemma~\ref{lem:shuffle107}. In general any derivative
of the lowest type with a nonzero period against this column must
have order at least \(\sum_i d_i\) at this place.
\end{theorem}
\begin{proof}
Apply the two compact comparisons used in Theorem~\ref{thm:real108},
with the same sign-reversed form in the local multiplicity theorem.
The rank-\((k+1)\) positive compact comparison has \(i\) row
entries above \(u_i\), whereas \eqref{eq:src109} has \(i-1\).
The ranks differ by \(2s\), so the distinguished label at every
column generator changes sign. This gives exactly the negative
definite column in the comparison convention. For the row
calibration the rank-\((r-1)\) compact comparison again has
\(j-1\) column entries above \(\lambda_j\). There are actually
\(j\) such entries for \(j<k\) and \(k\) for \(j\ge k\).
The ratio of the required packet sign to the positive compact
one is \((-1)^{\min(j,k)-(j-1)}\). Upon returning to the
original Hermitian form, its negative gives plus in each of the
first \(k\) slots and alternating signs starting minus thereafter.
This is precisely the given discrete series. Xue's theorem
therefore gives the asserted multiplicity for each character.

The first \(k\) positive positions are \(j_i=i\). Consequently
the first \(k\) entries of the positive compact highest weight
are \(A_i=\lambda_i-k/2+i\) by
equation \ref{eq:mu106}.
For the compact column \(\nu_i=u_i-(k+1)/2+i\), this says
\(\nu_i-A_i=d_i\) for \(1\le i\le k\). Interlacing and
half-integrality imply \(d_i\ge0\). Compact branching requires
\(\nu_i\le A_i\) at each index. If all these entries are equal,
Lemma~\ref{lem:shuffle107} with \(I=\{1,\ldots,k\}\) gives
an integral embedding. Finally Proposition \ref{prop:jets106},
shows in its proof that after \(j\) derivatives each resulting
compact highest weight \(A'\) has
\(\sum_i\max(A'_i-A_i,0)\le j\). Since a contributing type
requires \(A'_i\ge\nu_i\) for \(1\le i\le k\), the sum of the
\(d_i\) is a lower bound.
\end{proof}

\begin{corollary}[Vanishing for the coded family]\label{cor:code109}
Suppose \(v_{i-1}-v_i=4\Delta\) and
\[
       u_i=v_{i-1}-\tfrac12-\Delta c_i,\qquad
       c_i\in\{0,1\},\quad 2\le i\le k-1 ,
\]
with \(\Delta>0\) common to the family and the map from words to
bit strings injective. These are the conditions in
\srcloc{Development/AbelianVarietyTheory/Definitions/AutomorphicMembership/ExistsSStageUnconditional}{16}.
If there are at least two words, every period on the lowest
compact row type vanishes for every word. Any differential word
giving a nonzero global value has total order at least
\[
       \sum_\tau (d_{\tau,1}+d_{\tau,k})
                       + f(k-2)(3\Delta-1).
\]
\end{corollary}
\begin{proof}
At least one bit is one somewhere. At that entry \(u_i\) is
in \(\ZZ+1/2\) and \(v_{i-1}\in\ZZ\), so \(\Delta\in\ZZ\)
and \(\Delta\ge1\). At every place and for every word, the
interior deficits are \(d_i=4\Delta-1-\Delta c_i\).
All of them are at least \(3\Delta-1\ge2\). Since the ranks
in Theorem~\ref{thm:rangenew50} have \(k>2\), the asserted
vanishing follows from Theorem~\ref{thm:src109}. A mixed
differential word of nonzero period must supply the separate
orders at all the places. Adding the endpoint and interior
contributions proves the bound.
\end{proof}

\begin{proposition}[One small real translate for all local models]\label{prop:translate109}
Fix the places and words above and a terminal generator. At each
place choose a nonzero lowest compact row vector \(w_\tau\), a
nonzero column vector for each word, and nonzero local Bessel
functionals. In every neighborhood of the identity in \(G_\infty\)
one translate of \(\bigotimes_\tau w_\tau\) has nonzero values
for all of these products of local functionals on the chosen
columns. This is possible with both matrix operator norms
\(\|g_\tau\|,\|g_\tau^{-1}\|\le2\) at every place. For any
automorphic row form on the original lowest type satisfying the
parameter hypotheses of Theorem~\ref{thm:raw97}, that theorem's
exponential upper bound also holds after any such bounded real
translate, with an absolute change in its constant.
\end{proposition}
\begin{proof}
Fix one local functional \(\ell\) and one nonzero dual column
vector \(z\). The subspace of column vectors annihilating the
entire smooth row is \(H_\tau\)-stable. By irreducibility of the
column and nonvanishing of \(\ell\), that subspace is zero.
Thus \(v\mapsto\ell(v,z)\) is nonzero. The compact-finite
Harish-Chandra module of a discrete series is irreducible
and dense in its space of smooth vectors. A nonzero lowest
type vector \(w_\tau\) generates this module under the complex
enveloping algebra. Indeed its generated submodule contains
its compact span, since the compact group is connected, and
conjugation on the Lie algebra then makes the generated
submodule compact-stable. Irreducibility gives the claim.
Some iterated Lie derivative on \(w_\tau\) therefore has
nonzero evaluation at \(z\).

The function \(g\mapsto\ell(\sigma_\lambda(g)w_\tau,z)\)
is real analytic near the identity. To see this in the smooth
topology, use ordered exponential coordinates. The derivative
bound of Lemma \ref{lem:cas96}, bounds each word of length
\(n+t\) on \(w_\tau\) by \(\{Cr(L+n+t+1)\}^{n+t}\|w_\tau\|\).
For fixed \(L,t\), dividing by \(n!\) leaves at most an exponential
in \(n\), with the base bounded by a fixed multiple of \(r\) in
the limit of \(n\)-th roots (use \(n!\ge(n/e)^n\)).
At total degree \(n\) the multinomial formula trades the product
of the individual factorials in the Taylor series for \(n!\)
and the \(n\)-th power of the sum of absolute coordinate sizes.
The series thus converges for a sufficiently small such sum in
every Hilbert derivative seminorm of fixed order \(t\).
These Taylor series represent the group translates: along
each one-parameter group this follows from Taylor's formula
with integral remainder and unitarity, and one successively
applies this identity in the ordered coordinates using absolute
convergence. The continuous functional may be applied to the
series. A nonzero iterated derivative at the identity precludes
identical vanishing there. The nonzero set of each such coefficient
is consequently dense and open in a connected coordinate ball.
Intersecting the finitely many sets at each place gives the
common translate in an arbitrarily small product neighborhood.

For the asserted bound write \(R(g)\) for right translation on
automorphic forms. Move each real derivative past \(R(g)\) by
acting on its Lie algebra argument by \(\operatorname{Ad}(g^{-1})\).
Its Frobenius norm increases by a factor at most four. By
multilinearity a word of length \(j\) in the proof of Theorem \ref{thm:decay96}, costs at most \(4^j\) in addition to the old
estimate. Right translation is unitary and leaves the finite
level fixed. In the proof of Theorem~\ref{thm:raw97} the added
factor is absorbed by reducing the fixed coefficient of
\(\sqrt T/r\) in the chosen differentiation order. The factor
from the at most \(fr^2\) further restriction derivatives has
logarithm \(O(fr^2)\), which is absorbed as in that proof.
\end{proof}

Here the nonzero values are evaluations of the individual real
Bessel functionals on smooth vectors. A global theta lift still
requires its specified finite tests and its arithmetic normalization.
The preceding construction gives no arithmetic lattice for a
general small real translate.

\subsection{Full periodization below the auxiliary height bound}\label{nav:VII-20}
On the common fields and intervals put \(Q_0=g_1g_2\).
Use the isotropic height \(H\) of
Section \ref{nav:VII-03}.
It is the geometric mean over the real places of the largest
Gram--Schmidt length in the ordered list \(g_\tau^{-1}u_i\).
In particular it is a function on \(B_\infty\backslash G_\infty\).
Write \(\mathcal U=\{H<\sqrt{Q_0}\}\). This open set contains
the real maximal compact group in the signed coordinates. Indeed
on that group the \(u_i\) remain Euclidean orthogonal of length
\(\sqrt{2}\) at every place. More generally it contains every
point with \(\prod_\tau\|g_\tau^{-1}\|_{\mathrm{op}}^{1/f}
<\sqrt{Q_0/2}\), by using the original unmodified isotropic
vectors in the length minimization.

\begin{theorem}[The full automorphic norm on the height range]
\label{thm:low136}
Use \(K_\sigma\) with either or both terminal frequencies, or
use the refined column level realized by \(K^0\) with frequency
\(\varpi\) as in Theorem \ref{thm:midlevel131}. Suppose the real
functions in the period seeds vanish almost everywhere outside
\(\mathcal U\). On the family of Section \ref{nav:I-28}, for all sufficiently
large \(x\), summation over the full rational unitary group
defines an isometry in the natural seed norms,
\begin{equation}
 \left\|\sum_{\delta\in B(F_+)\backslash G(F_+)}
                         \Phi(\delta g)\right\|^2_{
                  L^2(G(F_+)\backslash G(\mathbb A_{F_+}))}
     =\|\Phi\|^2_{L^2(B(F_+)\backslash G(\mathbb A_{F_+}))}.
 \label{eq:low136}
\end{equation}
This holds for finite linear combinations of the selected types
and columns, including the refined zero-sum columns. For smooth
real sections compactly supported modulo \(B_\infty\) the sum is
ordinary. For \(L^2\) sections supported in \(\mathcal U\) it is
the isometric extension of these sums. Thus different selected
words give orthogonal automorphic functions in this construction.
For a given word and frequency with real section \(W\) of norm
one, its squared norm in class coordinates is
\[
    \frac{\operatorname{vol}_f(K)\,v_N\,\operatorname{vol}(H_\infty)}
         {\operatorname{vol}_f(N_f\cap K)\,d(\rho_c)}
         \sum_j\|\ell_j\|^2 .
\]
\end{theorem}
\begin{proof}
First take smooth real sections of compact support inside
\(\mathcal U\). Their scalar norms, with the original finite
support, descend to compactly supported envelopes modulo \(B(F_+)\).
The compact covering argument in Theorem~\ref{thm:maximal130}
shows that periodizations are locally finite and that the
product of two periodized envelopes is integrable. We may
therefore expand the mixed Gram pairing, unfold its first
rational index and then group the remaining terms, with absolute
integrability. Two finite cosets can meet only if the relative
rational element belongs at finite places to \(B_fKB_f\).

Group together the pairs with the same underlying isotropic
\(s\)-space. Their total integral is the Gram pairing of the
corresponding sums over the stabilizer of that space. Indeed
partition the index set of each sum into fibres over
\(Q(F_+)\backslash G(F_+)\), with \(Q\) the stabilizer of \(L\).
The terms under consideration are exactly the products with the
same fibre index. Unfolding that index integrates the two
intermediate sums over \(Q(F_+)\backslash G(\mathbb A_{F_+})\).
Theorem~\ref{thm:maximal130} evaluates this as the input Gram
pairing, also at \(K^0\) by Theorem \ref{thm:midnorm131}.

Consider a pair of different underlying isotropic spaces.
For each finite cell the two unipotent characters either
disagree or agree on their common real subgroup. In the first
case the orbital integral is zero by Lemma \ref{lem:zero133}.
In the second case Theorem \ref{thm:height135}, gives
\(H(g)H(\gamma^{-1}g)\ge Q_0\) for every real argument \(g\).
The ranks tend to infinity on these intervals, so \(s\ge2\)
as required there. If the two sections are both nonzero at this
argument, both heights are strictly below \(\sqrt{Q_0}\),
contradicting that inequality. Thus the integrand vanishes in
every such cell. This argument treats a word paired with itself
as well as distinct words. Only the already evaluated
same-subspace contribution survives, proving \eqref{eq:low136}.

Smooth compact sections inside \(\mathcal U\) are dense in the
space of square-integrable sections supported there. In local
trivializations first truncate to compact subsets of the open set
and then approximate by smooth functions using a partition of
unity. The associated bundles have finite-dimensional fibres
and the quotient Haar measure is a smooth density, so this
procedure gives norm convergence. The finite class description
in Section \ref{nav:VI-15}, carries the approximations to the
seed norm. The identity for differences shows that their
automorphic images are Cauchy. Their limits in the complete
automorphic \(L^2\) space are independent of the approximations
and give the asserted isometry on the closure. Polarization
gives the mixed identity. Finally, finite class and compact
column integration in Theorem \ref{thm:identity127}, give the displayed
constant and the input orthogonality of distinct types.
\end{proof}

\begin{corollary}[The truncated part of a specified real section]
\label{cor:cut136}
For a norm-one real section \(W\) of a selected word and
frequency, set
\[
    \eta(W)=\int_{\{H\ge\sqrt{Q_0}\}}\|W(g)\|^2\,d\dot g
         \quad (d\dot g\text{ on } B_\infty\backslash G_\infty).
\]
Let \(\Phi^{<}\) be its seed with real section
\(\mathbf1_{\mathcal U}W\). For any column \((\ell_j)\), the
squared automorphic norm of the extension in
\eqref{eq:low136} for \(\Phi^{<}\) equals
\((1-\eta(W))c_{\rho,K}\sum_j\|\ell_j\|^2\). At \(K^0\) the
squared class-coordinate constant is the same as at \(K_\sigma\).
Their squared constants for the column probability norms have
ratio \((T^k-1)^f\).
\end{corollary}
\begin{proof}
The cutoff is left \(B_\infty\)-invariant, so it preserves
the transformation law and scales the squared real norm to
\(1-\eta(W)\). Apply Theorem~\ref{thm:low136} and Theorem \ref{thm:midnorm131}. The zero-section case gives zero on both sides.
\end{proof}

The construction evaluates a full automorphic pairing for a
specified height truncation. For the originally prescribed theta
vectors one still needs a bound for the contribution from
\(H\ge\sqrt{Q_0}\) after periodization, including its pairing
with the truncated part. A cutoff can also change the real
representation generated by a section. The integral theta
comparison requires the prescribed projection and the
arithmetic normalization in addition to the complex norm.

\subsection{The prescribed compact row type and real translations}\label{nav:VII-21}
Fix a row in \eqref{eq:src109} at one of the choices in
Theorem~\ref{thm:common111}, and let \(K_\infty^{\rm max}\)
be its real maximal compact group in the signed coordinates.
Write \(\sigma\) for the product of the prescribed real discrete
series, \(U\) for the tensor product of their lowest compact
types, and \(\mathcal C_c\) for the space of finite class
covectors of type \(\rho_c\), with its ordinary sum inner product.
We use the refined norm-\(T\) level \(K^0\).

\begin{theorem}[Automorphic copies of the prescribed compact type]
\label{thm:copies137}
There is a nonnegative smooth cutoff \(\theta\) on
\(B_\infty\backslash G_\infty\), right invariant under
\(K_\infty^{\rm max}\) and compactly supported in \(\{H<2\}\),
with the following property. Using the real Bessel models of
the prescribed discrete series, the corresponding cutoff sums
give a \(K_\infty^{\rm max}\)-equivariant injection
\[
    \bigoplus_{c\in\mathcal A} U\otimes\mathcal C_c
                  \ \longrightarrow\ L^2([G])^{K^0}.
\]
The images for distinct words are orthogonal. On the summand for
\(c\) the map is an isometry after a single positive scalar
normalization, also on its refined zero-sum columns.
\end{theorem}
\begin{proof}
For each \(c\), the product of the nonzero local Bessel
functionals of Theorem~\ref{thm:src109} gives a smooth
column-valued function
\[
     W_{c,v}(g)=\Lambda_c(\sigma(g)v),\qquad  v\in U.
\]
Here the functional is viewed as a linear map on smooth row
vectors with values in \(V_{\rho_c}\), by duality with the
finite-dimensional dual column. It transforms under the left
unipotent and column groups by the factors in equation \ref{eq:real127}
with terminal vector \(\varpi e\). Right translation by a member
of \(K_\infty^{\rm max}\) replaces \(v\) by its compact translate.
At the identity \(H=\sqrt2\), and \(H\) is unchanged by right
translation by the maximal compact group since Euclidean norms
are unchanged. Choose a relatively compact neighborhood of the
identity coset and its compact orbit inside \(\{H<2\}\). Take a
nonnegative smooth bump supported in this open set, equal to one
near the orbit, and average it over the compact group. The result
is again compactly supported there and equals one near the orbit.

For every \(c\) choose a nonzero pure tensor \(v\in U\).
Proposition~\ref{prop:translate109} gives points arbitrarily
close to the identity at which \(W_{c,v}\) is nonzero.
Continuity gives positive integral for \(\|\theta W_{c,v}\|^2\).
The Hermitian form on \(U\) obtained by integrating these
cutoff sections is invariant under the compact group, by right
invariance of the cutoff and quotient measure. Irreducibility
of \(U\) and Schur's lemma make it \(a_c\) times the given row
inner product for some \(a_c>0\).

Form the period seed from \(\theta W_{c,v}\) and any class
covector tuple. The map is linear in both arguments and extends
to their tensor product. Formula \eqref{eq:low136}, polarized
on such tensors, gives \(a_c c_{\rho_c,K^0}\) times the tensor
product inner product. To see the two factors directly, first
average the pointwise rank-one operator over the left compact
column group as in Theorem \ref{thm:identity127}. This gives the
inner product of the two real sections divided by \(d(\rho_c)\)
times the identity on the column representation. Finite class
integration then pairs the class tuples. Distinct words have
zero inner product by Theorem~\ref{thm:low136}. Multiplication
on each summand by \((a_c c_{\rho_c,K^0})^{-1/2}\) proves the
assertion. The construction is right equivariant for the
maximal compact group since the cutoff is invariant and
ordinary periodization commutes with right translations.
\end{proof}

For a real group element \(a=(a_\tau)\) put
\[
 \ell(a)=\frac1f\sum_\tau\log
       \max\{\|a_\tau\|_{\rm op},\|a_\tau^{-1}\|_{\rm op}\}.
\]
It is nonnegative, symmetric in \(a,a^{-1}\), and subadditive
under multiplication by the matrix norm inequality. Write
\(R_B,R_G\) for right translations on the seed and automorphic
quotients, and \(P\) for the isometry of Theorem~\ref{thm:low136}.

\begin{theorem}[Agreement of translated matrix coefficients]
\label{thm:matrix137}
Let \(\Phi,\Psi\) be seeds as in Theorem~\ref{thm:low136},
supported at heights \(H<R,H<S\), respectively, with
\(R,S\le\sqrt{Q_0}\) and \(RS<Q_0\). Put \(h=\log(Q_0/(RS))\).
For every \(a\) with \(\ell(a)<h\),
\[
     \langle R_G(a)P\Phi,P\Psi\rangle
                       =\langle R_B(a)\Phi,\Psi\rangle .
\]
For an arbitrary finite complex Borel measure \(\mu\) on
\(G_\infty\), let \(R(\mu)=\int R(a)\,d\mu(a)\).
The difference of these two pairings with \(a\) replaced by
\(\mu\) has absolute value at most
\begin{equation}
     2\,|\mu|\{\ell\ge h\}\,\|\Phi\|\,\|\Psi\|.
     \label{eq:matrix137}
\end{equation}
\end{theorem}
\begin{proof}
Right translation changes the height by bounded factors,
\[
          e^{-\ell(a)}H(g)\le H(ga)\le e^{\ell(a)}H(g).
\]
Indeed apply \(a_\tau^{-1}\) to a set of isotropic vectors
minimizing the maximal length at \(g_\tau\). Take the geometric
mean and use the same upper bound with \(a^{-1}\) for the lower
bound. The support of \(R_B(a)\Phi\) is therefore in
\(H<e^{\ell(a)}R\). First take compact smooth input sections.
Their periodizations are ordinary and commute with right
translation. The proof of Theorem~\ref{thm:low136} for a mixed
pair only uses that the product of the two supporting height
bounds is less than or equal to \(Q_0\). It applies to
\(R_B(a)\Phi,\Psi\), since \(e^{\ell(a)}RS<Q_0\).
Their same-subspace pairing again unfolds to the seed pairing.
This proves equality even if the first translated height bound
exceeds \(\sqrt{Q_0}\).

For general seeds approximate each by compactly supported
sections inside its original height bound. Unitarity of the
two right regular representations and isometry of the original
periodizations allow passage to the limit on both sides.
For every \(a\) the absolute value of each pairing is at most
\(\|\Phi\|\|\Psi\|\) by Cauchy--Schwarz. Integrating their
difference against \(\mu\), the part where \(\ell(a)<h\) is zero.
Integration against total variation on the complement gives
\eqref{eq:matrix137}.
\end{proof}

\begin{corollary}[A uniform one-parameter spectral estimate]
\label{cor:heat137}
In either right regular representation let \(D_X\) denote the
self-adjoint generator given by \(R(\exp(vX))=\exp(ivD_X)\).
Put \(L_X=f^{-1}\sum_\tau\|X_\tau\|_{\rm op}\).
For \(t>0\) the two pairings of Theorem~\ref{thm:matrix137}
with the corresponding \(e^{-tD_X^2}\) differ by at most
\[
    2\exp\!\left(-\frac{h^2}{4tL_X^2}\right)\|\Phi\|\|\Psi\|
                    \qquad (L_X>0).
\]
For \(L_X=0\) they agree. For the copies in
Theorem~\ref{thm:copies137}, \(L_X\le1\) and \(t=1\), this
error is at most \(2e^{-c x^{2\eps}}\|\Phi\|\|\Psi\|\) with an
absolute \(c>0\) throughout the common intervals.
\end{corollary}
\begin{proof}
The scalar Fourier transform of the Gaussian probability
\((4\pi t)^{-1/2}e^{-v^2/(4t)}\,dv\) is \(e^{-t\xi^2}\).
The spectral theorem thus realizes the heat operator by
integrating this one-parameter unitary group against that
probability. The matrix exponential series gives
\(\ell(\exp(vX))\le |v|L_X\). The Gaussian mass of
\(\{|v|\ge h/L_X\}\) is at most \(\exp(-h^2/(4tL_X^2))\).
Indeed after rescaling both halves use
\(e^{-(u+z)^2}\le e^{-z^2}e^{-u^2}\) for \(u,z\ge0\).
Equation~\eqref{eq:matrix137} proves the first bound.
For \(L_X=0\) the generator is zero. For the displayed copies
take \(R=S=2\). The auxiliary primes are at least
\((k\exp(x^\eps))^8\), so \(h=\log(Q_0/4)\ge 15x^\eps\)
for all sufficiently large \(x\), which gives the last estimate.
\end{proof}

The positive constants \(a_c\) come from cutoff real models.
The construction realizes the compact row representation
in automorphic \(L^2\) and controls the stated translation
averages. Identifying the specified discrete-series projection
there and controlling its integral theta periods require further
estimates.

\subsection{Full automorphic norms for real sections with a height moment}\label{nav:VII-22}
Take the fields and intervals of Theorem~\ref{thm:invar148},
and retain the finite columns, levels and frequencies of
Theorem~\ref{thm:low136}. In this paragraph \(r=\dim_E V\)
denotes the row rank, \(s\) the length of its isotropic flag,
and \(f=[F_+:\QQ]\). Put
\[
 Q_0=g_1g_2,\qquad R_0=\sqrt{Q_0},\qquad D_0=2frs,
 \qquad m=2D_0,\qquad C_*=\frac{4}{\sqrt3}\,17^{rs}.
\]
Write \(P\) for summation from the seed quotient to the full
automorphic quotient, with the extension constructed in
Theorem \ref{thm:tail149}.
For a seed \(\Phi\) put \(\Phi^{<}=\mathbf1_{H<R_0}\Phi\).

\begin{proposition}[A real moment sufficient for the complex norm]
\label{prop:mom150}
If \(\|H^m\Phi\|_{\rm seed}\le A^m\|\Phi\|_{\rm seed}\), then
the full automorphic sum exists in \(L^2\) and
\[
  \left|\|P\Phi\|-\|\Phi\|_{\rm seed}\right|
          \le (1+C_*)(A/R_0)^m\|\Phi\|_{\rm seed}.
\]
For any orthogonal projection \(\Pi\) on automorphic \(L^2\),
\(\|\Pi P\Phi-\Pi P\Phi^{<}\|
\le C_*(A/R_0)^m\|\Phi\|_{\rm seed}\).
If \(A\le R_0\exp(-\eta x^\eps)\) for fixed \(\eta>0\),
both relative errors, normalized by \(\|\Phi\|_{\rm seed}\),
are at most \(\exp(-\eta D_0 x^\eps)\) for all sufficiently
large \(x\).
\end{proposition}
\begin{proof}
Take \(\sigma=D_0\) in Theorem \ref{thm:tail149}. Its weighted
tail is at most \((A/R_0)^m\|\Phi\|_{\rm seed}\), as is the
ordinary seed norm above \(R_0\). Summation is isometric on the
lower part. The reverse triangle inequality applied to the
two decompositions proves the first assertion. The operator
norm of an orthogonal projection is at most one, which gives
the second. Finally \(\log(1+C_*)\le C D_0\) with an absolute
\(C\), since \(f\ge1\). This proves the uniform error bound.
\end{proof}

There is also a family for which the contribution after
removing the height cutoff can be estimated directly. The
initial sections can be the compact row copies of
Theorem~\ref{thm:copies137}. For real \(a\) use the symmetric
length \(\ell(a)\) and right translations \(R_B(a),R_G(a)\)
of Theorem~\ref{thm:matrix137}.

\begin{theorem}[Full periodization of translation averages]
\label{thm:smooth150}
Let an initial seed \(\Phi\) be supported in \(H<R<R_0\).
Let \(\mu\) be a complex Borel measure on \(G_\infty\) with
\(\int e^{m\ell(a)}\,d|\mu|(a)<\infty\).
Set \(W=R_B(\mu)\Phi\) and
\(\delta=|\mu|\{\ell\ge\log(R_0/R)\}\). Then
\[
         P W=R_G(\mu)P\Phi,\qquad
         \|P(W-W^{<})\|\le2\delta\|\Phi\|_{\rm seed},
         \qquad W^{<}=\mathbf1_{H<R_0}W .
\]
In addition,
\[
          \big|\|P W\|-\|W\|_{\rm seed}\big|
                     \le 2\delta\|\Phi\|_{\rm seed}.
\]
\end{theorem}
\begin{proof}
Theorem~\ref{thm:matrix137}, gives
\(H(ga^{-1})\le H(g)e^{\ell(a)}\).
Thus the seed \(R_B(a)\Phi\) is supported at heights below
\(R e^{\ell(a)}\). Minkowski's integral inequality and
unitarity of right translation give
\[
   \|H^m R_B(\mu)\Phi\|_{\rm seed}
          \le R^m\|\Phi\|_{\rm seed}\int e^{m\ell(a)}d|\mu|(a).
\]
In particular the weighted domain of \(P\) contains \(W\).
For compactly supported \(\mu\) and bounded compactly
supported seeds, ordinary summation commutes with its right
translations and their integral. Indeed on compact test sets
only finitely many rational indices occur, uniformly for \(a\)
in a compact set, by the compact lifting argument of
Theorem~\ref{thm:maximal130}. For a square-integrable initial
seed in \(H<R\), approximate it there by bounded compact
sections. Its translates for \(a\) in the support of such a
\(\mu\) lie in one bounded height range. Theorem \ref{thm:sectors149},
gives continuity on that range, proving commutation in the
limit. Finally restrict a general \(\mu\) to increasing compact
sets exhausting \(G_\infty\). The omitted integrals converge
to zero in the weighted seed norm by the displayed inequality
and in the ordinary seed norm by their total variations.
Theorem \ref{thm:tail149}, allows passage to the limit on the
automorphic side. This proves \(PW=R_G(\mu)P\Phi\).

Split \(\mu=\mu_0+\mu_1\) at
\(\ell=\log(R_0/R)\), giving the boundary to \(\mu_1\), and
put \(W_i=R_B(\mu_i)\Phi\). The support of \(W_0\) lies in
\(H<R_0\) up to a null set, so \(\|PW_0\|=\|W_0\|_{\rm seed}\).
The seed norm of \(W_1\) is at most
\(\delta\|\Phi\|_{\rm seed}\); its automorphic norm satisfies
the same bound by commutation and \(\|P\Phi\|=\|\Phi\|_{\rm seed}\).
Since \(W-W^{<}=W_1-W_1^{<}\) and the isometry applies to
\(W_1^{<}\), its periodization satisfies the first error bound.
Comparing the two full norms with the equal norms of \(W_0\)
proves the last inequality.
\end{proof}

\begin{corollary}[Gaussian averages preserving the prescribed compact row]
\label{cor:gauss150}
On the common fields choose \(\Phi\) in the compact row copies
of Theorem~\ref{thm:copies137}, with real support \(H<2\).
Let \(X=(X_\tau)\) be in the real Lie algebra and put
\(L_X=f^{-1}\sum_\tau\|X_\tau\|_{\rm op}\). For \(t>0\)
average right translates by
\[
          k\exp(uX)k^{-1},\qquad
       k\in K_\infty^{\rm max},\qquad
       (4\pi t)^{-1/2}e^{-u^2/(4t)}\,du ,
\]
using Haar probability for \(k\). For the resulting section
\(W\), Theorem~\ref{thm:smooth150} applies to the entire
Gaussian. If \(L_X>0\), each of its two error bounds is at most
\[
 2\exp\!\left(-\frac{\log^2(R_0/2)}{4tL_X^2}\right)
           \|\Phi\|_{\rm seed}.
\]
For \(t=1\) and \(L_X\le1\) this is at most
\(2e^{-c x^{2\eps}}\|\Phi\|_{\rm seed}\) with absolute
\(c>0\) at all sufficiently large \(x\). The averaged
automorphic map on the original copies is injective and
equivariant for \(K_\infty^{\rm max}\). The same tail bound
holds after any orthogonal automorphic projection.
\end{corollary}
\begin{proof}
Compact conjugation preserves each operator norm, and the
exponential series gives
\(\ell(k\exp(uX)k^{-1})\le |u| L_X\).
The Gaussian has all the required exponential moments.
The mass beyond \(\ell=\log(R_0/2)\) is bounded by the
two-sided Gaussian tail in the proof of
Corollary~\ref{cor:heat137}, giving the displayed estimate.
The auxiliary levels in that corollary satisfy
\(\log(R_0/2)\ge7x^\eps\) eventually. If \(L_X=0\),
the translation is the identity and the tail is zero.
For fixed \(k\), the operator on either quotient is unitarily
conjugate to \(e^{-tD_X^2}\). Its quadratic form on any
nonzero vector is strictly positive, since the spectral
integrand \(e^{-t\lambda^2}\) is positive at every real
\(\lambda\). Its compact average is therefore injective, and
it commutes with the compact group by invariance of Haar
measure. Compose it with the injection of
Theorem~\ref{thm:copies137} and use the exact commutation in
Theorem~\ref{thm:smooth150}. Applying a contraction to the
tail gives the last assertion.
\end{proof}

These estimates on the automorphic quotient hold on the fields
where \eqref{eq:invar148} bounds the quadratic invariants.
For the specific theta sections of the iterated lift, the
moment in Proposition~\ref{prop:mom150}, or a representation
by translates with the required tail, still has to be evaluated.
The scalar comparing the projected complex functional with
the integral theta lattice and the remaining strict orders in
\eqref{eq:tail114} are also required for \eqref{eq:window}.

\subsection{The prescribed discrete series at the prime terminal frequency}\label{nav:VII-23}
Use the fields of Theorem~\ref{thm:invar148} and the actual
common terminal prime \(T=\varpi\bar\varpi\) of
Theorem~\ref{thm:common111}. Write \(m_0=2(s+1)\) for the rank
parameter of the chain, \(r=k+m_0-1\) for the row rank,
\(f=[F_+:\QQ]\), \(D_0=2frs\), and \(R_0=\sqrt{g_1g_2}\).
We use the original frequency \(\varpi\) at the refined level.

\begin{proposition}[The moment required at frequency one]
\label{prop:scale152}
Suppose a nonzero pure tensor of real Bessel sections of the
prescribed row has square-integrable factors \(W_{1,\tau}\)
at frequency one. Normalize them to norm one. Let
\[
        A_\tau=
             \|h_\tau^{4rs}W_{1,\tau}\|_2^{\,1/(4rs)}
             \quad\text{and}\quad A_1=\Big(\prod_\tau A_\tau\Big)^{1/f}.
\]
For the transported norm-one tensor at frequency \(\varpi\),
the height moment satisfies
\[
        \|H^{2D_0}W_\varpi\|_2^{\,1/(2D_0)}=\sqrt T\,A_1 .
\]
Its mass in \(H<R_0\) equals the mass of the original tensor
in \(H<R_0/\sqrt T\). On these intervals,
\[
  \log(R_0/\sqrt T)\ \le\ -(8-o(1))m_0x^\eps .
\]
In particular the sufficient bound of
Proposition~\ref{prop:mom150} with fixed \(\eta\) becomes
\(A_1\le(R_0/\sqrt T)e^{-\eta x^\eps}\) in unit frequency.
\end{proposition}
\begin{proof}
Use Proposition \ref{prop:scale151}
on each real factor. This transports the Bessel functional by
left translation; for fixed input vectors any other nonzero
choice of the functional differs by a scalar by
Theorem~\ref{thm:src109}. Normalizing a nonzero factor in
\(L^2\) removes the absolute value of that scalar. Since
\(H^{2D_0}=\prod_\tau h_\tau^{4rs}\), product measure and
Tonelli's theorem give the asserted moment, also if it is
infinite. The same change of variables gives the mass formula
for the indicator. Both \(g_i\) lie between \((kb)^8\) and
\(2(kb)^8\) by Proposition~\ref{prop:index90}, so
\(\log R_0=(8+o(1))x^\eps\).
The proof of Theorem~\ref{thm:common111} gives
\(\log T\ge(16-o(1))m_0 x^\eps\), and \(m_0\to\infty\).
These inequalities prove the stated bound.
\end{proof}

Use Euclidean distance in the signed coordinates, as for \(H\),
and put
\[
 H_{\rm end}(g)=
 \left(\prod_\tau
  \operatorname{dist}(g_\tau^{-1}u_s,g_\tau^{-1}L_{s-1})\,
  \operatorname{dist}(g_\tau^{-1}e,g_\tau^{-1}L)\right)^{1/f}.
\]
These two distances concern the last simple coordinate of the
unipotent character. Let \(\Pi\) be orthogonal projection on
automorphic \(L^2\) onto the isotypic space of the prescribed
real row \(\sigma\) of \eqref{eq:src109}, followed by projection
onto its lowest compact type. These projections commute.

\begin{theorem}[A bound for the projection at the terminal prime]
\label{thm:project152}
There is an absolute \(c>0\). For all sufficiently large \(x\)
in the common intervals, \(0<U\le T^{1/4}\), and every seed
\(\Phi\) of frequency \(\varpi\) supported in
\[
       H_{\rm end}\le U,\qquad H<R_0\exp(T^{1/8}),
\]
one has
\begin{equation}
       \|\Pi P\Phi\|\ \le\
       \exp\!\left(-\frac{c\sqrt T}{r^2 U}\right)\|\Phi\|_{\rm seed}.
       \label{eq:project152}
\end{equation}
Without an upper restriction on \(H\), every seed supported in
\(H_{\rm end}\le U\) in the weighted domain of
Proposition~\ref{prop:mom150} satisfies
\[
 \|\Pi P\Phi\|\le
       \exp\!\left(-\frac{c\sqrt T}{r^2 U}\right)
       \left(\|\Phi^{<}\|_{\rm seed}
        +C_*\big\|\mathbf1_{H\ge R_0}(H/R_0)^{2D_0}\Phi\big\|_{\rm seed}
       \right),\qquad C_*=\frac{4}{\sqrt3}\,17^{rs}.
\]
The constants are uniform in the selected words and finite
class columns.
\end{theorem}
\begin{proof}
In Theorem \ref{thm:project151}, the compact Lie and Casimir bound
for this row is \(L_*\le\exp(C_\psi x^\eps)\) with a fixed
constant, as computed in the proof of Theorem~\ref{thm:raw97}.
All ranks and \(f\) are polynomially bounded in \(x\).
Together with \(\log T\ge(16-o(1))m_0x^\eps\), this shows
that \(r^2(L_*+f+1)=o(T^{1/4})\).
Let \(C_1\ge1\) exceed the absolute constant in Theorem \ref{thm:project151}, and take
\[
          j=\left\lfloor
                    \frac{\sqrt T}{8C_1 f r^2 U}\right\rfloor .
\]
Uniformly for \(U\le T^{1/4}\) and large \(x\), this integer
is at least \(\sqrt T/(16C_1 f r^2 U)\) and the base
\(C_1 r^2(L_*+fj+1)U/\sqrt T\) is at most one half.
Thus the projected operator at height \(R\) is bounded by
\[
           B(R)\exp\!\left(-\frac{c_1\sqrt T}{r^2 U}\right)
\]
with absolute \(c_1>0\), using \(B(R)\) of that theorem.
For \(R_0<R\le R_0\exp(T^{1/8})\),
\(\log B(R)\le rs\log17+2frs\,T^{1/8}
=o(T^{1/4}/r^2)\), again because the ranks are polynomial.
For \(R\le R_0\) the prefactor is one. This proves
\eqref{eq:project152} after decreasing \(c_1\).

For the unbounded assertion split off \(H<R_0\). On each
of the upper shells in Theorem \ref{thm:tail149}, the same
constant \(\exp(-c_1\sqrt T/(r^2 U))\) multiplies the
operator bound \(B(R)\) at its upper endpoint. The
Cauchy--Schwarz summation in the proof of that theorem with
\(\sigma=D_0\) gives exactly the second term displayed here.
The projection is bounded, so the sum can be taken in
automorphic \(L^2\). In particular this reasoning uses no
assumption on an untruncated \(L^2\) realization of a
particular local Bessel functional.
\end{proof}

\begin{corollary}[The projection of the compact row copies and their averages]
\label{cor:compact152}
For the large terminal frequency the row copies in
Theorem~\ref{thm:copies137} can be chosen so that, for every
initial seed \(\Phi\) in their span and every compactly
conjugated Gaussian average \(W\) of Corollary~\ref{cor:gauss150},
\[
     \max\{\|\Pi P\Phi\|,\|\Pi PW\|\}
             \le \exp(-c'\sqrt T/r^2)\|\Phi\|_{\rm seed}
\]
on the same fields, with an absolute \(c'>0\).
\end{corollary}
\begin{proof}
On the maximal compact group, the indicated Euclidean
distances are \(\sqrt2\) and one at each place, so
\(H_{\rm end}=H=\sqrt2\). Both functions are continuous
and right invariant under that compact group. Choose the bump
in Theorem~\ref{thm:copies137} to be one near the compact
orbit of the identity and supported where both functions are
less than two. This retains the nonzero row copies by the
same local nonvanishing and Schur orthogonality argument.
Apply \eqref{eq:project152} with \(U=2\). The convolution on
the automorphic space in Corollary~\ref{cor:gauss150} is a
contraction. It commutes with the \(\sigma\)-isotypic projection
since that subspace is reducing for real right translations.
Its compact conjugation average commutes also with the
compact-type projector. The exact identity
\(PW=R_G(\mu)P\Phi\) therefore gives the estimate for \(W\).
\end{proof}

Theorem~\ref{thm:project152} bounds projection onto the
required real representation of the automorphic group.
For the integral theta functional one still needs the
contribution from larger terminal coordinates and its
normalization on the specified lattice. The cohomology
estimate \eqref{eq:invar148} on these fields still leaves
the remaining summands in \eqref{eq:tail114}.

\subsection{Which part of the real period can carry the discrete series}\label{nav:VII-24}
Continue on the fields and intervals of
Theorems~\ref{thm:invar148} and~\ref{thm:project152}. Use
their prime \(T\), \(m_0=2(s+1)\), \(R_0=\sqrt{g_1g_2}\),
and \(b=\exp(x^\eps)\). At each real place let \(a_{i,\tau}\)
be the successive Euclidean Gram--Schmidt lengths of the list
\[
 g_\tau^{-1}u_1,\ldots,g_\tau^{-1}u_s,\quad
 g_\tau^{-1}(\varpi_\tau e).
\]
Set \(h_i=(\prod_\tau a_{i,\tau})^{1/f}\) and let \(M\) be
the geometric mean of the local maximum consecutive ratios
\(a_{i+1,\tau}/a_{i,\tau}\). The global height \(H\) uses
only the first \(s\) entries. Let \(\Pi\) denote the same
automorphic discrete-series and lowest-compact-type projection
as in Theorem~\ref{thm:project152}.

\begin{theorem}[The projection below the full height cutoff]
\label{thm:alllow154}
At frequency \(\varpi\), for every \(0<R\le R_0\) and every
seed supported in \(H<R\),
\begin{equation}
        \|\Pi P\Phi\|\le
          \exp\!\left(-\frac{c\sqrt T}{r^2R}\right)\|\Phi\|_{\rm seed}
         \label{eq:alllow154}
\end{equation}
with absolute \(c>0\), for all sufficiently large \(x\).
There is no condition on the negative vector modulo the
isotropic space. In particular this applies to the entire
lower part of any square-integrable section, including each
prescribed real Bessel section for which that part is in the
seed Hilbert space.
\end{theorem}
\begin{proof}
By Corollary \ref{cor:ang153},
\(M\ge\sqrt T/H\). Apply that corollary with
\(V=\sqrt T/R\). Its factor \(B(R)\) is exactly one.
The real row parameter \(L_*\), the dimensions, and
\(R_0\le2(kb)^8\) satisfy
\(r^2(L_*+f+1)R_0=o(\sqrt T)\), by the size estimates in
Theorem~\ref{thm:project152}. Choose
\(j=\lfloor V/(8C_1 f r^2)\rfloor\) with \(C_1\) at least
the absolute constant in Corollary \ref{cor:ang153}.
Then \(fj\ge V/(16C_1 r^2)\) and
\(C_1r^2(L_*+fj+1)/V\le1/2\). This proves
\eqref{eq:alllow154} uniformly in \(R\). The calculation uses
the dual vector to the final quotient line at each place.
Its norm decreases when the distance of \(g_\tau^{-1}e\)
from \(g_\tau^{-1}L\) increases.
\end{proof}

\begin{theorem}[Reduction to the region of large isotropic vectors]
\label{thm:cone154}
Let \(\mathcal D=\{M<b^2\}\). For any seed supported in
\(H<R_0\exp(b^{3/2})\),
\[
  \|\Pi P\Phi-\Pi P(\mathbf1_{\mathcal D}\Phi)\|
               \le e^{-c b^2/r^2}\|\Phi\|_{\rm seed}
\]
with an absolute \(c>0\). In \(\mathcal D\) all \(s\)
isotropic lengths satisfy
\begin{equation}
        h_i>\sqrt T\, b^{-2(s-i+1)}
                 \quad(1\le i\le s).
        \label{eq:cone154}
\end{equation}
The same error bound holds for seeds in the weighted domain
without a height cutoff if \(\|\Phi\|_{\rm seed}\) on the right
is replaced by
\[
 \|\Phi^{<}\|_{\rm seed}
  +C_*\big\|\mathbf1_{H\ge R_0}(H/R_0)^{2D_0}\Phi\big\|_{\rm seed}.
\]
\end{theorem}
\begin{proof}
We need the dependence on the row weights, since \(b^2\) is
now the ratio threshold. In \eqref{eq:src109} the \(v_i\)
are polynomially bounded and \(|b_{j,\tau}|\le2b\) for all
sufficiently large \(x\), by the construction preceding
Lemma~\ref{lem:spread111}. The compact highest weights differ
from these parameters by at most \(r\) in each entry.
The formulas in the proof of Theorem~\ref{thm:raw97} therefore
give \(L_*\ll\sqrt r(2b+r)\): the compact Lie norm is bounded
by \(\sqrt r\) times the largest absolute compact weight, and
the Casimir is the difference of the two sums of squares for
\(\lambda_\tau\) and \(\rho\). The ranks and \(f\) are
polynomial in \(x\). Hence \(r^2(L_*+f+1)=b^{1+o(1)}\) as
an upper bound. Here and below \(\eps>0\) is fixed.

Apply Corollary \ref{cor:ang153}, to the part with \(M\ge b^2\).
With \(V=b^2\), take \(j=\lfloor V/(8C_1fr^2)\rfloor\)
as above. After optimizing the derivative order this gives
\(B(R)e^{-c_1 b^2/r^2}\) at every bounded height \(R\).
When \(R\le R_0\exp(b^{3/2})\), the logarithm of \(B(R)\)
is at most \(rs\log17+2frs\,b^{3/2}=o(b^2/r^2)\).
Absorbing this factor proves the first bound. Summing the
corresponding estimates on all shells with the weights of
Theorem \ref{thm:tail149}, exactly as in
Theorem~\ref{thm:project152}, proves the assertion for
unbounded supports. Lastly, on \(\mathcal D\) use
\(M^{s-i+1}h_i\ge\sqrt T\) from Corollary \ref{cor:ang153}.
This proves \eqref{eq:cone154} simultaneously for every
index, without counting a union over the places.
\end{proof}

\begin{corollary}[Separation from the compact real seeds]
\label{cor:far154}
On \(\mathcal D\) one has
\[
          \log h_1\ge(14-o(1))s x^\eps,\qquad
          \log h_s\ge(8-o(1))m_0x^\eps-2x^\eps .
\]
In particular every one of these \(s\) geometric means
eventually exceeds \(R_0\). For the Gaussian averages in
Corollary~\ref{cor:gauss150} at the prime frequency,
an initial seed supported at \(H<R\le R_0\) satisfies
\(\|\Pi R_G(\mu)P\Phi\|\le
e^{-c\sqrt T/(r^2R)}\|\Phi\|_{\rm seed}\).
\end{corollary}
\begin{proof}
Use \(\log T\ge(16-o(1))m_0x^\eps\) and \(m_0=2(s+1)\)
in \eqref{eq:cone154}. The least of its lower bounds occurs
at \(i=1\); its logarithm is at least
\((8-o(1))\,2(s+1)x^\eps-2s x^\eps\).
Meanwhile \(\log R_0=(8+o(1))x^\eps\) and \(s\to\infty\).
The Gaussian convolution on the automorphic space is a
contraction commuting with both projections, as in
Corollary~\ref{cor:compact152}. Apply Theorem~\ref{thm:alllow154}.
\end{proof}

The lower-height isometry of Theorem~\ref{thm:low136} and
\eqref{eq:alllow154} now apply to exactly the same restricted
seeds. The first computes their full norm and the second
bounds the portion in the required real row. The period on
\(\mathcal D\), its integral scale, and the outstanding strict
groups in \eqref{eq:tail114} still enter the prime-sum estimate
\eqref{eq:window}.

\subsection{Where the normalized real Bessel functions have their mass}\label{nav:VII-25}
For a selected word and translate, let
\(J_\varpi:\bigotimes_\tau\sigma_\tau\to
L^2(B_\infty\backslash G_\infty,E_\rho)\) be the isometric
embedding of Corollary \ref{cor:rel155}, with terminal frequency \(\varpi_\tau\) at
every real place. This space consists of sections with the
specified Bessel character and compact column type. Its
measure is the product of the compatible quotient measures
in that corollary. We write \(J_1\) for the corresponding
embedding at frequency one. Use the same fields, intervals
and prime \(T\) as in Theorem~\ref{thm:invar148} and
Theorem~\ref{thm:cone154}.

\begin{theorem}[Simultaneous concentration at the real places]
\label{thm:localmass156}
Let \(v\) be any unit vector in the product lowest compact
type of the row and let \(W=J_\varpi v\). Set \(V_x=b r^4\).
For the individual lengths and ratios of
Theorem~\ref{thm:cone154}, put
\[
    \Omega_x=\bigcap_{\tau}
       \left\{\max_{1\le i\le s}\frac{a_{i+1,\tau}}{a_{i,\tau}}
                                                     <V_x\right\}.
\]
For an absolute \(c>0\) and all sufficiently large \(x\),
\begin{equation}
   \int_{\Omega_x}\|W(g)\|^2\,d\dot g\ \ge\ 1-e^{-c b r^2}.
   \label{eq:mass156}
\end{equation}
At every point of this region, simultaneously at every real
place and every \(1\le i\le s\),
\[
        a_{i,\tau}>\sqrt T\,V_x^{-(s-i+1)} .
\]
In particular \(\log a_{1,\tau}\ge(15-o(1))s x^\eps\) there.
All these assertions concern the entire real quotient.
\end{theorem}
\begin{proof}
Fix a place \(\tau\). The isometry of \(J_\varpi\) transfers
the norm of every right Lie derivative of \(W\) to the row
Hilbert space. Use the pointwise inequality of Theorem \ref{thm:energy153} at just this place, leaving the other coordinates
as parameters. There are \(r^{2j}\) words of order \(j\).
By Lemma \ref{lem:cas96} each derivative on the lowest type
has norm at most \(\{Cr(L_*+j+1)\}^j\). This bound also holds
after tensoring with the identity at the other places: expand
the vector in an orthonormal basis there and sum squares.
With \(M_\tau=\max_i a_{i+1,\tau}/a_{i,\tau}\) it follows that
\[
        \|\mathbf1_{M_\tau\ge V_x}W\|_2
         \ \le\ \left(\frac{C_1r^2(L_*+j+1)}{V_x}\right)^j .
\]
Here \(L_*\ll\sqrt r(2b+r)\) by the proof of
Theorem~\ref{thm:cone154}. Take \(j=\lfloor V_x/(8C_1r^2)\rfloor\),
enlarging \(C_1\) to be at least one. Since \(r\to\infty\)
and \(r\) is polynomial in \(x\), for all large \(x\) the base
is at most \(1/2\) and \(j\gg b r^2\). The complement of
\(\Omega_x\) is a union of \(f\) such events. Sum their
squared \(L^2\)-norms, using \(f\) polynomial in \(x\) and
\(\|W\|=1\). This proves \eqref{eq:mass156}.

At one place the consecutive ratios telescope from \(a_{i,\tau}\)
to \(a_{s+1,\tau}=\sqrt T\, Z_\tau\ge\sqrt T\).
This proves the individual bounds. Finally
\(\log V_x=x^\eps+4\log r=(1+o(1))x^\eps\), whereas
\(\log\sqrt T\ge(8-o(1))\,2(s+1)x^\eps\). Subtract
\(s\log V_x\) to obtain the claim at \(i=1\).
\end{proof}

\begin{corollary}[A lower bound for the height moment]
\label{cor:mom156}
For every pure product of nonzero lowest-type vectors, normalize
its unit-frequency local Bessel functions to norm one and
define \(A_1\) as in Proposition~\ref{prop:scale152}. Then
\[
       A_1\ \ge\
        \frac{(1-e^{-c b r^2})^{1/(4D_0)}}{b r^4}.
\]
Consequently none of these products satisfies the sufficient
small-moment condition \(A_1\le (R_0/\sqrt T)e^{-\eta x^\eps}\)
of Proposition~\ref{prop:scale152}, for any fixed \(\eta>0\)
and large \(x\) on the prescribed intervals.
\end{corollary}
\begin{proof}
The real functions in the statement are square-integrable and
nonzero by Corollary \ref{cor:rel155}. At each place, changing
the terminal frequency by the isometry of Proposition \ref{prop:scale151} gives another isometric Bessel embedding.
Uniqueness of the functional and equality of norms give the
same embedding up to a unit scalar. Thus the product at the
prime frequency has the distribution in
Proposition~\ref{prop:scale152}. On \(\Omega_x\) one has
\(H\ge h_s>\sqrt T/V_x\) and at least the mass in
\eqref{eq:mass156} lies there. Integrating \(H^{4D_0}\)
against \(\|W\|^2\) gives
\[
   (\sqrt T A_1)^{4D_0}
     \ \ge\ (\sqrt T/V_x)^{4D_0}(1-e^{-c b r^2}).
\]
This remains true for an infinite positive moment.
Finally \(\log R_0=(8+o(1))x^\eps\), while
\(\log\sqrt T\ge(8-o(1))m_0 x^\eps\) with \(m_0\to\infty\).
It follows that \(V_x R_0/\sqrt T\to0\), and the factor
\((1-e^{-c b r^2})^{1/(4D_0)}\) tends to one. These
estimates prove the stated incompatibility.
\end{proof}

The norm in \eqref{eq:mass156} is the real norm with formal
degree normalization. Passing to the global theta period still
requires its finite tests, summation on the automorphic quotient,
and its scale in the prescribed integral lattice. The estimate
here locates the real contribution at each place on the same
fields where Theorem~\ref{thm:invar148} controls the indicated
part of strict cohomology.

\subsection{The complex projected period on the full real quotient}\label{nav:VII-26}
Fix \(x\) on the common intervals of Theorem~\ref{thm:invar148},
one translate, its full finite family of words and the prime
of Theorem~\ref{thm:common111}. Use the actual refined level
\(K^0\) and terminal frequency \(\varpi\).
Write \(D_c\) for the space of finite class covectors of word
\(c\) with the inner product induced by its seed. The common
real row has lowest compact type \(U\); the automorphic copies
of this type inside that discrete series form \(U\otimes\mathcal M\).
Write \(J_U\) for the product Bessel embeddings with the
normalization in Corollary \ref{cor:rel155}, tensored with
their finite class covectors. These embeddings give actual
real vectors to which the projected period can be applied.

\begin{theorem}[The full projected period and its finite Gram matrix]
\label{thm:global158}
For every square-integrable seed in these word spaces the limit
of the projected automorphic sums over height truncations exists
in \(U\otimes\mathcal M\). Call it \(P_U\). Put \(D=\bigoplus_c
D_c\). There is a bounded map \(A:\mathcal M\to D\) at this
field and level, and a positive semidefinite matrix \(G_D=AA^*\)
on \(D\), such that the projected period on normalized real
vectors is \(P_U J_U=1_U\otimes A^*\). For any right
compact-invariant measurable region \(E_0\) in the real period
quotient the Gram operator of \(P_U\mathbf1_{E_0}J_U\) is
\begin{equation}
       1_U\otimes T_{E_0} G_D T_{E_0},\qquad
       T_{E_0}|_{D_c}=t_c(E_0)1_{D_c},
       \quad t_c(E_0)=\int_{E_0}\|(J_cv)(g)\|^2\,d\dot g ,
       \label{eq:gram158}
\end{equation}
where \(v\) is any unit vector of \(U\).
In particular, for each \(c\) using \(E_0\) instead of the
full quotient multiplies every projected vector of that word
by the same scalar \(t_c(E_0)\).
\end{theorem}
\begin{proof}
The level, irreducible column types and fixed row are precisely
the data in Section \ref{nav:VII-08}. Real sections times the finite covectors have the
seed inner product by the identity-term calculation in Theorem \ref{thm:identity127}. Theorem \ref{thm:adjoint157} extends \(\Pi P\) to
this Hilbert space. Boundedness and
\(\mathbf1_{H<R}\Phi\to\Phi\) in the seed norm prove the limit.
Corollary \ref{cor:mass157} identifies its restriction to
\(\mathbf1_{E_0}J_U\) with \(1_U\otimes A^*T_{E_0}\).
Taking the adjoint and composing proves \eqref{eq:gram158}.
The class spaces \(D_c\) are finite dimensional for the fixed
number field and finite level. The formula includes zero
eigenvalues, with their full multiplicities.
\end{proof}

\begin{corollary}[The contribution from the region of large real lengths]
\label{cor:recovery158}
Set \(\delta_x'=e^{-c b r^2}\) with the constant in
Theorem~\ref{thm:localmass156}. For its region \(\Omega_x\),
the norms of the differences between the full and restricted
projected period operators on \(J_U\), and between their Gram
operators, are at most
\[
           \delta_x'\,\|A\|,\qquad 2\delta_x'\,\|G_D\|
\]
respectively, with no upper height restriction. On each
individual word the restricted norm lies between \(1-\delta_x'\)
and one times its full norm for every input, including the
refined zero-sum columns. There also exists a finite common
height \(R_x\) for which the same statements on
\(\Omega_x\cap\{H<R_x\}\) hold with \(2\delta_x'\) in place
of \(\delta_x'\).
\end{corollary}
\begin{proof}
The ratios defining \(\Omega_x\) use Euclidean lengths, so
right translation by the maximal compact group preserves them.
Theorem~\ref{thm:localmass156} gives
\(1-\delta_x'\le t_c(\Omega_x)\le1\) for every word. Thus
\(\|1-T_{\Omega_x}\|\le\delta_x'\). Omitting the tensor
identity on \(U\), the operator difference is \(A^*(1-T)\).
For the Gram difference use
\(G_D-TG_D T=(1-T)G_D+TG_D(1-T)\) and \(\|T\|\le1\).
For a single word the outcome is exactly multiplied by its
scalar \(t_c\). Finally monotone convergence gives
\(t_c(\Omega_x\cap\{H<R\})\uparrow t_c(\Omega_x)\). There
are finitely many words and hence a single finite \(R_x\)
with every error below \(\delta_x'\). This proves the last
assertion, without an estimate for the size of \(R_x\).
\end{proof}

\begin{proposition}[Extraction at the separation height]
\label{prop:cut158}
Put \(t_c^0=t_c(\{H<R_0\})\). For each selected word \(t_c^0>0\).
On \(U\otimes D_c\) the full projected period equals its
truncation at \(R_0\) divided by \(t_c^0\), and
\[
 \bigl\|A^*|_{D_c}\bigr\|
       \ \le\ \frac{\exp[-c_0\sqrt T/(r^2R_0)]}{\sqrt{t_c^0}} .
\]
All scalars here use the actual local Bessel functions and
the original finite theta-period fields.
\end{proposition}
\begin{proof}
The height is \(\sqrt2\) at the identity and \(R_0>\sqrt2\)
eventually. Proposition~\ref{prop:translate109} gives, in
every neighborhood of the identity, a point at which the
Bessel coefficient of a nonzero pure product lowest-type
vector is nonzero. Scaling its local functional to the norm
of Theorem \ref{thm:rel155} preserves this property. Smoothness
therefore gives positive mass in \(H<R_0\). Compact
irreducibility makes this mass independent of the unit vector
of \(U\). By Theorem~\ref{thm:global158} the truncated
operator on this word equals \(t_c^0\) times the full
operator. Its seed input \(\mathbf1_{H<R_0}J_U\) has norm
\(\sqrt{t_c^0}\) as an operator on that summand. Apply
Theorem~\ref{thm:alllow154} and divide by \(t_c^0\).
\end{proof}

Thus at fixed arithmetic data the projected complex period has
a well-defined value for the full square-integrable real
vectors. The matrix \(G_D\) still depends on those arithmetic
data. Its size, its nonzero values on the prescribed integral
theta classes, and their lattice normalization remain to be
controlled as \(x\) varies. These requirements and the
remaining strict orders in \eqref{eq:tail114} enter the
deduction of \eqref{eq:window}.

\subsection{An exact comparison of two projected Bessel periods}\label{nav:VII-27}
Use the fields of Theorem~\ref{thm:invar148} and the prime
\(T=\varpi\bar\varpi\) of Theorem~\ref{thm:common111}.
Let \(a=a_\varpi\) be the rational isometry scaling the two
isotropic strings by \(\varpi,\bar\varpi^{-1}\). Put
\[
   K_0=K^0,\qquad K_1=a_fK_0a_f^{-1},\qquad L=K_0\cap K_1,
       \qquad d_{x,T}=T^{\,fs(r-s)/2},\quad I={r\brack s}_T^f .
\]
The two column levels at \(K_0,K_1\) both give the originally
prescribed refined class covectors. The intersection admits both
frequencies by Proposition \ref{prop:index159}. At \(K_0\) a unit-frequency Bessel
coefficient against a \(K_0\)-invariant row vanishes at a
column argument \(h\in H(\mathbb A_{F_+})\). Indeed Theorem \ref{thm:midlevel131} gives \(n_0\in N_f\cap K_0\) with
\(\chi_1(n_0)\ne1\). In the defining integral over \([N]\)
translate on the right by \(hn_0h^{-1}\). The row value
is fixed and the character is multiplied by the nontrivial
factor, proving vanishing. This is the reason for including
the change of finite level in the comparison.

Let \(\mathcal P_{\varpi,0}\) and \(\mathcal P_{1,1}\) be the
projected period maps at \(K_0,K_1\) on the real vectors
normalized by formal degree, using the same ordinary finite
class coordinates at both levels. Choose phases by transporting
the real vectors as in Proposition~\ref{prop:scale152}.
For a compact real support, define
\(\mathcal Q_{\beta,i}=\Pi e_i P\Phi_{\beta,L}\) on the
same inputs, where \(e_i\) averages over \(K_i\) and the
input specifies the real section and the class tuple. Use
compatible square-integrable real vectors by extension from
compact support as in the theorem below. Every output belongs
to the automorphic space of level \(L\).

\begin{theorem}[A single row isometry for both frequencies]
\label{thm:twofreq160}
For all selected words simultaneously and on their refined
zero-sum columns one has
\begin{align}
  \mathcal P_{\varpi,0}
    &=d_{x,T}^{-1}\,R(a_f^{-1})\,\mathcal P_{1,1},\notag\\
  \mathcal Q_{\varpi,0}
    &= d_{x,T}\,R(a_f^{-1})\,\mathcal Q_{1,1}.
     \label{eq:twofreq160}
\end{align}
The maps in the second equality extend to the prescribed
square-integrable real sections. In the common level-\(L\)
row space one may replace \(R(a_f^{-1})\) on their
finite-dimensional images from the lowest type by a single
unitary automorphism of that space, independent of the word.
With ordinary class coordinates the Gram matrices of the
two common-level maps satisfy
\[
      \mathcal Q_{\varpi,0}^{*}\mathcal Q_{\varpi,0}
           =d_{x,T}^{\,2}\,\mathcal Q_{1,1}^{*}\mathcal Q_{1,1}.
\]
At the respective coarse levels the two Gram matrices agree
when the inputs are given their isometric seed
normalizations.
\end{theorem}
\begin{proof}
Theorem \ref{thm:conj159} proves the first equality, with
\(d_{x,T}=\sqrt J\). Ordinary summation commutes with the
finite right translations for compact seeds and the archimedean
projection commutes with them on automorphic \(L^2\).
Consequently Theorem \ref{thm:common159} gives
\(\mathcal Q_{\varpi,0}=(J/I)\mathcal P_{\varpi,0}\) and
\(\mathcal Q_{1,1}=I^{-1}\mathcal P_{1,1}\).
The coarse-level maps are bounded at each fixed field and
level by Theorem~\ref{thm:global158} and its transport
in Theorem \ref{thm:conj159}. This extends each finite-averaged
map by continuity and proves \eqref{eq:twofreq160} for the
full real sections. Finite translation is unitary and the
same \(a_f\) occurs in every word. For the normalized real
vectors of the lowest type, each image has finite dimension
because both this type and the class spaces are finite.
Inside the sum of the two image spaces, their orthogonal
complements have equal dimension. Extend the isometry between
the images by a unitary between these complements, then by
the identity on the orthogonal complement of the sum in the
common row space. This proves the assertion about a single
automorphism. Taking inner products proves the Gram identity.

Finally the squared class norm of a real unit seed in ordinary
class coordinates is \(c_{\rho,K_i}\|\ell\|^2\). The ratio
\(c_{\rho,K_1}/c_{\rho,K_0}=d_{x,T}^2\) cancels the square
of the first scalar in \eqref{eq:twofreq160} when the inputs
are identified by their seed isometries.
\end{proof}

\begin{corollary}[The arithmetic cost of the common finite level]
\label{cor:integral160}
The scalar \(d_{x,T}\) is a positive rational integer and the
finite averaging maps have denominator dividing the positive
integer \(I\). At every \(x\) in the common intervals,
uniformly for all these words and every fixed \(A>0\),
\[
       \log(d_{x,T}I)=o(b/x^A),\qquad
       1-\frac{f}{T-1}\le
          \frac{c_{\rho,L}}{c_{\rho,K_0}}\le1 .
\]
\end{corollary}
\begin{proof}
The degree \(f\) is even on these intervals, as in the
proof of Theorem~\ref{thm:commonnorm128}. Thus \(d_{x,T}\)
is integral. At each split place the Gaussian binomial counts
subspaces and hence \(I\) is integral. The unaveraged finite
sum \(Ie_i\) has coefficient one on each coset. For its size
use
\[
   T^{s(r-s)}\ \le\ {r\brack s}_T\ \le\ 4T^{s(r-s)} .
\]
Indeed \(\prod_{j=1}^s (1-T^{-j})\ge1/4\): the first factor
is at least \(1/2\) and the remaining product is at least
\(1-\sum_{j\ge2}2^{-j}=1/2\). The factors in the numerator
of the normalized Gaussian binomial are at most one.
It follows that \(\log(d_{x,T} I)\le
\tfrac32 fs(r-s)\log T+f\log 4\). Each quantity on the
right is polynomially bounded in \(x\) by the proof of
Corollary~\ref{cor:cost128}, whereas \(b=\exp(x^\eps)\).
The bound on the class norm is Proposition \ref{prop:index159}.
\end{proof}

These identities apply to the two Bessel projections on the
same finite class lattice, with the stated finite averages.
For the theta lift the source specifies an additional equality
on coefficient spaces,
\[
     B_u(j_{N,u}\phi,j_E x)=\iota\bigl(B_{k,u}(\phi,x)\bigr).
\]
Here \(B_u\) is the complex theta form, \(B_{k,u}\) is the
form over the coefficient field, \(j_{N,u},j_E\) are its
two complex realization maps, and \(\iota\) is the chosen
embedding of the coefficient field. This is the condition at
\srcloc{Formalism/LocalRingCardinality/RhoDetPrimeEqMuTildeMulChiPow}{364}.
The comparison still requires its arithmetic scale, a nonzero
value on the designated theta classes, and the remaining
strict orders in \eqref{eq:tail114} on the same fields.

\section{Unfolding and comparison along the row}\label{nav:g-VII-2}

\subsection{Cuspidality at the original period level}\label{nav:VII-35}
\begin{theorem}[The projected first moment on the matching cuspidal row]
\label{thm:range239}
Choose the unit pure real vector \(u=\bigotimes_\tau u_\tau\)
in the lowest type of the common row. Keep its level \(K^0\) and
frequency \(\varpi\). For each word write \(S_cb=S_0(u\otimes b)\)
in the class norm of Theorem~\ref{thm:dyad234}. In
Theorem~\ref{thm:norm177} choose this real line, terminal depth
zero and the same row character. Include all exceptional places
of the row in its finite set, and take the other principal
subgroups there inside the kernels of the source characters.
Let \(P\) be projection to the whole cuspidal space
\(\mathcal H\) with these data. Then \(PS_c=S_c\) on \(D_c\).

Let \(J\) be either of the terminal spaces of
Theorem~\ref{thm:dyad234}, possibly restricted to the retained
classes of Corollary~\ref{cor:trace235}, and suppose \(d_J=\dim J>0\).
Use the random signed column of the corresponding first-moment
formula. For the matching complex theta operator \(\Theta\) put
\(\alpha=C_xP(\mathbf a)>0\) as in Theorem~\ref{thm:norm177}.
Then, writing \(O_{\gamma,c}\) for the nonparabolic orbital operators,
\[
 \begin{split}
 M_J&=1+d_J^{-1}\sum_\gamma
             \operatorname{Re}\operatorname{tr}(P_J O_{\gamma,c})
       =\mathbb E\|S_cb\|^2,\\
 e^{-\Delta_x}M_J&\le
      \alpha^{-1}\mathbb E\|\Theta S_cb\|^2\le e^{\Delta_x}M_J .
 \end{split}
\]
Some signed column in this distribution has nonzero forward
lift if and only if the real orbital sum is strictly greater than
\(-d_J\).
\end{theorem}
\begin{proof}
The ordinary periodization of Theorem~\ref{thm:ordinary163}
is square integrable. Right convolution by the real matrix unit
at \(u\) fixes it. On each archimedean copy of the row this matrix
unit projects to the chosen line. The cuspidality argument of
Theorem~\ref{thm:global167} applies to every finite class vector
at depth zero. In detail, a nonzero vector fixed by this
matrix unit generates a copy of the tempered discrete series.
On an arithmetic component the determinant of the level
stabilizer consists of norm-one units in \(E\), hence of roots
of unity. Thus a finite-index subgroup lies in the special
unitary group. The central real torus is compact, and after
pullback along the finite central isogeny one can take its
fixed character and apply Wallach's constant-term theorem to
the special unitary arithmetic quotient. Its proper parabolic
constant terms vanish. Passing back through the finite cover
preserves this vanishing, as does the passage to the adelic
constant terms described in Theorem~\ref{thm:global167}.
\href{https://mathweb.ucsd.edu/~nwallach/tempered-cuspidal.pdf}{Wallach,
Theorem 4.3}.

The finite right transformations of the seed pass through its
absolutely convergent rational sum. At \(v\mid T\) the first
test uses \(L_0=K^0_v\) with exactly this row character.
Away from \(T\) the seed is fixed by the chosen principal
subgroups at exceptional places and by the spherical compacts
elsewhere. These properties put every \(S_cb\) in \(\mathcal H\).
Taking a terminal character or a difference on the column changes
\(b\) inside the same domain. The row transformation law therefore
continues to hold for the signed columns and for their linear
combinations supported on retained classes.

The projection defect in Theorem~\ref{thm:dyad234} now vanishes.
Its averaging formula and absolute convergence give the first
identity, also on the retained classes by Corollary~\ref{cor:trace235}.
The norm comparison on \(\mathcal H\) gives the two inequalities
after averaging, since the terminal index at depth zero is one.
Each value being averaged is nonnegative, and the distribution
has finite support. This proves the last assertion.
\end{proof}

\subsection{The specified compact test and its Petersson norm}\label{nav:VII-36}
On the fields of Theorem~\ref{thm:column242} take a standard
definite pair of ranks \(r<s\), \(rs\) even. Pair its lift with
the full target or restrict it to the standard hyperplane. Denote
the group of integration on the target by \(H_*\), of rank
\(q=s\) or \(s-1\) respectively. In either case use the space
\(V\), polynomial \(Q\), and integer \(H\) of
Lemma~\ref{lem:branch243}. Write \(\mathcal I_V\) for its ordered
minor indices, \(k=r\sum_\tau b_\tau\), and \(\rho\) for the
target action on \(V\), including its vacuum character. On the
source compact group all these products transform by one character
\(\eta\). Indeed right multiplication of a full-column minor by a
column change multiplies it by the determinant of that change,
and the source vacuum character is fixed.

Keep the source function \(a\), its finite level \(J\) of
character \(\chi\), and the target level \(J_*\) of character
\(\chi_*\). Use the finite test \(\phi_{uv}\) of
Theorem~\ref{thm:edge242}, with \(v,j_*\) embedded in the full
target. In particular the average defining it includes both
\(\chi\) and \(\overline{\chi_*}\). Write
\(\kappa=\operatorname{vol}(U_{r,\infty}\times J)\),
\(\kappa_*=\operatorname{vol}(H_{*,\infty}\times J_*)\), and
\(\gamma_u,\gamma_v'\) for the two finite stabilizer orders.
With real Haar probability set
\[
 a_u=\int a(h,u)\eta(h)\,dh,\qquad
 L_v(P)=\sum_u\frac{a_u}{\gamma_u}\,
                          P(\partial)G_{\phi_{uv}}(0),\qquad
 \mathcal N_V=\sum_v\frac{1}{\gamma'_v}
                         \sum_{T\in\mathcal I_V}|L_v(P_T)|^2 .
\]
Use \(L(P)\) for the actual theta lift of the input with polynomial
\(P\) in the dimensionless coordinates of Theorem~\ref{thm:vac242},
restricted to \(H_*\) and projected to the stated finite character.
The Petersson pairing below is linear in its first argument.

\begin{theorem}[The period against a prescribed compact coefficient]
\label{thm:fixed243}
Let \(b\) be a target form of character \(\chi_*\) with
\(b(h,v)=B_v(\rho(h)Q)\) at its finite representatives, for
specified linear functionals \(B_v\) on \(V\). Then
\begin{align*}
 \frac{d_V H\,c^{k/2}}{\kappa\kappa_* C}
       \langle L(Q),b\rangle
   &=\sum_v\frac{1}{\gamma'_v}\sum_{T\in\mathcal I_V}
                       L_v(P_T)\overline{B_v(P_T)},\\
 \frac{d_V H\,\nu^D c^k}{\kappa^2\kappa_*}\,
                   \|L(Q)\|^2&=\mathcal N_V .
\end{align*}
Here \(C,c,\nu,D\) are as in Theorem~\ref{thm:vac242}.
In the full-target case the arithmetic quantity on the right of
the second identity is exactly \eqref{eq:pair55} when all
\(a_u=1\) and its \(F_{uv}\) are the actual finite sections.
In either case any nonzero vector in \(V\) gives a nonzero
projected lift if and only if \(\mathcal N_V>0\).
A signed sum with a common \(Q\), splitting and degree is
included by taking the signed sum of its \(L_v\)'s before the
two displayed operations.
\end{theorem}
\begin{proof}
The source compact action by \(\eta\) and the finite projections
give the following equality at every compact target argument,
by the class calculation of Theorem~\ref{thm:edge242}.
\[
       L(P)(h,v)=\kappa Cc^{-k/2}\,L_v(\rho(h)P).
\]
It is valid on the hyperplane because one can evaluate the original
kernel at \(h\oplus1,v\oplus1\) before integrating. The tensor
product of the irreducibles in Lemma~\ref{lem:branch243} is
irreducible for the product of the compact groups. For instance
the commuting operators are scalars by applying Schur's lemma
to each factor on the corresponding matrix blocks.

For any \(P\), the average of the rank-one operators
\((\rho(h)P)(\rho(h)P)^*\) is scalar by irreducibility and
has trace \(\|P\|^2_{\rm F}\). It equals
\(\|P\|^2_{\rm F} I_V/d_V\). On the other hand the average is
\(\|P\|^2_{\rm F}/(d_V C_{\boldsymbol b})\) times
\(\sum_{T\in\mathcal I_V}P_T P_T^*\) by the frame identity.
Evaluate between the two functionals \(L_v,B_v\) at \(P=Q\).
The ratio \(\|Q\|^2_{\rm F}/C_{\boldsymbol b}\) is \(1/H\).
Multiply by the class masses and the first formula follows.
For the norm use \(L_v\) as both functionals and retain
\(\kappa^2|C|^2c^{-k}\) from the pointwise expression. The
identity \(|C|^2=\nu^{-D}\) gives the second formula. For any
other \(P\ne0\) its squared norm on the left changes by the
positive factor \(\|P\|^2_{\rm F}/\|Q\|^2_{\rm F}\).

For a single terminal coefficient \(B_v(\rho(h)Q')\) with a
different \(Q'\), the same integration replaces
\(\|Q\|_{\rm F}^2\) in the first calculation by
\(\sum_j Q_j\overline{Q'_j}\) in a Fischer orthonormal basis.
This follows by polarization. Finite sums give the period against
any specified compact-finite test of this type, with its original
coefficients. Components of inequivalent compact type pair to zero
by Schur's lemma.
\end{proof}

For a general test in this compact isotypic component one can recover
the covector in the first identity without dividing its highest-weight
coordinates. Let \(\theta_Q\) be the character of the real diagonal
torus on \(Q\), including the fixed vacuum twist. Right integration
against \(\overline{\theta_Q}\) on that torus is an orthogonal
projection which fixes \(L(Q)\), by the pointwise formula. At one
place each indexed minor product is a torus eigenvector. To have the
weight of \(Q_l\), every minor must contain all the first \(r-1\)
rows. Its one remaining row must be \(r\) or \(s\), with exactly
\(l\) occurrences of \(s\). This forces the product to equal \(Q_l\).
For the full target all rows must be the first \(r\) and the same
assertion holds with \(A_0^b\). Thus the torus projector takes each
ordered product either to zero or to the distinguished product with
coefficient one, also after tensoring over places. Consequently a
class function written as an integral sum of matrix coefficients
\(B_i(\rho(h)R_i)\), with inputs in the module generated by the
\(P_T\)'s and with covectors in its integral dual, yields exactly
\(B_v(\rho(h)Q)\) with \(B_v\) still integral on that module.
Components of other compact types have zero pairing with this lift.
On the source an irreducible compact type contributes to \(a_u\)
only if it is the inverse of the one-dimensional character \(\eta\).
This follows by averaging its representation after twisting by \(\eta\).

\subsection{Character matching in the first seesaw}\label{nav:VII-37}
Fix the real tests in Theorem~\ref{thm:seesaw245}. The remaining
finite period is a linear functional \(\ell\) on
\(D\otimes\overline H\), where \(D=\mathcal V_Y\) and
\(H=\mathcal V_{-L}\). Use their finite oscillator Hilbert norms.
Let \(J_S,J_Y\) be the given fixing levels of the two forms, acting
on their integral Hermitian lattices. The level on \(Y\) acts
through its extension by the identity on \(L\). These actions
preserve the depth conditions: they preserve the rank-one
generators and integral homomorphisms in Proposition~\ref{prop:depth245},
commute with the selected idempotents, and conjugate a Weyl operator
to the operator of the transformed displacement. Thus \(J_S\) acts
on both \(D,H\), while \(J_Y\) acts on \(D\) and commutes with
\(J_S\). Each compact group acts through a finite quotient, since
the pointwise stabilizer of a basis in a smooth finite-dimensional
representation is open. Include the kernels of the level characters
in choosing these quotients. Set
\[
        \alpha=(\chi_S\delta_S)^{-1}\big|_{J_S},\qquad
        \beta=(\chi_Y\delta_Y)^{-1}\big|_{J_Y}.
\]
The splitting characters appear here because the finite actions
on the two small spaces are their own oscillator actions.

\begin{proposition}[The finite components detected by the period]
\label{prop:characters245}
Let \(D_\beta\) be the \(\beta\)-space for \(J_Y\). For each
irreducible unitary representation \(\pi\) of the finite quotient of
\(J_S\), let \(d_\pi=\dim\pi\), and let \(m_\pi,n_\pi\) be the
multiplicities of \(\alpha\otimes\pi\) in \(D_\beta\) and of \(\pi\)
in \(H\), respectively. Choose orthonormal bases \(p_{ia}\) in the
former copies and \(q_{ib}\) in the latter, identifying their
representation coordinates by the character twist. Here \(i\) ranges
from one to \(d_\pi\), while \(a,b\) index copies. Put
\[
          c_{\pi ab}=\ell(p_{1a}\otimes\overline{q_{1b}}).
\]
The only nonzero entries in these adapted bases can occur within
the indicated pairs of types, and within such a pair they are
\(\ell(p_{ia}\otimes\overline{q_{jb}})=\delta_{ij}c_{\pi ab}\).
In particular
\[
 \|\ell\|^2=\sum_\pi d_\pi\sum_{a,b}|c_{\pi ab}|^2.
\]
Let \(e\) denote the joint character projector on these finite
spaces. For any finite test \(z\) with \(\ell(z)\ne0\), one of the
matched individual pairs satisfies
\[
    |c_{\pi ab}|\ \ge\
       \frac{|\ell(z)|}{\|ez\|\,(\sum_\pi d_\pi m_\pi n_\pi)^{1/2}} .
\]
There is also a bound for the integral product basis \(b_i,\bar b'_j\)
of Proposition~\ref{prop:depth245}. Write \(w_i=\|b_i\|^2\) and
\(w'_j=\|b'_j\|^2\). Then
\[
 \max_{i,j}|\ell(e(b_i\otimes\bar b'_j))|
 \ \ge\ \frac{|\ell(z)|}
      {\|ez\|\,(\sum_{i,j}(w_iw'_j)^{-1})^{1/2}} .
\]
Each weight is the finite position volume of the corresponding
surviving orbit of cosets. Both inequalities concern the same
forms and the same finite tests as the original integral.
\end{proposition}
\begin{proof}
By change of variables, \(\ell\) transforms under the two small
actions by \(\alpha\) on the common group and \(\beta\) on \(J_Y\).
Consequently it kills the orthogonal complement of \(D_\beta\)
in its first variable. On any pair of irreducible copies it
gives a sesquilinear form \(B(p,q)=\ell(p\otimes\bar q)\).
Write \(B(p,q)=\langle Cp,q\rangle\), with first variable linear.
Its transformation law says
\(C\rho(j)=\alpha(j)\pi(j)C\) on a pair \(\rho,\pi\).
Schur's lemma makes this operator zero unless
\(\rho\simeq\alpha\otimes\pi\), and scalar under the specified
identification in that case. This proves the assertion about
entries. Summing their squared moduli in the orthonormal product
basis gives the squared norm of the functional. Since
\(\ell(z)=\ell(ez)\), Cauchy--Schwarz gives the first lower bound.

The integral orbit bases are orthogonal because their supports
are disjoint, and their entries on those supports have modulus one.
Thus \(b_i\otimes\bar b'_j/\sqrt{w_iw'_j}\) form another
orthonormal basis. Evaluate the same squared norm on it and bound
every unnormalized value by their maximum. Orthogonal projection
leaves the value of the functional unchanged by
Theorem~\ref{thm:seesaw245}. This proves the second bound.

If a row form is itself a signed sum with its stipulated level
character, the argument applies to that sum as a single input.
The coefficients \(c_{\pi ab}\) are then computed after taking
the signs. To apply the lower bounds to the desired arithmetic
comparison one still needs the size of an actual nonzero
\(\ell(z)\) and the normalization of the iterated integral map.
\end{proof}

\subsection{Unfolding the Bessel period on a full frame}\label{nav:VII-43}
Keep \(S=C\perp\mathbb H^r\), \(T=(D\perp L_n)\perp\mathbb H^r\),
\(X=X_*\), and the line conventions of
Proposition~\ref{prop:linkedline265}. Write \(N_S,N_T\) for the
unipotent radicals of the ordered isotropic flags and \(G=U(S)\).
Use brackets for rational adelic quotients. In the scaled
Witt coordinates an element of \(N_T\) has an upper unitriangular
Levi part \(v\), a block \(b:T_0\to E^r\) in the part fixing
the frame, and a skew-Hermitian central coordinate. In that
order of factors \(n_T=n_B(b)n_Z(z)m(v)\). The Bessel character
is the product of \(\chi(v)\) and \(\psi_E(e_r^{\,t}b\zeta_n)\).
Write \(v(n)\) for the analogous upper block of \(n\in N_S\).
The scalar \(\chi\) is as in Proposition~\ref{prop:flagunmarked264}.

Let \(A_D,B\) be compact-finite cusp forms on \(U(D),G\).
Use a horizontal Schwartz test with finite part as in
Theorem~\ref{thm:flagfinite264}, allowing finite linear
combinations in its residual variable. Write
\(\mathcal E(g;\Phi)\) for the adelic integral over
\(U_r(\mathbb A_E)\) in the two positions, evaluated at
\((X,p_*)\) as in that theorem. On decomposable data this is
the product of the real and finite integrals. Take compatible
Haar measures, giving the rational additive and unipotent
quotients volume one. The real factors of these measures may
absorb the constants from the finite conventions of
Theorem~\ref{thm:flagfinite264}.

\begin{proposition}[An orbital formula for the global Bessel period]\leavevmode
\label{prop:besorbit265}
The Bessel period on \(U(T)\) of the plain theta lift of \(B\),
paired with \(A_D\), is
\[
 P_T(\Phi)=\int_{R(\mathbb A_F)\backslash G(\mathbb A_F)}
           \mathcal J(g;\mathcal E(g;\Phi))\,d\dot g ,
 \qquad R=N_S\rtimes U(C).
\]
Here \(\mathcal J(g;f)\) is the integral on
\([U(D)]\), \([U(C)]\), and \([N_S]\) of
\[
 A_D(h)B(nk g)\,\overline{\chi(v(n))}\,
      \Theta^0\big(\mathcal W_n\,\omega^0(h,k)f\big).
\]
In this expression \(\Theta^0\) is the rational theta sum on
\(\operatorname{Hom}(-T_0,C)\), \(\omega^0\) is its residual
action, and \(\mathcal W_n\) is the displacement and central
phase on the new line of Proposition~\ref{prop:linkedline265}.
It is the identity on the \(D\)-part. The functional is jointly
continuous in \(g\) and the real Schwartz vector at fixed finite
level. In particular, if \(f\) is the product of a descending
core test and the line test transferred from \(q\), then
\(\mathcal J(1;f)\) is the old Jacobi period with test \(q\),
paired with the descent of \(A_D\).
\end{proposition}
\begin{proof}
Start by integrating \(A_D(h)B(g)\) times the theta sum at
\((n_T h,g)\) over \([U(D)]\), \([G]\), and \([N_T]\), with
the conjugate Bessel character. In the first polarization the
summation positions are all rational \(r\)-tuples \(x\) in
\(S\). On this fiber a shear in the skew-Hermitian center
acts by the character of the Hermitian matrix \(x^*h_Sx\).
Indeed its phase is half the alternating trace pairing of
the retained columns with their product by the shear.
The pairing of Hermitian and skew-Hermitian matrices obtained
by taking the trace after multiplication by \(\xi_0\) is
nondegenerate. Diagonal entries are detected by diagonal
matrices, and a pair of off-diagonal entries is detected by
varying an arbitrary element of \(E\) and its negative
conjugate in the transposed positions. Integration over the
rational central quotient therefore retains exactly
\(x^*h_Sx=0\). Every retained frame has full rank. At a single
marked place the first-position box has independent columns
modulo the maximal ideal in normalized lattice bases.
The actions on these columns of \(g\) and of the upper
unitriangular factor preserve rank. The other unipotent
factors act on the fiber. A singular rational frame can
therefore give no term.

Witt extension over \(E/F\) is transitive on the full ordered
isotropic frames. It also extends each isometry \(X\mapsto Xv\)
for \(v\in U_r(E)\). The stabilizer up to these transformations
is \(R\). Indeed preserving \(X\) up to \(v\) fixes the ordered flag,
acts trivially on its successive one-dimensional quotients,
and permits an arbitrary isometry of \(C\) in the remaining
Levi factor. Thus representatives of the frames are
\(\gamma^{-1}Xv_0\), with \(\gamma\in R(F)\backslash G(F)\)
and \(v_0\in U_r(E)\), after choosing representatives for the
first quotient. The latter coordinate is unique for each such
choice because \(X\) is injective. Unfolding \(\gamma\) gives
\(R(F)\backslash G(\mathbb A_F)\). For \(v_0\), change
\(b\mapsto v_0 b\) in the noncentral block and conjugate the
central matrix correspondingly. These rational substitutions
preserve the quotients and their measures, and the last row
of \(v_0\) fixes the Bessel character. The remaining argument
of the first position is \(Xv_0v\). The sum over \(v_0\) and
the integral on \([U_r]\) give integration over the entire
adelic upper unitriangular group with weight
\(\overline{\chi(v)}\). The multiplication law in the radical
permits this order of integration because the central quotient
has already been integrated and the upper group preserves
the product additive measure of the block \(b\).

On the fiber of \(X\), the remaining action of \(b\) adds
\(Xb\) to the residual map into the isotropic columns of \(S\).
In the second polarization this is a modulation. Write \(y\)
for the scaled opposite position. Its physical row is obtained
by taking \(-\xi_0^{-1}y^*h_{T_0}\) in the dual frame.
Its character under that displacement is
\(\psi_E(-\operatorname{tr}(yb))\). In fact pairing the
isotropic displacement with the physical opposite position
cancels \(\xi_0\) against the conjugate of \(-\xi_0^{-1}\),
and Fourier evaluation uses the negative of this trace. Integrating
in \(b\) against the conjugate Bessel character therefore
selects \(y=-\zeta_n e_r^{\,t}\), by ordinary character
orthogonality on the rational additive quotient. Its unscaled
row is just \(p_*\) of Proposition~\ref{prop:positions264}.
The remaining rational sum is \(\Theta^0\). This establishes
the two-position integral before disintegration.
\srcloc{Development/AbelianVarietyTheory/Definitions/TwistedInnerProducts/WittDecompSymplInclIsotropicInclAnisotropic}{268--287}.

Disintegrate \(R(F)\backslash G(\mathbb A_F)\) by its \(R\)-orbits.
The group \(U(C)\) preserves Haar measure on \(N_S\). Indeed on
the free block its determinant over \(F\) is a power of the
norm of a unitary determinant, which equals one. It fixes the
central and upper matrix coordinates. Thus \(R\) is unimodular
and the \(R\)-fiber has the product integral on \([N_S]\) and
\([U(C)]\). A left translation by \(nk\in R\) sends the first
position \(Xu\) to \(Xv(n)^{-1}u\). The change of variable
in \(u\) supplies \(\overline{\chi(v(n))}\). On the opposite
position the last coordinate is fixed. Its residual action is
precisely \(\mathcal W_n\omega^0(1,k)\) by the block and
central calculations of Proposition~\ref{prop:linkedline265}.
The \(h\)-action fixes the new line and commutes with this
displacement. The split core characters are \(\gamma_i\) on both
sides, and the remaining character on \(U(C)\) is the line factor
computed in that proposition. This gives the displayed integrand.
Changing the representative of the outer coset just changes
variables on the \(R\)-fiber.

We give the convergence details since \(U(C)\) can be
indefinite. Before unfolding, the absolute rational theta
sums are bounded by a polynomial in the heights of \(g,h\)
on Siegel sets, uniformly on a compact fundamental set for
the unipotent quotient. One obtains the bound by summing a
Schwartz majorant in integral position coordinates, or after
the fixed partial Fourier transformations, as in
Theorem~\ref{thm:seesaw245}. Rapid decrease of \(A_D,B\)
then gives absolute convergence for unfolding. Restricting
the rational position sum by additive orthogonality preserves
this majorant. In unfolding the upper group one unfolds the
sum of the absolute values before applying cancellation.

For the integral on the \(R\)-fiber the real Schwartz
seminorms are uniformly bounded as the input varies in a
bounded set of tests and \(g\) varies in a compact set.
The restrictions \(B(nk g)\) decrease rapidly in the chamber
variables for \(U(C)\), uniformly for \(n\) in a compact
fundamental set. To verify this, on a reduction set for \(k\)
take a rational isotropic flag of \(C\) adapted to its
positive chamber. Wedge the first \(j\) flag vectors with
the wedge of all \(r\) distinguished isotropic vectors of \(S\),
obtaining a nonzero rational vector \(w_j\).
The element \(n^{-1}\) fixes \(w_j\), since it changes the
vectors of \(C\) only by isotropic columns and acts with
determinant one on those columns. In an adelic product norm,
\(\|(nk g)^{-1}w_j\|\) is consequently bounded above by a
constant times the inverse of the product of the first \(j\)
expanding chamber coordinates, with their field multiplicities.
This follows by applying the inverse Siegel decomposition to
the wedge. The unipotent factor fixes its flag wedge and the
remaining compact factors are bounded. The positive chamber
coordinates and hence the height in a Siegel set are bounded by
fixed powers of the maximum of these expanding products and one.
The same assertion is immediate for a compact chamber.
For every rational \(\delta\in G(F)\) the product norm of
\(\delta w_j\) is bounded below by a positive constant, by
applying the product formula to a nonzero rational coordinate.
Since \(\delta w_j=(\delta nk g)((nk g)^{-1}w_j)\) in the
exterior representation, every representative \(\delta nk g\)
has height bounded below by a positive power of the expanding
product in question, up to a fixed factor. On a Siegel
representative of \(nk g\) the rapid decrease of \(B\) thus
dominates every polynomial in the chamber variables of \(k\).
Finitely many reduction sets suffice. Together with the
rapid decrease of \(A_D\) this supplies an integrable
majorant for the fiber integral and proves continuity by
dominated convergence.

At the identity a core test times the linked line test
factors the rational theta sum. The integral in \(h\) gives
the actual descending form on \(U(C)\). The line factor is
the conjugate Jacobi theta at \(n,k\) by
Proposition~\ref{prop:linkedline265}. The displayed simple
root character supplies its flag character. The remaining
integral is the asserted old Jacobi period.
\end{proof}

\subsection{Realizing group convolution in the frame integral}\label{nav:VII-44}
Fix one horizontal edge \(i\ge2\) of
Theorem~\ref{thm:fullperiod266}. Write \(G\) for the product
of its real source groups and \(R\) for the product of the
real flag stabilizers in Propositions~\ref{prop:besorbit265}
and~\ref{prop:jacorbit266}. Give these groups and \(R\backslash G\)
the compatible measures of those propositions. Write \(B\) for
the incoming row, and \(\beta\) for the specified residual
real Schwartz vector. It includes the descending core test
and the line test of the transfer. Absorb the prescribed
finite residual vector into the functional \(J_B(g;w)\), so
that \(w\) is its real variable. Write \(\pi(g)B\) for right
translation. For \(b\in C_c^\infty(G)\) put
\(\pi(b)B=\int_G b(g)\pi(g)B\,dg\).

\begin{lemma}[Real test for a group convolution]\leavevmode
\label{lem:transfer270}
Fix a compact set \(Q\subset G\). To each \(b\) supported in
\(Q\) one can assign a real horizontal Schwartz test \(\Phi_b\),
linearly in \(b\). The outgoing period on this test and the full
finite partition equals
\[
                c_{i,f}v_{i,f} J_{\pi(b)B}(1;\beta).
\]
Here \(c_{i,f},v_{i,f}\) are the original positive base scalar
and finite orbit measure. For every integer \(h\ge0\),
\[
       \|(1+\mathcal N_i)^h\Phi_b\|_2
                     \le A_{i,h}\|b\|_{C^{2h}(Q)} .
\]
The number operator \(\mathcal N_i\) may be taken in any
fixed normalized real oscillator coordinates. The constants
depend on those coordinates, \(Q\), the local frame charts,
measures and \(\beta\). They are independent of both cusp
forms. On the last edge the same assertion holds in the raw
terminal convention, taking the absolute value of the period.
\end{lemma}
\begin{proof}
Write \(\mathcal E_\Phi(g)\) for the real upper-group integral
with evaluation at the opposite position. In the two
unfoldings the change \(g\mapsto r g\), \(r\in R\), acts
on this residual vector by an invertible operator \(\sigma(r)\).
Indeed \(r^{-1}X=Xu_r^{-1}\) in the first position. A left
change by \(u_r\) in the upper-group integral supplies the
simple-root character. In the opposite position \(r\) fixes
the selected row and acts on its fiber by the core oscillator
and, on a Bessel edge, by its line displacement. These are
precisely the factors in the two orbital formulas. Composition
gives \(\mathcal E_\Phi(rg)=\sigma(r)\mathcal E_\Phi(g)\).
Every residual vector occurs at the identity by the chart
construction of Lemma~\ref{lem:realflag265} and extension off
the opposite position. Thus composition gives a representation
on the whole residual space. It is smooth there in the
Schwartz topology by the oscillator action. Right translation
in the unimodular \(R\)-fiber of the unfolded integral gives
\[
 J_B(rg;\sigma(r)w)=J_B(g;w),\qquad
 J_B(g;w)=J_{\pi(g)B}(1;w).
\]
The second equality follows directly from the factor \(B(nkg)\)
in each formula.

Prescribe the equivariant section
\[
           F_b(g)=\int_R b(rg)\sigma(r)^{-1}\beta\,dr .
\]
Its support modulo \(R\) is contained in the image of \(Q\).
Right invariance of \(dr\) gives \(F_b(rg)=\sigma(r)F_b(g)\).
We construct \(\Phi_b\) with \(\mathcal E_{\Phi_b}=F_b\).
Cover the image of \(Q\) by finitely many relatively compact
section charts and choose a partition of unity there. In one
chart write \(s(y)\) for the section and \(\tau(y)\) for its
partition function. Let \(j(w)\) multiply \(w\) by a fixed
Schwartz function of the opposite position whose value at
\(p_*\) is one. In the residual variables before evaluation,
\(\Omega(g)\) denotes the full action after the first
polarization, as in Lemma~\ref{lem:realflag265}. On the
orbit prescribe the vector-valued function
\[
 \Phi_b(s(y)^{-1}Xu,\cdot)=
   \tau(y)\chi(u)\alpha_0(u)\,
          \Omega(s(y))^{-1}j\big(F_b(s(y))\big),
\]
where \(\alpha_0\) is a fixed smooth upper-group bump of
integral one. The orbit coordinates are smooth and the
support in the first position is compact in the full-rank
locus. Extend in transverse coordinates by a fixed bump of
value one on the orbit, using a finite family of submanifold
charts as in that lemma. Sum the resulting chart extensions.
Substitution in the upper integral gives \(F_b(s(y))\) on
each section after summing the partition, and equivariance
gives the same equality on \(G\).

For \(y\) in the closed chart supports, every relevant \(r\)
lies in the compact set \(R\cap Qs(\operatorname{supp}\tau)^{-1}\).
Derivatives of \(F_b(s(y))\) of order \(l\) therefore cost
at most \(C_l\|b\|_{C^l(Q)}\) in each fixed residual Schwartz
seminorm. This follows by differentiating \(b(rs(y))\)
under the integral and bounding \(\sigma(r)^{-1}\beta\) on
that compact set. Take the \(C^l\) norms using a fixed
left-invariant frame and derivatives through order \(l\).
Derivatives of the other chart factors and the residual
oscillator action are bounded on their compact parameters
in the required seminorms. In position coordinates
\((1+\mathcal N_i)^h\) is a sum of words of length at most
\(2h\) in positions and derivatives. The first coordinates
stay on a fixed compact set and the residual seminorms
bound the weighted \(L^2\) norms. A fixed metaplectic
coordinate change takes each such word to a combination of
words of the same length. This proves the estimate.

Finally \(J_B(g;F_b(g))=\int_R b(rg)J_B(rg;\beta)\,dr\).
Integration over the quotient gives \(J_{\pi(b)B}(1;\beta)\).
The \(g\)-integration is now over a compact set and the two
orbital formulas provide an integrable majorant on each
remaining fiber. Their finite orbital integrations give
the stated factor. Lemma~\ref{lem:topchar266} takes the last
test to the raw character by a fixed metaplectic action and
a permutation of marked terms with scalars of absolute value
one. It preserves the bound after changing \(A_{i,h}\).
\end{proof}

\subsection{A finite character on a proper block of the definite core}\label{nav:VII-45}
Let \(H=U(C_m)\) have rank \(N=k+m(r+6)\), and choose \(T\) as
in Lemma~\ref{lem:flag272}. Each \(H(F_v)\), \(v\mid T\), is
\(\mathrm{GL}_N(\mathbb Q_T)\). Choose the split factor and write its
core lattice as \(L=O^k\oplus O\oplus O^{N-k-1}\), \(O=\mathbb Z_T\).
Write the coordinates \(w,y,u\) in this order. Fix a nonzero
residue covector with integral lift \(\ell_v\). Define
\[
 a_v(w,y,u)=(w,\ell_v(w)+T y,u),\qquad
 z(w,y,u)=(zw,y,u)\quad(z^{T-1}=1).
\]
These scalars form a copy \(Z_v\) of \(\mathbb F_T^\times\).
Set \(a=(a_v)_{v\mid T}\), extended by one elsewhere,
\(K=\prod_{v\mid T}\mathrm{GL}(L)\), and \(Z=\prod_{v\mid T}Z_v\).
For \(\lambda=\prod_v\lambda_v\) with every \(\lambda_v\ne1\), write
\(S_\lambda=\sum_{z\in Z}\overline{\lambda(z)}R(z)\) and
\(D_\lambda=S_\lambda R(a)\).

\begin{proposition}[Compression on the definite automorphic quotient]\leavevmode
\label{prop:corehecke272}
Let \(e_K\) be Haar probability averaging over \(K\).
For \(K\)-fixed \(v,w\) on \([H]\),
\[
 \langle R(a)v,D_\lambda w\rangle=\langle v,L_Tw\rangle,\qquad
 \langle D_\lambda v,D_\lambda w\rangle
                  =|Z|\langle v,L_Tw\rangle .
\]
Here \(L_T=\prod_{v\mid T}(1-M_v)\) and \(M_v\) is convolution
on fixed vectors with the probability on the double coset of
\(\operatorname{diag}(T,T^{-1},1,\ldots,1)\).
On every real isotypic subspace nontrivial on at least one real
special unitary factor, the spectrum on the fixed vectors satisfies
\[
          (1-c_T)^f I\ \le L_T\le (1+c_T)^f I,\qquad
                  c_T=\frac{3T-1}{T(T+1)} .
\]
\end{proposition}
\begin{proof}
The image \(a_v L\) is given by \(y=\ell_v(w)\bmod T\).
For \(z\ne1\), its translate by \(z\) has equation
\(y=z^{-1}\ell_v(w)\bmod T\). These are distinct hyperplanes
modulo \(T\) with sum \(L\) upon taking inverse images. Their
intersection has index \(T\) in each lattice. Consequently the
elementary divisors of \(a_v^{-1}za_v\) in \(L\) have one exponent
\(1\), one exponent \(-1\), and all others zero. Compact compression
therefore gives \(M_v\) for every \(z\ne1\) and the identity
for \(z=1\). The coefficients on \(z\ne1\) sum to \(-1\).
Multiplication at distinct places proves
\(e_K R(a)^\ast S_\lambda R(a)e_K=L_T e_K\).
The identity \(S_\lambda^\ast S_\lambda=|Z|S_\lambda\) gives the
second pairing.

We verify the absence of local derived-group invariants needed
for the spectral estimate. The special unitary group \(H^1\) of
the core is simply connected and absolutely almost simple over
\(F\), and \(H^1(F_v)=\mathrm{SL}_N(\mathbb Q_T)\) is noncompact.
Strong approximation gives density of \(H^1(F)\) in the adeles
of \(H^1\) away from \(v\).
\href{https://doi.org/10.1070/IM1969v003n06ABEH000838}{V.~Platonov,
The problem of strong approximation} and
\href{https://library.slmath.org/books/Book61/files/70rapi.pdf}{A.~Rapinchuk,
Strong approximation for algebraic groups, Theorem 2.3}.
If a smooth automorphic function is right invariant under
\(H^1(F_v)\), left rational invariance and this density give
invariance under \(H^1(F_\infty)\). Indeed for \(g\in H(\mathbb A_F)\)
and \(s\in H^1(F_\infty)\), approximate the components of
\(gsg^{-1}\) away from \(v\) by rational elements of \(H^1\).
At \(v\) correct each rational element by a local element of
\(H^1(F_v)\). Conjugating this correction through \(g_v\)
preserves the derived group. The corrected translates have
the value at \(g\) and converge to \(gs\). Smoothing proves
the same assertion for invariant vectors in \(L^2\), since
the local invariant subspace is stable under every translation.
An irreducible real type nontrivial on the derived group has
zero invariant subspace there. The local matrix coefficient
argument of Lemma~\ref{lem:coefficient199} now applies. Its
three-coordinate root subgroup proof is valid for \(N\ge3\)
in every unitary representation without \(\mathrm{SL}_N\)-invariants.
It gives \(\|M_v\|\le c_T\). Each \(M_v\) is self-adjoint since
its double coset equals its inverse. The factors at different
places commute, which proves the spectral bounds.
\end{proof}

\subsection{A projected Fourier coefficient on the original row group}\label{nav:VII-46}
In Theorem~\ref{thm:graph274} write \(G_0=U(S_0)\),
\(H_0=U(W_0)\), of ranks \(r=k+m-1\) and \(k\). Write
\(F_0 A(h)=\int_{[N_0]} A(nh)\overline{\chi_\varpi(n)}\,dn\).
On \(G_0\) use the original terminal graph \(\Lambda_{\ell_{0,v}}\)
and the trivial compact row character at \(v\mid T\).
For \(\lambda\) nontrivial at every such \(v\), put
\(A_\lambda=\sum_{z\in Z}\overline{\lambda(z)}R_0(z)A\).
Let \(w\ge0\) be integral and denote the right-hand side of
Theorem~\ref{thm:bottom271} with \(w_m=w\) by \(\mathcal D_w\).
Use \(N,d_\rho,d_0,a_0,B,L_0\) of Theorem~\ref{thm:corepair272}
and set
\[
 \mathcal C_w=
 \frac{\sqrt{|Z|d_\rho}\,(1-c_T)^f\,\mathcal D_w}
 {2(1+c_T)^{f/2}N^{d_0}L_0^2 B^{6a_0}} .
\]

\begin{theorem}[Descent to the original column rank and terminal graph]\leavevmode
\label{thm:small274}
For each unit bottom row of Theorem~\ref{thm:bottom271} and each
such \(\lambda\), there is a smooth left \(G_0(F)\)-invariant
row \(A\) on \(G_0(\mathbb A_F)\), fixed by all the stated
terminal graph stabilizers, such that
\[
          |F_0 A_\lambda(1)|\ge (1+D_m)^w\mathcal C_w .
\]
Here \(D_m\) is the final Hermite degree of a horizontal lift
from that bottom row. In particular the projection to the
specified central character of \(H_0\) is nonzero on the
original field, frequency and terminal graph. The row \(A\)
is a restriction of a translated horizontal lift.
\end{theorem}
\begin{proof}
Choose the raw lift \(B_m\) and its integral coefficient from
Theorem~\ref{thm:bottom271}. The raw period has modulus at least
\((1+D_m)^w\mathcal D_w\). Choose the coordinates on the core
as in Theorem~\ref{thm:graph274}. Put
\(B_\lambda=\sum_z\overline{\lambda(z)}R(za)B_m\) on the
large row and \(P=F_{\chi_+}B_\lambda\) on its definite core.
The compression calculation of Theorem~\ref{thm:corepair272}
works directly with \(\chi_+\), before its rational change
of frequency. In detail, projection of \(F_{\chi_+}B_m\)
to the original finite level and conjugate real type has
norm at least the raw period modulus divided by its column
norm. Applying \(L_T\) costs at most the lower spectral
factor \((1-c_T)^f\), and extraction against the integral
list gives a column \(Q\) for which
\[
 |\langle D_\lambda Q,P\rangle|
 \ge\frac{|Z|(1-c_T)^f(1+D_m)^w\mathcal D_w}
              {\sqrt{h_m}\,N^{d_0}L_0 B^{4a_0}} .
\]
This is the same two-sided compression identity, since
\(F_{\chi_+}\) commutes with translations on \(H\).
The upper frame bounds give
\(\|Q\|\le L_0 B^{2a_0}/\sqrt{h_m d_\rho}\), and
Proposition~\ref{prop:corehecke272} gives
\(\|D_\lambda Q\|\le\sqrt{|Z|}(1+c_T)^{f/2}\|Q\|\).
Cauchy--Schwarz in probability measure shows
\(\sup_{[H]}|P|\ge 2(1+D_m)^w\mathcal C_w\).

The continuous \(P\) attains this supremum on the compact
quotient. Choose a maximizing representative \(g\in H(\mathbb A_F)\)
whose components at \(v\mid T\) are all one. Indeed \(P\) is
left \(H(F)\)-invariant by rational conjugation in its integral,
and right invariant under an open subgroup at those places
by smoothness of \(B_\lambda\). The weak approximation argument
of Theorem~\ref{thm:graph274} supplies a single rational left
translate putting the components in that open subgroup.
Right translation inside it then sets them all to one. The same
argument sets \(g_v=1\) at any further prescribed finite set of
split places, including the split primes marking the theta tests.
Proposition~\ref{prop:fiber274} supplies an archimedean
\(n(x)\) for which the small integral of
\(B_\lambda(n_0 n(x)g)\) has modulus at least \(|P(g)|/2\).

Define \(A(y)=B_m(y n(x)g a)\) for \(y\in G_0(\mathbb A_F)\)
using the embedded subgroup. Restriction preserves smoothness
and left rational invariance. The terms \(n(x)\) and \(g\)
are trivial at \(v\mid T\). At those places the last test of
\(B_m\) is spherical for the large lattice, so the graph
stabilizers of Theorem~\ref{thm:graph274} fix \(A\). The
scalars \(z\) commute with both \(n(x)\) and \(g\), being
supported at the terminal places. Hence
\(A_\lambda(y)=B_\lambda(y n(x)g)\). Applying \(F_0\) at
one gives exactly the small integral just bounded. Finally
\(F_0 A_\lambda=|Z|p_{\lambda,H_0}F_0A\) by right equivariance
for \(H_0\), as in Corollary~\ref{cor:original200}.
\end{proof}

\subsection{A prescribed real projection after restriction}\label{nav:VII-47}
Keep \(B_m,\mathcal C_w,D_m\) and the common fields of
Theorem~\ref{thm:small274}. Write \(K=H(F_\infty)\) for the
large definite group. Let \(K_J\subset H_0(F_\infty)\) be the
product of the small definite groups at any subset \(J\) of
the real places, with identity elsewhere. Let \(\rho\) be
the real type of \(D_\lambda Q\) in the proof of that theorem.
For a type \(\nu\) of \(K_J\), let \(m_\nu\) be its multiplicity
in \(\rho|_{K_J}\) and \(d_\nu=\dim\nu\). Put
\(q_\nu=m_\nu d_\nu/d_\rho\). All compact group measures
in this subsection are probabilities.

\begin{theorem}[Extraction in an occurring real type]\leavevmode
\label{thm:extract275}
If \(q_\nu>0\), there is a smooth left rationally invariant row
\(A\) on the original \(G_0\), of real \(K_J\)-type \(\nu\)
and fixed by the original terminal graph stabilizers, with
\[
       |F_0 A_\lambda(1)|
          \ \ge\ \tfrac12\sqrt{q_\nu}\,(1+D_m)^w\mathcal C_w .
\]
For \(w=\lceil d_{\rm osc}/2\rceil\), its coefficient on
\(H_0(F_\infty)\) at the identity finite component has
\(L^2\) norm at least \(\sqrt{q_\nu}\mathcal C_w/2\).
The construction permits compact real projections of the lift
and of its restriction.
\end{theorem}
\begin{proof}
Apply the orthogonal real projector \(e_\rho\) for \(K\) to
\(B_m\), obtaining \(\widetilde B\). This is convolution with
\(d_\rho\overline{\chi_\rho}\) and preserves the final Hermite
degree. It commutes with the terminal finite operators and
with the large Fourier map, which is \(H\)-equivariant.
Write \(P=F_{\chi_+}D_{\lambda,G}\widetilde B\).
In the proof of Theorem~\ref{thm:small274}, the test
\(D_\lambda Q\) already has type \(\rho\). The pairing with
it is unchanged, and dividing its lower bound by
\(\|D_\lambda Q\|\) gives
\(\|P\|_{L^2([H])}\ge2(1+D_m)^w\mathcal C_w\).

At a finite level fixing \(P\), the compact quotient is
a finite union of components \(\Gamma_c\backslash K\)
for finite \(\Gamma_c\). On normalized measures its squared
norm is a weighted mean of the squared Haar norms of the
functions \(p_c(s)=P(s,c)\) on \(K\). This follows by
invariance of the conditional measure on each compact
orbit. Choose a finite representative \(c\) with
\(\|p_c\|_2\ge\|P\|_2\). Its function \(p=p_c\) is a sum
of matrix coefficients of \(\rho\) by Peter--Weyl.
For \(b,s\in K\) set
\[
        I(b,s)=d_\nu\int_{K_J}\overline{\chi_\nu(h)}
                                      p(bhs)\,dh .
\]
Then \(\int_{K^2}|I(b,s)|^2\,db\,ds=q_\nu\|p\|_2^2\).
Indeed write \(p(s)=\operatorname{tr}(M\rho(s))\).
The integral inserts the orthogonal projection \(P_\nu\)
onto the \(\nu\)-isotypic space between \(\rho(b)\) and
\(\rho(s)\). Schur orthogonality in \(s\) gives
\(d_\rho^{-1}\|M\rho(b)P_\nu\|_{\rm HS}^2\).
The average of \(\rho(b)P_\nu\rho(b)^*\) is scalar, and
its trace makes that scalar \(q_\nu\). Another application
of orthogonality gives the identity. Thus some \(b,s\)
give modulus at least \(\sqrt{q_\nu}\|p\|_2\).

We realize the left rotation \(b\) while restoring the
terminal components to one. Choose \(\gamma\in H(F)\)
with \(\gamma_\infty^{-1}\) arbitrarily close to \(b\)
and \(\gamma_v c_v\) inside the right level of \(P\) at
each \(v\mid T\). For weak approximation use the Cayley
chart in Theorem~\ref{thm:graph274}. At a real compact
factor the matrices with no eigenvalue \(-1\) are dense,
as is seen by a small scalar rotation. At the split finite
places the same chart is dense. Approximation in its
rational vector-space coordinates gives both requirements
together. Take \(g_\infty=s\), \(g_f=\gamma_f c\) and
set \(g_v=1\) at \(v\mid T\) by right level invariance.
For \(h\in K_J\), rational invariance gives
\(P(hg)=p(\gamma_\infty^{-1}hs)\). By continuity the integral
of this function against \(d_\nu\overline{\chi_\nu(h)}\)
has modulus at least \(\sqrt{q_\nu}\|P\|_2/2\).

Apply Proposition~\ref{prop:fiber274} at each \(hg\) and
integrate in \(h\). The integrands are uniformly bounded
since \([N]\) and \(K_J\) are compact and \(g\) is fixed.
Indeed all factors \(n_0 n(x)\) inside the row belong to
the large unipotent group. Dominated convergence applies.
Every \(h\) commutes with the real \(n(x)\) in that
proposition. The triangle inequality supplies an \(x\)
for which the character integral of the small Fourier
coefficients has modulus at least \(\sqrt{q_\nu}\|P\|_2/4\).
Start with \(y\mapsto\widetilde B(y n(x)g a)\) on \(G_0\).
Project its right translates by \(K_J\) to \(\nu\).
The resulting \(A\) is smooth, retains the terminal graph
invariance and gives exactly the selected expression after
the finite character sum and \(F_0\), as these operations
commute with the compact real projector. This proves the
pointwise estimate. The orbit under \(H_0(F_\infty)\)
still spans at most \((1+D_m)^{d_{\rm osc}}\) dimensions.
The compact projections preserve the real degree space, and
the last projection is inside the existing finite orbit
span on the small group. The evaluation bound of
Corollary~\ref{cor:real274} gives the norm estimate.
\end{proof}

\subsection{Unfolding the first nontrivial Fourier--Jacobi period}\label{nav:VIII-41}
Keep \(S=C\perp Ez\), \(T=D\perp(Ee+Ee')\) and the oscillator
of Proposition~\ref{prop:shell263}. Put \(G=U(S)\),
\(H=U(C)\), \(U=U(D)\). Let \(N\) be the unipotent radical of
the stabilizer of \(Ee\) in \(U(T)\), and denote rational adelic
quotients by brackets. In the mixed position \(x=X(e)\) write
a Schwartz test as \(a(x)\beta\), where \(\beta\) belongs
to the residual oscillator on \(\operatorname{Hom}(-D,S)\).
On this test the action of \(g\in G(\mathbb A_F)\) has the form
\(a(g^{-1}x)\Omega(g)\beta\). This formula defines the residual
action \(\Omega\) with the splitting factor of the actual kernel.
It follows by conjugating position translations and dual
modulations, since \(G\) stabilizes both hyperbolic summands
and preserves the measure of the position space.

Write \(\Theta_q(n,k)\) for the theta series on the line part
\(\operatorname{Hom}(-D,Ez)\) with Schwartz test \(q\), evaluated
at \(n\in N(\mathbb A_F)\), \(k\in U(\mathbb A_F)\). Here the
Heisenberg action is the one on the fiber \(x=z\). On \(U\)
the line uses the splitting attached to the odd line
\(\eta^{-1}\). In the Hermitian direct sum the residual theta
uses the core kernel on \(C\) and this line kernel on \(Ez\).
Give the compact unipotent quotient \([N]\) Haar probability
and use dual adelic measures on Heisenberg position and frequency
spaces, with rational additive quotients of volume one.

\begin{proposition}[The orbit formula at the prescribed character]\leavevmode
\label{prop:unfold263}
Let \(B\) on \([G]\) and \(A\) on \([U]\) be smooth compact-finite
cusp forms. Write \(\beta=b\otimes q\) in the direct sum, with
\(q\ne0\). The Fourier--Jacobi period of the plain horizontal
lift is
\[
 \begin{split}
 P(a,\beta)&=\int_{[U]} A(k)\int_{[N]}
       \overline{\Theta_q(n,k)}
       \int_{[G]} K_{T,S}(nk,g;a\otimes\beta)\,B(g)\,dg\,dn\,dk\\
 &=\int_{H(\mathbb A_F)\backslash G(\mathbb A_F)}
              a(g^{-1}z)\,\mathcal R(g)\,d\dot g .
 \end{split}
\]
The quotient measure disintegrates the two group measures.
With \(b_g\) the contraction of \(\Omega(g)\beta\) against
\(q\) in the line variable, linear in \(\beta\), one has
\[
 \mathcal R(g)=\int_{[H]}\int_{[U]} B(hg)\,A(k)\,
                       K_{D,C}(k,h;b_g)\,dk\,dh .
\]
The function \(\mathcal R\) is continuous. At the base point
it satisfies
\[
 \mathcal R(1)=\|q\|_2^2\int_{[H]} B(h)A_2(h)\,dh,\qquad
 A_2(h)=\int_{[U]}A(k)K_{D,C}(k,h;b)\,dk .
\]
\end{proposition}
\begin{proof}
The center of \(N\) shears \(e'\) by a trace-zero multiple
of \(e\). Parameterize it by \(c\in\mathbb A_F\) so that it
acts on the position \(x\) by \(\psi(c h_S(x,x))\).
Indeed a shear by \(\alpha e\), \(\bar\alpha=-\alpha\), adds
\(\alpha x\) to the dual column. In the Schr\"odinger formula
its phase is half the alternating cross pairing with \(x\).
Thus \(c\) is \(\operatorname{Tr}_{E/F}(\xi_0\alpha)/2\)
up to the common sign of the dual coordinate. This is an
invertible \(F\)-linear change of the center parameter. The
line test has character \(\psi(-c)\) in these coordinates.
Integrating the product with its conjugate over \(F\backslash
\mathbb A_F\) therefore selects precisely the rational positions
\(x\) of norm \(-1\). Witt extension over \(E/F\) makes these
one \(G(F)\)-orbit, with stabilizer \(H(F)\). Unfolding the
sum gives the domain \(H(F)\backslash G(\mathbb A_F)\) with
representative position \(z\). The rational frame scalings
in Proposition~\ref{prop:shell263} apply to both kernels, so the
character selected by this calculation is the grid character.

On this fiber the remaining action of \(N\) is the full
Heisenberg action in the line variable. To check this, modulo
the center an element of \(N\) is determined by its functional
on \(D\) taking values in \(Ee\). Under composition with \(X\)
it adds this functional times \(z\) to the residual columns.
These columns run through \(\operatorname{Hom}(-D,Ez)\);
the \(C\)-columns are fixed. The commutator on the line is
the alternating trace pairing, with the central action just
computed. The map is an isomorphism on their vector quotients
and on their one-dimensional centers, hence it identifies the
two Heisenberg groups over \(F\). The Jacobi action is given
by the ordinary Weyl operators. In coordinates with central
term one half of the commutator, the noncentral generators
are obtained by adding a functional to the \(D\)-columns
and inserting its Hermitian adjoint in the \(e'\)-column;
the quadratic correction in that column is the usual half
self-pairing. This gives the translation and modulation on
the line with the same phase in the two kernels.

On the subgroup \(H\times U\) the residual oscillator factors
as the core oscillator and the line oscillator just specified.
The symplectic direct sum gives this factorization up to a
character of the two groups. At a split place the scalar on the
chosen components for \(h\in H,\ k\in U\) is
\[
  \frac{\gamma_2(\det h)}
       {(\gamma_2\eta^{-1})(\det k)}
  =\frac{\gamma_2(\det h)}{\gamma_2(\det k)}\,\eta(\det k).
\]
The first factor on the right is the core splitting and the
second the line splitting. Hyperbolic substitution in the
\(x\)-variable has Jacobian one. Rational changes of polarization
preserve the adelic theta functional by Poisson summation, so
the character ratio is trivial on rational points. At almost
all nonsplit finite places the integral models match on the
two hyperspecial subgroups. A character trivial there is trivial
on the whole local group by the Cartan and hyperbolic-plane
argument of Proposition~\ref{prop:cancelsquare256}. The
remaining local factors are trivial by weak approximation for
these unitary groups, as in that proposition. This proves the
factorization on the adelic groups, including its phases.

For line tests \(v,q\), integration of the two Heisenberg theta
series on \([N]\) gives their inner product \(\langle v,q\rangle\).
In rationally dual position and frequency coordinates, integration
over the frequency quotient makes the two rational summation
indices equal. The sum of the resulting translated position
integrals over a fundamental domain is the full position integral
of \(v\bar q\). The common central phase cancels and the center
quotient has volume one. This computation applies for every
\(k\), since its line action is unitary on both tests. It also
applies to a general residual Schwartz vector by integration
in its line variable, or by Schwartz-continuous extension from
finite sums of tensors. The resulting core test is \(b_g\).
Disintegrating \(H(F)\backslash G(\mathbb A_F)\) over
\(H(\mathbb A_F)\backslash G(\mathbb A_F)\) now gives the
formula. Equivariance makes the inner expression invariant
under left translation of \(g\) by \(H(\mathbb A_F)\).

Here are convergence details for the two indefinite groups.
Uniformly on a compact fundamental set of \(N\), each absolute
theta sum grows at most polynomially in the heights on Siegel
sets for \(U\) and \(G\). This follows by bounding the rational
lattice sum by finitely many real Schwartz seminorms after
fixing the finite supports, as in Theorem~\ref{thm:seesaw245}.
It also majorizes the sum restricted to the norm shell by using
the absolute sum before the central integration. Rapid decrease
of \(A,B\) therefore justifies unfolding and Fubini. For continuity
of the inner expression, compact variation of \(g\) keeps \(b_g\)
in a bounded set of Schwartz tests. The restrictions \(B(hg)\)
decrease rapidly on Siegel sets of \(H\), uniformly in such \(g\).
Indeed extend an isotropic flag defining one of these sets in
\(C\) to a full rational isotropic flag in \(S\). The split torus
coordinates embed with additional entries one after the original
positive entries. The original lower bounds place them in a
Siegel chamber of \(G\) with a fixed lower bound. The bounded
unipotent coordinates embed in the larger parabolic. Bounded
anisotropic Levi coordinates commuting with the split torus
can be moved into the right compact set. Thus each of the
finitely many reduction sets lands in a Siegel set of \(G\)
with a fixed compact factor, also after the indicated translations.
Rapid decrease there supplies an integrable majorant for any
polynomial growth of the core kernel. Dominated convergence
proves continuity. Finally \(b_1=\|q\|_2^2 b\), which proves
the value at the identity.
\href{https://arxiv.org/abs/1106.2149}{S.~Miller and W.~Schmid,
Theorem A and Theorem 3.1}. For the corresponding norm-orbit unfolding in the
symplectic theta correspondence see
\href{https://archive.math.arizona.edu/xuehang/refinedFJ.pdf}
{H.~Xue, Lemma 8.3.1}.
\end{proof}

\subsection{The Bessel character in the terminal coordinates}\label{nav:VIII-42}
Let \(m=2s+2\ge4\), so \(S_m=V_{m,0}\) is the direct sum
of \(C_m\) and a staircase of length \(2s+1\). Denote the
first \(s\) unscaled staircase columns by \(b_j\) and the
middle column by \(z\). Let \(N\) be the flag unipotent
radical and \(H=U(C_m)\), embedded by fixing all staircase
columns. For \(n\in N\) write \(c_j(n)\) for the coefficient
of \(b_j\) in \((n-1)b_{j+1}\), \(j<s\). Write \(c_s(n)\)
for the coefficient of \(b_s\) in \((n-1)z\). The character
in the supplied terminal Bessel datum and its bilinear period
are
\[
 \begin{split}
  \chi_+(n)&=\psi_E\!\left(\xi_0\sum_{j=1}^s c_j(n)\right),\\
  \mathcal B^+(e,B)&=\int_{[H]} e(h)
             \int_{[N]}B(nh)\overline{\chi_+(n)}\,dn\,dh .
 \end{split}
\]
These formulas read off the isotropic columns, the middle
column and the coefficients of the given datum in the raw
Hermitian extension. In a convention conjugating the column
form one takes that form to be \(\bar e\).
\srcloc{Development/AbelianVarietyTheory/Definitions/Divisors/CmHostImaginaryGeneratorNeZero}{189--218}.
\srcloc{Development/AbelianVarietyTheory/Definitions/Divisors/BplusOfEquivariant}{269--291}.

\begin{lemma}[Passage to the prescribed terminal character]\leavevmode
\label{lem:topchar266}
A nonzero terminal Bessel period in the scaled frames of
Theorem~\ref{thm:bestransfer265} yields a nonzero
\(\mathcal B^+(e,B')\) with the same \(e\) and the same
earlier horizontal tests. The last horizontal test of \(B'\)
can have one original marked coset at each marked place and
one product of real Hermite functions.
\end{lemma}
\begin{proof}
Scaling \(b_j\) by \(d^j\) preserves its flag. In the scaled
coordinates the coefficient in the \(j\)-th simple position
is \(d^{j+1}d^{-j}c_j=dc_j\). The last vector is \(d^{s+1}z\),
so the same ratio applies to the last coefficient in \(d^sb_s\).
Thus the scaled character is
\(\psi_E(d\sum_jc_j)=\overline{\chi_+(n)}\).
The embedded core fixes both the raw and the scaled columns,
and is exactly \(H\).

Define a rational unitary involution \(\delta\) by multiplying
\(b_j\) and its opposite column by \((-1)^{j-1}\), multiplying
\(z\) by \((-1)^s\), and fixing \(C_m\). Opposite columns
still pair to one. Conjugating \(n\) by \(\delta\) negates
each \(c_j(n)\), including \(c_s(n)\). The involution preserves
the rational unipotent quotient and its Haar probability and
commutes with \(H\). For the right translate
\(B'(g)=B(g\delta)\), rational left invariance gives
\(B'(nh)=B(\delta^{-1}n\delta\,h)\). Changing variables in
the displayed integral therefore identifies it with the
scaled-character period of \(e,B\), with the same measure
on \(H\). If the probability on the compact terminal core
is used instead, both periods acquire the same positive
normalizing factor.

Right translating a theta lift by \(\delta\) acts on its last
Schwartz function by the Weil action on \(S_m\), by the defining
theta kernel. At a marked split place this is integral matrix
substitution times a unit-modulus determinant scalar. The
substitution preserves the original matrix lattice and its
depth because each staircase summand is preserved. It therefore
sends a single coset to a single coset. At the unmarked places
\(\delta\) fixes the original vacuum by
Proposition~\ref{prop:rowfinite262}. The scalars can be put
in the real test. At infinity the Weil action gives a Schwartz
function. Its terminal period with the chosen finite data is
continuous in this function, by
Proposition~\ref{prop:besorbit265} and the same rational
conjugation. Expansion in the original Hermite basis selects
a single product with nonzero period. This affects only the
last edge.
\end{proof}

\subsection{Weighted selection of an individual Hermite test}\label{nav:VIII-43}
Fix a field and an edge \(i\ge2\) as in
Lemma~\ref{lem:transfer270}. Let \(d_i\) be the total real
position dimension of its horizontal oscillator, and
\(g_i=\dim_{\mathbb R}G_i(\mathbb R\otimes_{\mathbb Q}F)\),
where \(G_i=U(V_{i,0})\) is regarded here as a group over
\(F\). Write \(T_i\) for the number of terms in the finite
partition on the edge. Its coefficients have absolute value
one. Put \(t_i=1\) on an odd-index edge and \(t_i=\|q\|_2^2\)
on an even-index edge, with the full line norm and measures
of Proposition~\ref{prop:jacorbit266}. Then \(J_B(1;\beta)=t_iP_i\),
where \(P_i\) is the specified incoming period.

At each intermediate edge fix the polynomial Fock realization
for compatible maximal compact subgroups of the row pair.
Such coordinates can be fixed before the input forms are chosen.
Indeed the two compact groups have compact product image in
the real symplectic group. Average a positive inner product over
this product. The skew operator for the symplectic form in that
metric, divided by the positive square root of minus its square,
is an invariant compatible complex structure, with the sign
chosen for positivity. An orthonormal complex basis gives
normalized monomials, or equivalently product Hermites in its
Schr\"odinger coordinates. Use the same maximal compact of a row
on its two incident edges. The underlying Schwartz space and its
seminorms transform by a fixed metaplectic map when changing
coordinates. At the very last edge the original terminal Hermite
coordinates may be used for the expansion.

\begin{proposition}[A weighted one-edge period bound]\leavevmode
\label{prop:weighted270}
Suppose \(i\ge3\) and the preceding test is a normalized product
Hermite of total degree \(D\) in the compatible coordinates.
For each integer \(w\ge0\) there is a constant \(C_i(w)>0\)
depending only on the local real data and the specified row
type. If \(P_i\ne0\), one original finite choice \(\mu\) and
one normalized product Hermite \(h_\nu\) on this edge satisfy
\[
 \frac{|P_{i+1}(\mu,h_\nu)|}{(1+|\nu|)^w}
 \ \ge\
 \frac{c_{i,f}v_{i,f}t_i}{2T_i C_i(w)}
       \frac{|P_i|}{(1+D)^{E_i(w)}},\qquad
 E_i(w)=d_{i-1}+4(w+d_i+1)+4g_i+4 .
\]
For \(i=2\) take any finite list of compact types containing
the second row. With \(M_2,H_2\) of
Proposition~\ref{prop:kernel270}, the same inequality holds
with the degree factor \((1+D)^{E_i(w)}\) replaced by
\(M_2H_2^{\,4(w+d_2+1)+4g_2+4}\).
\end{proposition}
\begin{proof}
For \(i\ge3\) the compact orbit of the preceding test is
contained in the Fock polynomials of degrees at most \(D\).
The action of a unitary symplectic change of basis on this
space is substitution on the variables times its action on
the vacuum line. On each degree its differential is the sum
of the infinitesimal actions on at most \(D\) tensor factors
and the fixed scalar action on the vacuum. Thus each fixed
compact generator has operator norm at most \(C_i(1+D)\).
The dimension of this space is \(\binom{D+d_{i-1}}{d_{i-1}}\).
The theta map with the preceding row fixed intertwines the
outgoing compact action on its test with right translation
on its output. Therefore the compact types of the incoming
form on edge \(i\) occur in this finite-dimensional space.
Their distinct dimensions sum to at most its dimension, and
\(C_K(\tau)\le C_i'(1+D)^2\) on them, by summing the squared
generator norms. Theorem~\ref{thm:realrow257} fixes the
discrete series at every real place. Thus the scalar for
\(\Omega_G\) on all these forms is fixed even when there
are several automorphic copies. The parameters of
Proposition~\ref{prop:kernel270} satisfy
\[
       M\le C_i''(1+D)^{d_{i-1}},\qquad H\le C_i'''(1+D).
\]
At the second row its finite type list instead gives
\(M_2,H_2\) directly.

Take \(b\) from that proposition with \(\pi(b)B=B\), and
let \(a_\nu=\langle\Phi_b,h_\nu\rangle\) be the coefficients
of the test in Lemma~\ref{lem:transfer270}. Use the raw
transform before expansion on the last edge. Set \(j=w+d_i+1\).
Cauchy--Schwarz and Parseval give
\[
   \sum_\nu |a_\nu|(1+|\nu|)^w
     \le \left(\sum_\nu(1+|\nu|)^{-2(d_i+1)}\right)^{1/2}
                  \|(1+\mathcal N_i)^j\Phi_b\|_2 .
\]
At degree \(l\) there are \(\binom{l+d_i-1}{d_i-1}\) indices,
so the numerical series converges.
By Lemma~\ref{lem:transfer270} the last norm is at most
\(A_{i,j}\|b\|_{C^{2j}}\).
The compact convolution estimate now bounds the whole
weighted sum by \(C_i(w)(1+D)^{E_i(w)}\). At \(i=2\) it
gives \(C_2(w)M_2 H_2^{4j+4g_2+4}\). All the constants just
used involve the charts, fixed residual tests, coordinates,
measures, group and representation parameters.

The period evaluated on the full finite partition and the
real test has modulus \(c_{i,f}v_{i,f}t_i|P_i|\).
Its expansion in the individual \(P_{i+1}(\mu,h_\nu)\) is
absolutely convergent. In fact the period has finite
Schwartz order by Proposition~\ref{prop:order268}, the term
on \(h_\nu\) grows at most polynomially in \(|\nu|\), and
the test coefficients have summable moments of every order.
The same continuity holds after a fixed coordinate change.
If every weighted period were smaller than the stated
right-hand side, the triangle inequality and the weighted
coefficient sum over the \(T_i\) finite choices would give
at most \(c_{i,f}v_{i,f}t_i|P_i|/2\). This contradiction
proves the assertion. A selected nonzero form on an
intermediate edge is compact-finite and cuspidal by
Theorem~\ref{thm:realrow257}.
\end{proof}

\subsection{A character-projected period at the larger terminal rank}\label{nav:VIII-44}
Use the bottom vector and data of Theorem~\ref{thm:bottom271}
and take \(w_m=0\). Denote its positive lower bound by \(\mathcal D\).
Use \(T\), \(a\), \(t\), and \(Z\) from Lemma~\ref{lem:flag272}
and Proposition~\ref{prop:corehecke272}. On \(G=U(S_m)\) use
the same translations \(D_{\lambda,G}\) as \(D_\lambda\), via the
inclusion of \(H=U(C_m)\). Define
\(\mathcal P_\varpi(B,Q)=\langle Q,F_{\chi_\varpi}B\rangle_{[H]}\)
with inner product conjugate-linear in \(Q\). Put
\[
             \mathcal D'=
       \frac{(1-c_T)^f\,\mathcal D}
            {\sqrt{h_m}\,N^{d_0}L_0 B^{4a_0}} ,
\]
where \(B\) in this denominator is the lattice constant of
Lemma~\ref{lem:terminalclass259}.

\begin{theorem}[Projected terminal Bessel periods on a proper scalar block]\leavevmode
\label{thm:corepair272}
There is a single horizontal lift \(B_m\) from the prescribed
bottom row and a single column \(Q\) from the conjugate list in
Proposition~\ref{prop:corelattice272} such that simultaneously
for every character \(\lambda\) nontrivial at all \(v\mid T\),
\[
 \left|\mathcal P_\varpi\bigl(D_{\lambda,G}R(t)B_m,\,
                              D_\lambda Q\bigr)\right|
       \ \ge\ |Z|\,\mathcal D' .
\]
With the integral column \(T^{a_0}D_\lambda Q\) the lower
bound is multiplied by \(T^{a_0}\). With orthogonal character
projections on both translated arguments the lower bound is
\(\mathcal D'/|Z|\) before this clearing of denominators.
\end{theorem}
\begin{proof}
Theorem~\ref{thm:bottom271} gives \(B_m\) and a column \(e_0\)
with \(|\int_{[H]} e_0 F_{\chi_+}B_m|\ge\mathcal D\). The
two upper frame bounds in Proposition~\ref{prop:schurframe259}
bound each original vector and covector norm by the square
root of its upper frame bound. The class norm formula gives
\(\|e_0\|\le L_0 B^{2a_0}/\sqrt{h_m d_\rho}\).
Let \(Y\) be the orthogonal projection of \(F_{\chi_+}B_m\)
to the real type of \(\bar e_0\) and the full original finite
core level. This function has finite norm by smoothness and
compactness of the definite and unipotent quotients. Since
the projected space contains \(\bar e_0\),
\(\|Y\|\ge\mathcal D/\|e_0\|\).

The real type is nontrivial on the special unitary group.
At the terminal real place the parameter of
Theorem~\ref{thm:realcore253} contains the cloud entries
\(v_j-1/2\). Consecutive cloud entries have difference
\(4\Delta>1\), with \(k\ge3\), and every added string lies
above the cloud by the core gap. Subtracting the ordered
half sum of roots leaves an unequal pair of highest weights.
The conjugate type has the same nontriviality. Hence
\(\|L_TY\|\ge(1-c_T)^f\|Y\|\) by
Proposition~\ref{prop:corehecke272}. Apply
Proposition~\ref{prop:corelattice272} to \(L_TY\).
Substituting the bound for \(\|e_0\|\) gives one \(Q\)
with \(|\langle Q,L_TY\rangle|\ge\mathcal D'\).

The rational change of flag gives
\(F_{\chi_\varpi}R(t)B_m=F_{\chi_+}B_m\), as functions before
real projection. Right translation of either row by an element
of \(H\) passes inside its unipotent integral as right translation
of the resulting function. Thus the new period pairs
\(D_\lambda Q\) with \(D_\lambda F_{\chi_+}B_m\).
By Lemma~\ref{lem:flag272} the latter input before translation
is \(K\)-fixed. Compression between \(K\)-fixed inputs gives
\(|Z|\langle Q,L_T F_{\chi_+}B_m\rangle\) by
Proposition~\ref{prop:corehecke272}. The real projector and
the finite-level projectors away from \(T\) commute with \(L_T\).
They fix \(Q\), so this last pairing equals
\(|Z|\langle Q,L_TY\rangle\). In particular \(Q\) works for
all the characters. Dividing each sum \(S_\lambda\) by \(|Z|\)
gives its orthogonal projection. The column integrality follows
from Proposition~\ref{prop:corelattice272}. On the row the
finite sums act on the last theta test by equivariance. The
rational element \(t\) also translates its real factor.
\end{proof}

Here the ambient and column ranks are \(N+m-1,N\), with
\(N=k+m(r+6)\). The graph mixes a \(k\)-dimensional block
with one further direction \emph{in the definite core}. In
Corollary~\ref{cor:original200} the ranks are \(k+m-1,k\),
and the graph uses the negative middle line of the ambient
space. There the scalars are central in the entire local
column group. Translations within that group preserve every
central character, since they commute with each scalar.
In particular they preserve a zero projection to a prescribed
central character. A comparison with the period on that graph
and its integral row normalization is still required.

\subsection{A lower bound on the original definite real group}\label{nav:VIII-45}
Keep the original ranks, frequency and terminal graph of
Theorem~\ref{thm:small274}. Use on the last oscillator the
Hermite coordinates for a positive metric compatible with
the compact core \(H(F_\infty)\). Let \(d_{\rm osc}\) be
the total number of real position variables across the real
places, and put \(w=\lceil d_{\rm osc}/2\rceil\). Denote
by \(H_{0,\infty}\) the product of the definite real column
groups and give it Haar probability.

\begin{corollary}[A definite real norm with the prescribed finite character]\leavevmode
\label{cor:real274}
For the row provided by Theorem~\ref{thm:small274},
\[
    \left(\int_{H_{0,\infty}}
        |F_0 A_\lambda(h_\infty,1_f)|^2\,dh_\infty\right)^{1/2}
                                    \ge\mathcal C_w .
\]
The right side is independent of the chosen final Hermite
index and of the real translation in the unipotent averaging.
On the compact adelic column quotient \(F_0A_\lambda\) has a
nonzero pairing with an integral column in some real type
and the prescribed central character, at a sufficiently fine
finite level.
\end{corollary}
\begin{proof}
In signed orthonormal coordinates the compact group of a
definite summand preserves the positive metric on the
oscillator and its compatible complex structure. Its Weil
action preserves total Hermite energy. Equivalently, the
compact group acts linearly on creation operators, and acts
by a scalar on the vacuum. The span of the real compact
translates of the final test is consequently of dimension
at most \((1+D_m)^{d_{\rm osc}}\), by counting multi-indices
of degree at most \(D_m\). Theta equivariance gives the same
bound for the \(H(F_\infty)\)-orbit span of \(B_m\). Finite
translations and their sum in \(B_\lambda\) commute with
this real group. The element \(n(x)\) acts only between
the extra space \(D\) and the flag spaces and fixes \(W_0\).
It commutes with \(H_{0,\infty}\). Also \(g_\infty\in H(F_\infty)\).
Thus a right real column translation \(k\) of
\(B_\lambda(y n(x)g)\) is obtained from the translation
of \(B_\lambda\) by \(g_\infty^{-1}kg_\infty\) on the large
row. Taking the small Fourier coefficient and evaluating
on the real column group are linear. The same dimension
bound holds for the right translates of the function inside
the displayed norm.

In a right-translation-invariant finite-dimensional space
of continuous functions on a compact group, the sum of the
squares of an orthonormal basis at a point is constant. A
translation changes the basis unitarily, and integration
shows the constant is the dimension for Haar probability.
Cauchy--Schwarz therefore bounds the absolute value at the
identity by the square root of that dimension times the
norm. Apply this to the preceding span and use
Theorem~\ref{thm:small274}. The choice of \(w\) absorbs the
final Hermite degree and gives the assertion.

The nonzero real restriction lies on a finite adelic
neighborhood of positive measure, by the finite smooth
level of the row. The full function on \([H_0]\) is
compact-finite at infinity by the same orbit-span argument.
Take a finite level fixing it and refine inside a positive
principal congruence group at an odd split prime. As in
Proposition~\ref{prop:terminalcore255}, the class stabilizers
are trivial since they are finite and this congruence group
is torsion-free. Each occurring compact unitary type has
integral highest weights. Its Schur and determinant model
and the dual admit integral lattices spanning their spaces,
so their matrix coefficients span the type on each class.
Some class coefficient \(Q\) thus pairs nontrivially with
the function. The sum \(\sum_{z\in Z}\overline{\lambda(z)}
R(z)Q\) multiplies this pairing by \(|Z|\), because the
function has character \(\lambda\). One can use a common
finite refinement under the translates. The scalars preserve
the standard column lattices at \(T\), so pulling back
integral matrix coefficients uses no denominator in the
fixed-right-argument convention of
Proposition~\ref{prop:corelattice272}. This proves existence
at that level.
\end{proof}

The norm just estimated uses the real group on one finite
component. The lower estimate in Corollary~\ref{cor:original200}
requires a prescribed global cuspidal row, a prescribed real
column type, the original levels away from \(T\), and a
comparison in the specified integral normalization. The
restriction constructed here proves the original terminal
graph and frequency conditions. Its projection to those
global spaces and the required quantitative comparison remain
open.

\subsection{Passage from a Bessel period to a Jacobi period}\label{nav:IV-39}
Use \(i=2r+2\), \(\ell=r+1\), and the groups and measures
of Proposition~\ref{prop:jacorbit266}. Write \(s=\dim_E S\)
and \(t=\dim_E D\). Write \(c_{\ell,f}\) and \(v_f\) for the
base scalar and compact-orbit measure of
Theorem~\ref{thm:flagfinite264} at this edge.

\begin{theorem}[Transfer to the next Fourier--Jacobi period]\leavevmode
\label{thm:jactransfer266}
Fix a nonzero Bessel period \(P_i\) of a compact-finite cusp
form \(B_i\) at the prescribed level against the specified
descent \(A_i=\operatorname{desc}_b A_{i+1}\).
There is a fixed distinguished line test \(q_\ell\), chosen
before the two forms, and one horizontal test at the original
depth with one original coset at each marked place, such that
\[
   |P_{i+1}(q_\ell)|\ \ge\
       \frac{c_{\ell,f}v_f\,\|q_\ell\|_2^2}
            {2(N\mathfrak m)^{\ell s+rt}}\,|P_i| .
\]
The real cutoff inducing this estimate has quotient integral
one. For nonvanishing the horizontal real test may be a single
product Hermite function, with the same \(q_\ell\).
\end{theorem}
\begin{proof}
The middle ideal in \(S\) is \(I_\ell^{(2r+1)}=\mathfrak l^\ell\)
in the notation of Proposition~\ref{prop:positions264}. The
undeepened line summand of the residual map lattice from \(D\)
to \(S_0=C\perp L\) is therefore
\(\mathfrak A\mathfrak l^\ell
\operatorname{Hom}(\Lambda_D,\mathcal O_E z)\).
The Jacobi line rule multiplies this by
\(\mathfrak m/\bar{\mathfrak m}\), as in
Proposition~\ref{prop:linkedline265}. At an unmarked place
take its ordinary vacuum. At a marked place its vacuum is the
indicator of the zero coset of depth \(a_v\) in the original
line position lattice. Indeed the position lattice acquires
the factor \(\pi^{a_v}\) and the opposite lattice the dual
factor \(\pi^{-a_v}\). Translation and modulation then give
exactly this invariant line. Take a fixed nonzero Gaussian at
every real place. These choices define \(q_\ell\) and have
positive finite \(L^2\)-norm.

Use \(b\otimes q_\ell\) as residual test. Its finite part has
one original marked coset, by the direct sum \(C\perp L\).
For the remaining finite variables use the two position tests
of Proposition~\ref{prop:positions264}. Their integrated value
at the identity is \(c_{\ell,f}(b\otimes q_\ell)_f\), and their
finite support projects to the orbit of measure \(v_f\).
Represent this orbit by the level group. Right invariance of
the test and of \(B_i\) makes the finite part constant there
after disintegration, as in Theorem~\ref{thm:bestransfer265}.

At infinity choose a residual Schwartz test, in both the
opposite position and the residual oscillator, whose value
at the zero opposite position is \((b\otimes q_\ell)_\infty\).
On real local sections use the continuous function obtained
by acting on this test and inserting its evaluated value in
\(J_{q_\ell}\) of Proposition~\ref{prop:jacorbit266}. At the
base point the function has value \(\|q_\ell\|_2^2P_i\).
Choose a product neighborhood throughout which the difference
from this value is less than half its modulus.
Lemma~\ref{lem:realflag265} realizes a nonnegative quotient
bump of integral one supported in this neighborhood. Its
upper-group integral has the phase required by the flag
character. The orbital formula now gives a new period of
modulus at least \(c_{\ell,f}v_f\|q_\ell\|_2^2|P_i|/2\).

The marked expansion has exactly
\((N\mathfrak m)^{\ell s+rt}\) terms with coefficients of
absolute value one. The line is already a single original
coset at each of those places. The triangle inequality gives
the assertion for an individual term. Its period is a
continuous functional of the horizontal real test by
Proposition~\ref{prop:jacorbit266}. The Hermite expansion
and cuspidality argument of Theorem~\ref{thm:bestransfer265}
now give a compact-finite lift with nonzero Jacobi period.
As there, the numerical lower bound uses the localized smooth
test.
\end{proof}

\subsection{Local integrals for horizontal pairs of opposite parity}\label{nav:IX-40}
Use the original pair \(V_{i+1,0},V_{i,0}\) on a field of
Theorem~\ref{thm:first250}. Its ranks are \(N=u+s,\ u=r+is\),
where \(s=r+6\) is odd. Attach
\(\gamma_i=\chi_{m-i}\) to the even space and
\(\gamma_i\eta^{-1}\) to the odd space. Use the character of
Proposition~\ref{prop:parity262}, or any parity character with
the same infinite type, unramified off the marked set \(S_0\)
and trivial on its prescribed principal unit groups. Write
\(\Lambda_N,\Lambda_u\) for the actual Hermitian lattices.
These include the fractional ideals in the orthogonal extensions.
\srcloc{Development/AbelianVarietyTheory/Definitions/Divisors/BasePlaceNotMemChiGridBadPlaces}{168--188,401--439}.

\begin{proposition}[The horizontal finite tests at the original level]\leavevmode
\label{prop:rowfinite262}
On every horizontal edge of the bottom row there is a fixed
product of finite tests satisfying the original depth equations.
At each marked place its factor is one integral coset indicator,
and its local doubling form on vectors fixed by \(K_u(a_v)\)
is \(D_{i,v}\) times the source inner product, where
\[
 D_{i,v}=(Nv)^{-a_v(uN+u^2)}
                \prod_{h=1}^{u}\bigl(1-(Nv)^{-h}\bigr)^{-1}.
\]
At the other finite places its normalized doubling section gives
the spherical integral. The tests are invariant under the given
principal congruence groups on both sides at marked places
and the two lattice stabilizers elsewhere.
\end{proposition}
\begin{proof}
The full map lattice
\(\mathfrak A\operatorname{Hom}(\Lambda_N,\Lambda_u)\) is
self-dual for the trace pairing by
Proposition~\ref{prop:tracelat253}. At a marked place its two
halves are dual rectangular matrix lattices. Take the indicator
of \((I_u\ \ 0)+\varpi_v^{a_v}\operatorname{Mat}_{u,N}(\mathcal O_v)\)
in lattice bases on the chosen half. The marked characters have
exponents at most \(a_v\), by the level of \(\gamma_i\) and the
choice of \(\eta\). Their determinant factors are therefore
one on both congruence groups. Left translation by \(h\) gives
overlap with this rectangle precisely when \(h\in K_u(a_v)\).
Indeed the first \(u\) columns of both an intersecting matrix
and its translate are congruent to \(I_u\), so their quotient
is in \(K_u(a_v)\). In that subgroup the entire coset is
preserved. The integral is the additive volume \((Nv)^{-a_vuN}\)
times the group volume of \(K_u(a_v)\). Counting invertible
matrices modulo \(\varpi_v^{a_v}\) gives the stated product,
with the measures of Theorem~\ref{thm:markedcore253}. The same
theorem's translation and modulation calculation proves the
depth equations. Projectivity and the Chinese remainder theorem
lift all the chosen marked matrix residues at once, including
for the fractional lattices.

At an unmarked split place the position indicator is that of
the integral rectangle in the two lattice bases, dilated by
the local generator of \(\mathfrak A\) in the chosen component.
Dilation commutes with both substitutions and their determinant
characters. Its normalized coefficient is thus the usual split
unramified coefficient, also at the dyadic places.

At a nonsplit \(v\) the extension is unramified, the residue
characteristic is odd, and the change of variables in
Proposition~\ref{prop:tracelat253} gives a trace-zero unit and
a character of conductor \(\mathcal O_v\) for the original tensor
lattice. Both Hermitian lattices split integrally into hyperbolic
planes and, in odd rank, one unit line. The proof is successive
lifting of isotropic residue vectors and their partners as in
Proposition~\ref{prop:vacuumcore254}. Polarize the even space
by its two integral isotropic halves. In the resulting position
space take the lattice indicator and give its lattice volume one.
Translations and dual modulations fix it. On the odd space the
group acts by substitution with trivial determinant twist, since
the character attached to the even space is unramified with
trivial restriction to \(F_v^\times\), and a uniformizer is
common to the two fields.

On the even space the integral Siegel Levi preserves the indicator
by change of basis. Its character is trivial on unit determinants.
Integral unipotents multiply it by a phase of an integral
quadratic form on its support. A Weyl element interchanging
isotropic coordinates acts by partial Fourier transform with
unramified unit factors. The Fourier transform preserves the
corresponding self-dual lattice indicator. In the Weil formula
the index of a unimodular quadratic form here is one. Indeed
the integral of its quadratic phase over any dual-lattice coset
outside the lattice vanishes by orthogonality of its nontrivial
linear term, whereas the integral over the lattice is its volume
one. This evaluates the stabilizing Gaussian integral defining
the index. The other splitting factor on an integral Weyl
representative is the character of a unit. These elements
generate the lattice stabilizer. Reduction modulo the maximal
ideal gives the Siegel Bruhat decomposition over the residue field,
whose parabolic factors and partial exchanges lift integrally.
An element left congruent to the identity has an invertible
diagonal block, so block elimination factors it into integral
upper and lower unipotents and a Levi element. This proves
invariance for both compact groups with the specified splitting.
For the character and Weil-index factors in the Schr\"odinger
splitting see \href{https://doi.org/10.1007/BF02773003}
{S.~Kudla, Theorem 3.1} and
\href{https://msp.org/pjm/2001/199-1/pjm-v199-n1-p10-s.pdf}
{S.-Y.~Pan, Sections 2.1 and 3.2}.

The sign on the integration group can now be checked explicitly.
Put \(Q=Nw=(Nv)^2\). In a Cartan representative of \(U(u)\)
write the hyperbolic entries as \(\varpi^{t_j},\varpi^{-t_j}\),
\(t_j\ge0\), and set \(T=\sum_jt_j\). If \(N\) is even, the
position space consists of \(N/2\) copies of the small space.
Substitution on each gives normalized overlap \(Q^{-T}\).
If \(u\) is even, the small torus acts through its Siegel Levi
on \(u/2\) copies of the large space. On the \(j\)-th copy
dilation by \(\varpi^{t_j}\) has normalized overlap
\(Q^{-Nt_j/2}\), by the square root of the lattice index.
Its splitting character on the chosen uniformizer to power
\(t_j\) is \((-1)^{Nt_j}\) by restriction to \(F_v^\times\).
Thus in both cases the coefficient is
\[
                   (-1)^{NT}Q^{-NT/2}.
\]
In the doubled group the primitive exterior generator of the
diagonal \(u\)-plane has lattice norm \(Q^T\) after applying
\((h,1)^{-1}\). This follows by choosing one column from each
coordinate pair in the exterior product on a Cartan representative.
The Siegel Levi determinant consequently has valuation \(T\).
At the theta point the normalized unramified section evaluates
its unramified character on this determinant and raises its
absolute value to \(s/2+u/2=N/2\). This gives the displayed
coefficient. Bi-invariance on the small compact group, with
the unramified determinant character trivial there, extends the equality
to all \(h\). Flat extension multiplies both restrictions by
the same power of the Iwasawa height. The unramified doubling
calculation therefore applies with exactly the omitted places
\(S_0\).
\end{proof}

\subsection{Passage from a Jacobi period to a Bessel period}\label{nav:IX-41}
Let \(i=2r+1\), \(r\ge1\), \(i+1\le m\). Use the finite
and quotient measures of Theorem~\ref{thm:flagfinite264}
and Proposition~\ref{prop:besorbit265}. Denote by \(c_{r,f}\)
the positive scalar in the former theorem's base value,
and by \(v_f\) its support probability on the product of
compact quotient orbits. Put \(s=\dim_E S\), \(t=\dim_E T_0\),
and \(s_0=\dim_E C\).

\begin{theorem}[Transfer to the next Bessel period]\leavevmode
\label{thm:bestransfer265}
Take a compact-finite row cusp form \(B_i\) at the prescribed
level. Let \(A_i\) be the descent of \(A_{i+1}\) by the
specified core test of Theorem~\ref{thm:downcore254}. Suppose its
Jacobi period \(P_i(q)\) against \(B_i\) is nonzero, with \(q\)
in the distinguished finite line vacuum and a fixed real
Schwartz test. There is a horizontal test with one original
integral coset at every marked place, and the original vacuum
elsewhere at finite places, whose lift satisfies
\[
 |P_{i+1}|\ \ge\
   \frac{c_{r,f}\,v_f}{2(N\mathfrak m)^{\,r(s+t)+s_0}}\,
   |P_i(q)|.
\]
Here \(P_{i+1}\) is the Bessel period against \(A_{i+1}\) at
the grid's distinguished character. The real first-position
cutoff induces a bump of quotient integral one by
Lemma~\ref{lem:realflag265}. The other real tests use the
specified core and line vectors. For nonvanishing a single
product of real Hermite functions can be used on the entire
horizontal edge.

In particular, for \(m\ge4\) the datum from
Theorem~\ref{thm:third263} gives a nonzero fourth Bessel
period on the same field, interval, and finite levels,
retaining its original terminal integral coefficient and its
preceding three horizontal tests when those give a compact-finite
third row.
\end{theorem}
\begin{proof}
Form the residual test as the tensor product of the specified
descending core test and the conjugate transfer of \(q\) from
Proposition~\ref{prop:linkedline265}. At the finite places
take the full new line indicator and put its normalization
scalar into the real line test. Denote the resulting residual
product by \(\beta_f\otimes\beta_\infty\). At infinity choose
a Schwartz test \(\rho_\infty\) in the second-position and
core variables with value \(\beta_\infty\) at \(p_*\).
For instance multiply \(\beta_\infty\) by a product of real
bumps of value one at the required rows.

At the finite places take the two positions of
Proposition~\ref{prop:positions264}. Apply
Theorem~\ref{thm:flagfinite264} termwise to the marked
partition of the full new line. The coefficient is supported
on \(R_fK_f\) and its value at \(1\) is \(c_{r,f}\beta_f\).
On the image of \(K_f\) in the quotient use representatives
in \(K_f\). Both the full horizontal test and \(B_i\) are
right invariant there, by Proposition~\ref{prop:rowfinite262}
and the level hypothesis. Consequently integration of
Proposition~\ref{prop:besorbit265} in the finite variable
gives the factor \(v_f c_{r,f}\).

Use real sections \(s(y)\) through one. Before multiplication
by the first-position cutoff, the remaining scalar function of
\(y\) on these sections is
\[
   J(y)=\mathcal J\big((s(y),1);
          [\Omega(s(y))\rho_\infty](p_*)\otimes\beta_f\big).
\]
It is continuous by Proposition~\ref{prop:besorbit265},
since the real oscillator action and evaluation are continuous
on Schwartz space. The last assertion of that proposition gives
\(J(1)=P_i(q)\). Take a product neighborhood where
\(|J(y)-P_i(q)|<|P_i(q)|/2\). Lemma~\ref{lem:realflag265}
supplies a first-position test that multiplies \(J(y)\) by a
nonnegative cutoff \(a\) supported there, with quotient integral
one. Outside its support the integrand vanishes. The upper-group
bump in that lemma also has integral one, in the actual
disintegrated measures. Thus the period for the finite sum
is \(c_{r,f}v_f\int a(y)J(y)\,dy\), of modulus at least
\(c_{r,f}v_f|P_i(q)|/2\).

The two finite positions contribute \((N\mathfrak m)^{r(s+t)}\)
terms with unit-modulus coefficients. The full new line
contributes \((N\mathfrak m)^{s_0}\) more choices, each with
coefficient one. They occupy disjoint matrix blocks, and the core
test already has one marked coset. Linearity and the triangle
inequality give the displayed bound for an individual product
of original marked cosets. At that finite choice the period
is a continuous linear functional of the real test by the
Schwartz seminorm estimate in Proposition~\ref{prop:besorbit265}.
The real Hermite expansion converges in Schwartz space.
Nonvanishing on the selected test therefore implies nonvanishing
on an individual product Hermite function in the original
polarization. The corresponding output is a compact-finite
cusp form by Theorem~\ref{thm:realrow257}. The quantitative
inequality uses the smooth cutoff, whose derivatives can depend
on the two forms.

For \(i=3\), choose the compact-finite third row and nonzero
Jacobi period of Theorem~\ref{thm:third263}. Its \(q\) has exactly
the distinguished finite vacuum, its descent is the one just
used, and its level is the original one. The construction
changes only the test on the fourth horizontal lift. This gives
the stated fourth period.
\end{proof}

\section{The prescribed initial and terminal periods}\label{nav:g-VII-3}

\subsection{Nonvanishing at the finite level of the initial restriction}\label{nav:I-52}
Use \(G,H,\iota,z,\mathfrak m\) from Theorem~\ref{thm:short246}
and write \(J=\iota^{-1}K(\mathfrak m)\) on the finite adeles.
Take Haar product measures, with counting measure on rational
points. Let \(\phi\) be a bounded continuous automorphic function
on \(H\) invariant under \(J\). For a real-place function \(F_\infty\)
on \(G_\infty\) put \(\mathcal K=F_\infty{\bf1}_{K(\mathfrak m)}\). Define
the Poincar\'e series and its restricted period at \(g\in H(\mathbb A_F)\) by
\[
 P_g(t)=\sum_{\gamma\in G(F)}\mathcal K(\iota(g)^{-1}\gamma t),
 \qquad
 \mathcal P_g=\int_{H(F)\backslash H(\mathbb A_F)}
                      P_g(\iota h)\,\phi(h)\,dh .
\]
The initial restriction pairing in the source uses such a series
of a prescribed discrete-series coefficient. Its contribution from
the identity coset, divided by \(\operatorname{vol}J\), is
\[
        A_g=\int_{H_\infty}F_\infty(\iota k)\,\phi(g(k,1))\,dk .
\]

\begin{theorem}[The initial restriction period and its nonidentity terms]
\label{thm:restriction246}
Suppose \(F_\infty\) is measurable and for some \(B\ge0\) and
every \(t\in G_\infty\) satisfies
\[
 \int_{H_\infty}|F_\infty(t\iota k)|\,dk
              \le B\prod_\tau |t_\tau z|_\tau^{-(2n+5)} .
\]
Here the norms are in the adapted frames at the identity,
\(n=\dim_E X\), and \(f=[E^+:\QQ]\).
Then the restricted series is absolutely integrable and
\[
 \left|\frac{\mathcal P_g}{\operatorname{vol}J}-A_g\right|
    \le \|\phi\|_\infty B\, C_n^f
                  (b_g\sqrt{N\mathfrak b})^{-(2n+4)}
\]
with the constants in Theorem~\ref{thm:packing246}.

For every \(g\) with \(\phi(g)\ne0\) there also exists
\(F_\infty\in C_c^\infty(G_\infty)\) for which all nonidentity
terms vanish and \(\mathcal P_g\ne0\) at this same finite level.
If \(\phi\) belongs on the right to a specified isotypic component
\(\sigma\) of a maximal compact \(K_H\subset H_\infty\), the
function \(F_\infty\) can be taken of right type \(\bar\sigma\)
under \(\iota K_H\).
\end{theorem}
\begin{proof}
Group the sum first by the right cosets of \(H(F)\) in \(G(F)\).
Left invariance of \(\phi\) unfolds the integral for each such
coset to one integral over \(H(\mathbb A_F)\). For its
representative \(\gamma\), the finite variables satisfying the
support condition form either the empty set or one right coset
\(h_0J\): the product of the inverse of one allowed group element
with another is \(\iota(h_0^{-1}h_f)\in K(\mathfrak m)\).
Thus each nonempty finite integral of the absolute indicator
has volume \(\operatorname{vol}J\). The archimedean argument for
the coefficient is \(t=\iota(g_\infty)^{-1}\gamma_\infty\)
followed by \(\iota k\), so the integral of its absolute value
is at most \(B D_g(\gamma z-z)\).
By Theorem~\ref{thm:packing246} these bounds are summable, including
the identity term, whose weight is one. Applying Tonelli's theorem
first with absolute values justifies unfolding the full series on
the quotient and then multiplying by the bounded \(\phi\).
For the identity term change variables from \(h\) to \(g h'\);
the finite component of \(h'\) runs over \(J\), on which \(\phi\)
is right invariant. This gives \(\operatorname{vol}J\,A_g\).
Summing the other terms gives the displayed bound.

For the existence assertion, the diagonal image of
\(\mathfrak b\Lambda_g\) in the real vector space
\(\prod_\tau X_\tau\) is a lattice. After applying
\(\iota(g_\infty)^{-1}\) it is still discrete. Choose a neighborhood
\(U\) of zero meeting this transformed lattice only in zero. Choose
a nonnegative smooth bump on \(G_\infty\), positive at the identity,
supported in a sufficiently small compact neighborhood there such that
\(tz-z\in U\) throughout its support. If an unfolded term meets the
finite test, \(w=\gamma z-z\in\mathfrak b\Lambda_g\) by
Theorem~\ref{thm:short246}. Every argument
\(t=\iota(g_\infty)^{-1}\gamma_\infty\iota k\) of its real test
satisfies \(tz-z=\iota(g_\infty)^{-1}w\). Hence the only class
meeting both supports is the identity. Compact support makes
the original series locally finite. Its unfolding is absolutely
integrable on the restriction by the same support argument and
compactness of the intersection with the closed subgroup \(\iota H_\infty\).

By continuity shrink the bump so that
\(|\phi(g(k,1))-\phi(g)|<|\phi(g)|/2\) wherever
\(F_\infty(\iota k)\ne0\). The integral of the bump on \(H_\infty\)
is positive. The triangle inequality now gives
\(|A_g|\ge \tfrac12|\phi(g)|\int_{H_\infty} F_\infty(\iota k)\,dk>0\).
Finally apply the right \(\bar\sigma\)-isotypic projector on the
bump if \(\phi\) has the stated type. In the bilinear pairing the
projector on the bump moves to the \(\sigma\)-isotypic projector
on \(\phi\), by right translation and
\(\chi_\sigma(k^{-1})=\overline{\chi_\sigma(k)}\). It therefore
preserves the period. Right multiplication by \(\iota K_H\) fixes
\(z\), so this average still excludes every nonidentity term.
\end{proof}

The compactly supported construction gives a nonzero restricted
period for a given nonzero row and its finite level. For its use
in the theta chain, the core form must also lie in a specified
\(G_\infty\) discrete series and in the image of the descending
column lifts. The first bound quantifies the error for a proposed
coefficient in that series. Its identity action and the uniform
constants for that coefficient still enter the comparison.

\subsection{The first restriction on the original intervals}\label{nav:I-53}
Take the fields and intervals of Theorem~\ref{thm:region198} and the
first pair \(H=U(Y)\subset G=U(X)\) of Theorem~\ref{thm:short246}.
Here \(n=\dim_E X=k+m\), \(f=[E^+:\QQ]\), and \(b=\exp(x^\eps)\).
Use the original lattice, finite modulus \(\mathfrak m\) and its
restriction \(J=\iota^{-1}K(\mathfrak m)\) from
Theorem~\ref{thm:restriction246}. Write \(m_0=g_1g_2\), the product
of the two marked level primes, and fix \(C>0\). Put
\(Q=C f n^4\log b\). For a real-place function \(F_\infty\) on
\(G_\infty\) set \(V(h)=F_\infty(\iota h)\) on \(H_\infty\).

\begin{theorem}[The restriction bound for a prescribed real coefficient]
\label{thm:initial247}
Let \(u\) be a nonzero bounded smooth cuspidal row fixed on the right
by \(K_H(m_0)\) and \(J\). Suppose a measurable \(F_\infty\) and a
constant \(c\ne0\) satisfy the following three conditions, with the
measures of Theorem~\ref{thm:restriction246}.
\begin{align*}
 \int_{H_\infty}V(h)u(s h)\,dh &= c\,u(s) &&(s\in H(\mathbb A_F)),\\
 \int_{H_\infty}|F_\infty(t\iota h)|\,dh
  &\le |c| e^Q\prod_\tau |t_\tau z|_\tau^{-(2n+5)} &&(t\in G_\infty),\\
 \int_{H_\infty}|V(\exp(-A)h)-V(h)|\,dh
  &\le |c| e^Q\|A\|_1 &&(A\in\operatorname{Lie}(H_\infty)).
\end{align*}
For all sufficiently distant intervals there is a point \(g\) at
which the original restricted period satisfies
\[
     |\mathcal P_g|\ge \tfrac14\operatorname{vol}(J)|c|\,\|u\|_\infty .
\]
For fixed \(C\) the start of this range is uniform in the row,
its supremum, the coefficient and its scalar \(c\).
\end{theorem}
\begin{proof}
Put \(U=\|u\|_\infty\) and choose \(g\) with \(|u(g)|>U/2\).
Apply the first hypothesis at \(s\exp A\) and at \(s\). The change
of variable \(h\mapsto\exp(-A)h\) in the first integral and the third
hypothesis give
\[
                |u(s\exp A)-u(s)|\le U e^Q\|A\|_1 .
\]
The tail bound at \(t=1\) supplies the absolute convergence needed
here. Theorem~\ref{thm:cusp247} therefore applies with \(L=Ue^Q\),
\(q_*=n\), and \(d=2f(n-1)\).

Choose the cyclotomic modulus \(M\) in Lemma~\ref{lem:dimension239}.
On these intervals \(\log M=o(\log b)\) and \(\log m_0=O(\log b)\).
The ranks in Theorem~\ref{thm:rangenew50} give \(f\le n\) for large
\(x\), \(n\to\infty\), and \(\log n=O(\log x)\). The different
bound in Lemma~\ref{lem:lattices247} gives \({\rm rd}(E)\le M\).
In particular
\[
 \log\kappa=\log m_0+2\log M+(d+5)\log(2d),\qquad
               -\log b_g\ll_C f^2 n^5\log b .
\]
For the second inequality take logarithms in Theorem~\ref{thm:cusp247}.
The terms are at most \(f\log(4\sqrt{fM})\) plus
\(fn(\log(4\kappa)+Q+\log 2)\), with the displayed bound since
\(d\le2fn\), \(f\le n\) and \(\log b=x^\eps\).

Let \(p_0\) be the auxiliary prime at which the original grid
specifies its depth, and write \(p_1,p_2\) for its two type primes.
For the fixed positive depth constant \(C_{\rm dep}\) the exponent
at the selected prime above \(p_0\) is \(\lfloor S/\log p_0\rfloor+1\),
where
\[
    S=C_{\rm dep} f^2 n^6
         \log\big(b n m_0 p_1p_2 |\operatorname{disc}E|\big).
\]
These are the original parameters in
\srcloc{Development/AbelianVarietyTheory/Definitions/Divisors/LeastSplitAuxiliaryPrimeThetaChainSpec}{250--256}.
The ideal norm of that prime is at least \(p_0\). Hence
\(\log N\mathfrak m>S\ge C_{\rm dep}f^2 n^6\log b\).
The height loss above is at most half this lower bound once \(n\)
is large enough. We obtain \(b_g N\mathfrak m\ge(N\mathfrak m)^{1/2}\).

The identity term divided by \(\operatorname{vol}J\) is \(c\,u(g)\).
By Theorems~\ref{thm:packing246} and~\ref{thm:restriction246},
with \(\sqrt{N\mathfrak b}=N\mathfrak m\), the other terms together
have absolute value, after the same division, at most
\[
        U|c|\exp\{Q+f(6n+9)-(n+2)\log N\mathfrak m\} .
\]
By the depth lower bound the exponent tends uniformly to \(-\infty\).
For large \(x\) the error is at most \(U|c|/4\). Subtracting this
from \(|c\,u(g)|>U|c|/2\) gives the claimed estimate.
\end{proof}

The hypotheses specify the remaining real-place calculation at this
step. The source lists the corresponding local branching coefficient,
its integrated tail and its derivative in
\srcloc{Development/AutomorphicArithmetic/ThetaCorrespondence/RealBands/BranchDefs}{11--19}.
For the designated row one must obtain such a coefficient in its
prescribed discrete series. The restricted period must then be
related to the integral pairing through all later lifts and the
prescribed projection. The unrestricted cohomology sum in
\eqref{eq:region198} is a separate input on these same intervals.

\subsection{A single real estimate for the prescribed initial row}\label{nav:I-54}
On the fields and intervals of Theorem~\ref{thm:initial247} keep its
\(G,H,n,f,\mathfrak m,m_0,J\). Let the discrete series \(\sigma\)
of \(H_\infty\) have the spaced row parameters in
\eqref{eq:letters111} and the lists following it. Take an isometric
\(H_\infty\)-equivariant realization \(T\) of \(\sigma\) in cuspidal
\(L^2\) and a unit vector \(w\) in its lowest compact type \(S\).
Write \(u_w=T w\), using its smooth representative, and assume the
boundedness and the two finite invariances of Theorem~\ref{thm:initial247}.
Use the larger real representation \(\pi\) in the convention of
Theorem~\ref{thm:compact248}, with a multiplicity-one isometric
embedding \(j:\sigma\to\pi|_{H_\infty}\). Since \(F_v\) is a
conjugate coefficient, the representation generated on the right
by this kernel is contained in the contragredient of \(\pi\).
Thus for a prescribed core series one takes \(\pi\) to be its
contragredient in this convention.

\begin{theorem}[A normalized overlap at the original grid level]
\label{thm:ratio248}
At every \(x\) in every sufficiently distant interval, each point \(g\) with
\(|u_w(g)|>\|u_w\|_\infty/2\) satisfies
\(b_gN\mathfrak m\ge(N\mathfrak m)^{1/2}\), uniformly in \(T,w\).
At any such point, for every \(v\) with \(\mathcal T(v)<\infty\) and
\(\alpha=\langle v,jw\rangle\ne0\), choose \(F_{v'}\) as in
Theorem~\ref{thm:compact248}. Then the restricted period satisfies
\[
 \left|\frac{\mathcal P_g}{\operatorname{vol}J}
            -\frac{|\alpha|^2}{d_\sigma}u_w(g)\right|
 \le \|u_w\|_\infty d_S\mathcal T(v) C_n^f
                       (N\mathfrak m)^{-(n+2)} .
\]
In particular \( |\mathcal P_g| \) is at least
\(\operatorname{vol}(J)\|u_w\|_\infty |\alpha|^2/(4d_\sigma)\)
whenever
\[
    \frac{d_\sigma d_S\mathcal T(v)}{|\langle v,jw\rangle|^2}
                     \le \frac{(N\mathfrak m)^{n+2}}{4 C_n^f} .
\]
All quantities in this comparison belong to the same initial pair
and the same interval. Rescaling the Haar measure of \(H_\infty\)
changes the integral in \(\mathcal T\) and the formal degree by
reciprocal factors.
\end{theorem}
\begin{proof}
Let \(B\ge1\) bound the absolute Harish-Chandra coordinates of the
row across the real places. The construction following
\eqref{eq:letters111} gives \(\log(2+B)=O(\log b)\) uniformly.
Proposition~\ref{prop:gradient248} gives the Lipschitz hypothesis
of Theorem~\ref{thm:cusp247} with \(L/\|u_w\|_\infty\le
40\sqrt{n-1}(B+2n)\), independently of \(v\). Take \(q_*=n\),
\(d=2f(n-1)\). Taking logarithms of the cusp bound yields
\[
 -\log b_g\ll f n\{\log(2m_0 M^2 f(2+B))+fn\log(2fn)\}
\]
for every \(g\) above half the supremum. Indeed
\({\rm rd}(E)\le M\), \(\log\kappa=\log m_0+2\log M+(d+5)\log(2d)\),
and \(\log(1+L/\|u_w\|_\infty)\ll\log(2+B)+\log(2n)\).
Such points exist by the definition of the positive supremum.
Since \(f\le n\), \(\log n=O(\log x)\), \(\log m_0=O(\log b)\) and
\(\log M=o(\log b)\), the displayed bound is
\(o(f^2 n^6\log b)\). The depth computation in the proof of
Theorem~\ref{thm:initial247} now gives the stated height at \(g\).

Write \(V=F_{v'}|_{H_\infty}\). Since \(V\in L^1\) and \(T\) is
bounded and equivariant, the identity
\(\int V(h)u_w(\,\cdot\,h)\,dh=(|\alpha|^2/d_\sigma)u_w\) holds
in \(L^2\) by Theorem~\ref{thm:compact248}. The integral is also
pointwise continuous by dominated convergence and boundedness of
\(u_w\). Pairing against a compactly supported continuous function
and applying Fubini identifies it with the Hilbert integral. The
equality therefore holds everywhere, since quotient Haar measure
has full support. Use this identity for the main term of
Theorem~\ref{thm:restriction246} and the bound
\(\mathcal T(v')\le d_S\mathcal T(v)\) for the remaining terms.
The height estimate proves the first inequality of the theorem.
Under the last displayed hypothesis in its statement the error
is at most one quarter of \(\|u_w\|_\infty |\alpha|^2/d_\sigma\).
The lower bound follows from the choice of \(g\).
\end{proof}

There is a fixed-parameter existence argument for vectors in the
last comparison. Suppose the adjacent gaps of the larger series
exceed \(s+2D\), where \(s=2n+5\) and \(D>4n^2\) is an integer.
Since compact-finite vectors are dense, some such \(v\) has
\(\langle v,jw\rangle\ne0\). Lemma~\ref{lem:translate163} bounds its
coefficient by a constant times \(\prod_\tau\|g_\tau\|_{\rm op}^{-(s+2D)}\).
Here a compact-finite vector in the product is a finite sum of
tensors, by compact admissibility at each of the finitely many
places. To integrate a single real factor write \(t=b_1 h_0\),
where \(h_0\) fixes \(z\) and \(\|b_1\|_{\rm op}=\|b_1^{-1}\|_{\rm op}
\le2|tz|\). Indeed for a positive line the positive and negative
coordinate lengths of \(tz\) are \(\cosh l,\sinh l\). The boost on a
signed plane sends \(z\) to these coordinates and has norm \(e^l\).
A block unitary aligns them with \(tz\). For a negative line exchange
the signs, and in a definite space a compact isometry suffices.
The norm of \(b_1 h\) is at least \(|tz|\) and at least
\(\|h\|_{\rm op}/(2|tz|)\). Raising these bounds to powers \(s+D\)
and \(D\), respectively, gives the integrable majorant
\(2^D|tz|^{-s}\|h\|_{\rm op}^{-D}\) for the integrated power.
This proves \(\mathcal T(v)<\infty\) for the fixed parameters.

The compact dimension itself satisfies
\(\log d_S=O(f n^2\log(2+B+n))\) by the Weyl dimension formula.
Each numerator is at most \(2(B+n)+n\), and each denominator is a
positive integer. Theorems~\ref{thm:tensor249}, \ref{thm:transverse250} and~\ref{thm:direction249} now control the decay constants and the prescribed component together as the series vary. The source obtains a translate by estimating a second moment of a
branching coefficient in
\srcloc{Development/WeilBound/RepresentationTheoryGL/HolomorphicDiscreteSeries/SchurEigenProjections/HolomorphicMoments/MomentHigherEstimateHolds}{16--26}.
Its use to select a nonzero component appears in
\srcloc{Development/AutomorphicArithmetic/ThetaCorrespondence/FockModels/RealBranch}{140--169}.
The inequality above requires the component along the specified \(w\).
It concerns the first restriction. The later projected integral
and the cohomology sum in \eqref{eq:region198} remain to be bounded
on the same intervals.

\subsection{A quantitative period after two horizontal lifts}\label{nav:I-55}
On the original fields use the period and measures of
Theorem~\ref{thm:secondrow256}. Take \(m\ge2\), as holds on
the intervals in question. Write \(A_2(e)\) for descent from a
terminal class coefficient \(e=e_{c,a,b}\) of
Theorem~\ref{thm:individual259} to \(C_2\), with the given
tests on all edges above it. Let \(b_1\) be the prescribed
test on the second core edge. For a Hermite product \(h_\alpha\)
on the first edge, denote its two factors on \(Y\) and \(Ez\)
by \(p_\alpha\) and \(q_\alpha\), respectively. Use the
finite components of Theorem~\ref{thm:firstsquare255}.
First lift a prescribed smooth row \(u\) on \(Y\) to \(C_1\)
using \(p_\alpha\). Lift the resulting form to \(Z=C_2\perp Ez'\)
with the test \(b_1\otimes Uq_\alpha\) of
Proposition~\ref{prop:cancelsquare256}, replacing its marked
line component by the individual residue indicator \(\nu\).
Write \(B_{2,\alpha,\nu}\) for the resulting form, and set
\[
 J_{e,\alpha,\nu}=
      \int_{[U(C_2)]} A_2(e)(k)\,B_{2,\alpha,\nu}(k\oplus1)\,dk .
\]
The count of allowed \(\nu\) is
\(D_L=(\operatorname{Nm}\mathfrak m)^{\dim_E C_1}\).

\begin{theorem}[Second horizontal periods with individual integral tests]\leavevmode
\label{thm:secondbound260}
Let \(\ell(A)=\mathcal B(A,u)\) on \(H(v)\).
With the terminal constants of Theorem~\ref{thm:individual259}
and the first real test \(\phi\) of
Proposition~\ref{prop:projecttests260}, one has
\[
 \sum_e\sum_\alpha\frac{1}{D_L}\sum_\nu |J_{e,\alpha,\nu}|^2
 \ \ge\ \frac{e^{-\delta}C_{\rm fix}\|\ell\|_*^2}
               {h_m d_\rho B^{4a_0}\|\phi\|_2^2}.
\]
The sum over Hermite indices denotes the supremum of finite
partial sums. For every \(0<\eta<1\) and \(\|\ell\|_*>0\), one
terminal coefficient, one Hermite product and one marked
residue satisfy
\[
 |J_{e,\alpha,\nu}|\ \ge\
 \frac{(1-\eta)e^{-\delta/2}\sqrt{C_{\rm fix}}\,\|\ell\|_*}
 {H_*\sqrt{d_\rho}\,N^{d_0}B^{2a_0}\|\phi\|_2 C(v,\phi)}.
\]
At the two horizontal steps the forms occurring in this
nonzero restriction are cuspidal. For any nonzero row \(u\)
of the prescribed realization there is a pure compact-finite
\(v\) to which the individual lower bound applies.
\end{theorem}
\begin{proof}
For the single product \(h_\alpha\), factor the two kernels by
Proposition~\ref{prop:cancelsquare256}. The first restriction
of the core descent equals the period on \(C_2\) with the full
Fourier-transformed line test. In detail, the descent to \(C_1\)
is the integral of its second core kernel against \(A_2(e)\).
Upon inserting it into the restriction, first integrate the
\(Y\)-kernel against \(u\) to form the first horizontal lift.
This lift is cuspidal or zero by the real first occurrence in
Theorem~\ref{thm:secondrow256}. It and \(A_2(e)\) give rapid
decay at the two variables of the remaining integral, with the
intervening theta kernels of polynomial growth. The two orders
of integration and the replacement by the second horizontal
kernel are therefore justified as in that theorem. This
argument gives the identity for every product test, including
those with a vanishing lift. By
Proposition~\ref{prop:oppline256} are scalars
\(\epsilon_\nu\) of modulus one such that
\[
 \mathcal B(F_e(h_\alpha),u)
     =D_L^{-1}\sum_\nu \epsilon_\nu J_{e,\alpha,\nu}.
\]
The cosets are those of the actual negative line lattice and
the other marked factor is the specified rectangle. This is
the finite Fourier decomposition at the original levels.
The first finite averaging of Theorem~\ref{thm:firstsquare255}
also preserves the first period since the two input forms
have the required level invariances.

Put \(z_e=\mathcal B(P_vT^\dagger e,u)\). Apply
Proposition~\ref{prop:projecttests260} simultaneously to all
terminal coefficients. For every \(e\), \(z_e\) is the limit
of the finite sums of the last display with coefficients
\(c_{j,\alpha}\). Cauchy--Schwarz in the Hermite index and
the equal-weight average over \(\nu\) bound the squared
modulus of the \(j\)-th sum by
\[
 \|\Phi_j\|_2^2\sum_{\alpha\in\operatorname{supp}c_j}
                           D_L^{-1}\sum_\nu|J_{e,\alpha,\nu}|^2 .
\]
Using the upper limit for \(\|\Phi_j\|_2\) gives a lower
bound \(|z_e|^2/\|\phi\|_2^2\) for the full sum at this \(e\).
Summing and applying Theorem~\ref{thm:individual259} proves
the first inequality.

For the same finite sums the triangle inequality and the
absolute Hermite coefficient bound give
\[
 |z_e|\le C(v,\phi)\|\phi\|_2\sup_{\alpha,\nu}|J_{e,\alpha,\nu}|.
\]
Choose the single \(e\) of Theorem~\ref{thm:individual259}
and take an entry within the indicated factor \(1-\eta\) of
the resulting lower bound for the supremum. This choice works
also when that supremum is infinite. Since the period is
nonzero, both horizontal lifts are nonzero. Their cuspidality
then follows from real first occurrence at each corresponding
target and the tower property in Theorem~\ref{thm:realrow257}.
The existence of \(v\) for a prescribed smooth row follows from
Theorem~\ref{thm:moment258}.
\end{proof}

The estimate reaches the restriction on \(C_2\) with one
integral terminal input and one integral residue at the marked
places. The factor \(\sqrt{C_{\rm fix}}/\|\phi\|_2\) records
the normalization of the actual even core transforms. For a
general compact-finite \(v\), the weighted integral in
\(C(v,\phi)\) and the norm of its restriction functional still
enter the bound. The final arithmetic projection and the full
residual cohomology estimate on these intervals remain open
in the argument.

\subsection{Nonvanishing of the terminal Bessel period}\label{nav:I-58}
Fix a primitive Dirichlet character, one of the sufficiently
distant intervals of Theorem~\ref{thm:lowesttwo261}, and any
one of its admissible fields. Keep the original horizontal
spaces \(V_{i,0}\), core spaces \(C_i\), congruence levels
and the characters of Proposition~\ref{prop:parity262}.
Let \(m\ge4\) be the even chain length on this field.

\begin{theorem}[Nonvanishing at the prescribed terminal character]\leavevmode
\label{thm:fullperiod266}
Given a unit vector in the lowest compact type of a prescribed
old cuspidal row, one can choose a terminal integral class
coefficient \(e\), a descending core chain starting with \(e\),
and a horizontal chain starting with that row vector such that
\[
                    \mathcal B^+(e,B_m)\ne0 .
\]
Here \(\mathcal B^+\) is the global bilinear Bessel period
with the original raw terminal flag, middle vector, character
and embedded core of Lemma~\ref{lem:topchar266}. After the
two initial lifts of Theorem~\ref{thm:lowesttwo261}, each
horizontal edge can have a single original coset indicator
at every marked place and a single product of real Hermite
functions. Away from the marked places it uses the original
lattice vacua. The coefficient \(e\) and the entire descending
chain are the ones selected for the second restriction.
\end{theorem}
\begin{proof}
Theorems~\ref{thm:downcore254} and~\ref{thm:lowesttwo261}
give the terminal coefficient, descents \(A_i\), and initial
tests for a compact-finite \(B_2\) with
\(\int_{[U(C_2)]} A_2(h) B_2(h\oplus1)\,dh\ne0\).
For the staircase of length one the Bessel radical is
trivial, so this is the required even-index period at \(i=2\).
Theorem~\ref{thm:third263} yields a compact-finite cusp form
\(B_3\) with nonzero Jacobi period for the distinguished
integral line vacuum. Equivalently one may use
Theorem~\ref{thm:jactransfer266} at \(r=0\).

Suppose first that the current index \(i=2r+1\ge3\) is odd.
The induction hypothesis gives the nonzero Jacobi period for
the distinguished test. Since \(A_i\) is the specified
descent of \(A_{i+1}\), Theorem~\ref{thm:bestransfer265}
gives the next nonzero Bessel period, with the same terminal
coefficient. Take its individual Hermite test. Its output is
cuspidal and compact-finite at the prescribed level, as
required at the next step. If \(i=2r+2<m\) is even, the
input is that nonzero Bessel period, and
Theorem~\ref{thm:jactransfer266} gives the next nonzero Jacobi
period, now for \(q_{r+1}\), with a compact-finite cuspidal
output. At each step the previous tests are fixed. This
finite induction reaches \(i=m\), with the Bessel period
paired against \(A_m=e\). Lemma~\ref{lem:topchar266} puts it
in the required raw character convention, using only the
last horizontal test. All the constructions took place at
the field and finite levels fixed at the start.

The finite costs in the two transfer estimates can be tracked
along this chain. For clarity let \(M\) be the largest
\(E\)-dimension of any row or core space. At edge \(i\)
let \(r_i,\ell_i\) be the two flag lengths of
Proposition~\ref{prop:positions264}. The number of marked
terms is at most
\((N\mathfrak m)^{(\ell_i+r_i)M+M}\), including the full
line partition on the odd-index edges. The reciprocal
finite orbit measure is at most
\((N\mathfrak m)^{2\ell_iM}\) by
Theorem~\ref{thm:flagfinite264}. The reciprocal squared norm
of the finite zero-coset line on an even-index edge is at
most \((N\mathfrak m)^M\) when each undeepened marked
position lattice has volume one and the unmarked vacua
have norm one. This follows by integrating its indicator.
There are fewer than \(m\) edges under consideration and
\(r_i,\ell_i\le m/2\). Hence the product of these three
finite costs over the chain is at most
\((N\mathfrak m)^{3m^2M}\). The separate base scalars
\(c_{\ell_i,f}\) and their measure conventions are given
exactly in Theorem~\ref{thm:flagfinite264}.
\end{proof}

In the assertion just proved \(e\) is integral and the period is
a nonzero complex number. A uniform lower estimate for the
designated integral projection in Corollary~\ref{cor:original200}
requires a comparison in its fixed arithmetic normalization.
For example, the Hermite selection in each transfer proves
nonvanishing without bounding the index or the normalized
coefficient of the selected Hermite function. The uniform
strict cohomology sum in \eqref{eq:region198} is also open.
These estimates are needed to reach the prime-sum bound
\eqref{eq:window} for Main Theorem~\ref{thm:grh16}.

\subsection{Normalized Bessel functionals under a weight shift}\label{nav:VII-54}
Keep the Euclidean Witt frame, dual vector \(\omega\), evaluation
map \({\rm ev}\), and spectral projections of
Proposition~\ref{prop:cofactor286}. Give the dual defining space
the invariant Hermitian form dual to \(h\), so pairing against
\(\omega\) evaluates at \(U_1\). Give its tensor product with
the row the resulting invariant Hermitian form \([\, ,\,]\).
All positive Hilbert norms below belong to the discrete series
or the compact column. Use the same Haar measures for every
row, and write \(d_\mu\) for its formal degree.

\begin{proposition}[The coefficients for a single lowering step]\leavevmode
\label{prop:transfer286}
Start with a row of signs \eqref{eq:src109} and parameters
strictly interlacing the compact column in the order
\[
 u_1>\mu_1>u_2>\cdots>\mu_{k-1}>u_k>\mu_k>\cdots>\mu_r .
\]
Suppose the row gaps exceed two. Write
\(P_i=\iota_i p_i\), where \(\iota_i\) includes the shifted
unitary row and \(p_i\iota_i=1\). The maps can be normalized
so that \([\iota_i y,\iota_i y]=\epsilon_i\|y\|^2\) with
\(\epsilon_i\in\{1,-1\}\). Take \(\Lambda_\mu\) in the
normalization of Theorem~\ref{thm:rel155}. If
\(\Lambda_\mu\,{\rm ev}\,\iota_i\) is nonzero, it equals
\(\beta_i\Lambda_{\mu-e_i}\) in the same normalization, and
\[
 |\beta_i|^2\,\frac{d_\mu}{d_{\mu-e_i}}
 =\frac{(2\pi)^{2s}}{2}\,
 \frac{\prod_{a=1}^k|\mu_i-1/2-u_a|}
      {\prod_{j\ne i}|\mu_i-1-\mu_j|}.
\]
If \(\mu_i-1/2\) is a column parameter, the map is zero.
All nonzero terms retain the strict interlacing above. On the
original fields and intervals, after at most \(4r^2\)
successive lowerings, each displayed coefficient is bounded by
\((2\pi)^{2s}(1+2L)^2\), where \(L\) bounds the absolute row
and column entries and may be taken to satisfy \(\log(2+L)=
O(x^\varepsilon)\).
\end{proposition}
\begin{proof}
Different summands are orthogonal for the invariant tensor
form. For example, the split Casimir of
Proposition~\ref{prop:cofactor286} is symmetric for this form
and has distinct real eigenvalues. Its symmetry follows by
taking the difference of the three quadratic Casimirs for
invariant forms, each symmetric as the real infinitesimal
generators are skew adjoint. The tensor form is nondegenerate
on the compact-finite module, so its restriction to each
summand is nondegenerate. An invariant Hermitian form on an
irreducible unitarizable Harish--Chandra module is a real
scalar multiple of its Hilbert form, by Schur's lemma applied
to the maps into the conjugate compact-finite dual. Rescale
the inclusions to obtain the indicated signs.

Write \(L_i v=p_i(\omega\otimes v)\) and \(R_i y={\rm ev}\,\iota_i y\).
Both commute with \(B\). In the Hilbert forms they satisfy
\(\langle R_i y,v\rangle_\mu=\epsilon_i\langle y,L_i v\rangle_i\)
on compact-finite vectors. This follows by pairing \(\iota_i y\)
with \(\omega\otimes v\) in the invariant tensor form.
Continuity gives the equality on smooth vectors, since the
inclusions and projections act on smooth globalizations.
We use the convention that the Hilbert pairing is linear in
the first variable in this proof.

To track the modulus in passing from Hilbert vectors to Bessel
functionals, include a dual compact column vector and write
\(\ell_\nu(v\otimes z)=z(\Lambda_\nu v)\). In the proof of
Theorem~\ref{thm:rel155}, the integral of a matrix coefficient
over the period group against its character is precisely
\(\ell_\nu(a)\overline{\ell_\nu(b)}\). The compact column and
the measures are the same on the two sides. Extend \(R_i,L_i\)
by the identity on its dual. Since they commute with the period
group, the adjoint equality holds for the integrands and gives
\[
 \ell_\mu(R_i y)\overline{\ell_\mu(v)}
       =\epsilon_i\,\ell_i(y)\overline{\ell_i(L_i v)}
\]
whenever the shifted Bessel model exists. Here \(\ell_i\)
denotes the normalized functional on \(\mu-e_i\).
On compact-finite test vectors the integrals apply directly.
The maps preserve compact finiteness, since contraction or
insertion of a vector in a finite-dimensional compact
representation enlarges its compact span by a finite factor.
The matrix-coefficient integral is the one justified in
Theorem~\ref{thm:rel155}. Smooth density then gives the
display on all smooth inputs.

Multiplicity at most one for these real Bessel models
(the local theorem used in Theorem~\ref{thm:src109}) gives
\(\Lambda_i L_i=\kappa_i\Lambda_\mu\), since \(L_i\) is
\(B\)-equivariant. The last display with a nonzero
\(\ell_\mu(v)\) gives
\(\Lambda_\mu R_i=\epsilon_i\overline{\kappa_i}\Lambda_i\).
If the shifted model is absent, \(B\)-equivariance gives
\(\Lambda_\mu R_i=0\) directly. In the former case the signed
scalar of Proposition~\ref{prop:cofactor286} equals
\(\epsilon_i|\kappa_i|^2\), by composing the two maps. In
the latter case that scalar is zero. In particular a zero
of the numerator kills the entire outgoing map. When the map
is nonzero, \(|\beta_i|=|\kappa_i|\).
By Theorem~\ref{thm:norm91}, the ratio of formal degrees is
the ratio of absolute Vandermondes of the two rows, since
the measure factors cancel. Only the factors at position \(i\)
change. Multiplying by their ratio gives the asserted formula.

A row coordinate is integral and every column coordinate is
half-integral. A unit downward step can cross a column
coordinate only when \(\mu_i-1/2\) equals that coordinate.
Thus every contributing shift keeps the column order fixed.
The original adjacent row gaps exceed
\(8(f+1)(r+1)^3\) by the spacing in Section~\ref{sec:absolute163}.
After \(4r^2\) unit steps these gaps still exceed two
for sufficiently large \(x\). The original size bound follows
from \eqref{eq:letters111} and its endpoint formulas.
Increasing it by \(4r^2\), which is polynomial in \(x\), preserves
\(\log(2+L)=O(x^\varepsilon)\).

For the numerical bound set \(t=\mu_i-1/2\) and
\(a_j=\mu_j+1/2\). The numerator distances are \(|t-u_a|\),
and the denominator distances are \(|t-a_j|\) for \(j\ne i\).
After omitting a zero numerator assume first that \(i<k\).
We have \(u_{i+1}<t<u_i\). For \(2\le a\le i\) pair \(u_a\)
with \(a_{a-1}\ge u_a\). For \(i+1\le a<k\) pair \(u_a\)
with \(a_a\le u_a\), using half-integrality and
\(\mu_a<u_a\). Each paired denominator distance is at least
the numerator distance and the paired indices are distinct
from \(i\). Only the two endpoints remain. For \(i\ge k\)
we have \(t<u_k\) and can pair \(u_a\) with \(a_{a-1}\) for
\(2\le a\le k\), leaving only \(u_1\).
All unused denominator distances are at least one by the row
gaps. Each unpaired numerator distance is at most \(1+2L\).
This proves the bound, including equality cases in the
distance comparisons.
\end{proof}

\subsection{A terminal period bound from the second restriction}\label{nav:V-40}
Fix the original field, levels, unit lowest-type row,
terminal integral coefficient \(e\) and descending core
tests of Theorem~\ref{thm:lowesttwo261}. Keep the selected
second row \(B_2\) and its nonzero restriction \(P_2\).
Let \(m\ge4\) be the even row length. Fix the line tests on
the subsequent Jacobi edges as in Theorem~\ref{thm:fullperiod266}.
Fix their real coordinate and measure conventions as in
Proposition~\ref{prop:weighted270}. Let \(M_2,H_2\) be the
two compact spectral parameters for the initial \(B_2\).
For instance these can be taken from the compact orbit of
its second horizontal oscillator test, since the theta lift
intertwines that action. These parameters require only the
finite compact types and the specified infinitesimal character.

\begin{theorem}[A lower bound along the whole horizontal row]\leavevmode
\label{thm:period270}
Choose an integer \(w_m\ge0\). In descending order from
\(m-1\) to \(3\) put \(w_i=E_i(w_{i+1})\), and put
\(s_2=4(w_3+d_2+1)+4g_2+4\).
On these same data one can select one normalized product
Hermite and one original marked choice at each edge after
the second row. The resulting raw terminal Bessel period
obeys
\[
 \frac{|\mathcal B^+(e,B_m)|}{(1+D_m)^{w_m}}
 \ \ge\
 \frac{|P_2|}{M_2 H_2^{s_2}}
       \prod_{i=2}^{m-1}
       \frac{c_{i,f}v_{i,f}t_i}{2 T_i C_i(w_{i+1})}.
\]
Here \(D_m\) is the total degree of the final test in the
original terminal coordinates. All constants in the product
are local data of the specified edges. Taking \(w_m=0\)
gives a lower bound for the complex period in terms of the
initial restriction with no intermediate Hermite degrees
or periods on its right-hand side.
\end{theorem}
\begin{proof}
At edge two take the test given by
Proposition~\ref{prop:weighted270} with \(w=w_3\).
Denote by \(D_3\) the degree selected there. The result is
a lower bound for \(|P_3|/(1+D_3)^{w_3}\) by
\[
       \frac{|P_2|}{M_2H_2^{s_2}}\,
       \frac{c_{2,f}v_{2,f}t_2}{2T_2 C_2(w_3)} .
\]
In particular \(P_3\ne0\). Suppose one has reached row
\(i\) with its tests fixed and with this kind of bound for
\(|P_i|/(1+D_i)^{w_i}\). Apply the same proposition with
input degree \(D_i\) and weight \(w_{i+1}\). The degree
exponent required on the input is exactly \(w_i\).
The lower bound just obtained is multiplied by
\(c_{i,f}v_{i,f}t_i/(2T_i C_i(w_{i+1}))\), and now bounds
\(|P_{i+1}|/(1+D_{i+1})^{w_{i+1}}\). This selects one actual
test and preserves all earlier choices. Finite induction
proves the displayed result. The final raw character is
included before the last expansion by Lemma~\ref{lem:transfer270}.
Every step uses the finite partitions and unmarked vacua of
the original field. The descent ending in \(e\) was fixed
throughout.
\end{proof}

The right-hand side can be combined directly with the
quantitative estimate for \(P_2\) in
Theorem~\ref{thm:lowesttwo261}. Its factors \(C_i(w)\)
involve derivatives of real frame extensions and of fixed
oscillator vectors at the indicated orders. Their growth
as the field, ranks and length vary on the common interval
has yet to be bounded on the required scale. A separate
identity or estimate must relate this complex period to
the specified integral character projection in
Corollary~\ref{cor:original200}. The resulting comparison
and the upper bound for the cohomology mean in
Theorem~\ref{thm:average269} are both needed by the proposed
proof of Main Theorem~\ref{thm:grh16}.

\subsection{A terminal lower bound from the original bottom vector}\label{nav:V-41}
Take the original fields and levels of
Theorem~\ref{thm:lowesttwo261}, a unit lowest-type bottom row
\(u\), and a fixed pure lowest-type core vector \(v_0\).
Choose its polynomial even core tests as in that theorem.
Use \(B,H_*,N,d_0,a_0,d_\rho,\delta,C_{\rm fix}\) from
Theorem~\ref{thm:individual259}, in the terminal probability
normalization. The letters \(J,K_f,d_\sigma,d_\pi\) have the
measure conventions of Proposition~\ref{prop:poincarenorm261}.
That proposition provides a polynomial \(P_{\psi,\varepsilon}\)
for this family. Define the positive quantity
\[
 Z_0=
 \frac{e^{-\delta/2}\sqrt{C_{\rm fix}}}
      {H_*\sqrt{d_\rho}\,N^{d_0} B^{2a_0}}
 \frac{\operatorname{vol}(J)}{d_\sigma}
 \left(\frac{d_\pi}{\operatorname{vol}(K_f)}\right)^{1/2}
 \|u\|_\infty e^{-P_{\psi,\varepsilon}(x)} .
\]
In Lemma~\ref{lem:moment271} bound \(W(a_\tau)\) by
\(T_0^{C n^3}\) as in Proposition~\ref{prop:poincarenorm261},
giving \(\eta,S\) from prescribed data. Fix the subsequent
real core and line tests. Choose \(w_m\ge0\) integral,
use \(w_i,s_2\) of Theorem~\ref{thm:period270}, and set
\(q=d_1+s_2\).

\begin{theorem}[From a unit bottom row to the raw Bessel period]\leavevmode
\label{thm:bottom271}
There is a terminal integral coefficient \(e\) and a chain of
individual horizontal tests at the original finite levels for
which
\[
 \frac{|\mathcal B^+(e,B_m)|}{(1+D_m)^{w_m}}
 \ \ge\
 \frac{Z_0}{2S H_0^{s_2}(1+L_1+q/\eta)^q}
 \prod_{i=2}^{m-1}\frac{c_{i,f}v_{i,f}t_i}{2T_i C_i(w_{i+1})}.
\]
Here \(H_0,L_1\) are from Proposition~\ref{prop:spectrum271}.
In particular at \(w_m=0\) the right side involves only the
original bottom vector, field, group parameters, prescribed
local tests and measures. It is independent of every selected
Hermite index and of every individual intermediate period.
\end{theorem}
\begin{proof}
Choose the real translate \(a\) from
Proposition~\ref{prop:poincarenorm261}. With
\(\ell_a(F)=\mathcal B(R(a)F,u)\), that proposition and
Theorem~\ref{thm:individual259} give one \(e\) such that
\(\left|\ell_a(P_{v_0}T^\dagger e)\right|\ge Z_0\).
Keep this same \(e\). Lemma~\ref{lem:project271} identifies
this scalar with the first restriction on
\(\omega(a)\phi^0\). Expand that test in the original bottom
and line Hermites. For each product in its expansion the
two-kernel identity of Theorem~\ref{thm:secondbound260}
expresses the first restriction as
\(D_L^{-1}\sum_\nu\epsilon_\nu J_{e,\alpha,\nu}\), where
there are \(D_L\) terms and \(|\epsilon_\nu|=1\).
Schwartz convergence permits passage from finite Hermite
sums to the first restriction by
Proposition~\ref{prop:projecttests260}. If all individual
terms satisfied
\[
       |J_{e,\alpha,\nu}|<
              (Z_0/(2S))\,\exp(\eta|\alpha|),
\]
the triangle inequality and Lemma~\ref{lem:moment271} would
bound the modulus of every finite sum by \(Z_0/2\).
The limit contradicts its lower bound. Thus one pair
\(\alpha,\nu\) gives at least the displayed size. Set
\(P_2=J_{e,\alpha,\nu}\) for this choice. Its two earlier
horizontal tests and \(e\) are now fixed.

Put \(j=|\alpha|\). Proposition~\ref{prop:spectrum271} gives
\[
       M_2 H_2^{s_2}\le H_0^{s_2}(1+L_1+j)^q .
\]
For every real \(j\ge0\),
\((1+L_1+j)^q e^{-\eta j}\le(1+L_1+q/\eta)^q\).
Indeed the left side decreases for \(j\ge q/\eta\), and
below that point the power alone is at most the right side.
Hence \(|P_2|/(M_2H_2^{s_2})\) is at least the scalar before
the product in the theorem. Apply the proof of
Theorem~\ref{thm:period270} starting at this nonzero
restriction. That proof uses only its compact types and
nonvanishing, so it applies to the selected \(\alpha,\nu\).
It yields the asserted individual tests, measures and raw
terminal character. Finally \(\|u\|_\infty\) may itself be
replaced below by \(\operatorname{vol}([H])^{-1/2}\), since
\(u\) has unit \(L^2\) norm.
\end{proof}

The weights in this estimate satisfy \(w_i\ge4w_{i+1}+1\)
for \(3\le i<m\). Thus the derivative orders requested of
the fixed local data already grow at least geometrically
with the number of transfers. Establishing the needed
uniform scale, the upper bound for the full theta period
and the lower bound in the designated integral projection
remains open on the intervals of Section~\ref{sec:missing}.

\section{Restriction, levels and Fourier coefficients}\label{nav:g-VII-4}

\subsection{The initial degrees and the second row spectrum}\label{nav:IX-43}
Use the first two horizontal edges of
Theorem~\ref{thm:secondbound260}. A Hermite product on the
first core edge factors as \(h_\alpha=p_\alpha\otimes q_\alpha\)
in the bottom hyperplane and line. The second horizontal
test at infinity is \(b_1\otimes Uq_\alpha\), where \(b_1\)
is the fixed second core polynomial and \(U\) is the
opposite-line intertwiner. Let \(L_1\) be the total energy
of \(b_1\) across the real places. Use signed orthonormal
Hermitian bases respecting the two direct sums in
Proposition~\ref{prop:cancelsquare256}. Write \(d_1\) for
the total real position dimension on this second horizontal
edge, whose incoming index is \(1\).

\begin{proposition}[Compact spectral bounds before the remaining transfers]\leavevmode
\label{prop:spectrum271}
For every nonzero \(B_{2,\alpha,\nu}\) in that construction
the \(M_2,H_2\) of Proposition~\ref{prop:kernel270} can be
chosen to satisfy
\[
 M_2\le(1+L_1+|\alpha|)^{d_1},\qquad
 H_2\le H_0(1+L_1+|\alpha|).
\]
The constant \(H_0\) can be fixed from the real group, its
specified infinitesimal character, and its compact action on
the creation space and vacuum before choosing \(e,\alpha,\nu\)
or the bottom cusp vector.
\end{proposition}
\begin{proof}
At a real place the adjoint on the line summand sends a row
of complex coordinates to a conjugate column times signature
signs. The line generators have respective norms \(1,-1\).
Thus it is an isometry for the positive Hilbert--Schmidt
metrics, formed by making every signed basis vector of length
one. The sign switch identifies the alternating forms as in
Proposition~\ref{prop:cancelsquare256}. Both normalized
Heisenberg metrics use the same absolute trace scalar of the
additive character. An isometry carrying one symplectic form
to the other carries the positive complex structure to the
other, by the polar decomposition of its skew operator. Hence
the unitary change of real oscillator model intertwines the
two number operators. This also follows by transporting the
sum of squares of positions and frequencies, then subtracting
the same vacuum energy. The splitting scalar has no effect.
Thus the energy on the second test is at most \(L_1+|\alpha|\).
The compatible compact of the second row acts on this space
preserving the total number operator, since the signed bases
on the core and new line together are a signed orthonormal
basis for that row.

Set \(h=L_1+|\alpha|\). There are at most \((1+h)^{d_1}\)
multi-indices of degree at most \(h\), since each individual
coordinate is between zero and \(h\). Equivariance in the
second horizontal test puts every compact type of its output
among those of this polynomial space. This bounds the sum
of the distinct type dimensions. To see the dependence of
the eigenvalues explicitly, let \(A\) bound the operator norms
of the compact orthonormal Lie generators on the creation
space and let \(B\) bound the absolute values of their
vacuum scalars. Their norms through degree \(h\) are at most
\(B+Ah\) by the symmetric tensor action. Thus
\(C_K(\tau)\le(\dim K)(B+Ah)^2\) for each type in the list.
One may take
\[
       H_0=1+\sqrt{|\Omega_G(\pi_2)|}
                       +\sqrt{2\dim K}\,(A+B),
\]
with the metric and Casimir of Proposition~\ref{prop:kernel270}.
Theorem~\ref{thm:realrow257} fixes its scalar
\(\Omega_G(\pi_2)\) on the entire second row. This gives the
claimed \(H_2\) bound for every allowed finite residue.
\end{proof}

\subsection{Integral coefficient lattices on the graph correspondence}\label{nav:IX-44}
Use the coefficient representation \(W\) of
Proposition~\ref{prop:terminalcore255}. For \(b\in H(\mathbb A_{F,f})\)
put \(\Lambda_b=C_m(E)\cap b\widehat\Lambda\), where \(C_m(E)\)
here denotes the underlying \(E\)-vector space of the core.
Write \(W(\Lambda_b)\) for the tensor over the chosen embeddings
of its integral Young images and determinant powers. Enlarge the
coefficient field to include all finite character values. A common
choice for the right argument of a matrix coefficient is the fixed
lattice \(W^\circ=W(\Lambda_1)\). At a class represented by \(b\)
take the left covector in \(W(\Lambda_b)^\vee\).
This convention respects left rational translation, which acts
only on the left argument. It also defines a lattice at any finer
finite level contained in the stabilizer of \(\widehat\Lambda\).

\begin{proposition}[A common integral lattice for the translated columns]\leavevmode
\label{prop:corelattice272}
There is a list \(\mathcal Q\) of \(h_mN^{2d_0}\) integral
columns at the original level. For every vector \(Y\) of its
full real type and finite level, some \(Q\in\mathcal Q\) satisfies
\[
       |\langle Q,Y\rangle|\ge
       \frac{\|Y\|}{h_m\sqrt{d_\rho}\,N^{d_0} B^{2a_0}} .
\]
For the graph elements of Proposition~\ref{prop:corehecke272},
\(T^{a_0}D_\lambda Q\) is an integral column in the same
fixed-right-argument convention at the intersection of the
original level with all its conjugates by \(za\), \(z\in Z\).
Both assertions hold also for the conjugate coefficient representation.
\end{proposition}
\begin{proof}
Keep the \(N^{d_0}\) covectors of
Proposition~\ref{prop:schurframe259} for each class and use its
\(N^{d_0}\) vectors on the identity class for the right list at
every class. Their lower frame bounds are both \(B^{-2a_0}\).
The upper bounds are \(B^{2a_0}\) for the covectors and
\(L_0^2B^{2a_0}\) for the vectors. Form all products and extend
each class coefficient by zero on the other classes. The class
stabilizers are trivial by Proposition~\ref{prop:terminalcore255}.
With the probability on \([H]\), compact Schur orthogonality
divides the tensor-product frame operator by \(h_m d_\rho\).
Thus the sum of all squared pairings with \(Y\) is at least
\(B^{-4a_0}\|Y\|^2/(h_m d_\rho)\). Divide by the list length
and take square roots.

For \(u=za\), the elementary divisors of \(u\) at the chosen
split half are units and one \(T\), and those at the conjugate
half are their inverse duals. At all other places \(u=1\).
Consequently for every \(b\)
\[
          T\Lambda_b\subset\Lambda_{bu}\subset T^{-1}\Lambda_b .
\]
At an embedding of degree-\(d_\tau\) polynomial tensors, the
two inclusions with exponent \(d_\tau\) pass to the image of
the same integral Young operator. In the determinant line the
corresponding exponent is at most \(N\), by taking top
exterior powers. Dualization reverses inclusions and gives
the same two-sided estimate for negative determinant powers.
Tensoring over the embeddings gives
\(T^{a_0}W(\Lambda_b)\subset W(\Lambda_{bu})
\subset T^{-a_0}W(\Lambda_b)\), since
\(a_0=\sum_\tau(d_\tau+N|j_\tau|)\).

Write \(bu=\gamma t_c k_f\) with \(\gamma\in H(F)\), \(t_c\)
the old representative, and \(k_f\) in the old level. If
\(Q\) has support on \(c\), its translate at \(b\) has real
value \(l_c(\rho(\gamma_\infty^{-1}k)v)\), where \(k\)
is the real argument and \(v\) is its fixed right vector.
The left covector is integral on \(W(\Lambda_{bu})\), since
\(\Lambda_{bu}=\gamma\Lambda_{t_c}\). The displayed inclusion
therefore makes \(T^{a_0}\) times this covector integral on
\(W(\Lambda_b)\). If the support condition fails, the value
is zero. The common intersection fixes every translated form,
and character values are units in the enlarged coefficient
ring. Summing the translated columns proves integrality.
Complex conjugation preserves inclusions, norms and the
integer \(T\), proving the last assertion.
\end{proof}

\subsection{Restriction of a Fourier coefficient to a smaller row}\label{nav:IX-45}
In the embedding of Theorem~\ref{thm:graph274}, let \(I\) be
the isotropic flag space. The large middle space has the
orthogonal decomposition \(W=W_1\perp D\), with
\(W_1=W_0\perp Ez'\) belonging to the small row.
Set \(X_D=\operatorname{Hom}_E(D,I)\), regarded as an
additive \(F\)-group. For \(x\in X_D(F_\infty)\) let
\(n(x)\) be the real unipotent having this matrix as its
middle-to-flag block and zero skew-adjoint parameter. Its
finite components are one. Let \(\eta(x)=\chi_+(n(x))\).
Use probability measures on the compact rational unipotent
quotients. A box \(B_R\) is \([-R,R]^{\dim_{\mathbb R} X_D(F_\infty)}\)
in any fixed real linear coordinates, with Lebesgue volume \(|B_R|\).

\begin{proposition}[Averaging the small Fourier integrals]\leavevmode
\label{prop:fiber274}
For any smooth automorphic function \(B\) on the large group
and \(g\) in its adelic points, one has
\[
 \int_{[N]} B(ng)\overline{\chi_+(n)}\,dn
  =\lim_{R\to\infty}\frac1{|B_R|}\int_{B_R}\overline{\eta(x)}
      \int_{[N_0]} B(n_0 n(x)g)\overline{\chi_\varpi(n_0)}
                                            \,dn_0\,dx .
\]
In particular some real \(x\) gives a small integral of
modulus at least half the modulus of any nonzero left side.
\end{proposition}
\begin{proof}
In the ordering \(I,W,I^\vee\) the radical fixing \(I\)
pointwise has block matrices
\[
   \begin{pmatrix}
     1&X&Z-XX^\ast/2\\
     0&1&-X^\ast\\
     0&0&1
   \end{pmatrix},\qquad Z^\ast=-Z .
\]
The adjoint uses the Hermitian form on the middle space
and the pairing of the two outer spaces. The isometry condition
on the upper right entry is precisely that its sum with its
adjoint equal \(-XX^\ast\). The full flag radical also has
the upper unitriangular group on \(I\), acting dually on
\(I^\vee\). The subgroup \(L\) given by \(X=0\) and arbitrary
\(Z\) is normal and belongs to both radicals. Both characters
are trivial on \(L\). Modulo \(L\), the middle blocks add,
and conjugation by a flag matrix \(A\) sends \(X\) to \(AX\).
Hence \(\bar N=N/L=X_D\rtimes\bar N_0\), where
\(\bar N_0=N_0/L\). The same description holds on rational
points, by the displayed matrix representatives.

Average \(B(ng)\) on the left over \([L]\) and denote the
resulting function on \([\bar N]\) by \(b(\bar n)\).
Translations in \(L(\mathbb A_F)\) preserve its probability,
and rational conjugation preserves it by normality. This
proves that \(b\) is well defined. It is continuous and
bounded because all the unipotent quotients here are compact.
The quotient measure in \([\bar N]\) integrates first over
\([X_D]\) on the left of a representative of \([\bar N_0]\).
Indeed the semidirect product coordinates give a fundamental
set from fundamental sets for the two rational subgroups.
The flag matrices have determinant one in their action on
\(X_D\), preserving additive Haar measure. Normalization gives
the two probabilities. Averaging over \([L]\) in \(N_0\)
uses the same probability as in \(N\).

For fixed \(\bar n_0\), let \(A\) be its flag matrix. Then
\(\bar n_0\,\bar n(x)=(A_\infty x)\bar n_0\) in the quotient.
Since \(\chi_+\) is a character, its restriction \(\eta\)
to \(X_D\) is invariant under this conjugation, is trivial
on rational points and extends the displayed real character.
The continuous function
\(y\mapsto b(y\bar n_0)\overline{\eta(y)}\) on \([X_D]\)
therefore occurs on the right through \(y=A_\infty x\).
The box mean along this real subspace tends to its Haar
mean. In fact the characters of the compact additive
quotient are indexed by its rational dual under trace.
Each nontrivial one has nonzero real frequency on the full
archimedean space, also after the invertible transformation
\(A_\infty\). Integration in a coordinate of nonzero frequency
makes its box mean tend to zero. Finite sums of characters
are uniformly dense in the continuous functions by
Stone--Weierstrass. This proves the convergence for each
\(\bar n_0\). The bound \(\|b\|_\infty\) permits integration
over \([\bar N_0]\) by dominated convergence. The small
character is the restriction given in Theorem~\ref{thm:graph274}.
Restoring the common \(L\)-integration gives the identity.
Absolute values and the triangle inequality prove the final
assertion.
\end{proof}

\subsection{The degree required at the high real places}\label{nav:IX-46}
At a real place \(\tau\) write \(D_\tau\) for the degree of
the final Hermite test in the compatible coordinates of
Corollary~\ref{cor:real274}. For a product test the total
degree in Theorem~\ref{thm:period270} is
\(D_m=\sum_\tau D_\tau\). We consider the unrestricted
horizontal lift before projecting to the special large
compact type \(\rho\).

\begin{proposition}[A necessary Hermite degree for the full column]\leavevmode
\label{prop:degree275}
In the restrictions \(y\mapsto B_m(y n(x)g a)\) of
Theorem~\ref{thm:small274}, every occurring real column
highest weight \(\mu_\tau\) satisfies
\[
         \mu_{\tau,1}-\mu_{\tau,k}\ \le\ D_\tau .
\]
The bound survives the terminal character sum, the small
Fourier integral, compact projections on the large core
before restriction or on the column afterwards, and any
determinant twist. For the prescribed full word, every high
place must therefore have \(D_\tau\ge a+b_*\). Hence
\(D_m\ge h(a+b_*)\). A test in the span of all total degrees
less than this last threshold gives zero after the full
word projection for every such restriction.
\end{proposition}
\begin{proof}
Consider in the real column \(U(k)\) the circle
\(\operatorname{diag}(e^{it},1,\ldots,1,e^{-it})\). Its
weights on the defining representation and its dual are
\(1,0,-1\), with the zero weight omitted when \(k=2\).
On the complexified real symplectic space of the last
oscillator its weights are drawn from these three integers.
Indeed this space comes from the tensor of the last two
Hermitian spaces. The circle acts on the endpoint factor
through a copy of the defining representation, or its dual,
and trivially on the complementary coordinates. It is
trivial on the other tensor factor. Complexification of the
underlying real tensor adds the conjugate representation
and gives the same three possibilities.

For a positive compatible metric the creation space is an
invariant subspace of the complexified symplectic space.
The vacuum spans a line under the compact group. Its
character, including the chosen splitting, is trivial on the
displayed circle because it lies in \(SU(k)\), a connected
group with semisimple Lie algebra. The homogeneous Hermites
of degree \(D_\tau\) are spanned by length-\(D_\tau\)
products of creation operators applied to this vacuum.
Each adds a weight of absolute value at most one. Thus the
largest circle weight is at most \(D_\tau\).
The conclusion also holds on every conjugate of the circle
in the large definite \(U(N)\), by equivariance of the
compatible metric.

The map from the final oscillator test to the horizontal
lift intertwines the real endpoint action. In the small
restriction the real column commutes with \(n(x)\) and
acts through its conjugate by \(g_\infty\) inside the large
definite group. The terms at the finite places commute with
all these actions. Applying the small Fourier integral
intertwines the column translations. Each successive
compact projection mentioned in the statement stays in the
corresponding degree space or its image. A determinant
character is trivial on each of the circles. If an
irreducible column type occurs, its highest vector has
circle weight \(\mu_{\tau,1}-\mu_{\tau,k}\), which proves
the inequality. At the high places the word has width
\(a+b_*\) by Theorem~\ref{thm:word275}. A product test of
smaller total degree misses this width at some high place.
Linearity on homogeneous components proves the last assertion.
\end{proof}

The lower bound in Theorem~\ref{thm:word275} uses the compact
type paired with the integral large-core column. Its high-place
projection vanishes identically. A different component of the
full real theta lift can enter there only at the degrees just
estimated. A quantitative pairing with that component and
with the prescribed global cuspidal row at the original
levels is still required for the integral comparison of
Corollary~\ref{cor:original200}.

\subsection{The finite lattices in the restricted group}\label{nav:IX-47}
Fix the rational embedding of Theorem~\ref{thm:graph274} on a
field of Theorem~\ref{thm:conductor232}. Write \(S=S_m\),
\(C=C_m\), \(H=U(C)\), and let \(\Lambda_0,\Lambda_W\) be the
original signed row lattice and its definite column lattice.
At a finite place \(v\) of \(F=E^+\) let \(\pi_v\) be a
uniformizer of \(F_v\). All ideal inclusions below use the
integer ring of \(E\otimes_F F_v\), with both factors at a
split place. The raw lattice in the staircase is the direct
sum of the lines \(I_j b_j,I_j^{\rm op}b_j^\vee,I_z z\) and
the lattice in \(C\). Here \(b_j^\vee\) is the opposite vector
paired to \(b_j\) with value one. These are the actual local
fractional ideals of Proposition~\ref{prop:positions264}.

\begin{proposition}[A finite translate compatible with the column at all places]\leavevmode
\label{prop:lattices277}
There is a finite adele \(u\in H(\mathbb A_{F,f})\), equal to
one above \(T\), such that \(L_v=u_v\Lambda_{S,v}\) contains
\(\Lambda_{W,v}\) as an orthogonal summand at every finite place,
and \(v_0\) of Theorem~\ref{thm:graph274} belongs to
\(u_v\Lambda_{C,v}\) everywhere. For \(v\nmid T\) write
\(\pi=\pi_v\) and let \(e_v\) be the smallest nonnegative
integer satisfying the inclusions
\[
 \pi^{e_v}\sigma^{j-s-1}/2\in I_j,\qquad
 \pi^{e_v}\bar\sigma^{\,s+1-j}\in I_j^{\rm op}\quad(1\le j\le s),
 \qquad \pi^{e_v}\bar\varpi\in I_z .
\]
Then \(\pi^{e_v}\Lambda_{0,v}\subset L_v\) and
\(\pi^{e_v}\operatorname{pr}_{S_0}L_v\subset\Lambda_{0,v}\).
In particular \(e_v=0\) at every odd unmarked place where
\(\sigma\) and the staircase ideals are units. The choices of
\(u\) and \(e_v\) depend only on the embedded spaces and
lattices.
\end{proposition}
\begin{proof}
The core and small column lattices are Hermitian self-dual.
The extension \(E/F\) is unramified at every finite place by
the inertia argument of Lemma~\ref{lem:mass209}, and its dyadic
places split by Proposition~\ref{prop:tracelat253}. First make
a local model of the prescribed \(W_0,v_0\) inside
\(\Lambda_{C,v}\). At a split place a self-dual lattice is
a free module and its dual under evaluation. Embed the lattice
of \(W_0\) in the first \(k\) coordinates with their duals.
In the remaining coordinates take \((b,(T-1)b^\vee)\), with
\(b^\vee(b)=1\). This integral vector has norm \(T-1\), also
when \(v\) divides \(T-1\).

At an unramified field place of odd residue characteristic,
a self-dual lattice has an orthonormal basis. Its nonsingular
residue form has a vector of nonzero norm, since polarization
and nonsingularity rule out identically zero norm. Lift that
vector and scale the lift to norm one. The norm on units is
surjective by the residue norm and lifting successive digits
by the surjective residue trace. The unit line splits off by
orthogonal projection. Induction gives the basis. Embed the
column by the first \(k\) vectors. The complement has rank at
least two. In two orthonormal coordinates choose \(c\bar c=-1\).
The vectors \(l=w_1+c w_2\) and \(l'=(w_1-c w_2)/2\) in
these coordinates are isotropic and pair to one. The vector
\(l+(T-1)l'/2\) is integral of norm \(T-1\).

In both cases the resulting copy of \(W_0\perp Ev_0\) is
nondegenerate over the local field. Witt extension carries
it by a local unitary map to the given rational copy with the
specified coordinates. Apply this map to the core lattice.
At the terminal places the original core already has the two
properties by Theorem~\ref{thm:graph274}. At almost every
other place the initial rational embedding maps the column
lattice and \(v_0\) integrally into the original core.
Self-duality of the column then splits it off by orthogonal
projection. Thus the maps can be chosen to be the identity
outside a finite set and at all \(v\mid T\). They give \(u\).

In each original signed plane write \(u_j=e_++e_-\) and
\(u_j^\vee=(e_+-e_-)/2\). The integral generators are
\(u_j/2+u_j^\vee\) and \(u_j/2-u_j^\vee\). Under the rational
embedding the two flag vectors go to
\(\sigma^{j-s-1}b_j\) and
\(\bar\sigma^{s+1-j}b_j^\vee\). The middle generator goes to
\(\bar\varpi z+v_0\). The line ideals in the statement and the
integrality of \(v_0\) now put \(\pi^{e_v}\Lambda_{0,v}\)
inside \(L_v\). For \(x\in L_v\) and \(y\in\Lambda_{0,v}\),
pairing \(x\) with \(\pi^{e_v}y\) gives an integer. This is
the pairing of \(\pi^{e_v}\operatorname{pr}_{S_0}x\) with \(y\).
Self-duality of the signed lattice gives the other inclusion.
The formulas make \(e_v\) finite, with finite support.
\end{proof}

\subsection{The finite Fourier coefficient at every horizontal step}\label{nav:VI-41}
Fix \(2\le i<m\) on any of the fields and intervals of
Theorem~\ref{thm:third263}. Use the two position tests of
Proposition~\ref{prop:positions264}. Take a nonzero residual
product \(\beta_f\) on \(\operatorname{Hom}(-T_0,S_0)\) with
the original vacua away from the marked set and one original
integral coset per marked place. Write \(K_f\) for the
original level of \(U(S)\), including all its marked depths.
In the notation of Proposition~\ref{prop:flagunmarked264}
put \(R_f=R(\mathbb A_{F,f})\), where \(R\) preserves the
scaled ordered frame up to an upper unitriangular transformation.
Use the Haar measure giving \(O_{E_v}\) volume one in each
upper entry of \(U_\ell(E_v)\). Let \(d_E\) be the absolute
field discriminant and \(s=\dim_E S\).

\begin{theorem}[Support and value of the finite Fourier coefficient]\leavevmode
\label{thm:flagfinite264}
The restricted product of the local integrals
\(\mathcal E_v(g_v)\) for these tests has nonzero set exactly
\(R_fK_f\). Its value at \(1\) in the residual oscillator is
\[
  \mathcal E_f(1)
    = |d_E|^{\binom{\ell+1}{3}}\,
      (N\mathfrak j)^{\binom{\ell}{3}}\ \beta_f .
\]
At a marked \(v\), let \(p_v\) be the ordered-frame probability
in Lemma~\ref{lem:flagmarked264}. Give the compact orbit of the
full lattice stabilizer in \(R_v\backslash U(S)(F_v)\) probability
one. The measure of the projected level in that orbit is
\((Nv)^{a_v\ell(\ell-1)}p_v\). With these compact-orbit measures
at all finite places, the measure of the support in the quotient
is at least
\[
                (N\mathfrak m)^{-(2s\ell-2\ell^2+\ell)} .
\]
\end{theorem}
\begin{proof}
The unmarked support is Proposition~\ref{prop:flagunmarked264}.
At a marked place the full product test is fixed by \(K(a_v)\).
Indeed its split expansion in Proposition~\ref{prop:positions264}
consists of individual original integral cosets, each preserved
by matrix substitution at that depth. The determinant character
is trivial on this congruence group by
Proposition~\ref{prop:rowfinite262}. If the integral is nonzero
there is a \(u\) such that the columns of \(g^{-1}X_*u\), after
division by \(d^j\), lie in \(b_j+\pi^{a_v}L_{j,v}\).
Here the substitution in the first position is valid for the
whole group, independently of the residual action. The columns
have the original Gram matrix. Indeed the Gram matrix of \(X_*\)
is zero or is supported only at the last diagonal entry, and an
upper unitriangular matrix has last row equal to the last
coordinate functional. Lemma~\ref{lem:flagmarked264} supplies
a \(k\in K(a_v)\) with \(kX_*=g^{-1}X_*u\).
Hence \(gk\in R_v\). Conversely right invariance and the
left covariance proved in Proposition~\ref{prop:flagunmarked264}
apply at the marked places as well. It suffices to calculate
the base value.

At \(1\) the row evaluation leaves \(\beta_v\) unchanged.
In column \(j\), the coefficient of \(b_k\) for \(k<j\) is
\(d^{k-j}u_{kj}\). Thus the first indicators give, independently
in every entry,
\[
   u_{kj}\in\mathfrak j(dAB)^{j-k}O_{E_v}
      =\mathfrak D_{E/\QQ}^{-(j-k)}
         \mathfrak j^{\,1-(j-k)} O_{E_v}.
\]
For \(j=k+1\) this is the inverse different, so the entire
character is trivial there. The integral is the product of
the positive entry volumes times the residual vector. It is
the given vacuum at almost all places. Multiplying finite
volumes of a fractional ideal \(I\) gives \((NI)^{-1}\) in the
specified measures. There are \(\ell-m'\) entries with
\(j-k=m'\). Summing \(m'(\ell-m')\) and
\((m'-1)(\ell-m')\) for \(1\le m'<\ell\) gives
\(\binom{\ell+1}{3}\) and \(\binom{\ell}{3}\), respectively,
by the formulas for the first and second powers of successive
integers. This proves the base value and the product support.
In fact at almost all places the input group element lies in
the full stabilizer, so the factorizations into \(R_vK_v\)
assemble adelically.

It remains to pass from the ordered frame to its orbit modulo
unitriangular transformations. Write \(K\) for the full marked
lattice stabilizer and \(H\) for the pointwise stabilizer of
\(X_*\). The homomorphism \(R_v\to U_\ell(E_v)\),
\(gX_*=X_*u\), has kernel \(H\). Its image on \(R_v\cap K\)
is the subgroup of upper unitriangular matrices with entries in
\((dAB)^{j-k}O_{E_v}\), \(k<j\). Necessity follows by applying
the lattice condition to each of the two split components,
as in the column calculation. Conversely these conditions
make the normalized columns and the normalized rows in
Lemma~\ref{lem:flagmarked264} integral. Within each of the two
ordered spanning sets their change is triangular with unit
diagonal, so both still extend to bases. They have the same
prescribed pairing. The residue-field extension argument and
the integral right-inverse correction in that lemma now give
an element of \(K\) inducing \(u\). The same lemma identifies
the image of \(R_v\cap K(a_v)\) by multiplying each entry ideal
by \(\pi^{a_v}\). Both necessity and sufficiency here use
congruence of the entire ordered frame.

The index of the latter image in the former is
\((Nv)^{2a_v\binom{\ell}{2}}\) by product Haar measure, since
\(E_v=F_v\times F_v\). The preimages have kernels \(H\cap K\)
and \(H\cap K(a_v)\). Thus replacing \(H\) by \(R_v\) multiplies
the marked probability of the congruence orbit by this index.
Indeed that probability for a subgroup \(D\) is
\([D\cap K:D\cap K(a_v)]/[K:K(a_v)]\), by counting cosets
in \(K/K(a_v)\). This proves the asserted local factor.
The ordered-frame count in Lemma~\ref{lem:flagmarked264} is
at most \((Nv)^{a_v(2s\ell-\ell^2)}\), by replacing each
difference of powers by its leading power. Unmarked quotient
probabilities are one. Multiplication over the marked places
gives the stated lower bound.
\end{proof}

The base value here is an identity for the finite part of the
mixed Fourier coefficient. The global period also integrates
the real coefficient against the two automorphic forms. In
particular the lower bound in this theorem measures the finite
orbit. A bound for the integral projected global period still
requires control of that pairing in the prescribed integral
normalization.

\subsection{The prescribed column type at every real place}\label{nav:VII-48}
Give the compact quotients and definite real groups probability
measure. In Theorem~\ref{thm:raw276} put
\(M_B=|P_{\rm raw}|/\|e\|_{L^2([H])}>0\).
Write \(\rho\) for the real type of \(\bar e\).
Let \(\mu\) be the prescribed product of real word types,
\(d_\mu=\dim\mu\), and let \(\nu\) be its product with
\(\prod_\tau\det^{-j_\tau}\) as in
Proposition~\ref{prop:spacing276}, using \(\rho\) in that
proposition. Put \(m_\nu=[\rho|_{H_0(F_\infty)}:\nu]\ge1\).
As before \(Z=\prod_{v\mid T}\mathbb F_T^\times\) and
\(c_T=(3T-1)/(T(T+1))\).

\begin{theorem}[The projected coefficient in the full prescribed real type]\leavevmode
\label{thm:types276}
For each word and each \(\lambda=\prod_{v\mid T}\lambda_v\)
with all \(\lambda_v\ne1\), there is a smooth left rationally
invariant row \(A\) on \(G_0\), fixed by its prescribed terminal
graph stabilizers, whose central character sum at the original
frequency satisfies
\[
 \left(\int_{H_0(F_\infty)}
     |F_0A_\lambda(h,1_f)|^2\,dh\right)^{1/2}
 \ \ge\
 \frac{\sqrt{|Z|}(1-c_T)^{f/2}}{4}
       \sqrt{\frac{m_\nu}{d_\rho d_\mu}}\,M_B .
\]
The function under the norm has precisely real type \(\mu\).
The row is obtained from the modified horizontal lift by
real compact projection, translation, restriction and the
determinant twist of Proposition~\ref{prop:spacing276}.
\end{theorem}
\begin{proof}
The characters of the modified grid are unramified at \(T\).
Thus \(Y=e_\rho F_{\chi_+}B\), where \(e_\rho\) is the real
type projector on the large compact core, is fixed under
each original maximal compact subgroup there. Its pairing
against \(\bar e\) has modulus \(|P_{\rm raw}|\), so
\(\|Y\|\ge M_B\). The compact type \(\rho\) is nontrivial
on a real special unitary factor. In its parameter the
adjacent cloud entries have difference \(4\Delta>1\).
They differ by \(4\Delta-1>0\) in its highest weight.
The same holds after duality. Apply the finite translation
\(D_\lambda=S_\lambda R(a)\) defined in
Proposition~\ref{prop:corehecke272}, with the coordinate
choice of Theorem~\ref{thm:graph274}. That proposition gives
\[
 \|D_\lambda Y\|^2=|Z|\langle Y,L_TY\rangle
             \ \ge\ |Z|(1-c_T)^f M_B^2 .
\]
It applies to the whole real isotypic subspace, independently
of the size of the weights. The geometric embedding of
Theorem~\ref{thm:graph274} uses only the Hermitian spaces,
their lattices, and the additive character. They are the same
for the modified grid. The large character on the small
radical is therefore exactly \(\chi_\varpi\), and the
resulting row has the required terminal stabilizers.

Apply the compact averaging calculation in the proof of
Theorem~\ref{thm:extract275} to \(P=D_\lambda Y\) and the
full \(H_0(F_\infty)\). To recall the constants, choose a
finite component of real norm at least \(\|P\|\).
Schur orthogonality in the two real translates gives a
projected value of modulus at least
\(\sqrt{m_\nu d_\mu/d_\rho}\,\|P\|\).
Rational weak approximation realizes its left rotation
with at least half this modulus and sets \(g_v=1\)
for every \(v\mid T\). The averaging of the small Fourier
integrals in Proposition~\ref{prop:fiber274} then retains
at least half again at one real translate. The small real
projector commutes with the terminal translations. Finally
the determinant twist is trivial on the unipotent and on
the finite scalars. It converts \(\nu\) into \(\mu\) and
preserves the modulus at the identity. This gives a value of
\(|F_0A_\lambda(1)|\) at least
\(\sqrt{m_\nu d_\mu/d_\rho}\,\|P\|/4\).

On the real group at this finite component the function is
in the matrix coefficient space of \(\mu\), of dimension
\(d_\mu^2\). An orthonormal basis of that space has sum of
squared absolute values constant by translation and equal to
\(d_\mu^2\) by integration. Evaluation at one thus costs at
most \(d_\mu\) times the \(L^2\) norm. This proves the
inequality and uses no bound on the Hermite degree of the
selected test.
\end{proof}

The dimension estimate in Proposition~\ref{prop:spacing276}
bounds the logarithmic loss from the real projections by
\(O_{\psi,\varepsilon}(fN^2\log b)\). The positive quantity
\(M_B\) is the normalized raw period for the enlarged
auxiliary weights. Its size at the arithmetic scale of
Section~\ref{sec:missing} still requires an estimate.

\subsection{A bounded translation at one marked place}\label{nav:VII-49}
Let \(H=U(C)\) be the definite group of rank \(N\ge3\).
Choose one marked split place \(v_*\) of residue size \(Q\),
and on its first split factor let \(K_*=\mathrm{GL}(L_*)\)
for the core lattice. Fix a product finite level \(J\) with
\(J_{v_*}\subset K_*\). Let \(\omega>0\) be the smallest
mass of a finite class at level \(J\) on \([H]\) in probability
measure. Each class carries Haar probability on its real
orbit conditional on that class. Put \(K_\infty=H(F_\infty)\).
Let \(P\) be a continuous right \(J\)-invariant function of a single
irreducible real type \(\rho\), of dimension \(d\). In the
restriction to \(K_0=U(W_0)(F_\infty)\) let \(\nu\) have
multiplicity \(m_\nu>0\) and dimension \(d_\nu\).
Set \(\theta=m_\nu d_\nu/d\).

Choose a rational line in the perpendicular of \(W_0\perp Ev_0\)
in \(C\). The homomorphism \(D\) of the norm-one torus into
\(H\) acting by its argument on that line and by the identity
on its orthogonal complement is a section of the determinant.
It centralizes the entire embedded \(U(S_0)\).
Write \(c_Q=(3Q-1)/(Q(Q+1))\), and put
\[
    j=\left\lceil{\log(2d^3/\omega)\over-\log c_Q}\right\rceil .
\]

\begin{theorem}[Detection with bounded denominators at one split place]\leavevmode
\label{thm:mix277}
There are a finite norm-one idele \(\delta\), a real rotation
\(s\in K_\infty\), and a translation \(t\) supported at \(v_*\)
with \(t_*^{\pm1}L_*\subset\pi_*^{-j}L_*\), such that
\[
 \left|d_\nu\int_{K_0}\overline{\chi_\nu(h)}
              P(hs,D(\delta)t)\,dh\right|
       \ \ge\ \sqrt{\theta/2}\,\|P\|_{L^2([H])}.
\]
The integer \(j\) is independent of the choice of \(P\) in its
type and level.
\end{theorem}
\begin{proof}
For \(x\in[H]\) average the squared real projection over
its right argument. With all compact groups given probability
measure set
\[
 U(x)=\int_{K_\infty}
   \left|d_\nu\int_{K_0}\overline{\chi_\nu(h)}P(xhs)\,dh\right|^2ds .
\]
On a finite class \(c\), write its real function as
\(p_c(b)=\operatorname{tr}(M_c\rho(b))\).
Then \(\|M_c\|_{\rm HS}=\sqrt d\,\|p_c\|_2\), with the
norm on the real group. The character integral replaces
\(\rho(h)\) by the orthogonal projection of rank \(m_\nu d_\nu\).
Orthogonality in \(s,b\), as in Theorem~\ref{thm:extract275},
gives mean \(\theta\|p_c\|_2^2\) for \(U\) on this class.
The trace inequality bounds the integrand before squaring by
\(\sqrt{d\,m_\nu d_\nu}\|p_c\|_2\). Since each class has mass
at least \(\omega\), these statements imply
\[
 \int_{[H]}U=\theta\|P\|_2^2,\qquad
 \|U\|_\infty\le d^2\theta\|P\|_2^2/\omega,\qquad
 \|U\|_2\le d\theta\|P\|_2^2/\sqrt{\omega}.
\]
The last inequality uses \(U\ge0\). In the real variable \(U\)
is a sum of matrix coefficients of \(\rho\otimes\bar\rho\)
and is invariant under the real center.

Let \(M\) be right convolution with probability on
\(K_*\operatorname{diag}(\pi_*,\pi_*^{-1},1,\ldots,1)K_*\).
Let \(\Pi\) be projection onto the functions fixed by both
\(\mathrm{SL}_N(F_{v_*})\) and \(K_*\). These projections
commute by normality. On local special-linear invariants \(M\)
equals averaging over \(K_*\), since the middle diagonal has
determinant one. On their orthogonal complement the three-coordinate
estimate in Lemma~\ref{lem:coefficient199} gives norm at most
\(c_Q\). Its proof applies to any unitary representation on this
complement and uses the residue size \(Q\). It follows that
\(\|(M^j-\Pi)U\|_2\le c_Q^j\|U\|_2\).

Both images are fixed by \(J\) and have their real coefficients
in the same matrix coefficient space as \(U\). That space has
dimension at most \(d^4\). On a compact group an orthonormal
basis of a translation-invariant coefficient space has constant
sum of squared moduli, equal to its dimension by integration.
Thus point evaluation on any one class costs at most
\(d^2/\sqrt{\omega}\) times the global \(L^2\) norm. In particular
\[
  |(M^j-\Pi)U(x)|
       \le c_Q^j d^3\theta\|P\|_2^2/\omega
       \le \tfrac12\theta\|P\|_2^2\qquad(x\in[H]).
\]

The continuous function \(\Pi U\) depends at finite places
only on the determinant and is constant in its real argument.
Indeed strong approximation for the simply connected group
\(\mathrm{SU}(C)\) holds away from \(v_*\), where the group is
split and noncompact. For \(g\in H(\mathbb A_F)\) and
\(a\in\mathrm{SU}(C)(\mathbb A_F)\), approximate \(gag^{-1}\)
away from \(v_*\) by rational special unitary elements. At
\(v_*\) the difference can be corrected exactly by a right
local special unitary element, since that group is normal.
Rational left invariance, local right invariance and continuity
give \((\Pi U)(ga)=(\Pi U)(g)\).
This is the strong approximation theorem applied in
Proposition~\ref{prop:corehecke272}. At infinity every unitary
matrix is a scalar times a special unitary matrix, by taking
an \(N\)-th root of its determinant. The real central
invariance of \(U\) passes through \(\Pi\).

Projection preserves the mean since it fixes the constants.
Some \(x=(1,D(\delta))\) therefore has
\(\Pi U(x)\ge\theta\|P\|_2^2\). At that point \(M^jU\) is
at least half this number. It is the mean of \(U(xt)\) over
a \(j\)-fold product of the indicated compact double coset.
Each factor and its inverse takes \(L_*\) into
\(\pi_*^{-1}L_*\). Every element in the product satisfies
the asserted inclusions. Taking one value at least the
mean in the combined compact averages over \(t,s\) proves
the theorem.
\end{proof}

\subsection{Matching rectangular tests at a split finite place}\label{sec:matchingframes334}

In a square of the theta comparison the two spaces of matrices have
different row ranks. Fix a split finite place of the period field and
write \(k\) for the local field, \(O\) for its integers, \(\pi\) for a
uniformizer, and \(q=|O/\pi O|\). Normalize \(|\pi|=q^{-1}\).
Additive measure gives \(O\) mass one. Multiplicative Haar measure gives
\(K_m=\mathrm{GL}_m(O)\) mass one in \(\mathrm{GL}_m(k)\).
Write \(M_{m,n}\) for the space of \(m\)-by-\(n\) matrices. An integral
matrix with \(m\) rows is \emph{primitive} if it has rank \(m\) modulo
\(\pi\). Equivalently it has an integral right inverse. A test of
\emph{level \(e\)} is supported on \(M_{m,n}(O)\) and is constant on
cosets of \(\pi^e M_{m,n}(O)\).

Take \(a>b\ge1\) with \(n=a+b\) even, \(z=n/2\), and
\(\delta=z-b\). Let \(\alpha,\beta\) be the smooth unitary splitting
characters on \(k^\times\), and put \(\nu=\alpha/\beta\).
Set \(X_0=(I_a,0)\), \(Y_0=(0,I_b)\). For \(h\in\mathrm{GL}_n(k)\)
write \(X_h=X_0h^{-t}\) and \(Y_h=Y_0h\), where \(h^{-t}\) is the
transpose of the inverse. The two sections in the linear coordinates are
\[
\begin{aligned}
 A_\Psi(Y)&=\int_{\mathrm{GL}_b(k)}\Psi(gY)\nu(\det g)|\det g|^z\,d^\times g,\\
 S_\Psi(h)&=\alpha(\det h)|\det h|^{b/2} A_\Psi(Y_h),\\
 B_\Phi(h)&=\beta(\det h)|\det h|^{-a/2}
     \int_{\mathrm{GL}_a(k)}\Phi(gX_h)\nu(\det g)|\det g|^z\,d^\times g.
\end{aligned}
\]
We first construct \(\Phi\) for which the last integral has compact
support in \(g\). We treat spherical normalization separately below.

\Needspace{5\baselineskip}
\begin{lemma}[A finite determinant average]\label{lem:ramifiedsum334}
Suppose that \(\nu|_{O^\times}\) has positive conductor exponent \(c\). Thus
\(c\) is the least \(j\ge0\) for which \(\nu\) is trivial on \(U_j\),
where \(U_0=O^\times\) and \(U_j=1+\pi^jO\) for \(j>0\).
Let \(\Psi\) have level \(e\ge c\). Write \(R_e=O/\pi^eO\),
\(\gamma_m=\prod_{j=1}^m(1-q^{-j})\), and \(H_m(e)=q^{em^2}\gamma_m\).
For \(T\in M_b(R_e)\) let \(0\le t_1\le\cdots\le t_b\le e\)
be its Smith exponents, counting each zero pivot as \(e\), and
\(s(T)=\sum_i t_i\). At every primitive \(Y\),
\[
 A_\Psi(Y)=\frac{1}{H_b(e)}
 \sum_{\substack{T\in M_b(R_e)\\t_b(T)\le e-c}}
  \Psi(\widetilde T Y)\nu(\det\widetilde T)\,q^{-\delta s(T)}.
\]
Each summand is independent of the lift \(\widetilde T\).
\end{lemma}
\begin{proof}
The right inverse of \(Y\) shows that \(gY\) is integral precisely
when \(g\) is integral. On invertible matrices
\(d^\times g=\gamma_b^{-1}|\det g|^{-b}\,dg\). Indeed multiplication on
either side has additive Jacobian \(|\det g|^b\), and the proportion of
invertible matrices over the residue field is \(\gamma_b\), by choosing
independent columns in order. Singular matrices have additive measure
zero, by induction and Fubini applied to the nonzero determinant
polynomial. Since \(\delta>0\), the resulting integral over integral
matrices converges absolutely. The same conclusion for any full-rank
\(Y\) follows by changing its row basis to a basis of the saturated
integral row lattice.

Partition integral \(g\) by reduction modulo \(\pi^e\). Each class has
mass \(q^{-eb^2}\), and \(\Psi(gY)\) is fixed on it. Integral invertible
row and column operations diagonalize a representative modulo \(\pi^e\).
If \(t_b\le e-c\), every lift is invertible over \(k\), has determinant
valuation \(s(T)\), and its inverse has entries in \(\pi^{-t_b}O\).
Changing the lift changes its determinant by a factor in \(U_{e-t_b}\).
If \(t_b>e-c\), choose \(u\in U_{e-t_b}\) with \(\nu(u)\ne1\). In
diagonal coordinates multiplying the last row by \(u\) preserves the
whole class, its additive measure and \(|\det g|\). The integral over
that class equals \(\nu(u)\) times itself and vanishes. Substitution of
the class masses gives the formula, since \(H_b(e)=|\mathrm{GL}_b(R_e)|\).
\end{proof}

\Needspace{5\baselineskip}
\begin{theorem}[A matched section at every group element]\label{thm:matchingframe334}
For each \(\Psi\) of the preceding lemma there is a linear choice
\(\Phi=\mathcal T\Psi\) of level \(e\), supported on primitive
\(a\)-frames, such that \(B_\Phi(h)=S_\Psi(h)\) for all \(h\).
On primitive-supported tests transforming under left \(K_b\) by
\(\nu^{-1}\circ\det\), this construction is an isomorphism onto the
corresponding \(K_a\)-covariant tests. Over the ring generated by the
values of \(\nu\) on units, it permutes basis vectors indexed by direct
summands and multiplies each by a root of unity.
\end{theorem}
\begin{proof}
Right invariance of Haar measure gives
\(A_\Psi(CY)=\nu(\det C)^{-1}|\det C|^{-z}A_\Psi(Y)\).
Every full-rank \(a\)-frame is \(X_h\) for some \(h\). Set
\[
       F(X_h)=\nu(\det h)|\det h|^z A_\Psi(Y_h).
\]
If \(h'\) gives the same \(X_h\), then \(h'h^{-1}\) has block form
\(\left(\begin{smallmatrix}I_a&*\\0&D\end{smallmatrix}\right)\),
as is seen by fixing its first \(a\) columns. Its determinant is
\(\det D\) and \(Y_{h'}=D Y_h\). The preceding covariance makes \(F\)
well defined. Replacing \(h\) by \(\operatorname{diag}(C^{-t},I_b)h\)
gives \(F(CX)=\nu(\det C)^{-1}|\det C|^{-z}F(X)\).

Set \(\Phi(X)=F(X)\) on primitive frames and zero elsewhere. For such
\(X\), completing its rows to an integral basis gives \(h\in K_n\).
If the first rows change by multiples of \(\pi^e\), keep the remaining
rows of the completion unchanged and take the inverse transpose. This
gives \(h'\equiv h\bmod\pi^e\). The unit character and the finite formula
in Lemma~\ref{lem:ramifiedsum334} give equal values at the two frames.
Primitivity is decided modulo \(\pi\), so \(\Phi\) has level \(e\).
For arbitrary \(X_h\), the primitive bases of its row space are exactly
the integral bases of the intersection of that space with \(O^n\).
Consequently the elements \(g\) making \(gX_h\) primitive form one
coset \(K_a g_0\), of mass one. On it the twisted integrand equals
\(F(X_h)\). The coset also lies in a single double coset of \(K_a\),
so summation by elementary-divisor shells gives the same value.
Multiplying by the prefactor gives \(S_\Psi(h)\), since
\(\beta\nu=\alpha\) and \(z-a/2=b/2\).

For a primitive-supported covariant \(\Psi\) and primitive \(Y\),
only \(g\in K_b\) contributes, and \(A_\Psi(Y)=\Psi(Y)\).
On \(K_n\) the formula is therefore
\(\Phi(X_h)=\nu(\det h)\Psi(Y_h)\). Modulo \(\pi^e\) the row spans
of \(X_h\) and \(Y_h\) are orthogonal complements for the standard
perfect pairing. Every free direct summand admits an extended basis,
and passage to its orthogonal complement is a bijection. A covariant
function on either set of frames is determined by one value on a basis
of each summand. Changing this basis multiplies the value by a root
of unity, since \(\nu\) on units factors through \(R_e^\times\).
This proves the assertion over the specified ring, including invertibility.
\end{proof}

\Needspace{5\baselineskip}
\begin{proposition}[The two normalized spherical tests]\label{prop:sphericalframe334}
At a split place with unramified \(\alpha,\beta\), take the integral
lattice indicator in each rank \(m=a,b\). Replace \(z\) inside the
determinant integral by \(s\), multiply it by
\(c_m(s)=\prod_{j=1}^m(1-\nu(\pi)q^{m-j-s})\), and continue to \(s=z\).
The resulting \(B\) and \(S\) sections agree at every \(h\) and equal
one at the identity.
\end{proposition}
\begin{proof}
For additive probability on \(M_m(O)\) the generating function of
\(v(\det G)\) equals
\[
       \int u^{v(\det G)}\,dG=\prod_{j=1}^m
                  \frac{1-q^{-j}}{1-uq^{-j}}\qquad (|u|<q).
\]
The least valuation of the entries in the first column is \(r\) with
probability \((1-q^{-m})q^{-mr}\). Move the resulting primitive column
to the first basis vector by an integral invertible change of basis.
Conditionally on this column the other columns still have independent
additive Haar distribution. Expanding the determinant removes the first
row and column, adds \(r\) to its valuation, and proves the product by
induction. Substituting \(u=\nu(\pi)q^{m-s}\) and dividing by
\(\gamma_m\) gives \(c_m(s)^{-1}\) on a primitive frame for
\(\Re s>m-1\). This is the determinant integral of the matrix zeta
construction of Godement and Jacquet~\cite{XVII-XII-gj334}, with the present
Haar measures.

For arbitrary \(h\), choose \(k_0\in K_n\) whose bottom rows span the
saturated bottom row lattice of \(h\). Then
\(h=\left(\begin{smallmatrix}A&*\\0&D\end{smallmatrix}\right)k_0\),
so \(X_h=A^{-t}X_{k_0}\) and \(Y_h=D Y_{k_0}\).
The changes of basis in the two integrals give the normalized factors
\(\nu(\det A)|\det A|^s\) and \(\nu(\det D)^{-1}|\det D|^{-s}\).
Both are entire in \(s\). At \(s=z\), multiplying by the defining
prefactors gives equal sections, using \(|\det k_0|=1\),
\(\nu(\det k_0)=1\), and \(a+b=2z\).
\end{proof}


\subsection{A fixed functional on the matched frames}\label{sec:signedrow335}

A global class sum can prescribe different, even opposite, coefficients
on different local frames. We compute the local effect of matching on
each such functional. Retain the linear charts, \(c\ge1,\ e\ge c\),
and the level ring \(R_e\) of Section~\ref{sec:matchingframes334}.
Write \(\mathcal G\) for the rank-\(b\) direct summands of \(R_e^n\).
For each \(W\in\mathcal G\) choose \(k_W\in K_n\) such that the rows
of \(Y_W=Y_{k_W}\) give a basis of \(W\) after reduction. Let
\(X_W=X_{k_W}\). These large frames represent the complementary
summands. Evaluation at them gives coordinates on the primitive
covariant output of the transfer \(\mathcal T\).

\Needspace{5\baselineskip}
\begin{theorem}[Coefficients of a prescribed functional]\label{thm:signedrow335}
Fix scalars \(\lambda_W\) and put
\(\Lambda(\Phi)=\sum_{W\in\mathcal G}\lambda_W\Phi(X_W)\).
Write \(H=|\mathrm{GL}_b(R_e)|\) and
\(u=\nu(\pi)q^{-\delta}\). Define \(D_Z\) by
\(H\Lambda(\mathcal T\Psi)=\sum_{Z\in M_{b,n}(R_e)}D_Z\Psi(Z)\).
Give a rectangular matrix \(Z\) its \(b\) Smith exponents
\(0\le t_1\le\cdots\le t_b\le e\), with \(e\) for a zero pivot.
Then \(D_Z=0\) if \(t_b>e-c\). Otherwise, put \(s=\sum t_i\). For
each \(W\) containing the rows of \(Z\), write \(Z=T_{Z,W}Y_W\bmod\pi^e\).
There is a unique \(T_{Z,W}\in M_b(R_e)\), and
\[
 \begin{split}
 D_Z&=u^s\sum_{W\supseteq\operatorname{row}(Z)}
                     \epsilon_{W,Z}\lambda_W,\\
 \epsilon_{W,Z}&=\nu(\det k_W)\,
          \nu\bigl(\pi^{-s}\det\widetilde T_{Z,W}\bigr).
 \end{split}
\]
Each \(\epsilon_{W,Z}\) is a well-defined root of unity. For primitive
\(Z\) the sum has one term. Its coefficient is a unit multiple of
\(\lambda_{\operatorname{row}(Z)}\).
\end{theorem}
\begin{proof}
Insert the finite sum of Lemma~\ref{lem:ramifiedsum334} in the evaluation
of Theorem~\ref{thm:matchingframe334} at \(X_W\). The map
\(T\mapsto TY_W\bmod\pi^e\) is injective, since \(Y_W\) has an integral
right inverse. Its image consists of all matrices whose rows lie in \(W\),
by expansion in the chosen basis. An invertible column operation on
\(R_e^n\) takes \(Y_W\) to \((I_b,0)\), hence the Smith pivots of \(Z\)
are those of \(T\). The determinant of any contributing lift is
\(\pi^s\) times a unit. The independence of the unit character from
this lift was proved in Lemma~\ref{lem:ramifiedsum334}. Collecting the
coefficient of \(\Psi(Z)\) gives the displayed formula. A primitive
\(Z\) spans a direct summand with cardinality \(|R_e|^b\).
A rank-\(b\) direct summand containing it has the same cardinality,
so equals it. Every \(W\) occurs this way, for instance at \(Y_W\).
\end{proof}

\Needspace{5\baselineskip}
\begin{corollary}[A lower bound for every functional on the image]\label{cor:dualnorm335}
Give functionals the norm dual to the additive \(L^2\) norm, restricting
the large-side functional to the covariant output space. Then
\[
 \begin{split}
   \|\Lambda\mathcal T\|_*^2
    =\frac{p_a}{p_b}\|\Lambda\|_*^2+
       \frac{q^{ebn}}{H^2}
       \sum_{\substack{Z\ {\rm imprimitive}\\t_b(Z)\le e-c}} |D_Z|^2 .
 \end{split}
\]
Thus every nonzero functional on the matched frames survives pullback.
At \(e=c\) the squared norm ratio is \(p_a/p_b\) for every such
functional. In the original partial Fourier models the same lower
bound and the same minimal-level ratio hold, with the functional
restricted to the actual image.
\end{corollary}
\begin{proof}
Put \(G=|\mathcal G|\). The images of bases of an individual
\(W\) account for \(H\) matrices. Thus \(p_b=GHq^{-ebn}\).
On the large side all \(G\) complementary summands have the same
measure and changes of basis have character of absolute value one.
Consequently \(\|\Lambda\|_*^2=(G/p_a)\sum_W|\lambda_W|^2\).
Each small residue class has measure \(q^{-ebn}\), which makes its
evaluation functional have squared norm \(q^{ebn}\), and the evaluations
at different classes are orthogonal. The primitive terms in
\(q^{ebn}H^{-2}\sum_Z|D_Z|^2\) contribute
\(q^{ebn}H^{-1}\sum_W|\lambda_W|^2=(p_a/p_b)\|\Lambda\|_*^2\).
All other terms are as displayed. In particular the rows can have any
prescribed signs. The image is the full covariant space because the
transfer already maps the small primitive covariant subspace onto it.
The two changes of Fourier chart in Theorem~\ref{thm:mixedfourier334}
are scalar multiples of isometries with the product of their moduli,
including the scalar \(\sigma_P\), equal to one. Conjugating the
transfer by the isometries gives the stated bound on the transformed
image.
\end{proof}

\Needspace{5\baselineskip}
\begin{theorem}[Compact translates at the conductor level]\label{thm:compactrow335}
Let \(\mathcal H_e\) be the covariant space on primitive \(b\)-frames
at level \(e\), with its invariant Hilbert norm. Let
\(G_e=\mathrm{GL}_n(R_e)\) act by right changes of coordinates, and
let \(P_c\) be the average of the kernel of \(G_e\to G_c\). Its image
is the level-\(c\) frame space, of dimension
\[
         R=q^{(c-1)ab}\prod_{i=0}^{b-1}
                         \frac{q^n-q^i}{q^b-q^i}.
\]
This representation of \(G_c\) is irreducible. For every individual
\(f\in\mathcal H_e\) and every functional \(\Lambda\),
\[
 \frac1{|G_e|}\sum_{g\in G_e}|\Lambda(g f)|^2
       \ \ge\ R^{-1}\|\Lambda P_c\|_*^2\,\|P_cf\|^2 .
\]
At \(e=c\) this is an equality. At that level the character projection
for the full stabilizer of any specified \(W\), with the character by
which it acts on the basis line of \(W\), is the coordinate projection
onto that one line. The same assertions hold on the matched large
frame space with the induced action.
\end{theorem}
\begin{proof}
At level \(c\) there is one basis line for every summand \(W\).
The stabilizer acts on that line by the determinant character of its
action on \(W\), or its inverse according to the action on functions.
For distinct summands \(W,W'\) take \(x\in W'\) whose image modulo
\(W\) is nonzero. Since \(W\) is a direct summand, there is a linear
form \(l\) on \(R_c^n\) vanishing on \(W\) for which \(l(x)\) has
valuation \(t<c\). The endomorphism \(g=1+d\,x l\) fixes \(W\)
pointwise and preserves \(W'\). On both \(W'\) and the ambient module
its determinant is \(1+d\,l(x)\), by the rank-one determinant formula.
Choose \(d\in R_c\) so that this determinant is a unit with nontrivial
\(\nu\)-value. Such a unit exists in \(1+\pi^tR_c\) by the definition
of the conductor, using any unit of nontrivial character when \(t=0\).
Then \(g\) is invertible and acts by different phases on the two
basis lines. The \(W,W'\) entry of every commuting endomorphism
must vanish. Transitivity on the summands makes its diagonal entries
equal. Since an invariant subspace of a unitary representation has
a commuting orthogonal projection, the representation is irreducible.
For the stabilizer of a fixed \(W\) the same element \(g\) forces
the \(W'\) coordinate of any vector of its specified character to vanish
for each \(W'\ne W\). The \(W\) basis vector has that character.
Averaging the action against the inverse character is the orthogonal
projection onto this eigenspace and hence keeps exactly the \(W\)
coordinate. In particular this projector preserves the integral frame
lattice even at primes dividing the stabilizer order.

The congruence kernel acts transitively on the lifts of each primitive
frame modulo \(\pi^c\). Indeed extend a frame to an invertible matrix
and change its first rows to the prescribed lift, keeping the remaining
rows fixed. The change of coordinates is the identity modulo \(\pi^c\).
This proves the description of \(P_c\). At residue level one the number
of summands is the number of independent \(b\)-tuples divided by
the number of bases of a \(b\)-space. Every such summand has
\(q^{(c-1)ab}\) lifts, by writing a lift as a graph into a complementary
free \(a\)-space. This gives \(R\).

At level \(c\) the average of the operators
\(v\mapsto\langle v,g f\rangle\,g f\) (inner product linear in the
first variable) commutes with \(G_c\). Its trace is \(\|f\|^2\),
so its scalar is \(\|f\|^2/R\). Evaluating on the Riesz vector of
\(\Lambda\) proves the equality. For \(e>c\) fix a coset of the
congruence kernel and apply the squared-mean inequality to the values
\(\Lambda(g f)\) in this coset. Their mean is \(\Lambda(gP_cf)\).
The remaining average is over \(G_c\), to which the equality applies
with row \(\Lambda P_c\). The correspondence of complementary frames
transports the representation to the large side, up to the common
unitary determinant character in Theorem~\ref{thm:matchingframe334}.
Tensoring with that character preserves irreducibility and the squared
absolute values.
\end{proof}


\subsection{The cost of the prescribed terminal Hecke modification}
\label{sec:neighbornorm356}

Fix one ordinary word with the notation of
Theorem~\ref{thm:neighbormove356}. Define \(S_T'\) by taking
\(H_w^mS_w\) at every \(w\mid T\) and taking \(S_w\) elsewhere.
This is the terminal-prime prescription before the additional
scalar rescaling of the source. The local lattices define a
full global sublattice of \(S\). Indeed they have finite
index at finitely many places. Reduction modulo a common
integer annihilating the local quotients and the Chinese
remainder theorem realize their intersection with exactly
the specified localizations. Write \(h\) for the number of
high embeddings. At any finite place over a rational prime
\(p\), let \(v_w(p)=1\).

\begin{theorem}[Absolute index of the required terminal input lattice]\leavevmode
\label{thm:neighborindex356}
For the coded weights on the common period intervals,
\[
 \frac{\log [S:S_T']}{[J:\mathbb Q]d}
     =\frac{m\log T}{2}\!
       \left(h b_*+\frac{h(k-1)}2+
              \frac1k\sum_{\tau\ {\rm high}}|\lambda_\tau|\right)
     \ll_{\psi,\eps} b\,x^{-1/2-9\eps/8}.
\]
Here the index is the cardinality of the underlying finite
abelian quotient. Let \(\mathcal P\) be the joint signed
arithmetic period on any selection of words and translates
at this \(x\) taking values in one test space. In the following
inequality use the direct sums of their respective ordinary
lattices \(S\) and terminal modifications \(S'\), and their
total source dimension \(D\). Then
\[
       0\ \le\
       \frac{\log[\mathcal P(S):\mathcal P(S')]}{[J:\mathbb Q]D}
           \ \ll_{\psi,\eps}\ b\,x^{-1/2-9\eps/8}.
\]
The field, the terminal primes and all the signed maps here are
the ones used for these columns. No lower bound for their image
dimensions is required.
\end{theorem}
\begin{proof}
Apply the determinant formula of
Theorem~\ref{thm:neighbormove356} to the \(m\)-th power,
which multiplies its valuation by \(m\) even when the
elementary divisors of the power change their ordering.
The logarithm of the abelian index is the sum of these
valuations times \([J_w:\mathbb Q_T]\log T\), by Smith
normal form and the residue degrees. A given high \(\tau\)
occurs exactly when \(v_w(\tau(\pi))=1\). This value is
either zero or one since \(\pi\bar\pi=T\) is split.
Its average with weights \([J_w:\mathbb Q_T]/[J:\mathbb Q]\)
is \(1/2\). Indeed \(\tau(\pi)\) is a unit away from
\(T\) and each archimedean conjugate has magnitude
\(\sqrt T\). Taking its absolute norm to \(\mathbb Q\)
proves the assertion. This gives the exact index.

At a high embedding \(\lambda_1=a\), \(\lambda_k=-b_*\),
and every interior entry is polynomially bounded in \(x\)
by \eqref{eq:letters111}. Thus \(|\lambda_\tau|/k\le
a/k+C(x)\) for some fixed polynomial \(C(x)\). The
index contributions of \(b_*\) and of any fixed polynomial,
after multiplication by \(m f\log T\), are \(o(b/x^A)\)
for every fixed \(A>0\). The remaining contribution is
at most a constant times \(m f a\log T/k\), since \(h\le f\).
By the rank and prime estimates in the proof of
Theorem~\ref{thm:imagedet355}, this is
\[
       O_{\psi,\eps}\bigl(m b\,x^{-1-17\eps/16}\bigr)
                 =O_{\psi,\eps}\bigl(b\,x^{-1/2-9\eps/8}\bigr).
\]
The bounds are uniform in the selected columns. On their
direct sum the source indices multiply, so the normalized
logarithmic index is their dimension-weighted average.
Finally the map on the finite quotients
\(S/S'\longrightarrow\mathcal P(S)/\mathcal P(S')\)
is surjective. This proves the second inequality for the
actual signed map, including zero image.
\end{proof}

\begin{proposition}[The cost per image dimension at the terminal prime]\leavevmode
\label{prop:neighborrank356}
For one word take any map \(B\) on its ordinary initial
space, of image dimension \(1\le j\le d\). At \(w\mid T\)
let \(e_w(B)\) be the valuation of the determinant of
\(B(H_w^m S_w)\) in \(B(S_w)\). With \(s=|A_w|\),
\[
  0\le \frac{e_w(B)}{j}
     \le m\left[
        s\{b_*+k-1+C(x)\}
       +\frac{2(a+k)}{k-1}\{s\log2+\log(d/j)\}\right],
\]
where \(C(x)\) may be chosen as in
Theorem~\ref{thm:imagedet355}. In particular for
\(\log(d/j)\le x^{2+8\eps}\) the logarithmic cost per
image dimension is \(O_{\psi,\eps}(b\,x^{-1/2-9\eps/8})\).
This applies to the signed column on the prescribed tests.
\end{proposition}
\begin{proof}
In the reindexing of Theorem~\ref{thm:neighbormove356},
first use the reference coefficient lattice on the vector
markings. The product incidence and its transpose are
integral on the augmentation lattices, and their product
in both orders is \(T^{s(k-1)}I\). Hence passing from this
reference to the target by incidence shifts every increasing
source exponent upward by a number in \([0,s(k-1)]\).
This is exactly the min--max argument in
Theorem~\ref{thm:fourierslopes354}. Relative to the reference,
the exponents are sums of \(s\) independently sampled
coordinate weights, counted with their Weyl multiplicities.
The scalar in \(H_w\) adds \(s b_*\). With
\(t=(k-1)/(2(a+k))\), Lemma~\ref{lem:oneweight355} therefore
bounds the average of \(e^{t\gamma}\), for these single-step
exponents \(\gamma\), by
\[
        \exp\{t s(b_*+k-1+C(x))\}\,2^s .
\]
Their lower endpoint is nonnegative by integrality.
Apply the unit-minor argument in the proof of
Theorem~\ref{thm:imagedet355} to any rank \(j\) map on
this space. It bounds the change of its image determinant
per dimension by the mean of the largest \(j\) of these
exponents. Jensen's inequality bounds that mean by the
logarithm of the last display plus \(\log(d/j)\), all
divided by \(t\). Now use this single-step bound for
\(B,BH_w,\ldots,BH_w^{m-1}\). They all have the same
rank because \(H_w\) is invertible over the local field.
Their successive image indices telescope. This proves
the displayed inequality. Multiplying by \(\log T\) and
using \(s\le f\) and the same field parameters as in
Theorem~\ref{thm:imagedet355} gives its last assertion.
\end{proof}

In particular the prescribed terminal Hecke modification of
the input has a controlled index after projection. The
unmodified image can itself have a large index in the
designated test lattice, or have small rank. Its comparison
with the analytic Bessel pairing and the bounds at the other
coefficient primes required by Section~\ref{sec:missing}
must also be established.


\subsection{The determinant of the normalized input on a common terminal prime}
\label{sec:normalizedimage357}

For a single word write \(S\) for its ordinary initial lattice
of rank \(d\). Let \(S^\#_T\) be the lattice obtained from
\(S\) by the Hecke operation and the scalar normalization at
its terminal prime, with \(S\) at the other rational primes.
Use \(A\) and the common embedding of \(K_d\) of
Theorem~\ref{thm:terminalscalar357}. For two nested full
lattices \(L'\subset L\) in one \(J\)-space, let
\(\mathfrak d(L',L)\) be the integral ideal satisfying
\(\det L'=\mathfrak d(L',L)\det L\). Locally it is generated
by the determinant of any basis of \(L'\) in a basis of \(L\).

\begin{proposition}[The determinant ideal of the normalized source]\leavevmode
\label{prop:terminalideal357}
With the initial weights \(\lambda_\tau\), put
\[
 \theta=m\left(hb_*+\frac{h(k-1)}2+
                       \frac{1}{k}\sum_{\tau\ {\rm high}}|\lambda_\tau|\right).
\]
Then \(d\theta\) is a nonnegative integer, and
\[
   \mathfrak d(S'_T,S)=(\pi)^{d\theta},\qquad
   \mathfrak d(S^\#_T,S)=(\pi)^{d(A+\theta)}.
\]
In particular these equalities hold in the coefficient integers
even after a ramified extension above \(T\).
\end{proposition}
\begin{proof}
All coefficient embeddings in the type restrict to one embedding
of \(K_d\). When \(w\mid\pi\), every high embedding therefore
induces a place dividing the original \(\pi\); when
\(w\mid\bar\pi\), none does. Above \(\pi\) the induced places of
the high embeddings are pairwise distinct by complete splitting
of \(T\) in \(E_x\). Theorem~\ref{thm:neighbormove356} gives
\(v_w\det H_w^m=d\theta\) above \(\pi\) and zero above
\(\bar\pi\). Integrality of \(H_w\) gives nonnegativity.

The scalar part \(dmhb_*\) is integral, as is
\(dmh(k-1)/2\) since \(m\) is even. For the remaining
part, on each tensor of coefficient modules the mean of a
single diagonal coordinate at \(\tau\) is
\(|\lambda_\tau|/k\). Its product with the dimension of the
tensor is the sum of integer weights with multiplicity.
In the initial space each such tensor is repeated with the
integer multiplicities of the class and scalar-augmentation
bases from Theorem~\ref{thm:neighbormove356}. Hence
\(d|\lambda_\tau|/k\) is integral. Returning to the ordinary
discrete valuation gives \(e_w d\theta\) above \(\pi\),
exactly the order of \(\pi^{d\theta}\). Localization away
from \(T\) gives the unit ideal by definition. This proves
the first identity. Multiplication of a rank-\(d\) lattice
by \(\pi^A\) multiplies its determinant by \(\pi^{Ad}\).
Theorem~\ref{thm:terminalscalar357} now gives the second.
\end{proof}

\begin{theorem}[Removing the common terminal scalar after signed projection]\leavevmode
\label{thm:normalizedimage357}
At fixed \(x\) take any selected words and translates sharing
one generator \(\pi\), with signed arithmetic period maps to a
common \(J\)-space. Denote their joint map by \(P\), their
total source dimension by \(D\), and its image dimension by
\(j\). Use direct sums for the three input lattices. Then
\[
 \frac{\log[P(S):P(S^\#_T)]}{[J:\mathbb Q]}
       =\frac{Aj}{2}\log T+
                    \frac{\log[P(S):P(S'_T)]}{[J:\mathbb Q]} .
\]
The last summand lies between zero and
\(O_{\psi,\varepsilon}(D b\,x^{-1/2-9\varepsilon/8})\).
For the selected fixed-weight words of one translate, write
\(\delta=\sum_\tau\delta_\tau\) for the normal degree of
Corollary~\ref{cor:normaldegree347}, and \(P_0\) for the
same projected class sum with the normal-coordinate scalar
removed. On these words the joint lattice identity is
\[
                  P(S^\#_T)=\pi^A\bar\pi^{-\delta}\,P_0(S'_T).
\]
Consequently the absolute scalar on its \(j\)-th exterior
power at each complex embedding is \(T^{j(A-\delta)/2}\).
\end{theorem}
\begin{proof}
Theorem~\ref{thm:terminalscalar357} gives the same integral
scalar \(\pi^A\) on every summand of the joint domain. Linearity
gives \(P(S^\#_T)=\pi^A P(S'_T)\) after taking the sum of
the images. The relative index of the latter scalar multiplication
on a rank-\(j\) lattice is
\(|N_{J/\mathbb Q}\pi|^{Aj}\). This follows locally from the
determinant, also for a projective lattice without a global basis.
All archimedean conjugates of \(\pi\) have absolute value
\(\sqrt T\). Taking logarithms and using multiplicativity of
the index gives the first identity. The bounds for the residual
index are those of Theorem~\ref{thm:neighborindex356}.

At one translate the prescribed words all have the same
Hamming weight. Their last restriction includes
\(\prod_{\tau}\tau(\bar\pi)^{-\delta_\tau}=\bar\pi^{-\delta}\)
with the same \(\delta\) on every summand by
Corollary~\ref{cor:normaldegree347}. All remaining factors
of that restriction, the earlier theta maps and the tests
belong to \(P_0\). Pulling the common scalar through the
class sum proves the lattice identity. Its absolute value at
infinity follows from \(|\pi|=|\bar\pi|=\sqrt T\) there.
All identities include the zero image by taking the degree-zero
determinant to be one.
\end{proof}

The determinant just computed compares two inputs to the same
signed arithmetic period. For Section~\ref{sec:missing} one still
needs the size and rank of its image in the designated analytic
Bessel test lattice, and the bound on the residual strict groups
in Proposition~\ref{prop:res23}.

\clearpage
\section{References}\label{nav:bibliography-XVII}
\begingroup\def\section*#1{}\begin{thebibliography}{99}\small
\bibitem{XVII-XII-IV-sv163} W. Schmid and K. Vilonen, \emph{Two geometric character formulas for reductive Lie groups}, J. Amer. Math. Soc. \textbf{11} (1998), 799--867. \href{https://arxiv.org/pdf/math/9801081}{Sections 7--8, especially (8.3) and Theorem 8.5}.
\bibitem{XVII-XII-IV-molev106} A. I. Molev, \emph{Gelfand--Tsetlin bases for classical Lie algebras}, Handbook of Algebra, vol.~4 (2006), 109--170. \href{https://www.maths.usyd.edu.au/u/pubs/publist/preprints/2002/molev-18.pdf}{Theorem 2.1}.
\bibitem{XVII-XII-I-xue108} H. Xue, \emph{Bessel models for unitary groups and Schwartz homology}. \href{https://preprint.press.jhu.edu/ajm/sites/default/files/AJM-h-xue.pdf}{Theorem 1.1 and Section 1.2}.
\bibitem{XVII-XII-I-atobe108} H. Atobe, \emph{On the non-vanishing of theta liftings of tempered representations of \(U(p,q)\)}, Adv. Math. \textbf{363} (2020), 106984. \href{https://eprints.lib.hokudai.ac.jp/repo/huscap/all/84508/}{Appendix A}.
\bibitem{XVII-XII-I-ggp108} W. T. Gan, B. H. Gross and D. Prasad, \emph{Restrictions of representations of classical groups: examples}. \href{https://abel.math.harvard.edu/~gross/preprints/ggp-evidence4.pdf}{Section 2}.
\bibitem{XVII-XII-gj334} R. Godement and H. Jacquet, \emph{Zeta Functions of Simple Algebras}, Lecture Notes in Mathematics 260, Springer, 1972. \href{https://doi.org/10.1007/BFb0070263}{Monograph}.
\end{thebibliography}\endgroup
\end{document}
